% 附录 C — 文件定义（精译重写）
\biappendix{文件定义}{File Definitions}
\label{chap:files}

\section{引言}
\subsection*{Introduction}

本附录描述 RedHawk 支持的内部输入/输出文件格式与内容。
RedHawk 输入文件可在文件名前加设计名前缀；所有此类文件中，行首 \texttt{\#} 视为注释。

\begin{noteBox}
RedHawk 可直接读取 gzip 压缩格式（\texttt{*.gz}）的 LEF、DEF、STA、SPEF、VCD、FSDB 与 GDS 文件。
\end{noteBox}

\section{RedHawk 输入文件}
\subsection*{RedHawk Input Files}

\begin{enumerate}
\item 工艺文件（\texttt{*.tech}）
\item 全局开关配置 GSC 文件（\texttt{*.gsc}）
\item 全局系统需求 GSR 文件（\texttt{*.gsr}）
\item 焊盘实例/单元/位置文件（\texttt{*.pad}、\texttt{*.pcell}、\texttt{*.ploc}）
\item 库工艺文件（LEF 等，路径在 GSR 中引用）
\item 设计网表文件（DEF 等）
\item Synopsys 库文件（\texttt{*.lib}）
\end{enumerate}

下列文件信息须在 GSR 中引用路径。

\section{RedHawk 输出文件}
\subsection*{RedHawk Output Files}

分析各阶段均会生成结果文件与图形显示。详细报告说明见第~6~章及各主题章节。

\section{关键字语法约定}
\subsection*{Keyword Syntax Conventions}

RedHawk 命令、选项与关键字语法约定如下：
\begin{itemize}
\item \texttt{<x>} — 需指定的变量或取值
\item \texttt{[a|b|c]} — 必须选择 a、b、c 之一
\item \texttt{?x?} — 可选参数 x
\item \texttt{abcd} — 参数 a、b、c、d 均须指定
\item \texttt{<x> ...} — 可追加多个与 x 同类型的元素
\item \texttt{\{ \{a b c ...\} \{j k l ...\} ...\}} — 可追加多组相似元素集合
\item \texttt{+-*/} — 加减乘除算术运算符
\end{itemize}

\section{Apache 工艺文件（\texttt{*.tech}）}
\subsection*{Apache Technology File (*.tech)}

RedHawk 工艺文件为每个 IC 工艺提供金属层、过孔、介质、衬底及 bump 等参数。
需为每个工艺节点单独准备一份 \texttt{.tech} 文件。主要内容包括：
\begin{itemize}
\item 各金属层名称（与 LEF 一致）、厚度、电阻率、温度系数与 EM 限值
\item 过孔名称、单孔电阻与 via EM 限值
\item wire-bond 或 flip-chip bump 的 R/L/C
\item 介质层厚度、高度与介电常数
\item 衬底层厚度与电阻率
\end{itemize}

\subsection{单位与前缀}
\subsubsection*{Unit Prefix Conventions}

单位区分大小写。支持 t/g/M/k/m/u/n/p/f 等前缀（如 \texttt{1.3p = 1.3e-12}，数字与单位间无空格）。
长度换算：\texttt{1 mil = 0.001 inch = 2.54e-5 m}；\texttt{1 inch = 2.54e-2 m}；\texttt{1 micron = 1e-6 m}。

\subsection{工艺文件加解密}
\subsubsection*{Encrypting and Decrypting a Tech File}

\textbf{全文件加密：}使用 \texttt{bin} 目录下 \texttt{techEncrypt}/\texttt{techDecrypt}。
\begin{lstlisting}
techEncrypt <tech_filename>    # 生成 <tech_filename>.enc.tech
techDecrypt -i <encrypted_tech_file> -o <tech_file>
\end{lstlisting}

\textbf{部分加密：}在 \texttt{\#ENCRYPT\_START} 与 \texttt{\#ENCRYPT\_END} 注释行之间自动加密敏感 RC/高度/介电常数段。

\subsection{工艺文件关键字}
\subsubsection*{Technology File Keywords}

\begin{noteBox}
关键字与选项名不区分大小写。
\end{noteBox}

\subsubsection{\texttt{DIELECTRIC}}
\paragraph*{DIELECTRIC}

定义介质层厚度、介电常数及高度等参数；电容/电感提取时各层均须定义。（可选；默认：, dimensions are in microns）

\textbf{语法：}
\begin{lstlisting}
DIELECTRIC <dielectric_layer_name> {
constant <ER_value>
thickness <value>
[ Height <value> |
Above <dielectric_layer_name> ]
ER_VS_SI_SPACING {
(<spacing_1>, <ER1>)
...
(<spacing_n>, <ERn>)
}
}
\end{lstlisting}

\textbf{参数说明：}
\begin{itemize}
\item \texttt{dielectric <dielectric\_layer\_name>} — 定义介质层厚度、介电常数及高度等参数；电容/电感提取时各层均须定义。
\end{itemize}

\textbf{示例：}
\begin{lstlisting}
dielectric IMD3b {
thickness 0.22000
constant 2.9
ER_VS_SI_SPACING {
{ 0.2 2.8 }
{ 0.21 2.9 }
{ 0.23 2.99 }
{ 0.25 2.6 }
{ 0.26 3.0 }
}
above IMD3a
}
\end{lstlisting}

\subsubsection{\texttt{EM_PEAK_EQUATION_SOURCE_TECHFILE}}
\paragraph*{EM_PEAK_EQUATION_SOURCE_TECHFILE}

允许代工厂在 Peak EM 限值方程中内置自定义 Duty Ratio 因子。

\subsubsection{\texttt{EM_RULE_SET}}
\paragraph*{EM_RULE_SET}

在同一 tech 文件中定义多套 EM 规则集，供 EM_MODE 按角或分析模式选用。

\textbf{参数说明：}
在同一 tech 文件中定义多套 EM 规则集，供 EM\_MODE 按角或分析模式选用。

\subsubsection{\texttt{EM_TECH_FILE}}
\paragraph*{EM_TECH_FILE}

引用外部 EM 规则 tech 文件。

\textbf{语法：}
\begin{lstlisting}
Syntax
EM_TECH_FILE <EM rule set file>
\end{lstlisting}

\subsubsection{\texttt{HALF_NODE_SCALE_FACTOR}}
\paragraph*{HALF_NODE_SCALE_FACTOR}

半节点缩放因子，用于工艺尺寸缩放。（可选；默认：no scaling）

\textbf{语法：}
\begin{lstlisting}
HALF_NODE_SCALE_FACTOR <factor>
\end{lstlisting}

\textbf{示例：}
\begin{lstlisting}
HALF_NODE_SCALE_FACTOR 0.8
\end{lstlisting}

\subsubsection{\texttt{METAL}}
\paragraph*{METAL}

定义金属层电阻率、厚度、温度系数及 EM 限值等。（默认：is foundry-dependent；必需） 注意：定义金属层电阻率、厚度、温度系数及 EM 限值等。

\textbf{语法：}
\begin{lstlisting}
METAL <metal_layer_name> {
? MINWIDTH <drawn_wire_width_um> ?
? MINSPACE <wire_spacing_um> ?
WIDTH_ADJUST <adjustment>
? WIDTH-SPACE {
{<w1> <spacing1>}
{<w2> <spacing2>}
...
}?
Pitch <value of length>
Thickness <t_value>
Resistance <rpsq>
RHO <value_Ohm-microns>
CAP_DENSITY <value_fF/um^2> <associated_layer>
EM <max_wire_current_density>
? WIDTH_SI_VS_WIDTH_AND_SPACING ?
? T_EM <final_temp> ?
TNOM_EM <nom_temp>
EM_TEMP_RATING {
{ <temp1> <derating_factor1> }
{ <temp2> <derating_factor2> }
...
}
}
? EM_adjust <wire_width_adj_um> ?
\end{lstlisting}

\textbf{参数说明：}
定义金属层电阻率、厚度、温度系数及 EM 限值等。（默认：0) Coeff\_RT2 <R variation-second order c）

\textbf{示例：}
\begin{lstlisting}
POLYNOMIAL_BASED_EM_DC {
COND_RULE { L>10 AND w>0.18 AND (UPSTREAM) }
EM_POLYNOMIAL { 0.02*w IMAX sqrt(3*(w -0.01)^2)
COND_RULE { L>10 AND w>0.18 AND (DOWNSTREAM) }
EM_POLYNOMIAL { 0.04*w IMAX sqrt(4*(w -0.01)^2)
}
POLYNOMIAL_BASED_EM_PEAK {
COND_RULE { W/L > 2 AND W/L < 10 }
EM_POLYNOMIAL { 10 * ( w - 0.1 ) }
}
POLYNOMIAL_BASED_EM_RMS {
COND_RULE {L <= 10 }
EM_POLYNOMIAL { sqrt( 1 * delta_T *( w - 0.1 ) ^ 2 )
}
COND_RULE { L >= 20.0 AND w > 0.18 }
EM_POLYNOMIAL {1.3*w IMAX DERATE(min((w* 43.6/l),
( 0.06*l/25 )) )
EM_POLYNOMIAL { ANTI_DERATE ( 0.01*w ) IMAX (
DERATE ( MIN (( 0.001 + 0.00223086*w ),
( 0.032/w )))
}
}
METAL metal7
{
THICKNESS          0.7
RESISTANCE         0.025
EM                 2.4
EM_ADJUST          0.02
BLECH_JLC          250
ABOVE              diel
}
...
POLYNOMIAL_BASED_EM_RMS {
LENGTH_RANGES { 1e38 }
EM_POLYNOMIAL {sqrt(15.00*delta_T*(w - 0.015)^2*
(w - 0.015 + 0.155 )/(w - 0.015 + 0.055)) }
}
- Using EM_POLYNOMIAL format
...
POLYNOMIAL_BASED_EM_DC {
\end{lstlisting}

\begin{noteBox}
定义金属层电阻率、厚度、温度系数及 EM 限值等。
\end{noteBox}

\begin{noteBox}
定义金属层电阻率、厚度、温度系数及 EM 限值等。
\end{noteBox}

\begin{noteBox}
定义金属层电阻率、厚度、温度系数及 EM 限值等。
\end{noteBox}

\subsubsection{\texttt{RESISTIVE_ONLY}}
\paragraph*{RESISTIVE_ONLY}

仅提取电阻网络（忽略电容）。

\subsubsection{\texttt{CAPACITIVE_ONLY}}
\paragraph*{CAPACITIVE_ONLY}

仅提取电容（忽略电阻）。

\textbf{参数说明：}
\begin{itemize}
\item \texttt{RESISTIVE\_ONLY} — 仅提取电容（忽略电阻）。
\end{itemize}

\subsubsection{\texttt{POLYNOMIAL_BASED_THICKNESS_VARIATION}}
\paragraph*{POLYNOMIAL_BASED_THICKNESS_VARIATION}

基于多项式的金属厚度变化建模。 注意：基于多项式的金属厚度变化建模。

\textbf{参数说明：}
基于多项式的金属厚度变化建模。

\textbf{示例：}
\begin{lstlisting}
RHO_VS_SI_WIDTH_AND_THICKNESS {
METAL metal1
{
MINWIDTH 0.2
MINSPACE 0.3
PITCH 0.5
THICKNESS 0.4
RESISTANCE 0.091
EM 0.91
BLECH_JLC 10.0
HEIGHT 1
ABOVE DIEL1
RPSQ_VS_WIDTH_AND_SPACING {
WIDTHS { 0.42 0.84 1 2 }
VALUES {
{ 0.02 }
{ 0.03 }
{ 0.04 }
{ 0.041 }
}
}
Simple wire resistance is specified using the syntax:
Resistance <rpsq>
In the example below spacing is defined by columns left to right and
widths are defined by rows top to bottom:
ETCH_VS_WIDTH_AND_SPACING RESISTIVE_ONLY {
SPACINGS { 0.378 0.42 0.5 0.84 0.9 1 2 3 }
WIDTHS { 0.42 0.84 1.68 2.52 }
VALUES {
{ 0.0218 0.0217 0.0214 0.0206 0.0205 0.0203 0.0196 0.0194 }
{ 0.0233 0.0231 0.0228 0.0221 0.0220 0.0218 0.0209 0.0205 }
{ 0.0242 0.0241 0.0239 0.0233 0.0232 0.0230 0.0220 0.0215 }
{ 0.021 0.0245 0.0243 0.0238 0.0237 0.0236 0.0227 0.0221 }
}
}
metal M3 {
ETCH_VS_WIDTH_AND_SPACING CAPACITIVE_ONLY PERPENDICULAR_TO_REFERENCE
<width and spacing data>
ETCH_VS_WIDTH_AND_SPACING CAPACITIVE_ONLY PARALLEL_TO_REFERENCE
<width and spacing data>
\end{lstlisting}

\begin{noteBox}
基于多项式的金属厚度变化建模。
\end{noteBox}

\begin{noteBox}
基于多项式的金属厚度变化建模。
\end{noteBox}

\subsubsection{\texttt{SUBSTRATE}}
\paragraph*{SUBSTRATE}

定义衬底层厚度与电阻率。（可选；默认：none）

\textbf{语法：}
\begin{lstlisting}
SUBSTRATE [ BULK | EPI | PWELL | NWELL ]
{
? THERMAL_CONDUCTIVITY <value>?
? THICKNESS <value>?
? RESISTIVITY <value>?
? CAPACITANCE <value>?
}
\end{lstlisting}

\textbf{示例：}
\begin{lstlisting}
SUBSTRATE BULK
{
THERMAL_CONDUCTIVITY 120
THICKNESS 500
RESISTIVITY 0.012
}
\end{lstlisting}

\subsubsection{\texttt{UNITS}}
\paragraph*{UNITS}

声明 tech 文件中使用的单位制。（可选；默认：s: shown in table below） 注意：声明 tech 文件中使用的单位制。

\textbf{语法：}
\begin{lstlisting}
Syntax
units
{
capacitance           <value> (default: 1 pF; 1e-12 Farad)
inductance            <value> (default: 1 nH; 1e-9 Henry)
resistance            <value> (default: 1 Ohm)
length                <value> (default: 1 um; 1e-6 meter)
current               <value> (default: 1 mA; 1e-3 Ampere)
voltage               <value> (default: 1 Volt)
power                 <value> (default: 1 Watt)
time                  <value> (default: 1 ns; 1e-9 second)
frequency             <value> (default: 1 MHz; 1e+6 Hertz)
}
VIA <via_name> {
Resistance <resistance_value>
UpperLayer <metal_layer_above>
LowerLayer <metal_layer_below> ? <metal_layer_name2>?
Coeff_RT1 <R variation-first order coeff>
Coeff_RT2 <R variation-second order coeff>
TNOM <nominal_temp>
...
? RPV_VS_AREA {
{<min-area_1> <resistance-1>
...
<max-area_n> <resistance-n>
syntax and connected wire EM criteria following it in the syntax are
\end{lstlisting}

\textbf{参数说明：}
\begin{itemize}
\item \texttt{L/W} — 声明 tech 文件中使用的单位制。（默认：0) TNOM: specifies nominal temperature i）
\end{itemize}

\textbf{示例：}
\begin{lstlisting}
units
{
capacitance 1p
inductance 1n
resistance 1
length      1u
current     1m
voltage     1
power       1
time        1n
frequency   1M
}
VIA <via_name>
Specifies the via name and physical parameters defining the physical
characteristics of each via. Note that multiple vias can be defined under different via
names.
Note that Via parameter definitions are taken first from the Apache Tech file, but for
data not present in the Tech file, then Technology LEF file data is used. Required for
all vias.
\end{lstlisting}

\begin{noteBox}
声明 tech 文件中使用的单位制。
\end{noteBox}

\begin{noteBox}
声明 tech 文件中使用的单位制。
\end{noteBox}

\begin{noteBox}
声明 tech 文件中使用的单位制。
\end{noteBox}

\begin{noteBox}
声明 tech 文件中使用的单位制。
\end{noteBox}

\subsubsection{\texttt{VIAMODEL}}
\paragraph*{VIAMODEL}

定义过孔电阻、EM 限值及 via 模型参数。（必需）

\textbf{语法：}
\begin{lstlisting}
VIAMODEL <model_name>
{
<metal_layer_name> <X_loc_lower_left> <Y_loc_lower_left>
<X_loc_upper_right> <Y_loc_upper_right>
<via_layer_name>    <X_loc_lower_left> <Y_loc_lower_left>
<X_loc_upper_right> <Y_loc_upper_right>
<metal_layer_name> <X_loc_lower_left> <Y_loc_lower_left>
<X_loc_upper_right> <Y_loc_upper_right>
}
\end{lstlisting}

\textbf{示例：}
\begin{lstlisting}
VIAMODEL via1
{
metal2 -0.2 -0.2 0.2 0.2
\end{lstlisting}

\section{全局开关配置 GSC 文件}
\subsection*{Global Switching Configuration (GSC) File}

GSC 文件定义设计中模块、实例与电压域的开关状态，用于功耗计算与仿真。
可与 SIM2IPROF、AVM 或 APLMMX 的多状态定义配合，实现无向量（vectorless）多状态仿真。
在 GSR 中通过 \gsr{GSC\_FILES} 指定路径；亦可用 TCL \cmd{import gsc <file>} 动态读入。

\textbf{语法（GSR）：}
\begin{lstlisting}
GSC_FILES <gsc_FilePathName>
\end{lstlisting}

\textbf{语法（GSC 文件）：}
\begin{lstlisting}
[<blockName>|<instanceName>] ?<domain_name>?
[UNDECIDED|TOGGLE|HIGH|LOW|POWERUP|POWERDOWN|STANDBY|ENABLE|DISABLE|OFF|<custom_state>]
\end{lstlisting}

\begin{itemize}
\item 若 GSC 仅含单周期状态，RedHawk 在整个仿真时间内重复该状态
\item 若含多状态序列，触发多周期 GSC 流程
\item 多 GSC 文件中重复实例：后读文件覆盖先读值（警告 GSC-028）
\item 默认启用 \gsr{DYNAMIC\_GSC\_CHECK}，检查 GSC 中未覆盖的实例
\item 大文件可在首行加注释 \texttt{\#EXACT\_INSTANCE\_NAME\_MATCH} 加速解析
\end{itemize}

\subsection{GSC 状态关键字说明}
\subsubsection*{GSC State Keywords}

\begin{itemize}
\item \textbf{UNDECIDED}（默认）— 由 RedHawk 在仿真中决定实例/模块状态
\item \textbf{TOGGLE} — 实例在仿真中翻转；不推荐用于真实分析
\item \textbf{HIGH/LOW} — APL 表征顺序中的高/低翻转电荷条件（非必然对应功耗高低）
\item \textbf{POWERUP/POWERDOWN} — 电源门控模块上电/断电
\item \textbf{STANDBY} — 无输出翻转，时钟仍可能翻转
\item \textbf{DISABLE/ENABLE/OFF} — 实例 inactive/active/关断（OFF 时漏电为零）
\item \textbf{<custom\_state>} — SIM2IPROF/APLMMX 定义的自定义状态
\end{itemize}

\section{全局系统需求 GSR 文件}
\subsection*{Global System Requirements File (*.gsr)}

GSR 是 RedHawk 的\textbf{主配置文件}，汇总输入设计文件路径、分析类型、仿真参数与流程控制关键字。
可用 TCL \cmd{gsr get}/\cmd{gsr set} 查询或修改。关键字名不区分大小写。

本译本按原书 18 类列出全部 659 个 GSR 关键字；
每类以长表给出中文用途说明，语法与默认值见原书或在线 \cmd{help}。

\begin{noteBox}
GSR 关键字名不区分大小写；\gsr{DEFINE} 与 \gsr{INCLUDE} 可组织多文件 GSR 结构。
\end{noteBox}

\subsection{输入数据关键字}
\subsubsection*{Input Data Keywords}

本类共 101 个 GSR 关键字。以下逐条给出中文用途、语法、默认值与示例。

\subsubsection{\gsr{ADD_PLOC_FROM_DEF}}
\paragraph*{ADD_PLOC_FROM_DEF}

允许指定额外的 DEF 文件以定义 PLOC；RedHawk 仅读取其中 PIN 信息。（可选；默认：none）

\textbf{语法：}
\begin{lstlisting}
ADD_PLOC_FROM_DEF <DEF file path>
\end{lstlisting}

\subsubsection{\gsr{ADD_PLOC_FROM_TOP_DEF}}
\paragraph*{ADD_PLOC_FROM_TOP_DEF}

设为 1 时，使用顶层 DEF 的 PINS 段作为电源/地焊盘位置。（可选；默认：0 (no automatic copying from DEF)；兼容 DMP）

\textbf{语法：}
\begin{lstlisting}
ADD_PLOC_FROM_TOP_DEF [0|1] { <layerName1> <layerName2> ... }
ADD_PAD_LOC_FROM_TOP_DEF [0|1]{ <layerName1> <layerName2> ... }
\end{lstlisting}

\subsubsection{\gsr{ATE_USE_GDS_CELL}}
\paragraph*{ATE_USE_GDS_CELL}

GSR 关键字 ATE_USE_GDS_CELL：用于设置与 ATE、USE、GDS、单元 相关的 RedHawk 运行参数。（可选；默认：0；兼容 DMP）

\textbf{语法：}
\begin{lstlisting}
ATE_USE_GDS_CELL [1 | 0]
\end{lstlisting}

\subsubsection{\gsr{AUTO_PAD_CONNECTION_LAYERS}}
\paragraph*{AUTO_PAD_CONNECTION_LAYERS}

按指定金属层与捕捉范围自动为焊盘单元生成 P/G 连接 ploc。（可选；默认：none） 注意：按指定金属层与捕捉范围自动为焊盘单元生成 P/G 连接 ploc。

\textbf{语法：}
\begin{lstlisting}
AUTO_PAD_CONNECTION_LAYERS {
<layerA> <snap_range>
<layerB> <snap_range>
...
DEFAULT <default_snap_range>
REMOVE_FLOATING_PINS [ yes | no ]
}
\end{lstlisting}

\textbf{参数说明：}
按指定金属层与捕捉范围自动为焊盘单元生成 P/G 连接 ploc。（默认：range is set, the default value is used）

\textbf{示例：}
\begin{lstlisting}
AUTO_PAD_CONNECTION_LAYERS {
M1 10.0
M2
M3 15.0
DEFAULT 5.0
}
\end{lstlisting}

\begin{noteBox}
按指定金属层与捕捉范围自动为焊盘单元生成 P/G 连接 ploc。
\end{noteBox}

\subsubsection{\gsr{BLOCK_POWER_ASSIGNMENT}}
\paragraph*{BLOCK_POWER_ASSIGNMENT}

在早期设计中为未完成模块/区域分配功耗，用于评估电网响应。（默认：None；必需）

\textbf{语法：}
\begin{lstlisting}
<blockName> <type> <layerName> <netName>:<pinName>
<power/current> ...
When <type> is CELLTYPE, the syntax is:
<cellName> CELLTYPE <layerName> <pinName> <power/
current>
<powerW>/<GndCurrentA> : specifies either the power in Watts for power nets,
the ground current in Amps for ground nets,
'-1' : allows power calculation to determine the appropriate power/current based
on other information available, such as toggle rate, BPFS and APL
characterization. This can be used for either cell instances (blocks) or
regions. You must provide correct and complete power calculation data so
that power for the block can be calculated accurately. See section "Power
Calculation", page 4-34, for details. Use one line in the BPA specification for
every domain (net) in the block. Note that if one net in the block is assigned
‘-1’ for power/current in BPA, power/current for all domains in the block is
determined by the power calculation engine.
<cellName> : specifies the master cell name
BLOCK_POWER_ASSIGNMENT {
<BPA_instName> <cellName>
[ <layer name> | ALL | TOP | BOTTOM ] <netName>
[<powerW>|<GndCurrentA>]
<llx> <lly> <orientation>
...
?<BPA_instName> <cellName> OVERLAP_OK?
}
\end{lstlisting}

\textbf{参数说明：}
\begin{itemize}
\item \texttt{<instName>} — 在早期设计中为未完成模块/区域分配功耗，用于评估电网响应。（可选）
\end{itemize}

\subsubsection{\gsr{BLOCK_POWER_MASTER_CELL}}
\paragraph*{BLOCK_POWER_MASTER_CELL}

GSR 关键字 BLOCK_POWER_MASTER_CELL：用于设置与 模块、电源、主、单元 相关的 RedHawk 运行参数。（必需）

\textbf{参数说明：}
\begin{itemize}
\item \texttt{FULLCHIP} — GSR 关键字 BLOCK\_POWER\_MASTER\_CELL：用于设置与 模块、电源、主、单元 相关的 RedHawk 运行参数。
\end{itemize}

\textbf{示例：}
\begin{lstlisting}
cell TOP_region1 {
pin vdd1 {
type vdd
}
pin vdd2 {
type vdd
}
pin vss {
type vss
}
pgarc {
vdd1 vss
vdd2 vss
}
}
And the associated example GSR would be as follows:
BLOCK_POWER_ASSIGNMENT {
TOP/region1 REGION M1 VDD1:vdd1 0.1 100 100 300 450
TOP/region1 REGION M1 VDD2:vdd2 0.1 100 100 300 450
TOP/region1 REGION M1 VSS:vss 0.2 100 100 300 450
}
BPA region instances are listed in the file <design>.inst.pin. For static analysis,
BPA REGION instances are listed in files <design>.inst.arc and
<design>.inst.worst.
You can also exclude regions from power assignment in early analysis for designs
at a stage where place and route is complete for some hierarchical blocks, and you
want to use power calculation data for the completed blocks and use power
assignment for rest of the design. See section "Power Assignment to MMX Pin-
based Regions", page 4-65, for more details on MMX BPA use.
For power/current assigned to hierarchical elements in a design, power/current
assigned to child elements are included in the parent‘s assignment. So a parent
block’s power assignment is the sum of the power assigned specifically to its
children and also the power assigned to areas that are outside of its children. For
more details, see section "Early Analysis Flow", page 4-60.
c. Define a modified area within a block for power assignment
To define sub-regions inside an existing BPA instance (multiple specifications on
the same instance are allowed) for power assignment, the syntax is:
BLOCK_POWER_ASSIGNMENT (
<BPA_instName> BLOCK AREA <llx> <lly> <urx> <ury>
BLOCK_POWER_ASSIGNMENT {
\end{lstlisting}

\subsubsection{\gsr{BLOCK_POWER_ASSIGNMENT_FILE}}
\paragraph*{BLOCK_POWER_ASSIGNMENT_FILE}

GSR 关键字 BLOCK_POWER_ASSIGNMENT_FILE：用于设置与 模块、电源、分配、文件 相关的 RedHawk 运行参数。（可选；默认：none）

\textbf{语法：}
\begin{lstlisting}
BLOCK_POWER_ASSIGNMENT_FILE <file1>
...
or
BLOCK_POWER_ASSIGNMENT_FILE {
<file1>
...
}
\end{lstlisting}

\textbf{示例：}
\begin{lstlisting}
BLOCK_POWER_ASSIGNMENT_FILE {
bpa.txt
BPA_xyz.txt
}
The syntax and sample contents of a block power assignment file is as follows:
BLOCK_POWER_MASTER_CELL {
<BPMC_parameter_values>
...
}
BLOCK_POWER_ASSIGNMENT {
<BPA parameter_values>
\end{lstlisting}

\subsubsection{\gsr{BLOCK_POWER_MASTER_CELL}}
\paragraph*{BLOCK_POWER_MASTER_CELL}

GSR 关键字 BLOCK_POWER_MASTER_CELL：用于设置与 模块、电源、主、单元 相关的 RedHawk 运行参数。（可选；默认：none）

\textbf{语法：}
\begin{lstlisting}
BLOCK_POWER_MASTER_CELL {
<master_cell_name1> <BB_llx> <BB_lly> <BB_urx> <BB_ury> {
? AREA <llx> <lly> <urx> <ury> ?
...
? RECTILINEAR <x1 y1 x2 y3 x4 y5 ...>?
...
}
}
BLOCK_POWER_MASTER_CELL {
<master_cell_name1> <BB_llx> <BB_lly> <BB_urx> <BB_ury>
...
?<master_cell_name2> <bbox> {
<sub-area bbox>
...
}?
}
\end{lstlisting}

\textbf{示例：}
\begin{lstlisting}
BLOCK_POWER_MASTER_CELL {
rcell 0 0 400 200
myRegion 0 0 200 100 {
\end{lstlisting}

\subsubsection{\gsr{BLOCK_STA_FILES}}
\paragraph*{BLOCK_STA_FILES}

GSR 关键字 BLOCK_STA_FILES：用于设置与 模块、STA、文件 相关的 RedHawk 运行参数。（可选；默认：is 0)；兼容 DMP） 注意：GSR 关键字 BLOCK_STA_FILES：用于设置与 模块、STA、文件 相关的 RedHawk 运行参数。

\textbf{语法：}
\begin{lstlisting}
BLOCK_STA_FILES {
FREQ_OF_MISSING_INSTANCES <frequency_in_hertz>
[<top_block> | <Hierarchy_path/instance_name>]
<filename> ? INSTTYPE ?
...
<cell_name> <filename> CELLTYPE
...
}
\end{lstlisting}

\textbf{示例：}
\begin{lstlisting}
BLOCK_STA_FILES {
Top_level_block top.timing
InstA1 abc12.timing
InstB2 efc23.timing
InstC3 xyz56.timing
InstC3/instd1 ghi78.timing INSTTYPE
BlockZ1 blockz1.timing CELLTYPE
}
\end{lstlisting}

\begin{noteBox}
GSR 关键字 BLOCK\_STA\_FILES：用于设置与 模块、STA、文件 相关的 RedHawk 运行参数。
\end{noteBox}

\subsubsection{\gsr{BLOCK_VCD_FILES}}
\paragraph*{BLOCK_VCD_FILES}

GSR 关键字 BLOCK_VCD_FILES：用于设置与 模块、VCD、文件 相关的 RedHawk 运行参数。（可选；默认：None Global System Requirements File (*；兼容 DMP）

\textbf{语法：}
\begin{lstlisting}
BLOCK_VCD_FILES {
VCD_FILE {
<Hierarch_path/DEF_instance_name1> <VCD file1>
<VCD file1 options>
}
VCD_FILE {
<Hierarch_path/DEF_instance_name2> <VCD file2>
<VCD file2 options>
...
}
}
\end{lstlisting}

\textbf{参数说明：}
\begin{itemize}
\item \texttt{<Hierarch\_path/DEF\_instance\_name> <VCD file>} — GSR 关键字 BLOCK\_VCD\_FILES：用于设置与 模块、VCD、文件 相关的 RedHawk 运行参数。
\end{itemize}

\textbf{示例：}
\begin{lstlisting}
BLOCK_VCD_FILES {
VCD_FILE {
mydesign/block_A mydesign.vcd
FILE_TYPE VCD
FRONT_PATH "mydesign_vcd/blockA"
true_time 1
}
VCD_FILE {
top2 vcd.2
front_path "TOP/CHIP/top2/"
true_time 1
}
}
\end{lstlisting}

\subsubsection{\gsr{BPA_BY_CURRENT}}
\paragraph*{BPA_BY_CURRENT}

GSR 关键字 BPA_BY_CURRENT：用于设置与 BPA、BY、CURRENT 相关的 RedHawk 运行参数。（可选；默认：0 (off)）

\textbf{语法：}
\begin{lstlisting}
BPA_BY_CURRENT [ 0 | 1 ]
\end{lstlisting}

\textbf{示例：}
\begin{lstlisting}
BPA_BY_CURRENT 1
BLOCK_POWER_ASSIGNMENT {
adsU1/R1 REGION ALL vddcp 0.02 75.849 265.795 185.445 642.312
adsU1/R1 REGION ALL gnda 0.0111 75.849 265.795 185.445 642.312
}
In this example, a 0.02 Amp current is assigned to domain vddcp, and 0.0111 A to
domain gnda.
\end{lstlisting}

\subsubsection{\gsr{BPA_BY_LAYER}}
\paragraph*{BPA_BY_LAYER}

GSR 关键字 BPA_BY_LAYER：用于设置与 BPA、BY、层 相关的 RedHawk 运行参数。（可选；默认：power assigned using the BLOCK_POWER_ASS）

\textbf{语法：}
\begin{lstlisting}
BPA_BY_LAYER [ 0 | 1 ]
\end{lstlisting}

\subsubsection{\gsr{BPA_CONN_DISTANCE}}
\paragraph*{BPA_CONN_DISTANCE}

GSR 关键字 BPA_CONN_DISTANCE：用于设置与 BPA、CONN、DISTANCE 相关的 RedHawk 运行参数。（可选；默认：0）

\textbf{语法：}
\begin{lstlisting}
BPA_CONN_DISTANCE <distance_um>
\end{lstlisting}

\subsubsection{\gsr{BPA_CONN_MARGIN}}
\paragraph*{BPA_CONN_MARGIN}

GSR 关键字 BPA_CONN_MARGIN：用于设置与 BPA、CONN、MARGIN 相关的 RedHawk 运行参数。

\textbf{语法：}
\begin{lstlisting}
BPA_CONN_MARGIN <distance_um>
\end{lstlisting}

\textbf{参数说明：}
\begin{itemize}
\item \texttt{lie outside of the BPA boundary and cause disconnects. Optional; Default} — GSR 关键字 BPA\_CONN\_MARGIN：用于设置与 BPA、CONN、MARGIN 相关的 RedHawk 运行参数。
\end{itemize}

\subsubsection{\gsr{BPA_CURRENT_DENSITY}}
\paragraph*{BPA_CURRENT_DENSITY}

GSR 关键字 BPA_CURRENT_DENSITY：用于设置与 BPA、CURRENT、DENSITY 相关的 RedHawk 运行参数。（可选；默认：0；兼容 DMP）

\textbf{语法：}
\begin{lstlisting}
BPA_CURRENT_DENSITY <value>
\end{lstlisting}

\subsubsection{\gsr{CELL_GEOMETRIES_AS_MFILL}}
\paragraph*{CELL_GEOMETRIES_AS_MFILL}

GSR 关键字 CELL_GEOMETRIES_AS_MFILL：用于设置与 单元、GEOMETRIES、AS、MFILL 相关的 RedHawk 运行参数。（可选；默认：0；兼容 DMP）

\textbf{语法：}
\begin{lstlisting}
CELL_GEOMETRIES_AS_MFILL { OBS | SIGNAL | ALL/1 | NONE/0 ]
\end{lstlisting}

\textbf{参数说明：}
\begin{itemize}
\item \texttt{OBS} — GSR 关键字 CELL\_GEOMETRIES\_AS\_MFILL：用于设置与 单元、GEOMETRIES、AS、MFILL 相关的 RedHawk 运行参数。（默认：)）
\end{itemize}

\subsubsection{\gsr{CELL_PIN_FILE}}
\paragraph*{CELL_PIN_FILE}

GSR 关键字 CELL_PIN_FILE：用于设置与 单元、引脚、文件 相关的 RedHawk 运行参数。（可选；默认：None Global System Requirements File (*；兼容 DMP）

\textbf{语法：}
\begin{lstlisting}
CELL_PIN_FILE <filename>
The internal file syntax is as follows:
CELL [ STD_CELL | <cell_name> ]
PIN [ <Vdd_pin_name> | <Gnd_pin_name> ]
LAYER <metal_layer>
END PIN
END CELL
Pin file content example:
CELL STD_CELL
PIN VDD1
LAYER M1
END PIN
PIN VSS1
LAYER M1
END PIN
END CELL
\end{lstlisting}

\subsubsection{\gsr{CELL_TYPE_FILE}}
\paragraph*{CELL_TYPE_FILE}

GSR 关键字 CELL_TYPE_FILE：用于设置与 单元、TYPE、文件 相关的 RedHawk 运行参数。

\textbf{语法：}
\begin{lstlisting}
CELL_TYPE_FILE <filename>
\end{lstlisting}

\textbf{参数说明：}
GSR 关键字 CELL\_TYPE\_FILE：用于设置与 单元、TYPE、文件 相关的 RedHawk 运行参数。

\subsubsection{\gsr{CELL_WELL_CAP_FILE}}
\paragraph*{CELL_WELL_CAP_FILE}

GSR 关键字 CELL_WELL_CAP_FILE：用于设置与 单元、WELL、CAP、文件 相关的 RedHawk 运行参数。（可选；默认：none；兼容 DMP）

\textbf{语法：}
\begin{lstlisting}
CELL_WELL_CAP_FILE <filename>
\end{lstlisting}

\subsubsection{\gsr{CPA_FILES}}
\paragraph*{CPA_FILES}

GSR 关键字 CPA_FILES：用于设置与 CPA、文件 相关的 RedHawk 运行参数。（可选；默认：None Global System Requirements File (*；兼容 DMP）

\textbf{语法：}
\begin{lstlisting}
CPA_FILES {
PACKAGE <package_layout_filename> #mandatory option
MODEL <project_path> #mandatory option
\end{lstlisting}

\textbf{参数说明：}
\begin{itemize}
\item \texttt{PACKAGE} — GSR 关键字 CPA\_FILES：用于设置与 CPA、文件 相关的 RedHawk 运行参数。
\end{itemize}

\subsubsection{\gsr{DEF_FILES}}
\paragraph*{DEF_FILES}

GSR 关键字 DEF_FILES：用于设置与 DEF、文件 相关的 RedHawk 运行参数。（默认：None；兼容 DMP；必需） 注意：GSR 关键字 DEF_FILES：用于设置与 DEF、文件 相关的 RedHawk 运行参数。

\begin{noteBox}
GSR 关键字 DEF\_FILES：用于设置与 DEF、文件 相关的 RedHawk 运行参数。
\end{noteBox}

\subsubsection{\gsr{DEF_FILES}}
\paragraph*{DEF_FILES}

GSR 关键字 DEF_FILES：用于设置与 DEF、文件 相关的 RedHawk 运行参数。

\textbf{参数说明：}
\begin{itemize}
\item \texttt{<def\_FilePathName-1>.def} — GSR 关键字 DEF\_FILES：用于设置与 DEF、文件 相关的 RedHawk 运行参数。
\end{itemize}

\subsubsection{\gsr{DEF_FILES}}
\paragraph*{DEF_FILES}

GSR 关键字 DEF_FILES：用于设置与 DEF、文件 相关的 RedHawk 运行参数。

\subsubsection{\gsr{DEF_HONOR_HALF_NODE_SCALE_FACTOR}}
\paragraph*{DEF_HONOR_HALF_NODE_SCALE_FACTOR}

GSR 关键字 DEF_HONOR_HALF_NODE_SCALE_FACTOR：用于设置与 DEF、HONOR、半、节点、缩放、因子 相关的 RedHawk 运行参数。（默认：RedHawk calculates the width as (MINWIDT）

\textbf{语法：}
\begin{lstlisting}
DEF_HONOR_HALF_NODE_SCALE_FACTOR [ 0 | 1]
\end{lstlisting}

\textbf{参数说明：}
\begin{itemize}
\item \texttt{set this keyword to off. Optional} — 默认: On.（默认：On）
\end{itemize}

\subsubsection{\gsr{DEF_IGNORE_LAYERS}}
\paragraph*{DEF_IGNORE_LAYERS}

指定 层 in the DEF 文件 to be ignored during 设计 data input for RedHawk 分析. 兼容 DMP 分布式流程。 可选。默认：无。（可选；默认：None；兼容 DMP）

\textbf{语法：}
\begin{lstlisting}
DEF_IGNORE_LAYERS {
<layer_name1>
...
<layer_nameN>
}
\end{lstlisting}

\textbf{示例：}
\begin{lstlisting}
DEF_IGNORE_LAYERS {
m4
m6
}
\end{lstlisting}

\subsubsection{\gsr{DEF_IGNORE_NETS_WIDTH}}
\paragraph*{DEF_IGNORE_NETS_WIDTH}

GSR 关键字 DEF_IGNORE_NETS_WIDTH：用于设置与 DEF、忽略、网络、WIDTH 相关的 RedHawk 运行参数。（可选；默认：1 - On）

\textbf{语法：}
\begin{lstlisting}
DEF_IGNORE_NETS_WIDTH [1 |0 ]
\end{lstlisting}

\subsubsection{\gsr{DEF_IGNORE_PIN_LAYERS}}
\paragraph*{DEF_IGNORE_PIN_LAYERS}

GSR 关键字 DEF_IGNORE_PIN_LAYERS：用于设置与 DEF、忽略、引脚、金属层 相关的 RedHawk 运行参数。（可选；默认：None；兼容 DMP）

\textbf{语法：}
\begin{lstlisting}
DEF_IGNORE_PIN_LAYERS [
<layer_name1>
...
<layer_nameN>
}
\end{lstlisting}

\textbf{示例：}
\begin{lstlisting}
DEF_IGNORE_PIN_LAYERS {
metal3
metal5
}
\end{lstlisting}

\subsubsection{\gsr{DEF_IGNORE_SPECIFIC_LAYERS}}
\paragraph*{DEF_IGNORE_SPECIFIC_LAYERS}

GSR 关键字 DEF_IGNORE_SPECIFIC_LAYERS：用于设置与 DEF、忽略、SPECIFIC、金属层 相关的 RedHawk 运行参数。（可选；默认：None；兼容 DMP）

\textbf{语法：}
\begin{lstlisting}
DEF_IGNORE_SPECIFIC_LAYERS {
<def_block1_path> "<ignored_layer1> ..."
<def_block2_path> "<ignored_layer2> ..."
...
}
\end{lstlisting}

\textbf{参数说明：}
GSR 关键字 DEF\_IGNORE\_SPECIFIC\_LAYERS：用于设置与 DEF、忽略、SPECIFIC、金属层 相关的 RedHawk 运行参数。

\subsubsection{\gsr{DEF_IMPORT_REGION}}
\paragraph*{DEF_IMPORT_REGION}

GSR 关键字 DEF_IMPORT_REGION：用于设置与 DEF、导入、区域 相关的 RedHawk 运行参数。（可选；默认：None）

\textbf{语法：}
\begin{lstlisting}
DEF_IMPORT_REGION {
<Design name in the def> {
<low_left_x> <low_left_y> <up_right_x> <up_right_y>...
}
}
\end{lstlisting}

\subsubsection{\gsr{DEF_READ_ALL_IO_NETS}}
\paragraph*{DEF_READ_ALL_IO_NETS}

GSR 关键字 DEF_READ_ALL_IO_NETS：用于设置与 DEF、READ、ALL、IO、网络 相关的 RedHawk 运行参数。（可选；默认：0）

\textbf{语法：}
\begin{lstlisting}
DEF_READ_ALL_IO_NETS [0|1 ]
\end{lstlisting}

\subsubsection{\gsr{DEF_READ_CLOCK_ONLY}}
\paragraph*{DEF_READ_CLOCK_ONLY}

GSR 关键字 DEF_READ_CLOCK_ONLY：用于设置与 DEF、READ、CLOCK、仅 相关的 RedHawk 运行参数。（默认：0）

\textbf{语法：}
\begin{lstlisting}
DEF_READ_CLOCK_ONLY [ 0 | 1 ]
IMPORT_NETS {
"* [all]"
}
\end{lstlisting}

\subsubsection{\gsr{DEF_TLEF_MAP}}
\paragraph*{DEF_TLEF_MAP}

GSR 关键字 DEF_TLEF_MAP：用于设置与 DEF、TLEF、MAP 相关的 RedHawk 运行参数。（可选；默认：none；兼容 DMP）

\textbf{语法：}
\begin{lstlisting}
DEF_TLEF_MAP {
<def_file>   <tag>
...
}
\end{lstlisting}

\textbf{示例：}
\begin{lstlisting}
DEF_TLEF_MAP {
DesignDB/DEF/maia_code.def                HD1
DesignDB/DEF/maib_code.def                HD2
......
}
\end{lstlisting}

\subsubsection{\gsr{DEF_TRUE_PATH_EXTENSION}}
\paragraph*{DEF_TRUE_PATH_EXTENSION}

GSR 关键字 DEF_TRUE_PATH_EXTENSION：用于设置与 DEF、TRUE、PATH、EXTENSION 相关的 RedHawk 运行参数。

\textbf{语法：}
\begin{lstlisting}
DEF_TRUE_PATH_EXTENSION [ 0 | 1 ]
\end{lstlisting}

\textbf{参数说明：}
\begin{itemize}
\item \texttt{extenValue} — GSR 关键字 DEF\_TRUE\_PATH\_EXTENSION：用于设置与 DEF、TRUE、PATH、EXTENSION 相关的 RedHawk 运行参数。（可选；默认：0, Off；兼容 DMP）
\end{itemize}

\subsubsection{\gsr{DESIGN_IMPORT_REGION}}
\paragraph*{DESIGN_IMPORT_REGION}

GSR 关键字 DESIGN_IMPORT_REGION：用于设置与 DESIGN、导入、区域 相关的 RedHawk 运行参数。

\textbf{语法：}
\begin{lstlisting}
DESIGN_IMPORT_REGION {
<ll_x> <ll_y> <ur_x> <ur_y>
<ll_x> <ll_y> <ur_x> <ur_y>
...
}
\end{lstlisting}

\subsubsection{\gsr{FAST_DEF_READ}}
\paragraph*{FAST_DEF_READ}

GSR 关键字 FAST_DEF_READ：用于设置与 FAST、DEF、READ 相关的 RedHawk 运行参数。（可选；默认：0 for flat and 4 for DMP run）

\textbf{语法：}
\begin{lstlisting}
FAST_DEF_READ <num_proc>
\end{lstlisting}

\subsubsection{\gsr{FAST_DEF_READ_DIR}}
\paragraph*{FAST_DEF_READ_DIR}

GSR 关键字 FAST_DEF_READ_DIR：用于设置与 FAST、DEF、READ、DIR 相关的 RedHawk 运行参数。

\textbf{语法：}
\begin{lstlisting}
FAST_DEF_READ_DIR <path>
\end{lstlisting}

\subsubsection{\gsr{EXACT_PINPORT_VM}}
\paragraph*{EXACT_PINPORT_VM}

GSR 关键字 EXACT_PINPORT_VM：用于设置与 EXACT、PINPORT、VM 相关的 RedHawk 运行参数。（默认：0 in which tool picks a via model from t）

\subsubsection{\gsr{EXCLUDE_REGION}}
\paragraph*{EXCLUDE_REGION}

GSR 关键字 EXCLUDE_REGION：用于设置与 EXCLUDE、区域 相关的 RedHawk 运行参数。（可选；默认：None）

\textbf{语法：}
\begin{lstlisting}
EXCLUDE_REGION {
<cell_name/design_name> {
<ll_x> <ll_y> <ur_x> <ur_y>
...
}
\end{lstlisting}

\subsubsection{\gsr{GDS_CELLS}}
\paragraph*{GDS_CELLS}

GSR 关键字 GDS_CELLS：用于设置与 GDS、单元 相关的 RedHawk 运行参数。（可选；默认：none；兼容 DMP）

\subsubsection{\gsr{GDS_CELLS}}
\paragraph*{GDS_CELLS}

GSR 关键字 GDS_CELLS：用于设置与 GDS、单元 相关的 RedHawk 运行参数。

\textbf{参数说明：}
\begin{itemize}
\item \texttt{<GdsCellName-1> <path\_to\_cell-1>} — GSR 关键字 GDS\_CELLS：用于设置与 GDS、单元 相关的 RedHawk 运行参数。
\end{itemize}

\subsubsection{\gsr{GDS_CELLS}}
\paragraph*{GDS_CELLS}

GSR 关键字 GDS_CELLS：用于设置与 GDS、单元 相关的 RedHawk 运行参数。

\subsubsection{\gsr{GDS_CELLS_FILE}}
\paragraph*{GDS_CELLS_FILE}

GSR 关键字 GDS_CELLS_FILE：用于设置与 GDS、单元、文件 相关的 RedHawk 运行参数。（可选；默认：None；兼容 DMP）

\textbf{语法：}
\begin{lstlisting}
Syntax
GDS_CELLS_FILE <filename>
\end{lstlisting}

\subsubsection{\gsr{GSC_FILES}}
\paragraph*{GSC_FILES}

GSR 关键字 GSC_FILES：用于设置与 GSC、文件 相关的 RedHawk 运行参数。（默认：the OVERRIDE keyword is set to 0 and GSC）

\textbf{语法：}
\begin{lstlisting}
GSC_FILES <gsc_FilePathName>
[<blockName> | <instanceName> ] ? <domain_name>?
[ UNDECIDED | TOGGLE |HIGH | LOW | POWERUP | POWERDOWN |
STANDBY | ENABLE | DISABLE | OFF | <custom_state_name> ]
...
\end{lstlisting}

\textbf{参数说明：}
\begin{itemize}
\item \texttt{gsr set gsc ABC.gsc DMP compatible. Optional. Default} — GSR 关键字 GSC\_FILES：用于设置与 GSC、文件 相关的 RedHawk 运行参数。
\end{itemize}

\textbf{示例：}
\begin{lstlisting}
GSC_FILES DesignABC/gl_sw2.gsc
blockABC Powerup
block2C toggle
instB* standby
instC1 stateB2
ABC VDD1 High
\end{lstlisting}

\subsubsection{\gsr{HONOR_LEF_PIN_TYPE}}
\paragraph*{HONOR_LEF_PIN_TYPE}

GSR 关键字 HONOR_LEF_PIN_TYPE：用于设置与 HONOR、LEF、引脚、TYPE 相关的 RedHawk 运行参数。（可选；默认：)；兼容 DMP）

\textbf{语法：}
\begin{lstlisting}
Syntax
HONOR_LEF_PIN_TYPE [0|1]
\end{lstlisting}

\subsubsection{\gsr{HOOK_INTERNAL_PIN_POWER_EXTERNAL}}
\paragraph*{HOOK_INTERNAL_PIN_POWER_EXTERNAL}

GSR 关键字 HOOK_INTERNAL_PIN_POWER_EXTERNAL：用于设置与 HOOK、INTERNAL、引脚、电源、EXTERNAL 相关的 RedHawk 运行参数。（可选；默认：0）

\subsubsection{\gsr{IGNORE_INSTANCES_ON_TOP_DEF_REGIONS}}
\paragraph*{IGNORE_INSTANCES_ON_TOP_DEF_REGIONS}

在数据输入或分析流程中忽略与 实例、ON、顶层、DEF、REGIONS 相关的对象或检查。（可选；默认：None）

\textbf{语法：}
\begin{lstlisting}
IGNORE_INSTANCES_ON_TOP_DEF_REGIONS {
<ll_x> <ll_y> <ur_x> <ur_y>
...
}
\end{lstlisting}

\textbf{示例：}
\begin{lstlisting}
IGNORE_INSTANCES_ON_TOP_DEF_REGIONS {
0 0 100 100
150 150 200 250
}
\end{lstlisting}

\subsubsection{\gsr{IGNORE_INST_WITH_NO_MASTER}}
\paragraph*{IGNORE_INST_WITH_NO_MASTER}

在数据输入或分析流程中忽略与 INST、WITH、NO、主 相关的对象或检查。（可选；默认：) IGNORE_LEF_DEF_MISMATCH = 0 (default)；兼容 DMP）

\textbf{语法：}
\begin{lstlisting}
IGNORE_INST_WITH_NO_MASTER [0|1|2]
\end{lstlisting}

\subsubsection{\gsr{IGNORE_LEF_FOREIGN_OFFSET}}
\paragraph*{IGNORE_LEF_FOREIGN_OFFSET}

在数据输入或分析流程中忽略与 LEF、FOREIGN、OFFSET 相关的对象或检查。（可选；默认：1）

\textbf{语法：}
\begin{lstlisting}
IGNORE_LEF_FOREIGN_OFFSET [0|1]
\end{lstlisting}

\subsubsection{\gsr{IGNORE_PLOC_INTERNALNETS}}
\paragraph*{IGNORE_PLOC_INTERNALNETS}

在数据输入或分析流程中忽略与 焊盘位置、INTERNALNETS 相关的对象或检查。（可选；默认：(1), any plocs placed on internal nets a；兼容 DMP）

\textbf{语法：}
\begin{lstlisting}
IGNORE_PLOC_INTERNALNETS [ 0 |1 ]
\end{lstlisting}

\subsubsection{\gsr{IMPORT_NETS}}
\paragraph*{IMPORT_NETS}

GSR 关键字 IMPORT_NETS：用于设置与 导入、网络 相关的 RedHawk 运行参数。（可选；默认：none；兼容 DMP）

\textbf{语法：}
\begin{lstlisting}
IMPORT_NETS {
<net1_name>
<net2_name>
...
}
The keyword also supports the format "blocka/*clk*[all]”. Therefore, in the signal EM
flow, user can import all the clock nets with specified pattern within the specified
block or within the inner hierarchies inside the specified block.
IMPORT_NETS {
"<block_name>/<net pattern> [all]"
}
\end{lstlisting}

\textbf{参数说明：}
GSR 关键字 IMPORT\_NETS：用于设置与 导入、网络 相关的 RedHawk 运行参数。

\textbf{示例：}
\begin{lstlisting}
IMPORT_NETS {
" blocka/*clk* [all]"
}
In this case, "*clk*" nets, which are defined only within "blocka" or with in the blocks
instantiated under "blocka" will be imported. This format removes the redundancy of
specifying "*clk* with each and every block instantiated under "blocka" like:
IMPORT_NETS {
blocka/*clk*
blocka/blockb/*clk*
blocka/blockb/blockc/*clk*
}
\end{lstlisting}

\subsubsection{\gsr{IMPORT_NETS_FILE}}
\paragraph*{IMPORT_NETS_FILE}

GSR 关键字 IMPORT_NETS_FILE：用于设置与 导入、网络、文件 相关的 RedHawk 运行参数。（可选；默认：none）

\textbf{语法：}
\begin{lstlisting}
IMPORT_NETS_FILE <nets_file>
\end{lstlisting}

\subsubsection{\gsr{IMPORT_REGION}}
\paragraph*{IMPORT_REGION}

GSR 关键字 IMPORT_REGION：用于设置与 导入、区域 相关的 RedHawk 运行参数。（可选；默认：None；兼容 DMP）

\textbf{语法：}
\begin{lstlisting}
Syntax
IMPORT_REGION {
<cell_name/design_name> {
<ll_x> <ll_y> <ur_x> <ur_y>
...
}
...
}
\end{lstlisting}

\subsubsection{\gsr{IMPORT_REGION_CELL_OVERLAP}}
\paragraph*{IMPORT_REGION_CELL_OVERLAP}

GSR 关键字 IMPORT_REGION_CELL_OVERLAP：用于设置与 导入、区域、单元、OVERLAP 相关的 RedHawk 运行参数。（默认：partial）

\textbf{语法：}
\begin{lstlisting}
Syntax
IMPORT_REGION_CELL_OVERLAP [ partial | full ]
\end{lstlisting}

\subsubsection{\gsr{KEEP_POWER_DATA}}
\paragraph*{KEEP_POWER_DATA}

GSR 关键字 KEEP_POWER_DATA：用于设置与 KEEP、电源、DATA 相关的 RedHawk 运行参数。（可选；默认：RedHawk power data is retained）

\textbf{语法：}
\begin{lstlisting}
KEEP_POWER_DATA [ 0 |1 ]
\end{lstlisting}

\subsubsection{\gsr{LEF_CELL_IGNORE_PIN_LAYERS}}
\paragraph*{LEF_CELL_IGNORE_PIN_LAYERS}

GSR 关键字 LEF_CELL_IGNORE_PIN_LAYERS：用于设置与 LEF、单元、忽略、引脚、金属层 相关的 RedHawk 运行参数。（可选；默认：none）

\textbf{语法：}
\begin{lstlisting}
LEF_CELL_IGNORE_PIN_LAYERS {
<cell_name> <POWER | GROUND | SIGNAL | ALL | BOTH |
<pin_name> <layer1, layer2,...>
}
\end{lstlisting}

\textbf{参数说明：}
\begin{itemize}
\item \texttt{BOTH} — GSR 关键字 LEF\_CELL\_IGNORE\_PIN\_LAYERS：用于设置与 LEF、单元、忽略、引脚、金属层 相关的 RedHawk 运行参数。
\end{itemize}

\textbf{示例：}
\begin{lstlisting}
LEF_CELL_IGNORE_PIN_LAYERS {
cellA VDD2 MET2,MET1
cellA VDD* MET2
cellA VDD1 MET3
cell* VDD2 MET3,MET4
cell* VDD* MET1
cellB VDD2 MET5,MET3
}
\end{lstlisting}

\subsubsection{\gsr{LEFDEF_TECH_LAYER_MAP_FILE}}
\paragraph*{LEFDEF_TECH_LAYER_MAP_FILE}

GSR 关键字 LEFDEF_TECH_LAYER_MAP_FILE：用于设置与 LEFDEF、工艺、层、MAP、文件 相关的 RedHawk 运行参数。（可选；默认：none；兼容 DMP）

\textbf{语法：}
\begin{lstlisting}
LEFDEF_TECH_LAYER_MAP_FILE <layermap_file>
\end{lstlisting}

\subsubsection{\gsr{LEFDEF_TECH_LAYER_MAP}}
\paragraph*{LEFDEF_TECH_LAYER_MAP}

GSR 关键字 LEFDEF_TECH_LAYER_MAP：用于设置与 LEFDEF、工艺、层、MAP 相关的 RedHawk 运行参数。

\textbf{语法：}
\begin{lstlisting}
LEFDEF_TECH_LAYER_MAP {
<LEF/DEF_layer_name> <techlayer_name>
…
}
\end{lstlisting}

\textbf{示例：}
\begin{lstlisting}
LEFDEF_TECH_LAYER_MAP {
VIA0 VIA0_M0_STI
M0 M0_STI
}
\end{lstlisting}

\subsubsection{\gsr{LEF_FILES}}
\paragraph*{LEF_FILES}

GSR 关键字 LEF_FILES：用于设置与 LEF、文件 相关的 RedHawk 运行参数。（默认：None；兼容 DMP；必需）

\textbf{语法：}
\begin{lstlisting}
LEF_FILES {
<lef_FilePathName-1>.lef
...
<lef_fFilePathName-n>.lef
}
\end{lstlisting}

\textbf{参数说明：}
\begin{itemize}
\item \texttt{<lef\_FilePathName-1>.lef} — GSR 关键字 LEF\_FILES：用于设置与 LEF、文件 相关的 RedHawk 运行参数。
\end{itemize}

\textbf{示例：}
\begin{lstlisting}
LEF_FILES {
../design_data/ABCD.lef
}
\end{lstlisting}

\subsubsection{\gsr{LEF_IGNORE_PIN_LAYERS}}
\paragraph*{LEF_IGNORE_PIN_LAYERS}

指定 引脚 层 in the LEF 文件 to be ignored during 设计 data input for RedHawk 分析. 兼容 DMP 分布式流程。 可选。默认：无。（可选；默认：None；兼容 DMP）

\textbf{语法：}
\begin{lstlisting}
LEF_IGNORE_PIN_LAYERS {
<layer_name1>
...
<layer_nameN>
}
\end{lstlisting}

\textbf{示例：}
\begin{lstlisting}
LEF_IGNORE_PIN_LAYERS {
M4
M5
}
\end{lstlisting}

\subsubsection{\gsr{LIBERTY_DB}}
\paragraph*{LIBERTY_DB}

GSR 关键字 LIBERTY_DB：用于设置与 LIBERTY、DB 相关的 RedHawk 运行参数。（可选；默认：off）

\textbf{语法：}
\begin{lstlisting}
LIBERTY_DB [ 0 |1 | 2 ]
\end{lstlisting}

\textbf{参数说明：}
\begin{itemize}
\item \texttt{0 - off} — GSR 关键字 LIBERTY\_DB：用于设置与 LIBERTY、DB 相关的 RedHawk 运行参数。
\end{itemize}

\subsubsection{\gsr{LIB_FILES}}
\paragraph*{LIB_FILES}

GSR 关键字 LIB_FILES：用于设置与 LIB、文件 相关的 RedHawk 运行参数。（默认：None；兼容 DMP；必需）

\textbf{语法：}
\begin{lstlisting}
LIB_FILES {
[ <lib_filename> | <lib_file_dir> ] ? CUSTOM ?
...
}
\end{lstlisting}

\textbf{参数说明：}
\begin{itemize}
\item \texttt{<lib\_filename>} — GSR 关键字 LIB\_FILES：用于设置与 LIB、文件 相关的 RedHawk 运行参数。
\end{itemize}

\textbf{示例：}
\begin{lstlisting}
example, this allows including or excluding particular states from power
calculation using a custom LIB file. You can do this for a particular cell, or if
one or more cells is not specified, it applies globally for all cells. User has
LIB_FILES {
special/custom.lib CUSTOM
dirAB/memory/mem.lib
}
Sample contents of custom.lib file:
cell SDFFX1HVT {
exclude_states {
pin CK {
"EN * !(SE) * !(Q)"
}
}
}
cell cell_abc {
exclude_states {
pin * {
“RST&!CS”
}
}
}
\end{lstlisting}

\subsubsection{\gsr{LIB_HONOR_OUT_SIGNAL_SWING}}
\paragraph*{LIB_HONOR_OUT_SIGNAL_SWING}

GSR 关键字 LIB_HONOR_OUT_SIGNAL_SWING：用于设置与 LIB、HONOR、OUT、SIGNAL、SWING 相关的 RedHawk 运行参数。（默认：0）

\textbf{语法：}
\begin{lstlisting}
LIB_HONOR_OUT_SIGNAL_SWING [ 1| 0]
\end{lstlisting}

\subsubsection{\gsr{LIB_IGNORE_CELL_LEAKAGE}}
\paragraph*{LIB_IGNORE_CELL_LEAKAGE}

GSR 关键字 LIB_IGNORE_CELL_LEAKAGE：用于设置与 LIB、忽略、单元、LEAKAGE 相关的 RedHawk 运行参数。（默认：0）

\textbf{语法：}
\begin{lstlisting}
LIB_IGNORE_CELL_LEAKAGE [ 1| 0]
\end{lstlisting}

\subsubsection{\gsr{LIB_IGNORE_IO_VOLTAGE}}
\paragraph*{LIB_IGNORE_IO_VOLTAGE}

GSR 关键字 LIB_IGNORE_IO_VOLTAGE：用于设置与 LIB、忽略、IO、VOLTAGE 相关的 RedHawk 运行参数。（可选；默认：value 0 allows reading in these values f）

\textbf{语法：}
\begin{lstlisting}
LIB_IGNORE_IO_VOLTAGE [ 1| 0]
\end{lstlisting}

\subsubsection{\gsr{LIB_IGNORE_POWER_THRESHOLD}}
\paragraph*{LIB_IGNORE_POWER_THRESHOLD}

GSR 关键字 LIB_IGNORE_POWER_THRESHOLD：用于设置与 LIB、忽略、电源、THRESHOLD 相关的 RedHawk 运行参数。（可选；默认：None；兼容 DMP）

\subsubsection{\gsr{LIB_PIN_CAP}}
\paragraph*{LIB_PIN_CAP}

GSR 关键字 LIB_PIN_CAP：用于设置与 LIB、引脚、CAP 相关的 RedHawk 运行参数。（可选；默认：typical capacitance value；兼容 DMP）

\textbf{语法：}
\begin{lstlisting}
LIB_PIN_CAP [min|max|typ]
\end{lstlisting}

\subsubsection{\gsr{MACRO_POWER_FILES}}
\paragraph*{MACRO_POWER_FILES}

GSR 关键字 MACRO_POWER_FILES：用于设置与 宏、电源、文件 相关的 RedHawk 运行参数。（可选；默认：None）

\textbf{语法：}
\begin{lstlisting}
MACRO_POWER_FILES {
<memory_macro_filename>
...
}
\end{lstlisting}

\textbf{示例：}
\begin{lstlisting}
MACRO_POWER_FILES {
memory_macroMN.pwr
}
\end{lstlisting}

\subsubsection{\gsr{PACKAGE_SPICE_SUBCKT_INFO}}
\paragraph*{PACKAGE_SPICE_SUBCKT_INFO}

GSR 关键字 PACKAGE_SPICE_SUBCKT_INFO：用于设置与 PACKAGE、SPICE、SUBCKT、INFO 相关的 RedHawk 运行参数。（可选；默认：None）

\textbf{语法：}
\begin{lstlisting}
PACKAGE_SPICE_SUBCKT_INFO {
PATH <path_package_Spice_file>
? TOP <top subckt name> ?
}
Note that the PATH option must be specified. If TOP is not specified, then
the first subckt is treated as the top subckt.
\end{lstlisting}

\subsubsection{\gsr{PAD_FILES}}
\paragraph*{PAD_FILES}

GSR 关键字 PAD_FILES：用于设置与 焊盘、文件 相关的 RedHawk 运行参数。（可选；默认：none；兼容 DMP）

\textbf{参数说明：}
\begin{itemize}
\item \texttt{<padFilePathName-1>} — GSR 关键字 PAD\_FILES：用于设置与 焊盘、文件 相关的 RedHawk 运行参数。
\end{itemize}

\subsubsection{\gsr{PAD_FILES}}
\paragraph*{PAD_FILES}

GSR 关键字 PAD_FILES：用于设置与 焊盘、文件 相关的 RedHawk 运行参数。

\subsubsection{\gsr{PARTIAL_FLAT_SPEF}}
\paragraph*{PARTIAL_FLAT_SPEF}

GSR 关键字 PARTIAL_FLAT_SPEF：用于设置与 PARTIAL、FLAT、SPEF 相关的 RedHawk 运行参数。（可选；默认：0 (off)）

\textbf{语法：}
\begin{lstlisting}
PARTIAL_FLAT_SPEF [ 0 | 1 ]
\end{lstlisting}

\subsubsection{\gsr{PGNET_HONOR_DEF_TYPE}}
\paragraph*{PGNET_HONOR_DEF_TYPE}

GSR 关键字 PGNET_HONOR_DEF_TYPE：用于设置与 PGNET、HONOR、DEF、TYPE 相关的 RedHawk 运行参数。（可选；默认：1 (On)；兼容 DMP）

\textbf{语法：}
\begin{lstlisting}
PGNET_HONOR_DEF_TYPE [ 0 | 1 ]
\end{lstlisting}

\subsubsection{\gsr{PRINT_ONE_PLOC_PER_PADINST}}
\paragraph*{PRINT_ONE_PLOC_PER_PADINST}

GSR 关键字 PRINT_ONE_PLOC_PER_PADINST：用于设置与 PRINT、ONE、焊盘位置、PER、PADINST 相关的 RedHawk 运行参数。（可选；默认：0 (off) Global System Requirements File；兼容 DMP；必需）

\textbf{语法：}
\begin{lstlisting}
PRINT_ONE_PLOC_PER_PADINST [ 0|1 ]
RDL_CELL(S)
Redistribution layer geometries for flip chip designs are typically available only in
GDS format. To include RDL data in RedHawk analysis, two steps are required:
a. Conversion of RDL GDS to DEF form with the 'gds2def/gds2rh' utility
b. Instantiating this DEF from the top level design
Syntax
RDL_CELL(S) {
<inst_name1> <DEF_master_cell_name1> <path_RDL_master_DEF_file>
LOCATION <x_loc> <y_loc>
? OFFSET <x_offset> <y_offset>
? ORIENTATION <orientation> <dummy> ?
NET <Vdd_net_name> <power_pin_name>
NET <Vss_net_name> <ground_pin_name>
...
NEW_INSTANCE NEW_MASTER NEW_PATH
<inst2> <mastercell2> <DEF_path2>
<inst_nameB> <DEF_master_cell_nameB> <path_RDL_master_DEF_file>
LOCATION <x_loc> <y_loc>
? OFFSET <x_offset> <y_offset>
? ORIENTATION <orientation> <dummy> ?
NET <Vdd_net_name> <power_pin_name>
NET <Vss_net_name> <ground_pin_name>
...
NEW_INSTANCE NEW_MASTER NEW_PATH
\end{lstlisting}

\textbf{参数说明：}
\begin{itemize}
\item \texttt{<inst\_name>} — GSR 关键字 PRINT\_ONE\_PLOC\_PER\_PADINST：用于设置与 PRINT、ONE、焊盘位置、PER、PADINST 相关的 RedHawk 运行参数。（默认：0,0）
\end{itemize}

\textbf{示例：}
\begin{lstlisting}
Example
RDL_CELL {
rdl_inst chip_rdl /home/proj/data/def/chip_rdl.def
LOCATION 500 250
OFFSET -25 -25
ORIENTATION FS dummy
NET VDD VDD
NET VSS VSS
NEW_INSTANCE NEW_MASTER NEW_PATH
rdl_inst_A chip_rdl_C /home/proj/data/def/chip_rdl.def
}
\end{lstlisting}

\subsubsection{\gsr{READ_LEF_OBS}}
\paragraph*{READ_LEF_OBS}

GSR 关键字 READ_LEF_OBS：用于设置与 READ、LEF、OBS 相关的 RedHawk 运行参数。（可选；默认：off, 0）

\textbf{语法：}
\begin{lstlisting}
READ_LEF_OBS 1
\end{lstlisting}

\textbf{示例：}
\begin{lstlisting}
OBS section items in LEF included in power assignment:
\end{lstlisting}

\subsubsection{\gsr{OBS}}
\paragraph*{OBS}

GSR 关键字 OBS：用于设置与 OBS 相关的 RedHawk 运行参数。

\subsubsection{\gsr{END}}
\paragraph*{END}

GSR 关键字 END：用于设置与 END 相关的 RedHawk 运行参数。

\subsubsection{\gsr{REMOVE_PARENTLEF_GEOS}}
\paragraph*{REMOVE_PARENTLEF_GEOS}

GSR 关键字 REMOVE_PARENTLEF_GEOS：用于设置与 REMOVE、PARENTLEF、GEOS 相关的 RedHawk 运行参数。（可选；默认：off, 0）

\textbf{语法：}
\begin{lstlisting}
Syntax
REMOVE_PARENTLEF_GEOS [ 0 | 1 ]
\end{lstlisting}

\subsubsection{\gsr{REVERSE_DEF_READ_ORDER}}
\paragraph*{REVERSE_DEF_READ_ORDER}

GSR 关键字 REVERSE_DEF_READ_ORDER：用于设置与 REVERSE、DEF、READ、ORDER 相关的 RedHawk 运行参数。（可选；默认：the top DEFs are imported first, and the）

\textbf{语法：}
\begin{lstlisting}
Syntax
REVERSE_DEF_READ_ORDER [ 0 | 1 ]
\end{lstlisting}

\subsubsection{\gsr{STA_CRITICAL_PATH_FILE}}
\paragraph*{STA_CRITICAL_PATH_FILE}

GSR 关键字 STA_CRITICAL_PATH_FILE：用于设置与 STA、CRITICAL、PATH、文件 相关的 RedHawk 运行参数。（可选；默认：none）

\textbf{语法：}
\begin{lstlisting}
Syntax
\end{lstlisting}

\subsubsection{\gsr{STA_CRITICAL_PATH_FILE}}
\paragraph*{STA_CRITICAL_PATH_FILE}

GSR 关键字 STA_CRITICAL_PATH_FILE：用于设置与 STA、CRITICAL、PATH、文件 相关的 RedHawk 运行参数。

\textbf{示例：}
\begin{lstlisting}
Example
\end{lstlisting}

\subsubsection{\gsr{STA_CRITICAL_PATH_FILE}}
\paragraph*{STA_CRITICAL_PATH_FILE}

GSR 关键字 STA_CRITICAL_PATH_FILE：用于设置与 STA、CRITICAL、PATH、文件 相关的 RedHawk 运行参数。

\subsubsection{\gsr{STA_FILES}}
\paragraph*{STA_FILES}

GSR 关键字 STA_FILES：用于设置与 STA、文件 相关的 RedHawk 运行参数。（可选；默认：0), clockroots defined in the clock sect；兼容 DMP；必需） 注意：GSR 关键字 STA_FILES：用于设置与 STA、文件 相关的 RedHawk 运行参数。

\textbf{语法：}
\begin{lstlisting}
STA_FILES {
FREQ_OF_MISSING_INSTANCES <frequency in hertz>
APPLY_FOMI_TO_CLOCK_PINS [0|1]
EXTRACT_CLOCK_NETWORK [ 0| 1| 2]
<top_level_block_name> <sta_output_file>
}
\end{lstlisting}

\textbf{参数说明：}
\begin{itemize}
\item \texttt{FREQ\_OF\_MISSING\_INSTANCES <frequency in hertz>} — GSR 关键字 STA\_FILES：用于设置与 STA、文件 相关的 RedHawk 运行参数。（默认：0）
\end{itemize}

\begin{noteBox}
GSR 关键字 STA\_FILES：用于设置与 STA、文件 相关的 RedHawk 运行参数。
\end{noteBox}

\subsubsection{\gsr{STA_FILES}}
\paragraph*{STA_FILES}

GSR 关键字 STA_FILES：用于设置与 STA、文件 相关的 RedHawk 运行参数。

\subsubsection{\gsr{STACK_VIA_OFFSET_DISTANCE}}
\paragraph*{STACK_VIA_OFFSET_DISTANCE}

GSR 关键字 STACK_VIA_OFFSET_DISTANCE：用于设置与 STACK、过孔、OFFSET、DISTANCE 相关的 RedHawk 运行参数。

\textbf{语法：}
\begin{lstlisting}
STACK_VIA_OFFSET_DISTANCE {
Metal1 distance_in_um
Metal2 distance_in_um
...
}
Note that this keyword is mutually exclusive with GSR keyword
STACK_VIA_OFFSET_RESISTANCE. Only one of the two can be defined in a
GSR.
\end{lstlisting}

\subsubsection{\gsr{STACK_VIA_OFFSET_RESISTANCE}}
\paragraph*{STACK_VIA_OFFSET_RESISTANCE}

GSR 关键字 STACK_VIA_OFFSET_RESISTANCE：用于设置与 STACK、过孔、OFFSET、RESISTANCE 相关的 RedHawk 运行参数。

\textbf{语法：}
\begin{lstlisting}
STACK_VIA_OFFSET_REISTANCE {
Metal1 res_in_ohm
Metal2 res_in_ohm
...
}
Note that this keyword is mutually exclusive with GSR keyword
STACK_VIA_OFFSET_DISTANCE. Only one of the two can be defined in a GSR.
\end{lstlisting}

\subsubsection{\gsr{STANDARD_CELL_HEIGHT}}
\paragraph*{STANDARD_CELL_HEIGHT}

GSR 关键字 STANDARD_CELL_HEIGHT：用于设置与 STANDARD、单元、HEIGHT 相关的 RedHawk 运行参数。（可选；默认：RedHawk identifies standard cell heights；兼容 DMP）

\textbf{语法：}
\begin{lstlisting}
Syntax
STANDARD_CELL_HEIGHT {
<cell_height1>
<cell_height2>
...
}
\end{lstlisting}

\subsubsection{\gsr{TECH_CONNECTIVITY_FILE}}
\paragraph*{TECH_CONNECTIVITY_FILE}

GSR 关键字 TECH_CONNECTIVITY_FILE：用于设置与 工艺、CONNECTIVITY、文件 相关的 RedHawk 运行参数。（可选；默认：none）

\subsubsection{\gsr{TECH_FILES}}
\paragraph*{TECH_FILES}

GSR 关键字 TECH_FILES：用于设置与 工艺、文件 相关的 RedHawk 运行参数。（默认：none；兼容 DMP；必需） 注意：GSR 关键字 TECH_FILES：用于设置与 工艺、文件 相关的 RedHawk 运行参数。

\textbf{语法：}
\begin{lstlisting}
TECH_FILES <techFilePathName>.tech
\end{lstlisting}

\textbf{示例：}
\begin{lstlisting}
TECH_FILES abcDesign4.tech
\end{lstlisting}

\begin{noteBox}
GSR 关键字 TECH\_FILES：用于设置与 工艺、文件 相关的 RedHawk 运行参数。
\end{noteBox}

\subsubsection{\gsr{TECH_LEFS}}
\paragraph*{TECH_LEFS}

GSR 关键字 TECH_LEFS：用于设置与 工艺、LEFS 相关的 RedHawk 运行参数。（可选；默认：none；兼容 DMP）

\textbf{语法：}
\begin{lstlisting}
TECH_LEFS {
<tlef_file1>         <tag1>
<tlef_file2>         <tag2>
...
}
\end{lstlisting}

\textbf{示例：}
\begin{lstlisting}
TECH_LEFS {
abc1.tlef HD1
LEF_FILES {
../lef/global.tlef #global tech lef
abc1.tlef # Same block1 tech lef as used in TECH_LEFS keyword in the
above example
abc2.tlef #Same block2 tech lef as used in TECH_LEFS keyword in the
above example.
...
}
\end{lstlisting}

\subsubsection{\gsr{TEMPERATURE}}
\paragraph*{TEMPERATURE}

GSR 关键字 TEMPERATURE：用于设置与 温度 相关的 RedHawk 运行参数。（可选；默认：none；兼容 DMP）

\textbf{语法：}
\begin{lstlisting}
TEMPERATURE <actual_temperature_ºC>
\end{lstlisting}

\textbf{示例：}
\begin{lstlisting}
Temperature 100
\end{lstlisting}

\subsubsection{\gsr{TEMPERATURES}}
\paragraph*{TEMPERATURES}

GSR 关键字 TEMPERATURES：用于设置与 TEMPERATURES 相关的 RedHawk 运行参数。（可选；默认：none）

\textbf{语法：}
\begin{lstlisting}
TEMPERATURES {
<layer_1> <temp_1 ºC>
...
<layer_n> <temp_n ºC>
}
\end{lstlisting}

\textbf{示例：}
\begin{lstlisting}
TEMPERATURES {
M1 111
M2 112
M3 113
M4 114
...
}
\end{lstlisting}

\subsubsection{\gsr{THERMAL_PROFILE}}
\paragraph*{THERMAL_PROFILE}

GSR 关键字 THERMAL_PROFILE：用于设置与 THERMAL、PROFILE 相关的 RedHawk 运行参数。（可选；默认：none）

\textbf{语法：}
\begin{lstlisting}
THERMAL_PROFILE {
FILE <thermal_profile_file>
}
\end{lstlisting}

\subsubsection{\gsr{USE_DEF_ROW_HEIGHT}}
\paragraph*{USE_DEF_ROW_HEIGHT}

GSR 关键字 USE_DEF_ROW_HEIGHT：用于设置与 USE、DEF、ROW、HEIGHT 相关的 RedHawk 运行参数。（可选；默认：), selects standard cells with heights u；兼容 DMP）

\textbf{语法：}
\begin{lstlisting}
Syntax
USE_DEF_ROW_HEIGHT [ 0 | 1 ]
\end{lstlisting}

\subsubsection{\gsr{USE_DEF_VIAMODEL}}
\paragraph*{USE_DEF_VIAMODEL}

GSR 关键字 USE_DEF_VIAMODEL：用于设置与 USE、DEF、过孔模型 相关的 RedHawk 运行参数。（默认：1）

\textbf{语法：}
\begin{lstlisting}
USE_DEF_VIAMODEL [ 0 | 1]
\end{lstlisting}

\subsubsection{\gsr{USE_DEF_VIARULE}}
\paragraph*{USE_DEF_VIARULE}

GSR 关键字 USE_DEF_VIARULE：用于设置与 USE、DEF、VIARULE 相关的 RedHawk 运行参数。（默认：0）

\textbf{语法：}
\begin{lstlisting}
USE_DEF_VIARULE [ 0| 1 ]
\end{lstlisting}

\subsubsection{\gsr{USE_LEF_FOR_LOGICAL_CONNECTION}}
\paragraph*{USE_LEF_FOR_LOGICAL_CONNECTION}

GSR 关键字 USE_LEF_FOR_LOGICAL_CONNECTION：用于设置与 USE、LEF、FOR、LOGICAL、连接 相关的 RedHawk 运行参数。（可选；默认：modeled as gray box cells (GRAY_BOX_CELL）

\textbf{语法：}
\begin{lstlisting}
Syntax
USE_LEF_FOR_LOGICAL_CONNECTION {
<cellname1> white
<cellname2> white
...
\end{lstlisting}

\subsubsection{\gsr{USER_STA_FILE}}
\paragraph*{USER_STA_FILE}

GSR 关键字 USER_STA_FILE：用于设置与 USER、STA、文件 相关的 RedHawk 运行参数。（可选；默认：none；兼容 DMP） 注意：GSR 关键字 USER_STA_FILE：用于设置与 USER、STA、文件 相关的 RedHawk 运行参数。（必需）

\textbf{语法：}
\begin{lstlisting}
Syntax
USER_STA_FILE <filename>
The allowed format of information in the specified USER_STA file is different
depending on whether a signal/clock net or a switch is specified, as follows
a. Signal or clock timing window
<inst_name>/pin_name> TW <min_arriv_time> <max_arriv_time> ?<Freq>? ?[s|c ]?
b. Switch timing window
<inst_name/pin_name> SW <min_arriv_time> <max_arriv_time> ?<Freq>?
c. Slew
<inst_name> SL <rise_trans_time> <fall_trans_time>
<inst_name>/<pin_name> SL <rise_trans_time> <fall_trans_time>
...
\end{lstlisting}

\textbf{参数说明：}
\begin{itemize}
\item \texttt{<inst\_name>} — GSR 关键字 USER\_STA\_FILE：用于设置与 USER、STA、文件 相关的 RedHawk 运行参数。（可选）
\end{itemize}

\textbf{示例：}
\begin{lstlisting}
top/block/inst1                SL 0.1e-9 0.1e-9
\end{lstlisting}

\begin{noteBox}
GSR 关键字 USER\_STA\_FILE：用于设置与 USER、STA、文件 相关的 RedHawk 运行参数。
\end{noteBox}

\begin{noteBox}
GSR 关键字 USER\_STA\_FILE：用于设置与 USER、STA、文件 相关的 RedHawk 运行参数。（必需）
\end{noteBox}

\subsubsection{\gsr{USE_SIGNAL_LOAD_FROM_STA}}
\paragraph*{USE_SIGNAL_LOAD_FROM_STA}

GSR 关键字 USE_SIGNAL_LOAD_FROM_STA：用于设置与 USE、SIGNAL、LOAD、从、STA 相关的 RedHawk 运行参数。（可选；默认：), or if there is no signal load data in；兼容 DMP）

\textbf{语法：}
\begin{lstlisting}
USE_SIGNAL_LOAD_FROM_STA [ 0 | 1 ]
\end{lstlisting}

\subsubsection{\gsr{VCD_FILE}}
\paragraph*{VCD_FILE}

GSR 关键字 VCD_FILE：用于设置与 VCD、文件 相关的 RedHawk 运行参数。（可选；默认：None；兼容 DMP；必需）

\textbf{语法：}
\begin{lstlisting}
VCD_FILE {
<top_block_name> <VCD file path>
? FILE_TYPE [ VCD | FSDB | RTL_VCD | RTL_FSDB ]?
? DUMP_SAIF_FILE <filename>
? VCD_DRIVEN [0|1]?
? FRONT_PATH <“string”> ?
? SUBSTITUTE_PATH <“string”> ?
? SELECT_RANGE <start_time> <end_time> ?
? SELECT_TYPE [ WORST_POWER_CYCLE | WORST_DPDT_CYCLE |
[WORST_TILE_POWER_CYCLE | WORST_TILE_DPDT_CYCLE ]?
? START_TIME <start> ? ? END_TIME <end> ?
? TRUE_TIME [0|1] ?
? MAPPING <map_file> ?
? MAPPING_CASE_SENSITIVE [ 1 | 0]?
? PROPAGATE 2 ?
? STA_VCD_FREQ_RATIO <Ratio> ?
}
\end{lstlisting}

\textbf{参数说明：}
\begin{itemize}
\item \texttt{<top\_block\_name>} — GSR 关键字 VCD\_FILE：用于设置与 VCD、文件 相关的 RedHawk 运行参数。（默认：VCD) DUMP\_SAIF\_FILE : creates a specifie）
\end{itemize}

\subsubsection{\gsr{VCD_FILE}}
\paragraph*{VCD_FILE}

GSR 关键字 VCD_FILE：用于设置与 VCD、文件 相关的 RedHawk 运行参数。

\subsubsection{\gsr{VCD_CONSISTENT_SCENARIO_FILTER}}
\paragraph*{VCD_CONSISTENT_SCENARIO_FILTER}

GSR 关键字 VCD_CONSISTENT_SCENARIO_FILTER：用于设置与 VCD、CONSISTENT、SCENARIO、FILTER 相关的 RedHawk 运行参数。（默认：off for averaged mode）

\textbf{语法：}
\begin{lstlisting}
VCD_CONSISTENT_SCENARIO_FILTER [ 0 | 1 ]
\end{lstlisting}

\subsection{参数关键字}
\subsubsection*{Parameter Keywords}

本类共 40 个 GSR 关键字。以下逐条给出中文用途、语法、默认值与示例。

\subsubsection{\gsr{APACHE_DB_OVERWRITE}}
\paragraph*{APACHE_DB_OVERWRITE}

GSR 关键字 APACHE_DB_OVERWRITE：用于设置与 APACHE、DB、覆盖 相关的 RedHawk 运行参数。（可选；默认：1 (overwrite)；兼容 DMP）

\textbf{语法：}
\begin{lstlisting}
APACHE_DB_OVERWRITE [ 1 | 0 ]
\end{lstlisting}

\subsubsection{\gsr{APACHE_FILES}}
\paragraph*{APACHE_FILES}

GSR 关键字 APACHE_FILES：用于设置与 APACHE、文件 相关的 RedHawk 运行参数。（可选；默认：normal；兼容 DMP）

\textbf{语法：}
\begin{lstlisting}
APACHE_FILES [ normal | clean ]
\end{lstlisting}

\textbf{示例：}
\begin{lstlisting}
APACHE_FILES clean
\end{lstlisting}

\subsubsection{\gsr{CACHE_DIRECTORY}}
\paragraph*{CACHE_DIRECTORY}

GSR 关键字 CACHE_DIRECTORY：用于设置与 CACHE、DIRECTORY 相关的 RedHawk 运行参数。

\textbf{语法：}
\begin{lstlisting}
CACHE_DIR <path to cache directory>
\end{lstlisting}

\textbf{参数说明：}
\begin{itemize}
\item \texttt{Default} — GSR 关键字 CACHE\_DIRECTORY：用于设置与 CACHE、DIRECTORY 相关的 RedHawk 运行参数。
\end{itemize}

\textbf{示例：}
\begin{lstlisting}
CACHE_DIR /local/cache
\end{lstlisting}

\subsubsection{\gsr{CACHE_MODE}}
\paragraph*{CACHE_MODE}

GSR 关键字 CACHE_MODE：用于设置与 CACHE、模式 相关的 RedHawk 运行参数。（可选；默认：0 (off)；兼容 DMP）

\textbf{语法：}
\begin{lstlisting}
CACHE_MODE [ 0 | 1 ]
\end{lstlisting}

\subsubsection{\gsr{DEF_PG_NETS_FILE}}
\paragraph*{DEF_PG_NETS_FILE}

GSR 关键字 DEF_PG_NETS_FILE：用于设置与 DEF、PG、网络、文件 相关的 RedHawk 运行参数。（可选；默认：none；兼容 DMP）

\textbf{语法：}
\begin{lstlisting}
DEF_PG_NETS_FILE <filename>
\end{lstlisting}

\subsubsection{\gsr{DEF_SCALING_FACTOR}}
\paragraph*{DEF_SCALING_FACTOR}

GSR 关键字 DEF_SCALING_FACTOR：用于设置与 DEF、SCALING、因子 相关的 RedHawk 运行参数。（可选；默认：1；兼容 DMP） 注意：GSR 关键字 DEF_SCALING_FACTOR：用于设置与 DEF、SCALING、因子 相关的 RedHawk 运行参数。

\textbf{语法：}
\begin{lstlisting}
DEF_SCALING_FACTOR {
[ All <scale_factor> |
<DEF_top_block_path-1> <scale_factor>
...
<DEF_top_block_path-n> <scale_factor> ]
\end{lstlisting}

\textbf{参数说明：}
\begin{itemize}
\item \texttt{All <value>} — GSR 关键字 DEF\_SCALING\_FACTOR：用于设置与 DEF、SCALING、因子 相关的 RedHawk 运行参数。（默认：is 0), the x,y coordinates in the PAD\_FI）
\end{itemize}

\begin{noteBox}
GSR 关键字 DEF\_SCALING\_FACTOR：用于设置与 DEF、SCALING、因子 相关的 RedHawk 运行参数。
\end{noteBox}

\subsubsection{\gsr{DEF_SCALING_FACTOR}}
\paragraph*{DEF_SCALING_FACTOR}

GSR 关键字 DEF_SCALING_FACTOR：用于设置与 DEF、SCALING、因子 相关的 RedHawk 运行参数。

\subsubsection{\gsr{DEF_SCALING_FACTOR}}
\paragraph*{DEF_SCALING_FACTOR}

GSR 关键字 DEF_SCALING_FACTOR：用于设置与 DEF、SCALING、因子 相关的 RedHawk 运行参数。

\subsubsection{\gsr{EVA_PG_AWARE}}
\paragraph*{EVA_PG_AWARE}

GSR 关键字 EVA_PG_AWARE：用于设置与 EVA、PG、AWARE 相关的 RedHawk 运行参数。（可选；默认：0 (off, do not consider P/G weakness)；兼容 DMP）

\textbf{语法：}
\begin{lstlisting}
EVA_PG_AWARE [ 1 | 0 ] ?<PG_switch_ratio>?
\end{lstlisting}

\textbf{参数说明：}
\begin{itemize}
\item \texttt{<PG\_switch\_ratio>} — GSR 关键字 EVA\_PG\_AWARE：用于设置与 EVA、PG、AWARE 相关的 RedHawk 运行参数。（默认：ratio: 0）
\end{itemize}

\textbf{示例：}
\begin{lstlisting}
EVA_PG_AWARE 1 0.2
\end{lstlisting}

\subsubsection{\gsr{FREQUENCY}}
\paragraph*{FREQUENCY}

GSR 关键字 FREQUENCY：用于设置与 频率 相关的 RedHawk 运行参数。（默认：none；兼容 DMP；必需）

\textbf{语法：}
\begin{lstlisting}
FREQUENCY <value in Hertz>
\end{lstlisting}

\textbf{示例：}
\begin{lstlisting}
FREQUENCY 160e6 (or 160M, 160me)
Note that ‘me’ represents mega.
\end{lstlisting}

\subsubsection{\gsr{GENERATE_CPM}}
\paragraph*{GENERATE_CPM}

GSR 关键字 GENERATE_CPM：用于设置与 GENERATE、CPM 相关的 RedHawk 运行参数。（可选；默认：0；兼容 DMP）

\textbf{语法：}
\begin{lstlisting}
GENERATE_CPM [ 0 | 1 ]
\end{lstlisting}

\subsubsection{\gsr{GND_NETS}}
\paragraph*{GND_NETS}

GSR 关键字 GND_NETS：用于设置与 GND、网络 相关的 RedHawk 运行参数。（可选；默认：all the SPECIALNETS are designated as US；兼容 DMP） 注意：GSR 关键字 GND_NETS：用于设置与 GND、网络 相关的 RedHawk 运行参数。

\begin{noteBox}
GSR 关键字 GND\_NETS：用于设置与 GND、网络 相关的 RedHawk 运行参数。
\end{noteBox}

\subsubsection{\gsr{GND_NETS}}
\paragraph*{GND_NETS}

GSR 关键字 GND_NETS：用于设置与 GND、网络 相关的 RedHawk 运行参数。

\textbf{参数说明：}
\begin{itemize}
\item \texttt{<Vss\_domain\_net\_name>} — GSR 关键字 GND\_NETS：用于设置与 GND、网络 相关的 RedHawk 运行参数。
\end{itemize}

\subsubsection{\gsr{GND_NETS}}
\paragraph*{GND_NETS}

GSR 关键字 GND_NETS：用于设置与 GND、网络 相关的 RedHawk 运行参数。

\subsubsection{\gsr{LEF_SCALING_FACTOR}}
\paragraph*{LEF_SCALING_FACTOR}

GSR 关键字 LEF_SCALING_FACTOR：用于设置与 LEF、SCALING、因子 相关的 RedHawk 运行参数。（可选；默认：1；兼容 DMP）

\subsubsection{\gsr{LEF_SCALING_FACTOR}}
\paragraph*{LEF_SCALING_FACTOR}

GSR 关键字 LEF_SCALING_FACTOR：用于设置与 LEF、SCALING、因子 相关的 RedHawk 运行参数。

\textbf{参数说明：}
\begin{itemize}
\item \texttt{All <value>} — GSR 关键字 LEF\_SCALING\_FACTOR：用于设置与 LEF、SCALING、因子 相关的 RedHawk 运行参数。
\end{itemize}

\subsubsection{\gsr{LEF_SCALING_FACTOR}}
\paragraph*{LEF_SCALING_FACTOR}

GSR 关键字 LEF_SCALING_FACTOR：用于设置与 LEF、SCALING、因子 相关的 RedHawk 运行参数。

\subsubsection{\gsr{LICENSE_WAIT}}
\paragraph*{LICENSE_WAIT}

GSR 关键字 LICENSE_WAIT：用于设置与 LICENSE、WAIT 相关的 RedHawk 运行参数。（可选；默认：0 (off)；兼容 DMP） 注意：GSR 关键字 LICENSE_WAIT：用于设置与 LICENSE、WAIT 相关的 RedHawk 运行参数。

\textbf{语法：}
\begin{lstlisting}
LICENSE_WAIT [ 0 | 1 ]
\end{lstlisting}

\begin{noteBox}
GSR 关键字 LICENSE\_WAIT：用于设置与 LICENSE、WAIT 相关的 RedHawk 运行参数。
\end{noteBox}

\subsubsection{\gsr{MMX_RES_MAP_LIMIT}}
\paragraph*{MMX_RES_MAP_LIMIT}

GSR 关键字 MMX_RES_MAP_LIMIT：用于设置与 MMX、RES、MAP、LIMIT 相关的 RedHawk 运行参数。（可选；默认：, based on P/G grid weakness computation）

\textbf{语法：}
\begin{lstlisting}
MMX_RES_MAP_LIMIT <pin_limit>
\end{lstlisting}

\subsubsection{\gsr{MULTI_GND_PACKAGE_MODEL}}
\paragraph*{MULTI_GND_PACKAGE_MODEL}

GSR 关键字 MULTI_GND_PACKAGE_MODEL：用于设置与 MULTI、GND、PACKAGE、MODEL 相关的 RedHawk 运行参数。（可选；默认：, multi-ground domain designs have all g）

\subsubsection{\gsr{MULTI_THREADS}}
\paragraph*{MULTI_THREADS}

GSR 关键字 MULTI_THREADS：用于设置与 MULTI、THREADS 相关的 RedHawk 运行参数。（可选；默认：RedHawk all available CPUs are used；兼容 DMP）

\textbf{语法：}
\begin{lstlisting}
MULTI_THREADS [On=1|Off=0] ?<max_num_threads>?
\end{lstlisting}

\textbf{示例：}
\begin{lstlisting}
MULTI_THREADS 1 (uses all available CPUs- default)
MULTI_THREADS 0 (uses single-thread)
MULTI_THREADS 1 1 (uses single-thread)
MULTI_THREADS 1 4 (uses multi-thread with four CPUs)
\end{lstlisting}

\subsubsection{\gsr{PGPLOC_DEBUG}}
\paragraph*{PGPLOC_DEBUG}

GSR 关键字 PGPLOC_DEBUG：用于设置与 PGPLOC、DEBUG 相关的 RedHawk 运行参数。（可选；默认：0 (Off)；兼容 DMP）

\textbf{语法：}
\begin{lstlisting}
PGPLOC_DEBUG [ 0 | 1 ]
\end{lstlisting}

\subsubsection{\gsr{PRINT_EM_VIA_BOX}}
\paragraph*{PRINT_EM_VIA_BOX}

GSR 关键字 PRINT_EM_VIA_BOX：用于设置与 PRINT、电迁移、过孔、BOX 相关的 RedHawk 运行参数。（可选；默认：0；兼容 DMP）

\textbf{语法：}
\begin{lstlisting}
PRINT_EM_VIA_BOX [ 0 | 1 ]
\end{lstlisting}

\subsubsection{\gsr{PRINT_EM_VIA_INFO}}
\paragraph*{PRINT_EM_VIA_INFO}

GSR 关键字 PRINT_EM_VIA_INFO：用于设置与 PRINT、电迁移、过孔、INFO 相关的 RedHawk 运行参数。（可选；默认：0；兼容 DMP）

\textbf{语法：}
\begin{lstlisting}
PRINT_EM_VIA_INFO [ 0 | 1 ]
\end{lstlisting}

\subsubsection{\gsr{REPORT_DISCONN_MIN_LENGTH}}
\paragraph*{REPORT_DISCONN_MIN_LENGTH}

GSR 关键字 REPORT_DISCONN_MIN_LENGTH：用于设置与 REPORT、DISCONN、MIN、LENGTH 相关的 RedHawk 运行参数。（可选；默认：0）

\textbf{语法：}
\begin{lstlisting}
REPORT_DISCONN_MIN_LENGTH <min_length_microns>
\end{lstlisting}

\subsubsection{\gsr{REPORT_FLATTEN_LOG}}
\paragraph*{REPORT_FLATTEN_LOG}

GSR 关键字 REPORT_FLATTEN_LOG：用于设置与 REPORT、FLATTEN、日志 相关的 RedHawk 运行参数。（可选；默认：0；兼容 DMP）

\textbf{语法：}
\begin{lstlisting}
REPORT_FLATTEN_LOG [0 |1]
\end{lstlisting}

\subsubsection{\gsr{REPORT_PEAK_MEMORY}}
\paragraph*{REPORT_PEAK_MEMORY}

GSR 关键字 REPORT_PEAK_MEMORY：用于设置与 REPORT、峰值、MEMORY 相关的 RedHawk 运行参数。（可选；默认：memory use at the end of each stage is r；兼容 DMP）

\textbf{语法：}
\begin{lstlisting}
REPORT_PEAK_MEMORY [0|1]
\end{lstlisting}

\subsubsection{\gsr{REPORT_REDUCTION}}
\paragraph*{REPORT_REDUCTION}

GSR 关键字 REPORT_REDUCTION：用于设置与 REPORT、REDUCTION 相关的 RedHawk 运行参数。（默认：) normal/ 1: files that can be created f）

\textbf{语法：}
\begin{lstlisting}
REPORT_REDUCTION [ off/0 | normal/1 | max/2 ]
\end{lstlisting}

\textbf{参数说明：}
GSR 关键字 REPORT\_REDUCTION：用于设置与 REPORT、REDUCTION 相关的 RedHawk 运行参数。（可选；默认：file name is adsRpt/Dynamic/mcyc\_effvdd；兼容 DMP）

\subsubsection{\gsr{REPORT_PEAK_MEMORY}}
\paragraph*{REPORT_PEAK_MEMORY}

GSR 关键字 REPORT_PEAK_MEMORY：用于设置与 REPORT、峰值、MEMORY 相关的 RedHawk 运行参数。（可选；默认：memory use at the end of each stage is r）

\textbf{语法：}
\begin{lstlisting}
REPORT_PEAK_MEMORY [0|1]
\end{lstlisting}

\subsubsection{\gsr{SPLIT_VDD_EXTRACT_LP}}
\paragraph*{SPLIT_VDD_EXTRACT_LP}

GSR 关键字 SPLIT_VDD_EXTRACT_LP：用于设置与 SPLIT、VDD、提取、LP 相关的 RedHawk 运行参数。

\subsubsection{\gsr{SPLIT_VDD_EXTRACT_LP_FSIZE}}
\paragraph*{SPLIT_VDD_EXTRACT_LP_FSIZE}

GSR 关键字 SPLIT_VDD_EXTRACT_LP_FSIZE：用于设置与 SPLIT、VDD、提取、LP、FSIZE 相关的 RedHawk 运行参数。（可选；默认：5M node threshold at which point the fil）

\textbf{语法：}
\begin{lstlisting}
SPLIT_VDD_EXTRACT_LP <value>
The number of nodes in each split vdd file can be adjusted using the keyword
'SPLIT_VDD_EXTRACT_LP_FSIZE <value>'. Optional. Default: 1 M.
SPLIT_VDD_EXTRACT_LP_FSIZE <value>
\end{lstlisting}

\subsubsection{\gsr{STATIC_CONNECT_INST_FLOAT_GND}}
\paragraph*{STATIC_CONNECT_INST_FLOAT_GND}

GSR 关键字 STATIC_CONNECT_INST_FLOAT_GND：用于设置与 静态、CONNECT、INST、FLOAT、GND 相关的 RedHawk 运行参数。（可选；默认：, instances with floating grounds are a）

\textbf{语法：}
\begin{lstlisting}
STATIC_CONNECT_INST_FLOAT_GND [0 | 1]
\end{lstlisting}

\subsubsection{\gsr{STD_CELL_SINGLE_NOMINAL_VOLTAGE_ONLY}}
\paragraph*{STD_CELL_SINGLE_NOMINAL_VOLTAGE_ONLY}

GSR 关键字 STD_CELL_SINGLE_NOMINAL_VOLTAGE_ONLY：用于设置与 STD、单元、SINGLE、NOMINAL、VOLTAGE、仅 相关的 RedHawk 运行参数。（可选；默认：0）

\textbf{语法：}
\begin{lstlisting}
STD_CELL_SINGLE_NOMINAL_VOLTAGE_ONLY [ 0 | 1 ]
\end{lstlisting}

\subsubsection{\gsr{STD_INST_SINK_VIA_LAYERS}}
\paragraph*{STD_INST_SINK_VIA_LAYERS}

GSR 关键字 STD_INST_SINK_VIA_LAYERS：用于设置与 STD、INST、SINK、过孔、金属层 相关的 RedHawk 运行参数。（可选；默认：0）

\textbf{语法：}
\begin{lstlisting}
STD_INST_SINK_VIA_LAYERS {
<via_layers>
}
\end{lstlisting}

\subsubsection{\gsr{TEMPERATURE_DEVICE}}
\paragraph*{TEMPERATURE_DEVICE}

GSR 关键字 TEMPERATURE_DEVICE：用于设置与 温度、DEVICE 相关的 RedHawk 运行参数。（可选；默认：value is 5 degrees；兼容 DMP；必需）

\textbf{语法：}
\begin{lstlisting}
TEMPERATURE_DEVICE <T> ? <delta> ?
\end{lstlisting}

\subsubsection{\gsr{THERMAL_TEMP_RANGE}}
\paragraph*{THERMAL_TEMP_RANGE}

GSR 关键字 THERMAL_TEMP_RANGE：用于设置与 THERMAL、TEMP、RANGE 相关的 RedHawk 运行参数。（可选；默认：none）

\textbf{语法：}
\begin{lstlisting}
THERMAL_TEMP_RANGE <min_temp> <max_temp>
\end{lstlisting}

\subsubsection{\gsr{VDD_NETS}}
\paragraph*{VDD_NETS}

GSR 关键字 VDD_NETS：用于设置与 VDD、网络 相关的 RedHawk 运行参数。（默认：none；兼容 DMP；必需） 注意：GSR 关键字 VDD_NETS：用于设置与 VDD、网络 相关的 RedHawk 运行参数。

\textbf{语法：}
\begin{lstlisting}
VDD_NETS {
<Vdd_domain_net_name> <value_Volts> {
<equiv_Vdd_net_name1>
...
}
...
}
\end{lstlisting}

\textbf{参数说明：}
\begin{itemize}
\item \texttt{<Vdd\_domain\_net\_name>} — GSR 关键字 VDD\_NETS：用于设置与 VDD、网络 相关的 RedHawk 运行参数。
\end{itemize}

\textbf{示例：}
\begin{lstlisting}
VDD_NETS {
VDD 1.1 {
\end{lstlisting}

\begin{noteBox}
GSR 关键字 VDD\_NETS：用于设置与 VDD、网络 相关的 RedHawk 运行参数。
\end{noteBox}

\subsubsection{\gsr{VDD1CORE}}
\paragraph*{VDD1CORE}

GSR 关键字 VDD1CORE：用于设置与 VDD1CORE 相关的 RedHawk 运行参数。

\subsubsection{\gsr{VDDQ1}}
\paragraph*{VDDQ1}

GSR 关键字 VDDQ1：用于设置与 VDDQ1 相关的 RedHawk 运行参数。

\subsubsection{\gsr{VIA_IR_REPORT}}
\paragraph*{VIA_IR_REPORT}

GSR 关键字 VIA_IR_REPORT：用于设置与 过孔、IR、REPORT 相关的 RedHawk 运行参数。（可选；默认：output files are: If set to 2, the repor；兼容 DMP）

\textbf{语法：}
\begin{lstlisting}
VIA_IR_REPORT [ 0 | 1 ]
\end{lstlisting}

\subsection{自定义单元建模（CMM）关键字}
\subsubsection*{Custom Cell Modeling Keywords}

本类共 134 个 GSR 关键字。以下逐条给出中文用途、语法、默认值与示例。

\subsubsection{\gsr{CMM_CELLS}}
\paragraph*{CMM_CELLS}

GSR 关键字 CMM_CELLS：用于设置与 CMM、单元 相关的 RedHawk 运行参数。（可选；默认：“optimize” for cell_view and “original”；兼容 DMP）

\textbf{语法：}
\begin{lstlisting}
CMM_CELLS {
<cell_name> <model_cell_path> [original |optimize |raw ]
...
}
\end{lstlisting}

\subsubsection{\gsr{CMM_EXCLUDE_FILES}}
\paragraph*{CMM_EXCLUDE_FILES}

GSR 关键字 CMM_EXCLUDE_FILES：用于设置与 CMM、EXCLUDE、文件 相关的 RedHawk 运行参数。（可选；默认：, all available files in CMM models are）

\textbf{语法：}
\begin{lstlisting}
CMM_EXCLUDE_FILES {
<modelname1> <file_type1> <file_type2> ...
<modelname2> <file_type1> <file_type2> ...
...
}
\end{lstlisting}

\textbf{参数说明：}
\begin{itemize}
\item \texttt{<modelname1>} — GSR 关键字 CMM\_EXCLUDE\_FILES：用于设置与 CMM、EXCLUDE、文件 相关的 RedHawk 运行参数。
\end{itemize}

\textbf{示例：}
\begin{lstlisting}
CMM_EXCLUDE_FILES {
M1 ipf vcd spf sta ...
M2 bpfs gsc
}
So in the example, file types IPF, VCD, SPF and STA should not be imported for
CMM model M1, and BPFS and GSC information should not be imported for model
M2. Note that BLOCK_POWER_ASSIGNMENT (BPA) cannot be excluded from
the CMM.
Extra caution is needed when importing any instance power file (IPF), since any
instance not covered by IPF is assigned a power of 0, which may not be desirable.
So if any IPF file is imported, a warning is displayed that an IPF was imported for
some CMM models, but not for top level or other CMM.
\end{lstlisting}

\subsubsection{\gsr{CMM_EXPAND_PINS_AT_TOP}}
\paragraph*{CMM_EXPAND_PINS_AT_TOP}

GSR 关键字 CMM_EXPAND_PINS_AT_TOP：用于设置与 CMM、EXPAND、引脚、AT、顶层 相关的 RedHawk 运行参数。（可选；默认：0）

\textbf{语法：}
\begin{lstlisting}
CMM_EXPAND_PINS_AT_TOP [ 0 | 1 ]
\end{lstlisting}

\subsubsection{\gsr{CMM_INSTANCES}}
\paragraph*{CMM_INSTANCES}

GSR 关键字 CMM_INSTANCES：用于设置与 CMM、实例 相关的 RedHawk 运行参数。（可选；默认：none）

\textbf{语法：}
\begin{lstlisting}
CMM_INSTANCES {
<cmm_instance_path> {original | optimize}
...
}
\end{lstlisting}

\subsubsection{\gsr{CMM_LAYER_MAP_FILES}}
\paragraph*{CMM_LAYER_MAP_FILES}

GSR 关键字 CMM_LAYER_MAP_FILES：用于设置与 CMM、层、MAP、文件 相关的 RedHawk 运行参数。（可选；默认：none；兼容 DMP） 注意：GSR 关键字 CMM_LAYER_MAP_FILES：用于设置与 CMM、层、MAP、文件 相关的 RedHawk 运行参数。

\textbf{语法：}
\begin{lstlisting}
CMM_LAYER_MAP_FILES {
<CMM_cellname> [<filename> | AUTO | UNUSED]
...
}
\end{lstlisting}

\textbf{示例：}
\begin{lstlisting}
Assume that your CMM tech file has 12 layers, C1 to C12, and your top tech file has
16 layers, T1 to T16. In the layer map file, you would specify layers C1 to C4 as
“UNUSED”, then layers C5 to C12 map from T1 to T8 (these must be in stack order
in one-to-one mapping, so C5 <-> T1, C6<->T2, ..., C12<->T8. The combined tech
file for the top then is C1, C2, C3, C4, T1, T2, ...., T16. The layer properties for C1
to C4 layers in top tech are copied from the CMM tech file, the C5 to C12 layers in
the CMM tech are dropped, and T1 and T8 layer properties from top tech are used
for the layers in CMM. RedHawk does not check the property of layers C5 to C12 to
see if they match T1 to T8; RedHawk only checks that the temperatures used for
extraction for layer C5 to C12 matches the temperatures of layers T1 to T8.
For three CMM blocks in a top level run, CMM_LAYER_MAP_FILES could describe
the layers as follows:
CMM_LAYER_MAP_FILES {
dpL1DDATARAM ./myLayerMapFile
dpIDATARAM AUTO
# dpLSReqQueCam
\end{lstlisting}

\begin{noteBox}
GSR 关键字 CMM\_LAYER\_MAP\_FILES：用于设置与 CMM、层、MAP、文件 相关的 RedHawk 运行参数。
\end{noteBox}

\subsubsection{\gsr{CMM_POWER_OVERRIDE_IPF}}
\paragraph*{CMM_POWER_OVERRIDE_IPF}

GSR 关键字 CMM_POWER_OVERRIDE_IPF：用于设置与 CMM、电源、OVERRIDE、IPF 相关的 RedHawk 运行参数。（默认：0）

\textbf{语法：}
\begin{lstlisting}
CMM_POWER_OVERRIDE_IPF [ 0 | 1 ]
Power calculation keywords
\end{lstlisting}

\subsubsection{\gsr{APL_INTERPOLATION_METHOD}}
\paragraph*{APL_INTERPOLATION_METHOD}

GSR 关键字 APL_INTERPOLATION_METHOD：用于设置与 APL、INTERPOLATION、METHOD 相关的 RedHawk 运行参数。（默认：0；兼容 DMP）

\textbf{语法：}
\begin{lstlisting}
APL_INTERPOLATION_METHOD [0 | 4 ]
\end{lstlisting}

\subsubsection{\gsr{APL_MODE}}
\paragraph*{APL_MODE}

GSR 关键字 APL_MODE：用于设置与 APL、模式 相关的 RedHawk 运行参数。（可选；默认：(Accurate) mode, some grouping of APL sa）

\textbf{语法：}
\begin{lstlisting}
APL_MODE [ Accurate | Extended ]
\end{lstlisting}

\subsubsection{\gsr{BIASPIN_CONFIG_FILE}}
\paragraph*{BIASPIN_CONFIG_FILE}

GSR 关键字 BIASPIN_CONFIG_FILE：用于设置与 BIASPIN、CONFIG、文件 相关的 RedHawk 运行参数。（可选；默认：None；兼容 DMP）

\textbf{语法：}
\begin{lstlisting}
BIASPIN_CONFIG_FILE <filename>
The configuration file format is as follows:
<instance_name><biasPin1>=<biasV_1> [... <biasPinN>=<biasV_N>]
...
\end{lstlisting}

\subsubsection{\gsr{BLOCK_INSTANCE_POWER_FILE}}
\paragraph*{BLOCK_INSTANCE_POWER_FILE}

GSR 关键字 BLOCK_INSTANCE_POWER_FILE：用于设置与 模块、实例、电源、文件 相关的 RedHawk 运行参数。（兼容 DMP）

\textbf{语法：}
\begin{lstlisting}
BLOCK_INSTANCE_POWER_FILE {
FULLCHIP <block_name> <file_name>
CELLTYPE <block_cell_name> <file_name>
\end{lstlisting}

\textbf{参数说明：}
GSR 关键字 BLOCK\_INSTANCE\_POWER\_FILE：用于设置与 模块、实例、电源、文件 相关的 RedHawk 运行参数。

\subsubsection{\gsr{BLOCK_POWER_FOR_SCALING}}
\paragraph*{BLOCK_POWER_FOR_SCALING}

GSR 关键字 BLOCK_POWER_FOR_SCALING：用于设置与 模块、电源、FOR、SCALING 相关的 RedHawk 运行参数。（可选；默认：none；兼容 DMP）

\textbf{语法：}
\begin{lstlisting}
BLOCK_POWER_FOR_SCALING {
CELLTYPE <cell_name> [<pwr_W> ?<Vdd_pin>? |
<cell_current_A> ?<Vss_pin> ? ]
...
[FULLCHIP | BLOCK | <top_block_name>] <full_leaf_inst_path>
<pwr_W>
[FULLCHIP | BLOCK |<top_block_name>] <full_block_path>
<pwr_W> ? <domain> ? ? <pin> ?
...
[FULLCHIP | BLOCK | <top_block_name>]
<full_mVdd_inst_path>
[<pwr_W> ?<VDD_pin>? | <current_A> ?<Vss_pin>?]
...
CELLS [ comb | ff_latch | mem | clockinst | io ] <pwr_W>
...
STDCELL <cell_name> <power_W> <pin_name>
...
}
\end{lstlisting}

\subsubsection{\gsr{BLOCK_POWER_FOR_SCALING}}
\paragraph*{BLOCK_POWER_FOR_SCALING}

GSR 关键字 BLOCK_POWER_FOR_SCALING：用于设置与 模块、电源、FOR、SCALING 相关的 RedHawk 运行参数。

\textbf{参数说明：}
\begin{itemize}
\item \texttt{CELLTYPE <cell\_name>} — GSR 关键字 BLOCK\_POWER\_FOR\_SCALING：用于设置与 模块、电源、FOR、SCALING 相关的 RedHawk 运行参数。（可选）
\end{itemize}

\textbf{示例：}
\begin{lstlisting}
CELLTYPE cellA1 0.01
CELLTYPE cell3 0.001
FULLCHIP designABC 1.5
BLOCK block1 0.5
FULLCHIP block1/instance2 0.00322
BLOCK block1/sub_blockA 0.1
FULLCHIP block2/AN210J 0.001 VDD
FULLCHIP block2/AN210J 0.003 VDDC
FULLCHIP ia32/nand231 0.001 VDDQ
FULLCHIP ia32/nand342 0.005 VDD5
\end{lstlisting}

\subsubsection{\gsr{BLOCK_POWER_FOR_SCALING_FILES}}
\paragraph*{BLOCK_POWER_FOR_SCALING_FILES}

GSR 关键字 BLOCK_POWER_FOR_SCALING_FILES：用于设置与 模块、电源、FOR、SCALING、文件 相关的 RedHawk 运行参数。（可选；默认：None；兼容 DMP）

\subsubsection{\gsr{BLOCK_POWER_FOR_SCALING_FILES}}
\paragraph*{BLOCK_POWER_FOR_SCALING_FILES}

GSR 关键字 BLOCK_POWER_FOR_SCALING_FILES：用于设置与 模块、电源、FOR、SCALING、文件 相关的 RedHawk 运行参数。

\subsubsection{\gsr{BLOCK_POWER_FOR_SCALING_FILES}}
\paragraph*{BLOCK_POWER_FOR_SCALING_FILES}

GSR 关键字 BLOCK_POWER_FOR_SCALING_FILES：用于设置与 模块、电源、FOR、SCALING、文件 相关的 RedHawk 运行参数。

\subsubsection{\gsr{BLOCK_TOGGLE_FILES}}
\paragraph*{BLOCK_TOGGLE_FILES}

GSR 关键字 BLOCK_TOGGLE_FILES：用于设置与 模块、翻转、文件 相关的 RedHawk 运行参数。（可选；默认：None；兼容 DMP） 注意：GSR 关键字 BLOCK_TOGGLE_FILES：用于设置与 模块、翻转、文件 相关的 RedHawk 运行参数。

\begin{noteBox}
GSR 关键字 BLOCK\_TOGGLE\_FILES：用于设置与 模块、翻转、文件 相关的 RedHawk 运行参数。
\end{noteBox}

\subsubsection{\gsr{BLOCK_TOGGLE_FILES}}
\paragraph*{BLOCK_TOGGLE_FILES}

GSR 关键字 BLOCK_TOGGLE_FILES：用于设置与 模块、翻转、文件 相关的 RedHawk 运行参数。

\textbf{参数说明：}
\begin{itemize}
\item \texttt{<block\_name\_N>} — GSR 关键字 BLOCK\_TOGGLE\_FILES：用于设置与 模块、翻转、文件 相关的 RedHawk 运行参数。
\end{itemize}

\subsubsection{\gsr{BLOCK_TOGGLE_FILES}}
\paragraph*{BLOCK_TOGGLE_FILES}

GSR 关键字 BLOCK_TOGGLE_FILES：用于设置与 模块、翻转、文件 相关的 RedHawk 运行参数。

\subsubsection{\gsr{BLOCK_TOGGLE_RATE}}
\paragraph*{BLOCK_TOGGLE_RATE}

GSR 关键字 BLOCK_TOGGLE_RATE：用于设置与 模块、翻转、率 相关的 RedHawk 运行参数。（默认：toggle rate of the nets in a user-specif）

\textbf{示例：}
\begin{lstlisting}
example, ABC* matches all block instance and instance names starting with ABC.
DMP compatible. Optional. Default: clock rate: 2.0.
\end{lstlisting}

\subsubsection{\gsr{BLOCK_TOGGLE_RATE}}
\paragraph*{BLOCK_TOGGLE_RATE}

GSR 关键字 BLOCK_TOGGLE_RATE：用于设置与 模块、翻转、率 相关的 RedHawk 运行参数。

\subsubsection{\gsr{BLOCK_TOGGLE_RATE}}
\paragraph*{BLOCK_TOGGLE_RATE}

GSR 关键字 BLOCK_TOGGLE_RATE：用于设置与 模块、翻转、率 相关的 RedHawk 运行参数。

\subsubsection{\gsr{BLOCK_TOGGLE_RATE_FILES}}
\paragraph*{BLOCK_TOGGLE_RATE_FILES}

GSR 关键字 BLOCK_TOGGLE_RATE_FILES：用于设置与 模块、翻转、率、文件 相关的 RedHawk 运行参数。（可选；默认：None；兼容 DMP）

\textbf{语法：}
\begin{lstlisting}
BLOCK_TOGGLE_RATE_FILES <toggle_rate_filename>
\end{lstlisting}

\textbf{示例：}
\begin{lstlisting}
BLOCK_TOGGLE_RATE_FILES ABCDE.block_TR
\end{lstlisting}

\subsubsection{\gsr{BLOCK_SAIF_FILES}}
\paragraph*{BLOCK_SAIF_FILES}

GSR 关键字 BLOCK_SAIF_FILES：用于设置与 模块、SAIF、文件 相关的 RedHawk 运行参数。（可选；默认：none；兼容 DMP）

\textbf{语法：}
\begin{lstlisting}
BLOCK_SAIF_FILES {
SAIF_FILE {
\end{lstlisting}

\textbf{示例：}
\begin{lstlisting}
BLOCK_SAIF_FILES {
SAIF_FILE {
green_core3 green.saif
ROOT_INSTANCE green_tst/inst1/
SUBSTITUTE_PATH green_core3/
}
}
\end{lstlisting}

\subsubsection{\gsr{CELL_RC_FILES}}
\paragraph*{CELL_RC_FILES}

GSR 关键字 CELL_RC_FILES：用于设置与 单元、RC、文件 相关的 RedHawk 运行参数。（可选；默认：), RedHawk builds a c1-r-c2 equivalent c；兼容 DMP）

\subsubsection{\gsr{CELL_RC_FILE}}
\paragraph*{CELL_RC_FILE}

GSR 关键字 CELL_RC_FILE：用于设置与 单元、RC、文件 相关的 RedHawk 运行参数。

\subsubsection{\gsr{CELL_RC_FILE}}
\paragraph*{CELL_RC_FILE}

GSR 关键字 CELL_RC_FILE：用于设置与 单元、RC、文件 相关的 RedHawk 运行参数。

\subsubsection{\gsr{CELL_TOGGLE_RATE}}
\paragraph*{CELL_TOGGLE_RATE}

This keyword 允许指定 individual 翻转率s by leaf 单元 name. 兼容 DMP 分布式流程。 可选。默认：无。（可选；默认：None；兼容 DMP）

\textbf{语法：}
\begin{lstlisting}
CELL_TOGGLE_RATE {
<leaf_cell_name> <toggle_rate> ?<clock_toggle_rate>?
...
}
\end{lstlisting}

\subsubsection{\gsr{CELL_TOGGLE_RATE_FILE}}
\paragraph*{CELL_TOGGLE_RATE_FILE}

GSR 关键字 CELL_TOGGLE_RATE_FILE：用于设置与 单元、翻转、率、文件 相关的 RedHawk 运行参数。（默认：none；兼容 DMP；必需）

\textbf{语法：}
\begin{lstlisting}
CELL_TOGGLE_RATE_FILE <toggle_rate_file>
or
CELL_TOGGLE_RATE_FILE {
<toggle_rate_file1>
<toggle_rate_file2>
...
}
\end{lstlisting}

\subsubsection{\gsr{CLOCK_DOMAIN_TOGGLE_RATE}}
\paragraph*{CLOCK_DOMAIN_TOGGLE_RATE}

When set, 允许您 set a 默认 翻转率 for 实例 in certain clock domains, and 可选ly the clock 翻转率. 兼容 DMP 分布式流程。 默认: none.（默认：toggle rate for instances in certain clo；兼容 DMP）

\textbf{语法：}
\begin{lstlisting}
CLOCK_DOMAIN_TOGGLE_RATE <clock_domain_name>
<toggle_rate> ? <clock_toggle_rate> ?
\end{lstlisting}

\subsubsection{\gsr{CLOCK_DOMAIN_TOGGLE_RATE_FILE}}
\paragraph*{CLOCK_DOMAIN_TOGGLE_RATE_FILE}

GSR 关键字 CLOCK_DOMAIN_TOGGLE_RATE_FILE：用于设置与 CLOCK、DOMAIN、翻转、率、文件 相关的 RedHawk 运行参数。（默认：toggle rates for instances in certain cl；兼容 DMP）

\textbf{语法：}
\begin{lstlisting}
CLOCK_DOMAIN_TOGGLE_RATE_FILE                <filename>
The format of the file is shown below.
<clock_domain_name1> <toggle_rate1> ?<clock_toggle_rate>?
...
\end{lstlisting}

\subsubsection{\gsr{CYCLE_SELECTION_GRID_SIZE}}
\paragraph*{CYCLE_SELECTION_GRID_SIZE}

GSR 关键字 CYCLE_SELECTION_GRID_SIZE：用于设置与 CYCLE、SELECTION、GRID、SIZE 相关的 RedHawk 运行参数。（默认：none）

\textbf{语法：}
\begin{lstlisting}
CYCLE_SELECTION_GRID_SIZE <time_in_seconds>
\end{lstlisting}

\subsubsection{\gsr{GDS_BPFS_OVERRIDE_IPF}}
\paragraph*{GDS_BPFS_OVERRIDE_IPF}

GSR 关键字 GDS_BPFS_OVERRIDE_IPF：用于设置与 GDS、BPFS、OVERRIDE、IPF 相关的 RedHawk 运行参数。（默认：0）

\textbf{语法：}
\begin{lstlisting}
GDS_BPFS_OVERRIDE_IPF [ 0 | 1 ]
\end{lstlisting}

\subsubsection{\gsr{GSC_OVERRIDE_IPF}}
\paragraph*{GSC_OVERRIDE_IPF}

GSR 关键字 GSC_OVERRIDE_IPF：用于设置与 GSC、OVERRIDE、IPF 相关的 RedHawk 运行参数。（可选；默认：0；兼容 DMP）

\textbf{语法：}
\begin{lstlisting}
GSC_OVERRIDE_IPF [ 0 | 1 ]
\end{lstlisting}

\textbf{示例：}
\begin{lstlisting}
GSC_OVERRIDE_IPF 1
\end{lstlisting}

\subsubsection{\gsr{INSTANCE_POWER_FILES}}
\paragraph*{INSTANCE_POWER_FILES}

GSR 关键字 INSTANCE_POWER_FILES：用于设置与 实例、电源、文件 相关的 RedHawk 运行参数。（可选；默认：None；兼容 DMP）

\textbf{语法：}
\begin{lstlisting}
INSTANCE_POWER_FILES [ <inst_power_file_path> |
{ <inst1_power_file_path1>
<inst2_power_file_path2>
<inst3_power_file_path3> <inst3_hierpath>
<inst4_power_file_path4> <inst4_hierpath>
...
} ]
\end{lstlisting}

\textbf{参数说明：}
\begin{itemize}
\item \texttt{<instN\_power\_file\_pathN>} — GSR 关键字 INSTANCE\_POWER\_FILES：用于设置与 实例、电源、文件 相关的 RedHawk 运行参数。
\end{itemize}

\subsubsection{\gsr{INSTANCE_POWER_FILES}}
\paragraph*{INSTANCE_POWER_FILES}

GSR 关键字 INSTANCE_POWER_FILES：用于设置与 实例、电源、文件 相关的 RedHawk 运行参数。

\subsubsection{\gsr{INSTANCE_TOGGLE_RATE}}
\paragraph*{INSTANCE_TOGGLE_RATE}

GSR 关键字 INSTANCE_TOGGLE_RATE：用于设置与 实例、翻转、率 相关的 RedHawk 运行参数。（可选；默认：clock: 2；兼容 DMP）

\textbf{语法：}
\begin{lstlisting}
INSTANCE_TOGGLE_RATE {
<non-clock_inst_name> <output_TR> ?<clock_network_TR>?
<clock_net_inst_name> <output_TR> ?<clock_network_TR>?
...
}
\end{lstlisting}

\subsubsection{\gsr{INSTANCE_TOGGLE_RATE}}
\paragraph*{INSTANCE_TOGGLE_RATE}

GSR 关键字 INSTANCE_TOGGLE_RATE：用于设置与 实例、翻转、率 相关的 RedHawk 运行参数。

\subsubsection{\gsr{INSTANCE_TOGGLE_RATE_FILES}}
\paragraph*{INSTANCE_TOGGLE_RATE_FILES}

GSR 关键字 INSTANCE_TOGGLE_RATE_FILES：用于设置与 实例、翻转、率、文件 相关的 RedHawk 运行参数。（可选；默认：None；兼容 DMP）

\textbf{语法：}
\begin{lstlisting}
INSTANCE_TOGGLE_RATE_FILES <filename>
\end{lstlisting}

\subsubsection{\gsr{INTERCONNECT_GATE_CAP_RATIO}}
\paragraph*{INTERCONNECT_GATE_CAP_RATIO}

GSR 关键字 INTERCONNECT_GATE_CAP_RATIO：用于设置与 INTERCONNECT、GATE、CAP、RATIO 相关的 RedHawk 运行参数。（可选；默认：value of INTERCONNECT_GATE_CAP_RATIO；兼容 DMP）

\textbf{语法：}
\begin{lstlisting}
INTERCONNECT_GATE_CAP_RATIO <value>
\end{lstlisting}

\textbf{示例：}
\begin{lstlisting}
INTERCONNECT_GATE_CAP_RATIO 1.5
\end{lstlisting}

\subsubsection{\gsr{IPF_ERROR_THRESHOLD}}
\paragraph*{IPF_ERROR_THRESHOLD}

GSR 关键字 IPF_ERROR_THRESHOLD：用于设置与 IPF、错误、THRESHOLD 相关的 RedHawk 运行参数。（可选；默认：value, so that if at least the specified；兼容 DMP）

\textbf{语法：}
\begin{lstlisting}
IPF_ERROR_THRESHOLD <percent_instances>
\end{lstlisting}

\subsubsection{\gsr{ITR_OVERRIDE_BPFS}}
\paragraph*{ITR_OVERRIDE_BPFS}

GSR 关键字 ITR_OVERRIDE_BPFS：用于设置与 ITR、OVERRIDE、BPFS 相关的 RedHawk 运行参数。（可选；默认：0；兼容 DMP）

\textbf{语法：}
\begin{lstlisting}
Syntax
ITR_OVERRIDE_BPFS [0|1]
\end{lstlisting}

\subsubsection{\gsr{LIB2AVM_DELAY_METHOD}}
\paragraph*{LIB2AVM_DELAY_METHOD}

GSR 关键字 LIB2AVM_DELAY_METHOD：用于设置与 LIB2AVM、DELAY、METHOD 相关的 RedHawk 运行参数。（可选；默认：0）

\textbf{语法：}
\begin{lstlisting}
Syntax
LIB2AVM_DELAY_METHOD [0| 1]
\end{lstlisting}

\subsubsection{\gsr{LIB2AVM_MSTATE}}
\paragraph*{LIB2AVM_MSTATE}

GSR 关键字 LIB2AVM_MSTATE：用于设置与 LIB2AVM、MSTATE 相关的 RedHawk 运行参数。（可选；默认：0；兼容 DMP）

\textbf{语法：}
\begin{lstlisting}
Syntax
LIB2AVM_MSTATE [0| 1 | 2]
\end{lstlisting}

\subsubsection{\gsr{NET_LOAD_FILE}}
\paragraph*{NET_LOAD_FILE}

GSR 关键字 NET_LOAD_FILE：用于设置与 网络、LOAD、文件 相关的 RedHawk 运行参数。

\textbf{语法：}
\begin{lstlisting}
Syntax
NET_LOAD_FILE <filename>
\end{lstlisting}

\textbf{示例：}
\begin{lstlisting}
$OUTPUT_PIN_CAP_SCALE 1
ucpu0/clk_drvr_13/Y 1.2
ucpu1/clk_drvr_13/Y 1.2
DMP compatible. Optional. Default: None.
\end{lstlisting}

\subsubsection{\gsr{NET_TOGGLE_RATE}}
\paragraph*{NET_TOGGLE_RATE}

GSR 关键字 NET_TOGGLE_RATE：用于设置与 网络、翻转、率 相关的 RedHawk 运行参数。（可选；默认：s: TR: none；兼容 DMP）

\textbf{语法：}
\begin{lstlisting}
NET_TOGGLE_RATE {
<net_name1> <toggle_rate1> ?<duty_cycle1>?
<net_name2> <toggle_rate2> ?<duty_cycle2>?
...
}
\end{lstlisting}

\textbf{示例：}
\begin{lstlisting}
NET_TOGGLE_RATE {
net1 0.5 0.1
net2 0.3
}
\end{lstlisting}

\subsubsection{\gsr{NET_TOGGLE_RATE_FILE}}
\paragraph*{NET_TOGGLE_RATE_FILE}

GSR 关键字 NET_TOGGLE_RATE_FILE：用于设置与 网络、翻转、率、文件 相关的 RedHawk 运行参数。（可选；默认：None；兼容 DMP）

\textbf{语法：}
\begin{lstlisting}
Syntax
NET_TOGGLE_RATE_FILE <file_path>
\end{lstlisting}

\subsubsection{\gsr{POWER_ALLOW_MULTIPLE_STATE}}
\paragraph*{POWER_ALLOW_MULTIPLE_STATE}

GSR 关键字 POWER_ALLOW_MULTIPLE_STATE：用于设置与 电源、ALLOW、MULTIPLE、STATE 相关的 RedHawk 运行参数。（可选；默认：1；兼容 DMP）

\textbf{语法：}
\begin{lstlisting}
POWER_ALLOW_MULTIPLE_STATE [0 | 1 | 2 ]
\end{lstlisting}

\subsubsection{\gsr{POWER_ALLOW_PER_PIN_IPF}}
\paragraph*{POWER_ALLOW_PER_PIN_IPF}

GSR 关键字 POWER_ALLOW_PER_PIN_IPF：用于设置与 电源、ALLOW、PER、引脚、IPF 相关的 RedHawk 运行参数。（可选；默认：when the IPF (Instance_Power_File) is us；兼容 DMP）

\textbf{语法：}
\begin{lstlisting}
POWER_ALLOW_PER_PIN_IPF [ 1 | 0 ]
\end{lstlisting}

\subsubsection{\gsr{POWER_ANALYSIS_MODE}}
\paragraph*{POWER_ANALYSIS_MODE}

GSR 关键字 POWER_ANALYSIS_MODE：用于设置与 电源、ANALYSIS、模式 相关的 RedHawk 运行参数。（默认：) averaged : This method accounts for st）

\textbf{语法：}
\begin{lstlisting}
POWER_ANALYSIS_MODE [ averaged | averaged_fast |
averaged_sdp]
DMP compatible. Optional. Default: averaged_fast.
\end{lstlisting}

\subsubsection{\gsr{POWER_CYCLE_SELECT_BLACK_BOX}}
\paragraph*{POWER_CYCLE_SELECT_BLACK_BOX}

GSR 关键字 POWER_CYCLE_SELECT_BLACK_BOX：用于设置与 电源、CYCLE、SELECT、BLACK、BOX 相关的 RedHawk 运行参数。（可选；默认：none；兼容 DMP）

\textbf{语法：}
\begin{lstlisting}
POWER_CYCLE_SELECT_BLACK_BOX {
<x1> <y1> <x2> <y2>
<x3> <y3> <x4> <y4>
...
}
\end{lstlisting}

\subsubsection{\gsr{POWER_CYCLE_SELECT_MODE}}
\paragraph*{POWER_CYCLE_SELECT_MODE}

GSR 关键字 POWER_CYCLE_SELECT_MODE：用于设置与 电源、CYCLE、SELECT、模式 相关的 RedHawk 运行参数。（可选；默认：0；兼容 DMP）

\textbf{语法：}
\begin{lstlisting}
POWER_CYCLE_SELECT_MODE [0 | 1 ]
\end{lstlisting}

\textbf{参数说明：}
\begin{itemize}
\item \texttt{0} — GSR 关键字 POWER\_CYCLE\_SELECT\_MODE：用于设置与 电源、CYCLE、SELECT、模式 相关的 RedHawk 运行参数。（默认：) 1: runtime : high, memory use : high,）
\end{itemize}

\subsubsection{\gsr{POWER_CYCLE_SELECT_POWER_NETS}}
\paragraph*{POWER_CYCLE_SELECT_POWER_NETS}

GSR 关键字 POWER_CYCLE_SELECT_POWER_NETS：用于设置与 电源、CYCLE、SELECT、电源、网络 相关的 RedHawk 运行参数。（可选；默认：none；兼容 DMP）

\textbf{语法：}
\begin{lstlisting}
POWER_CYCLE_SELECT_POWER_NETS {
<net1>
<net2>
...
}
\end{lstlisting}

\subsubsection{\gsr{POWER_CYCLE_SELECT_REPORT_VDD}}
\paragraph*{POWER_CYCLE_SELECT_REPORT_VDD}

GSR 关键字 POWER_CYCLE_SELECT_REPORT_VDD：用于设置与 电源、CYCLE、SELECT、REPORT、VDD 相关的 RedHawk 运行参数。（可选；默认：none；兼容 DMP）

\textbf{语法：}
\begin{lstlisting}
POWER_CYCLE_SELECT_REPORT_VDD {
<Vdd_a>
<Vdd_s>
...
}
\end{lstlisting}

\subsubsection{\gsr{POWER_CYCLE_SELECT_WHITE_BOX}}
\paragraph*{POWER_CYCLE_SELECT_WHITE_BOX}

GSR 关键字 POWER_CYCLE_SELECT_WHITE_BOX：用于设置与 电源、CYCLE、SELECT、WHITE、BOX 相关的 RedHawk 运行参数。（默认：none）

\textbf{语法：}
\begin{lstlisting}
POWER_CYCLE_SELECT_WHITE_BOX {
<x1> <y1> <x2> <y2>
<x3> <y3> <x4> <y4>
...
}
\end{lstlisting}

\subsubsection{\gsr{POWER_DISABLE_SWITCH}}
\paragraph*{POWER_DISABLE_SWITCH}

GSR 关键字 POWER_DISABLE_SWITCH：用于设置与 电源、禁用、SWITCH 相关的 RedHawk 运行参数。（可选；默认：0；兼容 DMP）

\textbf{语法：}
\begin{lstlisting}
POWER_DISABLE_SWITCH [ 0 | 1 ]
\end{lstlisting}

\subsubsection{\gsr{POWER_DOMAIN_TOGGLE_RATE}}
\paragraph*{POWER_DOMAIN_TOGGLE_RATE}

GSR 关键字 POWER_DOMAIN_TOGGLE_RATE：用于设置与 电源、DOMAIN、翻转、率 相关的 RedHawk 运行参数。（可选；默认：toggle rate for instances in certain pow）

\textbf{语法：}
\begin{lstlisting}
POWER_DOMAIN_TOGGLE_RATE <power_domain_name>
<toggle_rate> ? <clock_toggle_rate> ?
\end{lstlisting}

\subsubsection{\gsr{POWER_DOMAIN_TOGGLE_RATE_FILE}}
\paragraph*{POWER_DOMAIN_TOGGLE_RATE_FILE}

GSR 关键字 POWER_DOMAIN_TOGGLE_RATE_FILE：用于设置与 电源、DOMAIN、翻转、率、文件 相关的 RedHawk 运行参数。（默认：toggle rates for instances in certain po）

\textbf{语法：}
\begin{lstlisting}
POWER_DOMAIN_TOGGLE_RATE              <filename>
The format of the file is shown below.
\end{lstlisting}

\subsubsection{\gsr{POWER_DRIVER_TOGGLE_RATE}}
\paragraph*{POWER_DRIVER_TOGGLE_RATE}

GSR 关键字 POWER_DRIVER_TOGGLE_RATE：用于设置与 电源、DRIVER、翻转、率 相关的 RedHawk 运行参数。（可选；默认：1 (On)）

\textbf{语法：}
\begin{lstlisting}
POWER_DRIVER_TOGGLE_RATE [ 0 | 1 ]
\end{lstlisting}

\subsubsection{\gsr{POWER_HIER_REPORT_LEVEL}}
\paragraph*{POWER_HIER_REPORT_LEVEL}

GSR 关键字 POWER_HIER_REPORT_LEVEL：用于设置与 电源、HIER、REPORT、LEVEL 相关的 RedHawk 运行参数。（可选；默认：)；兼容 DMP；必需）

\textbf{语法：}
\begin{lstlisting}
POWER_HIER_REPORT_LEVEL <level>
\end{lstlisting}

\textbf{示例：}
\begin{lstlisting}
POWER_HIER_REPORT_LEVEL 2
In this case, power for all the blocks to the 2nd hierarchy level are reported,
including the top level.
\end{lstlisting}

\subsubsection{\gsr{POWER_HONOR_LIB_K_FACTOR}}
\paragraph*{POWER_HONOR_LIB_K_FACTOR}

GSR 关键字 POWER_HONOR_LIB_K_FACTOR：用于设置与 电源、HONOR、LIB、K、因子 相关的 RedHawk 运行参数。

\textbf{语法：}
\begin{lstlisting}
POWER_HONOR_LIB_K_FACTOR [ 0 | 1 ]
\end{lstlisting}

\textbf{参数说明：}
GSR 关键字 POWER\_HONOR\_LIB\_K\_FACTOR：用于设置与 电源、HONOR、LIB、K、因子 相关的 RedHawk 运行参数。（可选；默认：0 (off, no scaling)）

\subsubsection{\gsr{POWER_IGNORE_ASYNCH_PIN}}
\paragraph*{POWER_IGNORE_ASYNCH_PIN}

GSR 关键字 POWER_IGNORE_ASYNCH_PIN：用于设置与 电源、忽略、ASYNCH、引脚 相关的 RedHawk 运行参数。（可选；默认：0 (off)）

\textbf{语法：}
\begin{lstlisting}
POWER_IGNORE_ASYNC_PIN [ 0 | 1 ]
\end{lstlisting}

\subsubsection{\gsr{POWER_INTERNAL_SCALING_FILE}}
\paragraph*{POWER_INTERNAL_SCALING_FILE}

GSR 关键字 POWER_INTERNAL_SCALING_FILE：用于设置与 电源、INTERNAL、SCALING、文件 相关的 RedHawk 运行参数。（默认：off）

\textbf{语法：}
\begin{lstlisting}
POWER_INTERNAL_SCALING_FILE <scaling_factor_filename>
Format of scaling factor file:
<cell_name1> <scaling_factor1> <power pin name>
<cell_name2> <scaling_factor2> <power pin name>
\end{lstlisting}

\subsubsection{\gsr{POWER_LEAKAGE_SCALING_FACTOR}}
\paragraph*{POWER_LEAKAGE_SCALING_FACTOR}

GSR 关键字 POWER_LEAKAGE_SCALING_FACTOR：用于设置与 电源、LEAKAGE、SCALING、因子 相关的 RedHawk 运行参数。（可选；默认：1；兼容 DMP）

\textbf{语法：}
\begin{lstlisting}
POWER_LEAKAGE_SCALING_FACTOR <factor>
\end{lstlisting}

\subsubsection{\gsr{POWER_LEAKAGE_SCALING_FILE}}
\paragraph*{POWER_LEAKAGE_SCALING_FILE}

GSR 关键字 POWER_LEAKAGE_SCALING_FILE：用于设置与 电源、LEAKAGE、SCALING、文件 相关的 RedHawk 运行参数。（可选；默认：none；兼容 DMP）

\textbf{语法：}
\begin{lstlisting}
POWER_LEAKAGE_SCALING_FILE <scaling_filename>
\end{lstlisting}

\subsubsection{\gsr{POWER_MCF_MULTI_CLOCK}}
\paragraph*{POWER_MCF_MULTI_CLOCK}

GSR 关键字 POWER_MCF_MULTI_CLOCK：用于设置与 电源、MCF、MULTI、CLOCK 相关的 RedHawk 运行参数。（可选；默认：0）

\textbf{语法：}
\begin{lstlisting}
POWER_MCF_MULTI_CLOCK [ 0 | 1 ]
\end{lstlisting}

\subsubsection{\gsr{POWER_MISSING_IPF_LEAK}}
\paragraph*{POWER_MISSING_IPF_LEAK}

GSR 关键字 POWER_MISSING_IPF_LEAK：用于设置与 电源、MISSING、IPF、LEAK 相关的 RedHawk 运行参数。（可选；默认：preserves leakage power values；兼容 DMP）

\textbf{语法：}
\begin{lstlisting}
POWER_MISSING_IPF_LEAK [ 0 | 1 ]
\end{lstlisting}

\subsubsection{\gsr{POWER_MISSING_IPF_POWER}}
\paragraph*{POWER_MISSING_IPF_POWER}

GSR 关键字 POWER_MISSING_IPF_POWER：用于设置与 电源、MISSING、IPF、电源 相关的 RedHawk 运行参数。（可选；默认：(0ff) sets missing power values to 0；兼容 DMP）

\textbf{语法：}
\begin{lstlisting}
POWER_MISSING_IPF_POWER [ 0 | 1 ]
\end{lstlisting}

\subsubsection{\gsr{POWER_MODE}}
\paragraph*{POWER_MODE}

GSR 关键字 POWER_MODE：用于设置与 电源、模式 相关的 RedHawk 运行参数。（可选；默认：Mixed；兼容 DMP）

\textbf{语法：}
\begin{lstlisting}
POWER_MODE [ APL | LIB | MIXED | APL_PEAK | APL_PEAK1 ]
\end{lstlisting}

\textbf{参数说明：}
\begin{itemize}
\item \texttt{APL} — GSR 关键字 POWER\_MODE：用于设置与 电源、模式 相关的 RedHawk 运行参数。
\end{itemize}

\textbf{示例：}
\begin{lstlisting}
POWER_MODE APL_PEAK
\end{lstlisting}

\subsubsection{\gsr{POWER_REPORT_BIAS_PIN}}
\paragraph*{POWER_REPORT_BIAS_PIN}

GSR 关键字 POWER_REPORT_BIAS_PIN：用于设置与 电源、REPORT、BIAS、引脚 相关的 RedHawk 运行参数。（可选；默认：), instance bias pin power is not includ）

\textbf{语法：}
\begin{lstlisting}
POWER_REPORT_BIAS_PIN            [ 0 | 1 ]
\end{lstlisting}

\subsubsection{\gsr{POWER_STATE_DEPENDENT_LEAKAGE}}
\paragraph*{POWER_STATE_DEPENDENT_LEAKAGE}

GSR 关键字 POWER_STATE_DEPENDENT_LEAKAGE：用于设置与 电源、STATE、DEPENDENT、LEAKAGE 相关的 RedHawk 运行参数。（可选；默认：0；兼容 DMP）

\textbf{语法：}
\begin{lstlisting}
POWER_STATE_DEPENDENT_LEAKAGE [0|1]
\end{lstlisting}

\subsubsection{\gsr{POWER_TRANSIENT_ANALYSIS}}
\paragraph*{POWER_TRANSIENT_ANALYSIS}

GSR 关键字 POWER_TRANSIENT_ANALYSIS：用于设置与 电源、TRANSIENT、ANALYSIS 相关的 RedHawk 运行参数。

\subsubsection{\gsr{GSC_FILES}}
\paragraph*{GSC_FILES}

GSR 关键字 GSC_FILES：用于设置与 GSC、文件 相关的 RedHawk 运行参数。

\subsubsection{\gsr{INSTANCE_POWER_FILES}}
\paragraph*{INSTANCE_POWER_FILES}

GSR 关键字 INSTANCE_POWER_FILES：用于设置与 实例、电源、文件 相关的 RedHawk 运行参数。

\subsubsection{\gsr{STATE_PROPAGATION}}
\paragraph*{STATE_PROPAGATION}

GSR 关键字 STATE_PROPAGATION：用于设置与 STATE、PROPAGATION 相关的 RedHawk 运行参数。

\subsubsection{\gsr{GSC_OVERRIDE_IPF}}
\paragraph*{GSC_OVERRIDE_IPF}

GSR 关键字 GSC_OVERRIDE_IPF：用于设置与 GSC、OVERRIDE、IPF 相关的 RedHawk 运行参数。

\subsubsection{\gsr{BLOCK_POWER_FOR_SCALING}}
\paragraph*{BLOCK_POWER_FOR_SCALING}

GSR 关键字 BLOCK_POWER_FOR_SCALING：用于设置与 模块、电源、FOR、SCALING 相关的 RedHawk 运行参数。

\subsubsection{\gsr{TOGGLE_RATE_RATIO_COMB_FF}}
\paragraph*{TOGGLE_RATE_RATIO_COMB_FF}

GSR 关键字 TOGGLE_RATE_RATIO_COMB_FF：用于设置与 翻转、率、RATIO、COMB、FF 相关的 RedHawk 运行参数。

\subsubsection{\gsr{INSTANCE_TOGGLE_RATE}}
\paragraph*{INSTANCE_TOGGLE_RATE}

GSR 关键字 INSTANCE_TOGGLE_RATE：用于设置与 实例、翻转、率 相关的 RedHawk 运行参数。（可选；默认：None；兼容 DMP）

\textbf{语法：}
\begin{lstlisting}
POWER_TRANSIENT_ANALYSIS {
<duration_frame1_sec> <config_file1>
<duration_frame2_sec> <config_file2>
...
}
\end{lstlisting}

\subsubsection{\gsr{POWER_USE_CCS}}
\paragraph*{POWER_USE_CCS}

GSR 关键字 POWER_USE_CCS：用于设置与 电源、USE、CCS 相关的 RedHawk 运行参数。（可选；默认：0）

\textbf{语法：}
\begin{lstlisting}
POWER_USE_CCS [ 0 | 1 ]
\end{lstlisting}

\subsubsection{\gsr{POWER_VCD_COVERED_THRESHOLD}}
\paragraph*{POWER_VCD_COVERED_THRESHOLD}

GSR 关键字 POWER_VCD_COVERED_THRESHOLD：用于设置与 电源、VCD、COVERED、THRESHOLD 相关的 RedHawk 运行参数。（默认：0）

\textbf{语法：}
\begin{lstlisting}
POWER_VCD_COVERED_THRESHOLD <value>
\end{lstlisting}

\subsubsection{\gsr{POWER_VCD_EVENT_SEQUENCE}}
\paragraph*{POWER_VCD_EVENT_SEQUENCE}

GSR 关键字 POWER_VCD_EVENT_SEQUENCE：用于设置与 电源、VCD、EVENT、SEQUENCE 相关的 RedHawk 运行参数。

\textbf{语法：}
\begin{lstlisting}
POWER_VCD_EVENT_SEQUENCE[ 0 | 1 | 2 ]
\end{lstlisting}

\textbf{参数说明：}
\begin{itemize}
\item \texttt{0} — GSR 关键字 POWER\_VCD\_EVENT\_SEQUENCE：用于设置与 电源、VCD、EVENT、SEQUENCE 相关的 RedHawk 运行参数。（默认：value, don't consider any sequence numbe）
\end{itemize}

\subsubsection{\gsr{POWER_VCD_HONOR_GLITCH_EVENT}}
\paragraph*{POWER_VCD_HONOR_GLITCH_EVENT}

GSR 关键字 POWER_VCD_HONOR_GLITCH_EVENT：用于设置与 电源、VCD、HONOR、GLITCH、EVENT 相关的 RedHawk 运行参数。

\textbf{语法：}
\begin{lstlisting}
POWER_VCD_HONOR_GLITCH_EVENT [ 0 | 1 ]
\end{lstlisting}

\subsubsection{\gsr{POWER_VCD_IGNORE_ERROR}}
\paragraph*{POWER_VCD_IGNORE_ERROR}

GSR 关键字 POWER_VCD_IGNORE_ERROR：用于设置与 电源、VCD、忽略、错误 相关的 RedHawk 运行参数。（可选；默认：RedHawk errors out when a block name is）

\textbf{语法：}
\begin{lstlisting}
POWER_VCD_IGNORE_ERROR [ 0 | 1 ]
\end{lstlisting}

\subsubsection{\gsr{POWER_VCD_LIMIT_TR}}
\paragraph*{POWER_VCD_LIMIT_TR}

GSR 关键字 POWER_VCD_LIMIT_TR：用于设置与 电源、VCD、LIMIT、TR 相关的 RedHawk 运行参数。（可选；默认：0）

\textbf{语法：}
\begin{lstlisting}
POWER_VCD_LIMIT_TR [0|1]
\end{lstlisting}

\subsubsection{\gsr{POWER_VCD_NUM_PROCESS}}
\paragraph*{POWER_VCD_NUM_PROCESS}

GSR 关键字 POWER_VCD_NUM_PROCESS：用于设置与 电源、VCD、NUM、PROCESS 相关的 RedHawk 运行参数。（可选；默认：, so RedHawk automatically decides how m；兼容 DMP）

\textbf{语法：}
\begin{lstlisting}
POWER_VCD_NUM_PROCESS <num_proc>
\end{lstlisting}

\subsubsection{\gsr{POWER_VCD_OVERRIDE_IPF}}
\paragraph*{POWER_VCD_OVERRIDE_IPF}

GSR 关键字 POWER_VCD_OVERRIDE_IPF：用于设置与 电源、VCD、OVERRIDE、IPF 相关的 RedHawk 运行参数。

\textbf{语法：}
\begin{lstlisting}
When set to 1, IPF data is ignored if the instance is covered by VCD
information. The default value of 0 maintains behavior prior to version 12.1, in
which, for Mixed-Mode analysis, when the design had VCD files as well as IPF for
both the block and the TOP-level instances, IPF power was considered where there
were differences in the switching scenario and power values. DMP compatible.
Optional. Default: 1.
POWER_VCD_OVERRIDE_IPF [ 0 | 1 ]
\end{lstlisting}

\subsubsection{\gsr{POWER_VCD_REUSE_EVENT}}
\paragraph*{POWER_VCD_REUSE_EVENT}

GSR 关键字 POWER_VCD_REUSE_EVENT：用于设置与 电源、VCD、REUSE、EVENT 相关的 RedHawk 运行参数。（可选；默认：1 (On)；兼容 DMP）

\textbf{语法：}
\begin{lstlisting}
POWER_VCD_REUSE_EVENT [ 0 | 1 | 2]
\end{lstlisting}

\textbf{参数说明：}
\begin{itemize}
\item \texttt{0} — GSR 关键字 POWER\_VCD\_REUSE\_EVENT：用于设置与 电源、VCD、REUSE、EVENT 相关的 RedHawk 运行参数。（默认：value, will check if global nets are reu）
\end{itemize}

\subsubsection{\gsr{POWER_VCD_TO_FSDB}}
\paragraph*{POWER_VCD_TO_FSDB}

GSR 关键字 POWER_VCD_TO_FSDB：用于设置与 电源、VCD、TO、FSDB 相关的 RedHawk 运行参数。（默认：1 (on)）

\textbf{语法：}
\begin{lstlisting}
POWER_VCD_TO_FSDB             [ 0 | 1 | 2 ]
\end{lstlisting}

\textbf{参数说明：}
\begin{itemize}
\item \texttt{0} — GSR 关键字 POWER\_VCD\_TO\_FSDB：用于设置与 电源、VCD、TO、FSDB 相关的 RedHawk 运行参数。
\end{itemize}

\subsubsection{\gsr{POWER_WORST_IO_PAD}}
\paragraph*{POWER_WORST_IO_PAD}

GSR 关键字 POWER_WORST_IO_PAD：用于设置与 电源、WORST、IO、焊盘 相关的 RedHawk 运行参数。（可选；默认：0）

\textbf{语法：}
\begin{lstlisting}
POWER_WORST_IO_PAD [ 0 | 1 ]
\end{lstlisting}

\subsubsection{\gsr{POWER_WORST_LEAKAGE}}
\paragraph*{POWER_WORST_LEAKAGE}

GSR 关键字 POWER_WORST_LEAKAGE：用于设置与 电源、WORST、LEAKAGE 相关的 RedHawk 运行参数。（可选；默认：0）

\textbf{语法：}
\begin{lstlisting}
POWER_WORST_LEAKAGE [ 0 | 1 ]
\end{lstlisting}

\subsubsection{\gsr{POWER_WORST_MBFF}}
\paragraph*{POWER_WORST_MBFF}

GSR 关键字 POWER_WORST_MBFF：用于设置与 电源、WORST、MBFF 相关的 RedHawk 运行参数。（可选；默认：the Power transient flow handles APL cur）

\textbf{语法：}
\begin{lstlisting}
POWER_WORST_MBFF [0| 1]
\end{lstlisting}

\subsubsection{\gsr{PRIMARY_OUTPUT_LOAD_CAPS}}
\paragraph*{PRIMARY_OUTPUT_LOAD_CAPS}

GSR 关键字 PRIMARY_OUTPUT_LOAD_CAPS：用于设置与 PRIMARY、OUTPUT、LOAD、CAPS 相关的 RedHawk 运行参数。（可选；默认：_LOAD is specified, the default Cload sp；兼容 DMP）

\subsubsection{\gsr{PRIMARY_OUTPUT_LOAD_CAPS}}
\paragraph*{PRIMARY_OUTPUT_LOAD_CAPS}

GSR 关键字 PRIMARY_OUTPUT_LOAD_CAPS：用于设置与 PRIMARY、OUTPUT、LOAD、CAPS 相关的 RedHawk 运行参数。（默认：_LOAD <output_cap_Farads> <Vdd_name> <ou）

\textbf{参数说明：}
\begin{itemize}
\item \texttt{IGNORE\_FANOUT\_CAP [ 0 | 1] <powerNet\_name>} — GSR 关键字 PRIMARY\_OUTPUT\_LOAD\_CAPS：用于设置与 PRIMARY、OUTPUT、LOAD、CAPS 相关的 RedHawk 运行参数。（默认：, this option is off and RedHawk only）
\end{itemize}

\subsubsection{\gsr{PRIMARY_OUTPUT_LOAD_CAPS}}
\paragraph*{PRIMARY_OUTPUT_LOAD_CAPS}

为 primary output 指定负载电容。可在 GSR 中按网名模式分配负载（例如 I/O 相关网与其他网使用不同电容值）；亦可用 \gsr{PRIMARY_OUTPUT_LOAD_CAPS_FILE} 从外部文件导入。

\begin{noteBox}
该关键字主要用于为 I/O 相关网分配负载电容；具体网名通配与电容数值见 GSR 语法示例或原书附录 C。
\end{noteBox}

\subsubsection{\gsr{PRIMARY_OUTPUT_LOAD_CAPS_FILE}}
\paragraph*{PRIMARY_OUTPUT_LOAD_CAPS_FILE}

GSR 关键字 PRIMARY_OUTPUT_LOAD_CAPS_FILE：用于设置与 PRIMARY、OUTPUT、LOAD、CAPS、文件 相关的 RedHawk 运行参数。

\textbf{语法：}
\begin{lstlisting}
PRIMARY_OUTPUT_LOAD_CAPS_FILE ./<file_pointer>
File contents:
<net_name> <load_cap_value_in_F>
\end{lstlisting}

\textbf{示例：}
\begin{lstlisting}
PRIMARY_OUTPUT_LOAD_CAPS_FILE polcf
PRIMARY_OUTPUT_LOAD_CAPS_FILE plocf.gz
\end{lstlisting}

\subsubsection{\gsr{PS_GENERATE_MCYC_SCENARIO}}
\paragraph*{PS_GENERATE_MCYC_SCENARIO}

GSR 关键字 PS_GENERATE_MCYC_SCENARIO：用于设置与 PS、GENERATE、MCYC、SCENARIO 相关的 RedHawk 运行参数。（默认：-1 (off)）

\textbf{语法：}
\begin{lstlisting}
PS_GENERATE_MCYC_SCENARIO [-1 | 1 | 2]
\end{lstlisting}

\subsubsection{\gsr{PS_GENERATE_VLESS_SCENARIO}}
\paragraph*{PS_GENERATE_VLESS_SCENARIO}

GSR 关键字 PS_GENERATE_VLESS_SCENARIO：用于设置与 PS、GENERATE、VLESS、SCENARIO 相关的 RedHawk 运行参数。（可选；默认：- 1；兼容 DMP）

\textbf{语法：}
\begin{lstlisting}
PS_GENERATE_VLESS_SCENARIO <integer>
\end{lstlisting}

\subsubsection{\gsr{PS_GLITCH_POWER_MODELING}}
\paragraph*{PS_GLITCH_POWER_MODELING}

GSR 关键字 PS_GLITCH_POWER_MODELING：用于设置与 PS、GLITCH、电源、MODELING 相关的 RedHawk 运行参数。（可选；默认：- 1）

\textbf{语法：}
\begin{lstlisting}
PS_GLITCH_POWER_MODELING <integer>
\end{lstlisting}

\subsubsection{\gsr{PS_LIB_EXTRACT_CCS}}
\paragraph*{PS_LIB_EXTRACT_CCS}

GSR 关键字 PS_LIB_EXTRACT_CCS：用于设置与 PS、LIB、提取、CCS 相关的 RedHawk 运行参数。（可选；默认：0）

\textbf{语法：}
\begin{lstlisting}
PS_LIB_EXTRACT_CCS [ 0 | 1 ]
\end{lstlisting}

\subsubsection{\gsr{PS_RTL_EP_MODE}}
\paragraph*{PS_RTL_EP_MODE}

GSR 关键字 PS_RTL_EP_MODE：用于设置与 PS、RTL、EP、模式 相关的 RedHawk 运行参数。（可选；默认：0）

\textbf{语法：}
\begin{lstlisting}
Syntax
PS_RTL_EP_MODE [ 0 | 1 ]
\end{lstlisting}

\subsubsection{\gsr{PS_RTL_EP_REPORT}}
\paragraph*{PS_RTL_EP_REPORT}

GSR 关键字 PS_RTL_EP_REPORT：用于设置与 PS、RTL、EP、REPORT 相关的 RedHawk 运行参数。（可选；默认：0, off；兼容 DMP）

\subsubsection{\gsr{PS_SP_GATED_CLOCK_LOGIC}}
\paragraph*{PS_SP_GATED_CLOCK_LOGIC}

GSR 关键字 PS_SP_GATED_CLOCK_LOGIC：用于设置与 PS、SP、GATED、CLOCK、LOGIC 相关的 RedHawk 运行参数。（可选；默认：1 (On)）

\textbf{语法：}
\begin{lstlisting}
Syntax
PS_SP_GATED_CLOCK_LOGIC [ 0 | 1 ]
\end{lstlisting}

\subsubsection{\gsr{PS_VLSG_POWER_DOMAIN_HANDLING}}
\paragraph*{PS_VLSG_POWER_DOMAIN_HANDLING}

GSR 关键字 PS_VLSG_POWER_DOMAIN_HANDLING：用于设置与 PS、VLSG、电源、DOMAIN、HANDLING 相关的 RedHawk 运行参数。

\textbf{语法：}
\begin{lstlisting}
Syntax
PS_VLSG_POWER_DOMAIN_HANDLING [ 0| 1 ]
\end{lstlisting}

\subsubsection{\gsr{PARA_CALC_POWER}}
\paragraph*{PARA_CALC_POWER}

GSR 关键字 PARA_CALC_POWER：用于设置与 PARA、CALC、电源 相关的 RedHawk 运行参数。（可选；默认：), performs power calculation and extrac；兼容 DMP）

\textbf{语法：}
\begin{lstlisting}
Syntax
PARA_CALC_POWER [ 0 | 1 ]
\end{lstlisting}

\subsubsection{\gsr{REPORT_MAXCAP_VIOLATION}}
\paragraph*{REPORT_MAXCAP_VIOLATION}

GSR 关键字 REPORT_MAXCAP_VIOLATION：用于设置与 REPORT、MAXCAP、VIOLATION 相关的 RedHawk 运行参数。（可选；默认：0 (off)）

\textbf{语法：}
\begin{lstlisting}
Syntax
REPORT_MAXCAP_VIOLATION [ 0 | n | -1 ]
\end{lstlisting}

\textbf{示例：}
\begin{lstlisting}
Example
REPORT_MAXCAP_VIOLATION 150
\end{lstlisting}

\subsubsection{\gsr{RTL_NAME_MAPPING}}
\paragraph*{RTL_NAME_MAPPING}

GSR 关键字 RTL_NAME_MAPPING：用于设置与 RTL、NAME、MAPPING 相关的 RedHawk 运行参数。（可选；默认：)）

\textbf{语法：}
\begin{lstlisting}
Syntax
RTL_NAME_MAPPING [ 0 | 1 | 2 | 3 ]
\end{lstlisting}

\subsubsection{\gsr{SAIF_FILE}}
\paragraph*{SAIF_FILE}

GSR 关键字 SAIF_FILE：用于设置与 SAIF、文件 相关的 RedHawk 运行参数。（可选；默认：none）

\textbf{语法：}
\begin{lstlisting}
SAIF_FILE {
<block_name> <SAIF_filename>
ROOT_INSTANCE <root_instance_name>
SUBSTITUTE_PATH <substitute_instance_path>
}
\end{lstlisting}

\textbf{参数说明：}
\begin{itemize}
\item \texttt{<block\_name>} — GSR 关键字 SAIF\_FILE：用于设置与 SAIF、文件 相关的 RedHawk 运行参数。
\end{itemize}

\subsubsection{\gsr{SCALE_CLOCK_POWER}}
\paragraph*{SCALE_CLOCK_POWER}

GSR 关键字 SCALE_CLOCK_POWER：用于设置与 缩放、CLOCK、电源 相关的 RedHawk 运行参数。（可选；默认：1；兼容 DMP）

\textbf{语法：}
\begin{lstlisting}
SCALE_CLOCK_POWER [0|1]
\end{lstlisting}

\subsubsection{\gsr{SCAN_CLK_DUTY_CYCLE}}
\paragraph*{SCAN_CLK_DUTY_CYCLE}

GSR 关键字 SCAN_CLK_DUTY_CYCLE：用于设置与 SCAN、CLK、DUTY、CYCLE 相关的 RedHawk 运行参数。（可选；默认：0；兼容 DMP）

\textbf{语法：}
\begin{lstlisting}
SCAN_CLK_DUTY_CYCLE <clock_duty_cycle>
\end{lstlisting}

\subsubsection{\gsr{SCAN_CONSTRAINT_FILES}}
\paragraph*{SCAN_CONSTRAINT_FILES}

GSR 关键字 SCAN_CONSTRAINT_FILES：用于设置与 SCAN、CONSTRAINT、文件 相关的 RedHawk 运行参数。（可选；默认：pattern of “01” is assumed；兼容 DMP）

\textbf{语法：}
\begin{lstlisting}
SCAN_CONSTRAINT_FILES {
<scan_chain_file_1> <pattern_1>
? <scan_chain_file_2> <pattern_2>?
...
? AUTO_TRACING <pattern>?
? AUTO_TRACING_BLOCK <hier_block_name>
[<pattern>|<patternfile>]?
? AUTO_TRACING_BLOCK_FILE <path_to_file> ?
? AUTO_TRACING_EXCLUDE_BLOCK <hier_block_name>
[<pattern>|<patternfile>]?
? AUTO_TRACING_EXCLUDE_BLOCK_FILE <path_to_file> ?
}
\end{lstlisting}

\textbf{示例：}
\begin{lstlisting}
Example
Multiple scan chains, each has its own pattern:
SCAN_CONSTRAINT_FILES {
scan_chain_1.spc
scan_chain_2.spc 101
\end{lstlisting}

\subsubsection{\gsr{SCANMODE}}
\paragraph*{SCANMODE}

GSR 关键字 SCANMODE：用于设置与 SCANMODE 相关的 RedHawk 运行参数。（可选；默认：0）

\textbf{语法：}
\begin{lstlisting}
SCANMODE[ 0 | 1 ]
\end{lstlisting}

\subsubsection{\gsr{SCAN_PI_CONSTRAINT_FILE}}
\paragraph*{SCAN_PI_CONSTRAINT_FILE}

GSR 关键字 SCAN_PI_CONSTRAINT_FILE：用于设置与 SCAN、PI、CONSTRAINT、文件 相关的 RedHawk 运行参数。（可选；默认：none）

\textbf{语法：}
\begin{lstlisting}
SCAN_PI_CONSTRAINT_FILE <pi_file>
\end{lstlisting}

\subsubsection{\gsr{SCAN_SHIFTIN_LSB}}
\paragraph*{SCAN_SHIFTIN_LSB}

GSR 关键字 SCAN_SHIFTIN_LSB：用于设置与 SCAN、SHIFTIN、LSB 相关的 RedHawk 运行参数。（默认：, it is shifted from LSB）

\textbf{语法：}
\begin{lstlisting}
Syntax
SCAN_SHIFTIN_LSB [ 1| 0 ]
\end{lstlisting}

\subsubsection{\gsr{SP_CG_CONSTRAINT_FILE}}
\paragraph*{SP_CG_CONSTRAINT_FILE}

GSR 关键字 SP_CG_CONSTRAINT_FILE：用于设置与 SP、CG、CONSTRAINT、文件 相关的 RedHawk 运行参数。（默认：none）

\textbf{语法：}
\begin{lstlisting}
Syntax
SP_CG_CONSTRAINT_FILE <CG_defin_file>
\end{lstlisting}

\subsubsection{\gsr{SP_CLOCK_GATING_ERROR_OUT_RATE}}
\paragraph*{SP_CLOCK_GATING_ERROR_OUT_RATE}

GSR 关键字 SP_CLOCK_GATING_ERROR_OUT_RATE：用于设置与 SP、CLOCK、GATING、错误、OUT、率 相关的 RedHawk 运行参数。（可选；默认：when more than 20% of the clock gates ha；兼容 DMP）

\textbf{语法：}
\begin{lstlisting}
Syntax
\end{lstlisting}

\subsubsection{\gsr{SP_CLOCK_GATING_RATIO}}
\paragraph*{SP_CLOCK_GATING_RATIO}

GSR 关键字 SP_CLOCK_GATING_RATIO：用于设置与 SP、CLOCK、GATING、RATIO 相关的 RedHawk 运行参数。（可选；默认：the clock gating ratio is automatically）

\textbf{语法：}
\begin{lstlisting}
Syntax
SP_CLOCK_GATING_RATIO <ratio>
\end{lstlisting}

\subsubsection{\gsr{SP_COMPLEX_CELL_HANDLING}}
\paragraph*{SP_COMPLEX_CELL_HANDLING}

GSR 关键字 SP_COMPLEX_CELL_HANDLING：用于设置与 SP、COMPLEX、单元、HANDLING 相关的 RedHawk 运行参数。（可选；默认：1；兼容 DMP）

\textbf{语法：}
\begin{lstlisting}
Syntax
SP_COMPLEX_CELL_HANDLING [1 | 0]
\end{lstlisting}

\subsubsection{\gsr{SP_LIMIT_TR}}
\paragraph*{SP_LIMIT_TR}

GSR 关键字 SP_LIMIT_TR：用于设置与 SP、LIMIT、TR 相关的 RedHawk 运行参数。（可选；默认：1）

\textbf{语法：}
\begin{lstlisting}
SP_LIMIT_TR [ 0 | 1 ]
\end{lstlisting}

\subsubsection{\gsr{STATE_PROPAGATION}}
\paragraph*{STATE_PROPAGATION}

GSR 关键字 STATE_PROPAGATION：用于设置与 STATE、PROPAGATION 相关的 RedHawk 运行参数。（默认：toggle rates in the design；兼容 DMP；必需）

\textbf{语法：}
\begin{lstlisting}
Syntax
STATE_PROPAGATION {
\end{lstlisting}

\textbf{参数说明：}
GSR 关键字 STATE\_PROPAGATION：用于设置与 STATE、PROPAGATION 相关的 RedHawk 运行参数。（默认：gate enable ratio is set at the output o）

\textbf{示例：}
\begin{lstlisting}
example, 0.6, is also on when FOP is set to 0.8.
If FOP is not set, any clock gate not otherwise specified has a default FOP
value of 1 (On).
GATED_CONTROL_FILE : specifies a file defining the On/Off status of individual
gated clock control cells. The format of the gated clock control file data is as
follows for both static and dynamic analyses:
<gated_clock_instance1> [ On | Off ]
...
GATED_CLOCK_COVERAGE : when set to 1, provides better design coverage
when using FLOP_ON_PERCENTAGE, and allows choosing different clock
gate turn-on scenarios for each frame so that overall analysis coverage is
increased. Note that POWER_TRANSIENT_ANALYSIS must also be turned
on.
\end{lstlisting}

\subsubsection{\gsr{STEINER_TREE_CAP}}
\paragraph*{STEINER_TREE_CAP}

GSR 关键字 STEINER_TREE_CAP：用于设置与 STEINER、TREE、CAP 相关的 RedHawk 运行参数。（可选；默认：none；兼容 DMP）

\textbf{语法：}
\begin{lstlisting}
STEINER_TREE_CAP <cap in pF per um>
\end{lstlisting}

\textbf{示例：}
\begin{lstlisting}
STEINER_TREE_CAP 2.0e-4
\end{lstlisting}

\subsubsection{\gsr{SUBSTRATE_PGARC}}
\paragraph*{SUBSTRATE_PGARC}

GSR 关键字 SUBSTRATE_PGARC：用于设置与 衬底、PGARC 相关的 RedHawk 运行参数。（默认：mode i）

\textbf{语法：}
\begin{lstlisting}
SUBSTRATE_PGARC [ 1 | 0 ]
\end{lstlisting}

\subsubsection{\gsr{TECHNOLOGY}}
\paragraph*{TECHNOLOGY}

GSR 关键字 TECHNOLOGY：用于设置与 TECHNOLOGY 相关的 RedHawk 运行参数。（可选）

\textbf{语法：}
\begin{lstlisting}
TECHNOLOGY <technology_node>
\end{lstlisting}

\subsubsection{\gsr{THERMAL_ANALYSIS}}
\paragraph*{THERMAL_ANALYSIS}

GSR 关键字 THERMAL_ANALYSIS：用于设置与 THERMAL、ANALYSIS 相关的 RedHawk 运行参数。（可选；默认：0）

\textbf{语法：}
\begin{lstlisting}
THERMAL_ANALYSIS [ 0 | 1 | 2 ]
\end{lstlisting}

\subsubsection{\gsr{THERMAL_MODEL}}
\paragraph*{THERMAL_MODEL}

GSR 关键字 THERMAL_MODEL：用于设置与 THERMAL、MODEL 相关的 RedHawk 运行参数。（可选；默认：)）

\textbf{语法：}
\begin{lstlisting}
THERMAL_MODEL [0 | 1 | 2 ]
\end{lstlisting}

\subsubsection{\gsr{TOGGLE_RATE}}
\paragraph*{TOGGLE_RATE}

GSR 关键字 TOGGLE_RATE：用于设置与 翻转、率 相关的 RedHawk 运行参数。（默认：signal toggle rate of the nets on the ch）

\textbf{语法：}
\begin{lstlisting}
TOGGLE_RATE <non_clock_TR> ?<clock_network_TR>?
\end{lstlisting}

\textbf{参数说明：}
\begin{itemize}
\item \texttt{<non\_clock\_TR>} — GSR 关键字 TOGGLE\_RATE：用于设置与 翻转、率 相关的 RedHawk 运行参数。（默认：2）
\end{itemize}

\textbf{示例：}
\begin{lstlisting}
example, a clock net generally has a toggle rate of 2.0 with respect to its clock
domain, since the net switches once from 0->1 and once 1->0 within a clock cycle.
Note that if there is no power consumption table in .lib the toggle rate is taken from
TOGGLE_RATE, and charge is scaled to meet power. DMP compatible. Optional.
Default: non-clock: 0.3; clock: 2.0.
TOGGLE_RATE 0.2 1.0
\end{lstlisting}

\subsubsection{\gsr{TOGGLE_RATE_RATIO_COMB_FF}}
\paragraph*{TOGGLE_RATE_RATIO_COMB_FF}

GSR 关键字 TOGGLE_RATE_RATIO_COMB_FF：用于设置与 翻转、率、RATIO、COMB、FF 相关的 RedHawk 运行参数。（可选；默认：1；兼容 DMP）

\textbf{语法：}
\begin{lstlisting}
TOGGLE_RATE_RATIO_COMB_FF <TR_ratio_combin_re_FF>
\end{lstlisting}

\textbf{示例：}
\begin{lstlisting}
TOGGLE_RATE_RATIO_COMB_FF 0.7
\end{lstlisting}

\subsubsection{\gsr{USE_FAST_DECAP_ALG}}
\paragraph*{USE_FAST_DECAP_ALG}

GSR 关键字 USE_FAST_DECAP_ALG：用于设置与 USE、FAST、去耦电容、ALG 相关的 RedHawk 运行参数。（可选；默认：1 On；兼容 DMP）

\textbf{语法：}
\begin{lstlisting}
USE_FAST_DECAP_ALG [ 0 | 1 ]
\end{lstlisting}

\subsubsection{\gsr{USE_LIB_MAX_CAP}}
\paragraph*{USE_LIB_MAX_CAP}

GSR 关键字 USE_LIB_MAX_CAP：用于设置与 USE、LIB、MAX、CAP 相关的 RedHawk 运行参数。（可选；默认：0 (Off)；兼容 DMP）

\textbf{语法：}
\begin{lstlisting}
USE_LIB_MAX_CAP [ 0 | 1 ]
\end{lstlisting}

\subsubsection{\gsr{USER_MCF_FILE}}
\paragraph*{USER_MCF_FILE}

GSR 关键字 USER_MCF_FILE：用于设置与 USER、MCF、文件 相关的 RedHawk 运行参数。（可选；默认：；兼容 DMP）

\textbf{语法：}
\begin{lstlisting}
USER_MCF_FILE <MCF_file>
\end{lstlisting}

\subsubsection{\gsr{VCD_PREPARE_SCENARIO}}
\paragraph*{VCD_PREPARE_SCENARIO}

GSR 关键字 VCD_PREPARE_SCENARIO：用于设置与 VCD、PREPARE、SCENARIO 相关的 RedHawk 运行参数。（可选；默认：1 (On)；兼容 DMP）

\textbf{语法：}
\begin{lstlisting}
Syntax
VCD_PREPARE_SCENARIO [ 0 | 1 ]
\end{lstlisting}

\subsubsection{\gsr{VCD_SCENARIO_COMPRESS}}
\paragraph*{VCD_SCENARIO_COMPRESS}

GSR 关键字 VCD_SCENARIO_COMPRESS：用于设置与 VCD、SCENARIO、COMPRESS 相关的 RedHawk 运行参数。（可选；默认：0 (off)；兼容 DMP）

\textbf{语法：}
\begin{lstlisting}
Syntax
VCD_SCENARIO_COMPRESS [ 0 | 1 ]
\end{lstlisting}

\subsubsection{\gsr{VCD_TIME_ALIGNMENT}}
\paragraph*{VCD_TIME_ALIGNMENT}

GSR 关键字 VCD_TIME_ALIGNMENT：用于设置与 VCD、TIME、ALIGNMENT 相关的 RedHawk 运行参数。（可选；默认：(1), VCD start times are aligned and cyc；兼容 DMP）

\textbf{语法：}
\begin{lstlisting}
VCD_TIME_ALIGNMENT [ 0 | 1 ]
\end{lstlisting}

\subsubsection{\gsr{VCD_X_LOGIC_STATE}}
\paragraph*{VCD_X_LOGIC_STATE}

GSR 关键字 VCD_X_LOGIC_STATE：用于设置与 VCD、X、LOGIC、STATE 相关的 RedHawk 运行参数。（默认：)；兼容 DMP）

\textbf{语法：}
\begin{lstlisting}
VCD_X_LOGIC_STATE [-1|0|1|2|3]
\end{lstlisting}

\subsubsection{\gsr{VECTORLESS_BLOCK}}
\paragraph*{VECTORLESS_BLOCK}

GSR 关键字 VECTORLESS_BLOCK：用于设置与 VECTORLESS、模块 相关的 RedHawk 运行参数。（可选；默认：none；兼容 DMP）

\textbf{语法：}
\begin{lstlisting}
Syntax
VECTORLESS_BLOCK {
<Block1>
<Block2>
...
}
\end{lstlisting}

\subsection{电迁移（EM）关键字}
\subsubsection*{Electromigration Keywords}

本类共 81 个 GSR 关键字。以下逐条给出中文用途、语法、默认值与示例。

\subsubsection{\gsr{ANALYZE_NETS}}
\paragraph*{ANALYZE_NETS}

GSR 关键字 ANALYZE_NETS：用于设置与 ANALYZE、网络 相关的 RedHawk 运行参数。（可选；默认：none；兼容 DMP）

\textbf{语法：}
\begin{lstlisting}
Syntax
ANALYZE_NETS {
<signal_net_name_1>
<signal_net_name_2>
...
}
\end{lstlisting}

\subsubsection{\gsr{ANALYZE_NETS_FILE}}
\paragraph*{ANALYZE_NETS_FILE}

GSR 关键字 ANALYZE_NETS_FILE：用于设置与 ANALYZE、网络、文件 相关的 RedHawk 运行参数。（可选；默认：none；兼容 DMP）

\textbf{语法：}
\begin{lstlisting}
Syntax
ANALYZE_NETS_FILE <file_with_signal_nets_names>
\end{lstlisting}

\subsubsection{\gsr{BRIDGE_WIRE_CONNECTION}}
\paragraph*{BRIDGE_WIRE_CONNECTION}

GSR 关键字 BRIDGE_WIRE_CONNECTION：用于设置与 BRIDGE、WIRE、连接 相关的 RedHawk 运行参数。

\textbf{语法：}
\begin{lstlisting}
BRIDGE_WIRE_CONNECTION 1
\end{lstlisting}

\subsubsection{\gsr{CONFIGURABLE_REPORT_FILE}}
\paragraph*{CONFIGURABLE_REPORT_FILE}

GSR 关键字 CONFIGURABLE_REPORT_FILE：用于设置与 CONFIGURABLE、REPORT、文件 相关的 RedHawk 运行参数。（默认：list of items that are included in the c）

\subsubsection{\gsr{EM_WORST_POWER}}
\paragraph*{EM_WORST_POWER}

GSR 关键字 EM_WORST_POWER：用于设置与 电迁移、WORST、电源 相关的 RedHawk 运行参数。（默认：Contents of Report file for signal EM:）

\subsubsection{\gsr{EM_WORST_SIGNAL}}
\paragraph*{EM_WORST_SIGNAL}

GSR 关键字 EM_WORST_SIGNAL：用于设置与 电迁移、WORST、SIGNAL 相关的 RedHawk 运行参数。（可选；默认：header of the ‘dump res_network’ when op；兼容 DMP）

\textbf{语法：}
\begin{lstlisting}
CONFIGURABLE_REPORT_FILE <config_filename>
\end{lstlisting}

\subsubsection{\gsr{DELTA_T_RMS_EM}}
\paragraph*{DELTA_T_RMS_EM}

GSR 关键字 DELTA_T_RMS_EM：用于设置与 DELTA、T、RMS、电迁移 相关的 RedHawk 运行参数。（可选；默认：none；兼容 DMP）

\textbf{语法：}
\begin{lstlisting}
DELTA_T_RMS_EM <temp_increase_deg C>
\end{lstlisting}

\subsubsection{\gsr{DETAILED_CONN_ISSUE_RPT}}
\paragraph*{DETAILED_CONN_ISSUE_RPT}

GSR 关键字 DETAILED_CONN_ISSUE_RPT：用于设置与 DETAILED、CONN、ISSUE、RPT 相关的 RedHawk 运行参数。（可选；默认：0）

\textbf{语法：}
\begin{lstlisting}
DETAILED_CONN_ISSUE_RPT [ 0 | 1 ]
\end{lstlisting}

\subsubsection{\gsr{EM_CHECK_2D}}
\paragraph*{EM_CHECK_2D}

GSR 关键字 EM_CHECK_2D：用于设置与 电迁移、检查、2D 相关的 RedHawk 运行参数。（可选；默认：1, On；兼容 DMP）

\textbf{语法：}
\begin{lstlisting}
EM_CHECK_2D [ 0 | 1 ]
\end{lstlisting}

\subsubsection{\gsr{EM_CUSTOM_CURRENT_FILE}}
\paragraph*{EM_CUSTOM_CURRENT_FILE}

GSR 关键字 EM_CUSTOM_CURRENT_FILE：用于设置与 电迁移、CUSTOM、CURRENT、文件 相关的 RedHawk 运行参数。

\textbf{语法：}
\begin{lstlisting}
EM_CUSTOM_CURRENT_FILE <filename>
\end{lstlisting}

\textbf{参数说明：}
\begin{itemize}
\item \texttt{net\_name} — GSR 关键字 EM\_CUSTOM\_CURRENT\_FILE：用于设置与 电迁移、CUSTOM、CURRENT、文件 相关的 RedHawk 运行参数。
\end{itemize}

\textbf{示例：}
\begin{lstlisting}
net1 INST ANA1 VREGO -1 -1 1e-3
net2 INST ANA1 SIGPIN -1 -1 2e-3
The example sets PEAK current to '1e-3' for net1 and '2e-3' for net2 at driving
points VREGO and SIGPIN, respectively. Other current values are unspecified.
Optional. Default: None
EM_CUSTOM_CURRENT_FILE abc2.ccf
\end{lstlisting}

\subsubsection{\gsr{EM_CCF_ONLY}}
\paragraph*{EM_CCF_ONLY}

GSR 关键字 EM_CCF_ONLY：用于设置与 电迁移、CCF、仅 相关的 RedHawk 运行参数。（可选；默认：0 (off)）

\textbf{语法：}
\begin{lstlisting}
EM_CCF_ONLY [ 1 | 0 ]
\end{lstlisting}

\subsubsection{\gsr{EM_DISABLE_SLICED_MINW_CHECK}}
\paragraph*{EM_DISABLE_SLICED_MINW_CHECK}

GSR 关键字 EM_DISABLE_SLICED_MINW_CHECK：用于设置与 电迁移、禁用、SLICED、MINW、检查 相关的 RedHawk 运行参数。（可选；默认：0）

\textbf{语法：}
\begin{lstlisting}
EM_DISABLE_SLICED_MINW_CHECK [1|0]
\end{lstlisting}

\subsubsection{\gsr{EM_DENSITY_ANALYSIS_LAYERS}}
\paragraph*{EM_DENSITY_ANALYSIS_LAYERS}

GSR 关键字 EM_DENSITY_ANALYSIS_LAYERS：用于设置与 电迁移、DENSITY、ANALYSIS、金属层 相关的 RedHawk 运行参数。

\textbf{语法：}
\begin{lstlisting}
EM_DENSITY_ANALYSIS_LAYERS {
layer1
layer2
...
}
\end{lstlisting}

\subsubsection{\gsr{EM_DUMP_PERCENTAGE}}
\paragraph*{EM_DUMP_PERCENTAGE}

GSR 关键字 EM_DUMP_PERCENTAGE：用于设置与 电迁移、DUMP、PERCENTAGE 相关的 RedHawk 运行参数。（可选；默认：None；兼容 DMP）

\textbf{语法：}
\begin{lstlisting}
EM_DUMP_PERCENTAGE [ <min_percent> | <min_percent>-<max_percent>]
\end{lstlisting}

\textbf{参数说明：}
GSR 关键字 EM\_DUMP\_PERCENTAGE：用于设置与 电迁移、DUMP、PERCENTAGE 相关的 RedHawk 运行参数。

\textbf{示例：}
\begin{lstlisting}
EM_DUMP_PERCENTAGE 0-20
For wires, the line format for EM data in the *.em file is as follows:
<layer> <segment_end_coords> <EM_Ratio> <Current_value> <netname> <width>
For vias, the line format for the EM data is:
<via_name> <x-y_coord> <EM_Ratio> <current_value> <netname>
<via_cut_bounding_box_coords> <current_direction>
Current values are in Amps, and coordinates and dimensions are in um. To include
values of <via_cut_bounding_box_coords> and <current_direction> in the report,
users must set the GSR keyword PRINT_EM_VIA_BOX to 1.
\end{lstlisting}

\subsubsection{\gsr{EM_LENGTH_FACTORS_DC_ONLY}}
\paragraph*{EM_LENGTH_FACTORS_DC_ONLY}

GSR 关键字 EM_LENGTH_FACTORS_DC_ONLY：用于设置与 电迁移、LENGTH、FACTORS、DC、仅 相关的 RedHawk 运行参数。（默认：1）

\textbf{语法：}
\begin{lstlisting}
EM_LENGTH_FACTORS_DC_ONLY [ 0 | 1]
\end{lstlisting}

\subsubsection{\gsr{EM_LENGTH_USE_MAX_LENGTH}}
\paragraph*{EM_LENGTH_USE_MAX_LENGTH}

GSR 关键字 EM_LENGTH_USE_MAX_LENGTH：用于设置与 电迁移、LENGTH、USE、MAX、LENGTH 相关的 RedHawk 运行参数。（可选；默认：)；兼容 DMP）

\textbf{语法：}
\begin{lstlisting}
EM_LENGTH_USE_MAX_LENGTH [0|1]
\end{lstlisting}

\subsubsection{\gsr{EM_MISSION_PROFILE}}
\paragraph*{EM_MISSION_PROFILE}

GSR 关键字 EM_MISSION_PROFILE：用于设置与 电迁移、MISSION、PROFILE 相关的 RedHawk 运行参数。（可选；默认：all coefficients are set to 1, in which）

\textbf{语法：}
\begin{lstlisting}
EM_MISSION_PROFILE <path_to_file>
Contents of a sample mission profile are shown below:
via VIA12 {
{ 0 5 }
{ 25 10 }
{ 105 100 }
{ 110 0.697 }
{ 125 0.25 }
{ 150 0.053 }
}
metal METAL1 {
{ 0 5 }
{ 25 11 }
{ 105 111 }
{ 110 0.697 }
{ 125 0.20 }
{ 150 0.053 }
}
In the above example, at a temperature of 125° C, the EM limit for METAL1 is
scaled by a factor of 0.20.
Two dimensional example mission profile:
metal M1 {
Temperature { 0 25 50 100 125 }
Life_Cycle { 1000 2000 3000 }
\end{lstlisting}

\subsubsection{\gsr{EM_MODE}}
\paragraph*{EM_MODE}

GSR 关键字 EM_MODE：用于设置与 电迁移、模式 相关的 RedHawk 运行参数。（可选；默认：peak；兼容 DMP）

\textbf{语法：}
\begin{lstlisting}
EM_MODE [ avg | rms | peak ]
\end{lstlisting}

\textbf{示例：}
\begin{lstlisting}
EM_MODE peak
\end{lstlisting}

\subsubsection{\gsr{EM_NET_INFO}}
\paragraph*{EM_NET_INFO}

GSR 关键字 EM_NET_INFO：用于设置与 电迁移、网络、INFO 相关的 RedHawk 运行参数。

\textbf{语法：}
\begin{lstlisting}
EM_NET_INFO <filename>
\end{lstlisting}

\textbf{参数说明：}
GSR 关键字 EM\_NET\_INFO：用于设置与 电迁移、网络、INFO 相关的 RedHawk 运行参数。（可选；默认：none；兼容 DMP）

\subsubsection{\gsr{EM_REPORT_MINWIDTH}}
\paragraph*{EM_REPORT_MINWIDTH}

GSR 关键字 EM_REPORT_MINWIDTH：用于设置与 电迁移、REPORT、MINWIDTH 相关的 RedHawk 运行参数。（可选；默认：minimum wire width for reporting EM viol）

\textbf{语法：}
\begin{lstlisting}
EM_REPORT_MINWIDTH <min_width_um>
\end{lstlisting}

\subsubsection{\gsr{EM_REPORT_MODE_ONLY}}
\paragraph*{EM_REPORT_MODE_ONLY}

GSR 关键字 EM_REPORT_MODE_ONLY：用于设置与 电迁移、REPORT、模式、仅 相关的 RedHawk 运行参数。（可选；默认：), or peak, but all three types of repor）

\textbf{语法：}
\begin{lstlisting}
EM_REPORT_MODE_ONLY [ 0 | 1 ]
\end{lstlisting}

\textbf{示例：}
\begin{lstlisting}
EM_REPORT_MODE_ONLY 1
\end{lstlisting}

\subsubsection{\gsr{EM_REPORT_PERCENTAGE}}
\paragraph*{EM_REPORT_PERCENTAGE}

GSR 关键字 EM_REPORT_PERCENTAGE：用于设置与 电迁移、REPORT、PERCENTAGE 相关的 RedHawk 运行参数。（可选；默认：100 (percent)；兼容 DMP）

\textbf{语法：}
\begin{lstlisting}
EM_REPORT_PERCENTAGE <percent_violation_thresh>
\end{lstlisting}

\textbf{示例：}
\begin{lstlisting}
EM_REPORT_PERCENTAGE 80
\end{lstlisting}

\subsubsection{\gsr{EM_REPORT_PERCENTAGE_BY_LAYER}}
\paragraph*{EM_REPORT_PERCENTAGE_BY_LAYER}

GSR 关键字 EM_REPORT_PERCENTAGE_BY_LAYER：用于设置与 电迁移、REPORT、PERCENTAGE、BY、层 相关的 RedHawk 运行参数。（可选；默认：none；兼容 DMP）

\textbf{语法：}
\begin{lstlisting}
EM_REPORT_PERCENTAGE_BY_LAYER ?[AVG| RMS| PEAK] ?{
<layerName1> <percent_of_limit1>
<layerName2> <percent_of_limit2>
...
}
EM_REPORT_DC_PERCENTAGE <percent_of_limit>
EM_REPORT_RMS_PERCENTAGE <percent_of_limit>
EM_REPORT_PEAK_PERCENTAGE <percent_of_limit>
\end{lstlisting}

\subsubsection{\gsr{EM_REPORT_LINE_NUMBER}}
\paragraph*{EM_REPORT_LINE_NUMBER}

GSR 关键字 EM_REPORT_LINE_NUMBER：用于设置与 电迁移、REPORT、LINE、NUMBER 相关的 RedHawk 运行参数。（可选；默认：1000；兼容 DMP）

\textbf{语法：}
\begin{lstlisting}
EM_REPORT_LINE_NUMBER [<Max_lines> | -1]
\end{lstlisting}

\textbf{示例：}
\begin{lstlisting}
EM_REPORT_LINE_NUMBER -1
\end{lstlisting}

\subsubsection{\gsr{EM_SCALE_DC}}
\paragraph*{EM_SCALE_DC}

GSR 关键字 EM_SCALE_DC：用于设置与 电迁移、缩放、DC 相关的 RedHawk 运行参数。

\subsubsection{\gsr{EM_SCALE_PEAK}}
\paragraph*{EM_SCALE_PEAK}

GSR 关键字 EM_SCALE_PEAK：用于设置与 电迁移、缩放、峰值 相关的 RedHawk 运行参数。

\subsubsection{\gsr{EM_SCALE_RMS}}
\paragraph*{EM_SCALE_RMS}

GSR 关键字 EM_SCALE_RMS：用于设置与 电迁移、缩放、RMS 相关的 RedHawk 运行参数。（可选；默认：1；兼容 DMP）

\textbf{语法：}
\begin{lstlisting}
EM_SCALE_DC <scale_value>
EM_SCALE_PEAK <scale_value>
EM_SCALE_RMS <scale_value>
\end{lstlisting}

\subsubsection{\gsr{EM_SLEW_NO_STA}}
\paragraph*{EM_SLEW_NO_STA}

GSR 关键字 EM_SLEW_NO_STA：用于设置与 电迁移、SLEW、NO、STA 相关的 RedHawk 运行参数。

\textbf{语法：}
\begin{lstlisting}
EM_SLEW_NO_STA [ ideal | derived | input_transition ]
\end{lstlisting}

\textbf{示例：}
\begin{lstlisting}
example, EM_SLEW_SIG_TIME and EM_SLEW_SIG_PERCENTAGE-- are
taken without any changes.
•    'input_transition' - means that the argument from GSR keyword
INPUT_TRANSITION (which has a default of 100ps) is used as the value for
the missing STA.data. Note the difference in this option compared to 'derived';
the 'derived' value is NOT considered a replacement for STA, but is taking the
time argument directly. Fundamentally, the 'input_transition' value still takes the
actual net load into account and computes from both the net load and the given
INPUT_TRANSITION number its own transition time in a non-linear fashion.
(Note that signal nets that are not driven at all can also be considered, but this
keyword is not applicable to them.) DMP compatible. Optional. Default: ideal.
EM_SLEW_NO_STA derived
\end{lstlisting}

\subsubsection{\gsr{EM_TECH_DC}}
\paragraph*{EM_TECH_DC}

GSR 关键字 EM_TECH_DC：用于设置与 电迁移、工艺、DC 相关的 RedHawk 运行参数。

\subsubsection{\gsr{EM_TECH_PEAK}}
\paragraph*{EM_TECH_PEAK}

GSR 关键字 EM_TECH_PEAK：用于设置与 电迁移、工艺、峰值 相关的 RedHawk 运行参数。

\subsubsection{\gsr{EM_TECH_RMS}}
\paragraph*{EM_TECH_RMS}

GSR 关键字 EM_TECH_RMS：用于设置与 电迁移、工艺、RMS 相关的 RedHawk 运行参数。（可选；默认：tech file；兼容 DMP）

\textbf{语法：}
\begin{lstlisting}
EM_TECH_DC <EM_ave_limit_filename>
EM_TECH_PEAK <EM_peak_limit_filename>
EM_TECH_RMS <EM_RMS_limit_filename>
\end{lstlisting}

\subsubsection{\gsr{EM_TECH_FILE}}
\paragraph*{EM_TECH_FILE}

引用外部 EM 规则 tech 文件。（可选；默认：tech file）

\textbf{语法：}
\begin{lstlisting}
EM_Tech_FILE <EM rule set file>
\end{lstlisting}

\subsubsection{\gsr{EM_TOPOLOGY_USE_ELECTRON_FLOW}}
\paragraph*{EM_TOPOLOGY_USE_ELECTRON_FLOW}

GSR 关键字 EM_TOPOLOGY_USE_ELECTRON_FLOW：用于设置与 电迁移、TOPOLOGY、USE、ELECTRON、FLOW 相关的 RedHawk 运行参数。（可选；默认：, RedHawk uses the width of the wire tha）

\textbf{语法：}
\begin{lstlisting}
EM_TOPOLOGY_USE_ELECTRON_FLOW [ 0 | 1 ]
\end{lstlisting}

\subsubsection{\gsr{ENABLE_AUTO_EM}}
\paragraph*{ENABLE_AUTO_EM}

启用 自动、电迁移 相关功能或检查。（可选；默认：you must execute a separate 'perform emc）

\textbf{语法：}
\begin{lstlisting}
Syntax
ENABLE_AUTO_EM [ 0 | 1 ]
\end{lstlisting}

\subsubsection{\gsr{ENABLE_BLECH}}
\paragraph*{ENABLE_BLECH}

启用 BLECH 相关功能或检查。（可选；默认：the Blech length filtering function is o；兼容 DMP）

\textbf{语法：}
\begin{lstlisting}
ENABLE_BLECH [ 0 | 1 ]
\end{lstlisting}

\subsubsection{\gsr{ENABLE_POLYNOMIAL_EM}}
\paragraph*{ENABLE_POLYNOMIAL_EM}

启用 多项式、电迁移 相关功能或检查。（可选；默认：0: ENABLE_BLECH           ENABLE_POLYNOM；兼容 DMP）

\textbf{语法：}
\begin{lstlisting}
ENABLE_POLYNOMIAL_EM [ 0| 1 ]
\end{lstlisting}

\subsubsection{\gsr{EQUAL_POTENTIAL_AROUND_PAD}}
\paragraph*{EQUAL_POTENTIAL_AROUND_PAD}

GSR 关键字 EQUAL_POTENTIAL_AROUND_PAD：用于设置与 EQUAL、POTENTIAL、AROUND、焊盘 相关的 RedHawk 运行参数。（默认：0；兼容 DMP）

\textbf{语法：}
\begin{lstlisting}
EQUAL_POTENTIAL_AROUND_PAD [ 0 | 1]
\end{lstlisting}

\subsubsection{\gsr{IGNORE_HALF_NODE_SCALE_FOR_EM}}
\paragraph*{IGNORE_HALF_NODE_SCALE_FOR_EM}

在数据输入或分析流程中忽略与 半、节点、缩放、FOR、电迁移 相关的对象或检查。（可选；默认：0 (Off)；兼容 DMP）

\textbf{语法：}
\begin{lstlisting}
IGNORE_HALF_NODE_SCALE_FOR_EM [ 0| 1 ]
\end{lstlisting}

\subsubsection{\gsr{IGNORE_INST_POSTPROCESS}}
\paragraph*{IGNORE_INST_POSTPROCESS}

在数据输入或分析流程中忽略与 INST、POSTPROCESS 相关的对象或检查。（可选；默认：0 (Off) Global System Requirements File）

\textbf{语法：}
\begin{lstlisting}
IGNORE_INST_POSTPROCESS [ 0| 1 ]
\end{lstlisting}

\subsubsection{\gsr{IGNORE_LEF_DEF_SCALE_FOR_EM}}
\paragraph*{IGNORE_LEF_DEF_SCALE_FOR_EM}

在数据输入或分析流程中忽略与 LEF、DEF、缩放、FOR、电迁移 相关的对象或检查。（可选；默认：0）

\textbf{语法：}
\begin{lstlisting}
IGNORE_LEF_DEF_SCALE_FOR_EM [ 0| 1 ]
\end{lstlisting}

\subsubsection{\gsr{NEW_MERGE_WIRE}}
\paragraph*{NEW_MERGE_WIRE}

GSR 关键字 NEW_MERGE_WIRE：用于设置与 NEW、MERGE、WIRE 相关的 RedHawk 运行参数。（可选；默认：0 - 0ff）

\textbf{语法：}
\begin{lstlisting}
NEW_MERGE_WIRE [ 0 | 1 ]
\end{lstlisting}

\subsubsection{\gsr{MERGE_ABUTTED_CUTS}}
\paragraph*{MERGE_ABUTTED_CUTS}

GSR 关键字 MERGE_ABUTTED_CUTS：用于设置与 MERGE、ABUTTED、CUTS 相关的 RedHawk 运行参数。（可选；默认：applies EM limits per via cut based on a）

\textbf{语法：}
\begin{lstlisting}
MERGE_ABUTTED_CUTS [0 | 1]
\end{lstlisting}

\subsubsection{\gsr{PS_LIB_EXTRACT_CCS}}
\paragraph*{PS_LIB_EXTRACT_CCS}

GSR 关键字 PS_LIB_EXTRACT_CCS：用于设置与 PS、LIB、提取、CCS 相关的 RedHawk 运行参数。（可选；默认：0）

\textbf{语法：}
\begin{lstlisting}
PS_LIB_EXTRACT_CCS [ 0 | 1 ]
\end{lstlisting}

\subsubsection{\gsr{SEM_ACCURACY}}
\paragraph*{SEM_ACCURACY}

GSR 关键字 SEM_ACCURACY：用于设置与 SEM、ACCURACY 相关的 RedHawk 运行参数。（可选；默认：none；必需）

\textbf{语法：}
\begin{lstlisting}
SEM_ACCURACY [clockmesh_moderate | ccs_dynamic ]
\end{lstlisting}

\subsubsection{\gsr{SEM_ANALYZE_NET_ONLY}}
\paragraph*{SEM_ANALYZE_NET_ONLY}

GSR 关键字 SEM_ANALYZE_NET_ONLY：用于设置与 SEM、ANALYZE、网络、仅 相关的 RedHawk 运行参数。（可选；默认：value ‘ALL’ analyzes all signal nets；兼容 DMP）

\textbf{语法：}
\begin{lstlisting}
SEM_ANALYZE_NET_ONLY [CLOCK|SIGNAL|LAYER|VIA|ALL]
\end{lstlisting}

\textbf{参数说明：}
\begin{itemize}
\item \texttt{CLOCK} — GSR 关键字 SEM\_ANALYZE\_NET\_ONLY：用于设置与 SEM、ANALYZE、网络、仅 相关的 RedHawk 运行参数。
\end{itemize}

\textbf{示例：}
\begin{lstlisting}
SEM_ANALYZE_NET_ONLY LAYER
SEM_CONNECT_NETS'
When set in the automatic signal-net connection flow, connects unconnected signal
nets for early analysis. Optional. Default: 0.
SEM_CONNECT_NETS [ 0 | 1 ]
\end{lstlisting}

\subsubsection{\gsr{SEM_DEFAULT_PARAMETERS}}
\paragraph*{SEM_DEFAULT_PARAMETERS}

GSR 关键字 SEM_DEFAULT_PARAMETERS：用于设置与 SEM、DEFAULT、PARAMETERS 相关的 RedHawk 运行参数。（可选；默认：values: none）

\textbf{语法：}
\begin{lstlisting}
SEM_DEFAULT_PARAMETERS {
PRIMARY_INPUT_PIN_SLEW <transition_time_sec>
PRIMARY_INPUT_PIN_SLEW_FILE <filename>
PRIMARY_OUTPUT_PIN_CAP <load_cap_Farads>
PRIMARY_OUTPUT_PIN_CAP_FILE <filename>
CONST_NET_SLEW <transition_time_sec>
CONST_NET_TOGGLE <toggle_rate>
}
\end{lstlisting}

\subsubsection{\gsr{SEM_DRV_CURRENT_FILE}}
\paragraph*{SEM_DRV_CURRENT_FILE}

GSR 关键字 SEM_DRV_CURRENT_FILE：用于设置与 SEM、DRV、CURRENT、文件 相关的 RedHawk 运行参数。

\textbf{语法：}
\begin{lstlisting}
SEM_DRV_CURRENT_FILE <drv current file>
\end{lstlisting}

\textbf{示例：}
\begin{lstlisting}
u_noram/u_dside/u_dbiu/U6101 Z 1e-11 -0.1 5e-11 0.1 5e-11
Optional. Default values: none.
\end{lstlisting}

\subsubsection{\gsr{SEM_DUTY_RATIO_ROOT}}
\paragraph*{SEM_DUTY_RATIO_ROOT}

GSR 关键字 SEM_DUTY_RATIO_ROOT：用于设置与 SEM、DUTY、RATIO、ROOT 相关的 RedHawk 运行参数。

\textbf{语法：}
\begin{lstlisting}
SEM_DUTY_RATIO_ROOT <root_value>
\end{lstlisting}

\textbf{示例：}
\begin{lstlisting}
SEM_DUTY_RATIO_ROOT 3
In this example, Ipeak_new = I_peak * r^(1/3)
\end{lstlisting}

\subsubsection{\gsr{SEM_ENABLE_SHORTS_REPORT}}
\paragraph*{SEM_ENABLE_SHORTS_REPORT}

GSR 关键字 SEM_ENABLE_SHORTS_REPORT：用于设置与 SEM、启用、SHORTS、REPORT 相关的 RedHawk 运行参数。（可选；默认：0）

\textbf{语法：}
\begin{lstlisting}
SEM_ENABLE_SHORTS_REPORT [ 0 | 1 ]
\end{lstlisting}

\subsubsection{\gsr{SEM_EXTRACT_LONG_WIRE}}
\paragraph*{SEM_EXTRACT_LONG_WIRE}

GSR 关键字 SEM_EXTRACT_LONG_WIRE：用于设置与 SEM、提取、LONG、WIRE 相关的 RedHawk 运行参数。（可选；默认：0）

\textbf{语法：}
\begin{lstlisting}
SEM_EXTRACT_LONG_WIRE [0 | 1]
\end{lstlisting}

\subsubsection{\gsr{SEM_GEO_OVERLAP_CHECK}}
\paragraph*{SEM_GEO_OVERLAP_CHECK}

GSR 关键字 SEM_GEO_OVERLAP_CHECK：用于设置与 SEM、GEO、OVERLAP、检查 相关的 RedHawk 运行参数。（可选；默认：0；兼容 DMP）

\textbf{语法：}
\begin{lstlisting}
SEM_GEO_OVERLAP_CHECK [0 | 1]
\end{lstlisting}

\subsubsection{\gsr{SEM_HIERARCHICAL_MODE}}
\paragraph*{SEM_HIERARCHICAL_MODE}

GSR 关键字 SEM_HIERARCHICAL_MODE：用于设置与 SEM、HIERARCHICAL、模式 相关的 RedHawk 运行参数。（可选；默认：, all signal nets are analyzed at the sa）

\textbf{语法：}
\begin{lstlisting}
SEM_HIERARCHICAL_MODE [ 0 | top_only | internal_only ]
\end{lstlisting}

\textbf{参数说明：}
\begin{itemize}
\item \texttt{0(default)} — GSR 关键字 SEM\_HIERARCHICAL\_MODE：用于设置与 SEM、HIERARCHICAL、模式 相关的 RedHawk 运行参数。
\end{itemize}

\subsubsection{\gsr{SEM_IGNORE_DISCONNECT}}
\paragraph*{SEM_IGNORE_DISCONNECT}

GSR 关键字 SEM_IGNORE_DISCONNECT：用于设置与 SEM、忽略、DISCONNECT 相关的 RedHawk 运行参数。（可选；默认：0）

\textbf{语法：}
\begin{lstlisting}
SEM_IGNORE_DISCONNECT [0 | 1]
\end{lstlisting}

\subsubsection{\gsr{SEM_IGNORE_NETS_MISSING_DATA}}
\paragraph*{SEM_IGNORE_NETS_MISSING_DATA}

GSR 关键字 SEM_IGNORE_NETS_MISSING_DATA：用于设置与 SEM、忽略、网络、MISSING、DATA 相关的 RedHawk 运行参数。（可选；默认：0；兼容 DMP）

\textbf{语法：}
\begin{lstlisting}
SEM_IGNORE_NETS_MISSING_DATA [ 0 | 1 ]
\end{lstlisting}

\subsubsection{\gsr{SEM_IMPORT_CONNECTED_NETS}}
\paragraph*{SEM_IMPORT_CONNECTED_NETS}

GSR 关键字 SEM_IMPORT_CONNECTED_NETS：用于设置与 SEM、导入、CONNECTED、网络 相关的 RedHawk 运行参数。（默认：0）

\textbf{语法：}
\begin{lstlisting}
SEM_ IMPORT_CONNECTED_NETS [ 0 | 1 ]
\end{lstlisting}

\subsubsection{\gsr{SEM_INOUT_PIN_AUTO_SELECTION}}
\paragraph*{SEM_INOUT_PIN_AUTO_SELECTION}

GSR 关键字 SEM_INOUT_PIN_AUTO_SELECTION：用于设置与 SEM、INOUT、引脚、自动、SELECTION 相关的 RedHawk 运行参数。（默认：0）

\textbf{语法：}
\begin{lstlisting}
SEM_INOUT_PIN_AUTO_SELECTION                [ 0 | 1 ]
\end{lstlisting}

\subsubsection{\gsr{SEM_KEEP_EMPTY_REPORT}}
\paragraph*{SEM_KEEP_EMPTY_REPORT}

GSR 关键字 SEM_KEEP_EMPTY_REPORT：用于设置与 SEM、KEEP、EMPTY、REPORT 相关的 RedHawk 运行参数。（可选；默认：the above files are not created）

\textbf{语法：}
\begin{lstlisting}
SEM_KEEP_EMPTY_REPORT [ 0 | 1 ]
\end{lstlisting}

\subsubsection{\gsr{SEM_MULTI_DRIVER_FILE}}
\paragraph*{SEM_MULTI_DRIVER_FILE}

GSR 关键字 SEM_MULTI_DRIVER_FILE：用于设置与 SEM、MULTI、DRIVER、文件 相关的 RedHawk 运行参数。（默认：off）

\textbf{语法：}
\begin{lstlisting}
SEM_MULTI_DRIVER_FILE <filename>
In <filename>,
<Signal net name> <Driver Instance name> <fsdb file> <fsdb
probe name> <start_time>
\end{lstlisting}

\textbf{示例：}
\begin{lstlisting}
net_1 drv_inst_1 filename.fsdb probename_1 1.1e-9
\end{lstlisting}

\subsubsection{\gsr{SEM_NET_CHECK_SHORT}}
\paragraph*{SEM_NET_CHECK_SHORT}

GSR 关键字 SEM_NET_CHECK_SHORT：用于设置与 SEM、网络、检查、短路 相关的 RedHawk 运行参数。（可选；默认：0）

\textbf{语法：}
\begin{lstlisting}
SEM_NET_CHECK_SHORT [ 0 | 1 ]
\end{lstlisting}

\subsubsection{\gsr{SEM_NET_FILTER_DRIVERS_FILE}}
\paragraph*{SEM_NET_FILTER_DRIVERS_FILE}

GSR 关键字 SEM_NET_FILTER_DRIVERS_FILE：用于设置与 SEM、网络、FILTER、DRIVERS、文件 相关的 RedHawk 运行参数。（可选；默认：location of generated file containing th）

\textbf{语法：}
\begin{lstlisting}
SEM_NET_FILTER_DRIVERS_FILE <filename>
\end{lstlisting}

\subsubsection{\gsr{SEM_NET_FILTER_EM_MODE}}
\paragraph*{SEM_NET_FILTER_EM_MODE}

GSR 关键字 SEM_NET_FILTER_EM_MODE：用于设置与 SEM、网络、FILTER、电迁移、模式 相关的 RedHawk 运行参数。（可选；默认：avg；兼容 DMP）

\textbf{语法：}
\begin{lstlisting}
SEM_NET_FILTER_EM_MODE <avg | peak>
\end{lstlisting}

\subsubsection{\gsr{SEM_NET_FILTER_MAX_PEAK}}
\paragraph*{SEM_NET_FILTER_MAX_PEAK}

GSR 关键字 SEM_NET_FILTER_MAX_PEAK：用于设置与 SEM、网络、FILTER、MAX、峰值 相关的 RedHawk 运行参数。（可选；默认：0）

\textbf{语法：}
\begin{lstlisting}
SEM_NET_FILTER_MAX_PEAK <I_max_peak>
\end{lstlisting}

\textbf{示例：}
\begin{lstlisting}
Example
SEM_NET_FILTER_MAX_PEAK 1e-03
\end{lstlisting}

\subsubsection{\gsr{SEM_NET_FREQ}}
\paragraph*{SEM_NET_FREQ}

GSR 关键字 SEM_NET_FREQ：用于设置与 SEM、网络、FREQ 相关的 RedHawk 运行参数。（默认：0 Global System Requirements File (*）

\textbf{语法：}
\begin{lstlisting}
SEM_NET_FREQ [ 0 | 1 ]
\end{lstlisting}

\subsubsection{\gsr{SEM_NET_INFO}}
\paragraph*{SEM_NET_INFO}

GSR 关键字 SEM_NET_INFO：用于设置与 SEM、网络、INFO 相关的 RedHawk 运行参数。（可选）

\textbf{语法：}
\begin{lstlisting}
SEM_NET_INFO <filename>
\end{lstlisting}

\textbf{示例：}
\begin{lstlisting}
<net_name> <trans_time_sec> <toggle_rate> <freq> <uni-dir_scale>
<bi-dir_scale> <extra-cap> <driver_cell> <pin> <Ceff>
<-> <voltage>
# slew_range 0.8
net1 1e-13         2     2e09       0.7     -    1e-14 - - - -
net_154* 10e-12 0.5 2e9 7 - 1e-14 - - - 5 0.85
...
In the example above, 'net1' is assigned a transition time of 1e-13 seconds, within
the slew range of 80%. It's also assigned a toggle rate of 2 and frequency of 2 GHz.
Its uni-directional currents are scaled by 0.7. No bi-directional scaling is used, but
an extra receiver side load pin cap of 1e-14 F is applied. Wild cards for <net_name>
are allowed. DMP compatible. Optional. Default: none.
\end{lstlisting}

\subsubsection{\gsr{SEM_NET_REPORT}}
\paragraph*{SEM_NET_REPORT}

GSR 关键字 SEM_NET_REPORT：用于设置与 SEM、网络、REPORT 相关的 RedHawk 运行参数。（可选；默认：0 (no report)）

\textbf{语法：}
\begin{lstlisting}
SEM_NET_REPORT [ 0 | 1 ]
\end{lstlisting}

\subsubsection{\gsr{SEM_NET_TOGGLE_RATE}}
\paragraph*{SEM_NET_TOGGLE_RATE}

GSR 关键字 SEM_NET_TOGGLE_RATE：用于设置与 SEM、网络、翻转、率 相关的 RedHawk 运行参数。（可选；默认：None；兼容 DMP）

\textbf{语法：}
\begin{lstlisting}
SEM_NET_TOGGLE_RATE {
<net_name1> <toggle_rate1>
<net_name2> <toggle_rate2>
...
}
\end{lstlisting}

\textbf{示例：}
\begin{lstlisting}
SEM_NET_TOGGLE_RATE {
net1 0.1
net2 0.2
...
}
\end{lstlisting}

\subsubsection{\gsr{SEM_NET_TOGGLE_RATE_FILE}}
\paragraph*{SEM_NET_TOGGLE_RATE_FILE}

GSR 关键字 SEM_NET_TOGGLE_RATE_FILE：用于设置与 SEM、网络、翻转、率、文件 相关的 RedHawk 运行参数。（可选；默认：None）

\textbf{语法：}
\begin{lstlisting}
SEM_NET_TOGGLE_RATE_FILE <filename>
\end{lstlisting}

\subsubsection{\gsr{SEM_NEW_SIGEM_INFO_RPT}}
\paragraph*{SEM_NEW_SIGEM_INFO_RPT}

GSR 关键字 SEM_NEW_SIGEM_INFO_RPT：用于设置与 SEM、NEW、SIGEM、INFO、RPT 相关的 RedHawk 运行参数。（可选；默认：0）

\textbf{语法：}
\begin{lstlisting}
SEM_NEW_SIGEM_INFO_RPT [ 0 | 1 ]
\end{lstlisting}

\subsubsection{\gsr{SEM_RECOVERY_FACTOR}}
\paragraph*{SEM_RECOVERY_FACTOR}

GSR 关键字 SEM_RECOVERY_FACTOR：用于设置与 SEM、RECOVERY、因子 相关的 RedHawk 运行参数。（默认：, RedHawk calculates the rectified avera）

\textbf{语法：}
\begin{lstlisting}
SEM_RECOVERY_FACTOR <R>
\end{lstlisting}

\textbf{参数说明：}
GSR 关键字 SEM\_RECOVERY\_FACTOR：用于设置与 SEM、RECOVERY、因子 相关的 RedHawk 运行参数。（可选；默认：none；兼容 DMP）

\subsubsection{\gsr{SEM_SLEW_OPTIMIZATION}}
\paragraph*{SEM_SLEW_OPTIMIZATION}

GSR 关键字 SEM_SLEW_OPTIMIZATION：用于设置与 SEM、SLEW、OPTIMIZATION 相关的 RedHawk 运行参数。（可选；默认：3）

\textbf{语法：}
\begin{lstlisting}
SEM_SLEW_OPTIMIZATION [ 1 | 2 | 3 ]
\end{lstlisting}

\subsubsection{\gsr{SEM_SPLIT_LONG_WIRE}}
\paragraph*{SEM_SPLIT_LONG_WIRE}

GSR 关键字 SEM_SPLIT_LONG_WIRE：用于设置与 SEM、SPLIT、LONG、WIRE 相关的 RedHawk 运行参数。（可选；默认：16 um）

\textbf{语法：}
\begin{lstlisting}
SEM_SPLIT_LONG_WIRE <PI_model_length>
\end{lstlisting}

\subsubsection{\gsr{SEM_TURBO_CLEAN_WIRE}}
\paragraph*{SEM_TURBO_CLEAN_WIRE}

GSR 关键字 SEM_TURBO_CLEAN_WIRE：用于设置与 SEM、TURBO、CLEAN、WIRE 相关的 RedHawk 运行参数。（可选；默认：0 (off)）

\textbf{语法：}
\begin{lstlisting}
SEM_TURBO_CLEAN_WIRE [ 0 | 1]
\end{lstlisting}

\subsubsection{\gsr{SEM_VECTORLESS_TIME_STEP}}
\paragraph*{SEM_VECTORLESS_TIME_STEP}

GSR 关键字 SEM_VECTORLESS_TIME_STEP：用于设置与 SEM、VECTORLESS、TIME、STEP 相关的 RedHawk 运行参数。（可选；默认：10ps）

\textbf{语法：}
\begin{lstlisting}
SEM_VECTORLESS_TIME_STEP <step_sec>
\end{lstlisting}

\subsubsection{\gsr{SIGNAL_PIN_FILE}}
\paragraph*{SIGNAL_PIN_FILE}

GSR 关键字 SIGNAL_PIN_FILE：用于设置与 SIGNAL、引脚、文件 相关的 RedHawk 运行参数。（可选）

\textbf{语法：}
\begin{lstlisting}
SIGNAL_PIN_FILE <signal_pin_file>
Signal_pin_file Syntax:
# comment
<pin_name> <x1><y1><x2><y2> <layer>
<direction:INPUT|OUTPUT|INOUT>
\end{lstlisting}

\textbf{示例：}
\begin{lstlisting}
#<pin_name> <x1> <y1> <x2> <y2> <layer>
<direction>
DATA0 15.05 4.65 15.15 4.75                 M2    INPUT
DOUT1 29.6 4.3 30.0 4.7                     M2    OUTPUT
Sample DEF with no pin geometry:
\end{lstlisting}

\subsubsection{\gsr{TEMPERATURE_EM}}
\paragraph*{TEMPERATURE_EM}

GSR 关键字 TEMPERATURE_EM：用于设置与 温度、电迁移 相关的 RedHawk 运行参数。（可选；默认：see section "Temperature Setting for Pow；兼容 DMP）

\textbf{语法：}
\begin{lstlisting}
TEMPERATURE_EM <temp ºC>
\end{lstlisting}

\subsubsection{\gsr{TEMPERATURE_REPORT_LINE_NUMBER}}
\paragraph*{TEMPERATURE_REPORT_LINE_NUMBER}

GSR 关键字 TEMPERATURE_REPORT_LINE_NUMBER：用于设置与 温度、REPORT、LINE、NUMBER 相关的 RedHawk 运行参数。（可选；默认：1000 lines）

\textbf{语法：}
\begin{lstlisting}
TEMPERATURE_REPORT_LINE_NUMBER <max_lines>
\end{lstlisting}

\subsubsection{\gsr{TEMPERATURES_EM}}
\paragraph*{TEMPERATURES_EM}

GSR 关键字 TEMPERATURES_EM：用于设置与 TEMPERATURES、电迁移 相关的 RedHawk 运行参数。（可选；默认：see section "Temperature Setting for Pow）

\textbf{语法：}
\begin{lstlisting}
TEMPERATURES_EM {
<layer_1> <temp_1 ºC>
...
<layer_n> <temp_n ºC>
}
\end{lstlisting}

\subsubsection{\gsr{THERMAL_COUPLING_PITCH}}
\paragraph*{THERMAL_COUPLING_PITCH}

GSR 关键字 THERMAL_COUPLING_PITCH：用于设置与 THERMAL、COUPLING、PITCH 相关的 RedHawk 运行参数。（默认：pitch multiplier is 15 for all metal lay）

\textbf{语法：}
\begin{lstlisting}
THERMAL_COUPLING_PITCH {
<ALL| metal layer> [spacing as a multiplier to metal
pitch]
<metal layer> [spacing as a multiplier to metal pitch]
<metal layer> [spacing as a multiplier to metal pitch]
…
}
\end{lstlisting}

\textbf{示例：}
\begin{lstlisting}
THERMAL_COUPLING_PITCH {
ALL 20 #assign 20 as multiplier for metal layer
METAL1 10 #metal 1 now gets 10 as pitch multiplier
fpoly 0
}
\end{lstlisting}

\subsubsection{\gsr{USE_DRAWN_WIDTH_FOR_EM}}
\paragraph*{USE_DRAWN_WIDTH_FOR_EM}

GSR 关键字 USE_DRAWN_WIDTH_FOR_EM：用于设置与 USE、DRAWN、WIDTH、FOR、电迁移 相关的 RedHawk 运行参数。（可选；默认：the default setting for this keyword is；兼容 DMP）

\textbf{语法：}
\begin{lstlisting}
USE_DRAWN_WIDTH_FOR_EM [ 0| 1 ]
USE_DRAWN_WIDTH_FOR_EM_     USE_DRAWN_WIDTH_FOR_EM_
LOOKUP 1                    LOOKUP 0
USE_DRAWN_WIDTH_FOR_    W (lookup) = W(D)
Not Supported
EM 1                    W= W(D)
USE_DRAWN_WIDTH_FOR_    W (lookup) = W(D)           W (lookup)=W(S)
EM 0                    W=W(S)                      W=W(S)
Some foundries require EM to be computed based on silicon width, but lookup
based on drawn width for 20 nm and lower nodes. This flow is enabled by setting
the two keywords as shown in the table above.
\end{lstlisting}

\subsubsection{\gsr{USE_DRAWN_WIDTH_FOR_EM_LOOKUP}}
\paragraph*{USE_DRAWN_WIDTH_FOR_EM_LOOKUP}

GSR 关键字 USE_DRAWN_WIDTH_FOR_EM_LOOKUP：用于设置与 USE、DRAWN、WIDTH、FOR、电迁移、LOOKUP 相关的 RedHawk 运行参数。（可选；默认：the default setting for this keyword is；兼容 DMP）

\textbf{语法：}
\begin{lstlisting}
USE_DRAWN_WIDTH_FOR_EM_LOOKUP [ 0 | 1 ]
\end{lstlisting}

\subsubsection{\gsr{VIA_COMPRESS}}
\paragraph*{VIA_COMPRESS}

GSR 关键字 VIA_COMPRESS：用于设置与 过孔、COMPRESS 相关的 RedHawk 运行参数。（可选；默认：1 (On)；兼容 DMP）

\textbf{语法：}
\begin{lstlisting}
VIA_COMPRESS [1 | 0 ]
\end{lstlisting}

\subsection{提取与网表关键字}
\subsubsection*{Extraction and Netlisting Keywords}

本类共 72 个 GSR 关键字。以下逐条给出中文用途、语法、默认值与示例。

\subsubsection{\gsr{ACCURATE_METAL_DENSITY}}
\paragraph*{ACCURATE_METAL_DENSITY}

GSR 关键字 ACCURATE_METAL_DENSITY：用于设置与 ACCURATE、金属、DENSITY 相关的 RedHawk 运行参数。（可选；默认：prior to Global System Requirements File；兼容 DMP）

\textbf{语法：}
\begin{lstlisting}
ACCURATE_METAL_DENSITY [1 | 0 ]
\end{lstlisting}

\subsubsection{\gsr{ALLOW_SEPARATED_METAL_PSEUDO_VIAS}}
\paragraph*{ALLOW_SEPARATED_METAL_PSEUDO_VIAS}

GSR 关键字 ALLOW_SEPARATED_METAL_PSEUDO_VIAS：用于设置与 ALLOW、SEPARATED、金属、PSEUDO、VIAS 相关的 RedHawk 运行参数。（可选；默认：(AUTO), when a particular foundry tech f）

\textbf{语法：}
\begin{lstlisting}
ALLOW_SEPARATED_METAL_PSEUDO_VIAS [ AUTO | 0 | 1 ]
\end{lstlisting}

\textbf{参数说明：}
GSR 关键字 ALLOW\_SEPARATED\_METAL\_PSEUDO\_VIAS：用于设置与 ALLOW、SEPARATED、金属、PSEUDO、VIAS 相关的 RedHawk 运行参数。（默认：) - as described above, depending on wha）

\subsubsection{\gsr{AUTO_INTERNAL_NET_EXTRACT}}
\paragraph*{AUTO_INTERNAL_NET_EXTRACT}

GSR 关键字 AUTO_INTERNAL_NET_EXTRACT：用于设置与 自动、INTERNAL、网络、提取 相关的 RedHawk 运行参数。（默认：1, On）

\textbf{语法：}
\begin{lstlisting}
AUTO_INTERNAL_NET_EXTRACT [ 0 | 1 ]
\end{lstlisting}

\subsubsection{\gsr{CELL_CURRENT_DIST_FILE}}
\paragraph*{CELL_CURRENT_DIST_FILE}

GSR 关键字 CELL_CURRENT_DIST_FILE：用于设置与 单元、CURRENT、DIST、文件 相关的 RedHawk 运行参数。（可选；默认：none；兼容 DMP）

\textbf{语法：}
\begin{lstlisting}
CELL_CURRENT_DIST_FILE <filename>
\end{lstlisting}

\subsubsection{\gsr{CELL_PIN_FILE}}
\paragraph*{CELL_PIN_FILE}

GSR 关键字 CELL_PIN_FILE：用于设置与 单元、引脚、文件 相关的 RedHawk 运行参数。（默认：Off (0) Global System Requirements File；兼容 DMP）

\textbf{语法：}
\begin{lstlisting}
CELL_PIN_FILE <file_name>
\end{lstlisting}

\subsubsection{\gsr{CEXTRACTION_USE_SPEF}}
\paragraph*{CEXTRACTION_USE_SPEF}

GSR 关键字 CEXTRACTION_USE_SPEF：用于设置与 CEXTRACTION、USE、SPEF 相关的 RedHawk 运行参数。（可选；默认：Off (0)）

\textbf{语法：}
\begin{lstlisting}
CEXTRACTION_USE_SPEF [ 0 | 1 ]
\end{lstlisting}

\subsubsection{\gsr{CEXTRACTION_SPEF_LAYER_MAP}}
\paragraph*{CEXTRACTION_SPEF_LAYER_MAP}

GSR 关键字 CEXTRACTION_SPEF_LAYER_MAP：用于设置与 CEXTRACTION、SPEF、层、MAP 相关的 RedHawk 运行参数。（可选；默认：none）

\textbf{语法：}
\begin{lstlisting}
CEXTRACTION_SPEF_LAYER_MAP <filename>
\end{lstlisting}

\subsubsection{\gsr{CLEAN_VIAS_AFTER_WIRES}}
\paragraph*{CLEAN_VIAS_AFTER_WIRES}

GSR 关键字 CLEAN_VIAS_AFTER_WIRES：用于设置与 CLEAN、VIAS、AFTER、WIRES 相关的 RedHawk 运行参数。（默认：0 (1 by default in v14）

\textbf{语法：}
\begin{lstlisting}
CLEAN_VIAS_AFTER_WIRES [ 0 | 1 ]
\end{lstlisting}

\subsubsection{\gsr{CMM_RIVETED_CONN}}
\paragraph*{CMM_RIVETED_CONN}

GSR 关键字 CMM_RIVETED_CONN：用于设置与 CMM、RIVETED、CONN 相关的 RedHawk 运行参数。（可选；默认：On (1)；兼容 DMP）

\textbf{语法：}
\begin{lstlisting}
CMM_RIVETED_CONN [ 0 | 1 ]
\end{lstlisting}

\subsubsection{\gsr{CONNECT_SWITCH_PINS}}
\paragraph*{CONNECT_SWITCH_PINS}

GSR 关键字 CONNECT_SWITCH_PINS：用于设置与 CONNECT、SWITCH、引脚 相关的 RedHawk 运行参数。（可选；默认：1；兼容 DMP）

\textbf{语法：}
\begin{lstlisting}
CONNECT_SWITCH_PINS [ 0 | 1 | 2 |3 ]
\end{lstlisting}

\textbf{参数说明：}
\begin{itemize}
\item \texttt{0} — GSR 关键字 CONNECT\_SWITCH\_PINS：用于设置与 CONNECT、SWITCH、引脚 相关的 RedHawk 运行参数。（默认：) 2 : shorts neither INT nor EXT switch）
\end{itemize}

\subsubsection{\gsr{COUPLEC}}
\paragraph*{COUPLEC}

GSR 关键字 COUPLEC：用于设置与 COUPLEC 相关的 RedHawk 运行参数。（可选；默认：0 (off)；兼容 DMP）

\textbf{语法：}
\begin{lstlisting}
COUPLEC [ 0 | 1 | 2]
\end{lstlisting}

\subsubsection{\gsr{DEF_IGNORE_SNET_SHIELD}}
\paragraph*{DEF_IGNORE_SNET_SHIELD}

GSR 关键字 DEF_IGNORE_SNET_SHIELD：用于设置与 DEF、忽略、SNET、SHIELD 相关的 RedHawk 运行参数。（可选；默认：1 (SHIELD SPECIALNETS ignored)；兼容 DMP）

\textbf{语法：}
\begin{lstlisting}
DEF_IGNORE_SNET_SHIELD [ 0 | 1 ]
\end{lstlisting}

\subsubsection{\gsr{DO_PININST_INTERNAL_CONN}}
\paragraph*{DO_PININST_INTERNAL_CONN}

GSR 关键字 DO_PININST_INTERNAL_CONN：用于设置与 DO、PININST、INTERNAL、CONN 相关的 RedHawk 运行参数。（可选；默认：0 (off)）

\textbf{语法：}
\begin{lstlisting}
DO_PININST_INTERNAL_CONN [ 0 | 1 ]
\end{lstlisting}

\subsubsection{\gsr{EXPAND_CELL_PIN_FILE}}
\paragraph*{EXPAND_CELL_PIN_FILE}

GSR 关键字 EXPAND_CELL_PIN_FILE：用于设置与 EXPAND、单元、引脚、文件 相关的 RedHawk 运行参数。

\textbf{语法：}
\begin{lstlisting}
EXPAND_CELL_PIN_FILE <expand_filename>
\end{lstlisting}

\subsubsection{\gsr{EXTRACTION_INC}}
\paragraph*{EXTRACTION_INC}

GSR 关键字 EXTRACTION_INC：用于设置与 提取、INC 相关的 RedHawk 运行参数。（可选；默认：0 (off)）

\textbf{语法：}
\begin{lstlisting}
EXTRACTION_INC [0 | 1]
\end{lstlisting}

\subsubsection{\gsr{EXTRACT_PIN_VOL_INSTS}}
\paragraph*{EXTRACT_PIN_VOL_INSTS}

GSR 关键字 EXTRACT_PIN_VOL_INSTS：用于设置与 提取、引脚、VOL、INSTS 相关的 RedHawk 运行参数。（可选；默认：None）

\textbf{语法：}
\begin{lstlisting}
EXTRACT_PIN_VOL_INSTS {
<hierarchical_path>/<block_inst1>
...
}
\end{lstlisting}

\textbf{示例：}
\begin{lstlisting}
EXTRACT_PIN_VOL_INSTS {
abcd3560/ifd_top_inst
abcd3560/vdec_pri_inst
abcd3560/ter_eq1
abcd3560/test_mips_inst
}
This example creates a probe for each power/ground pin specified in DEF files
describing these hierarchical blocks. To see the locations of these probes (block
pins) use the menu command View -> Connectivity -> Show Block Pins. The
names of the block pins that are created are reported, along with their measured
voltages, in the file adsRpt/fullchip_PV.rpt.
For each block that is specified in the GSR file, a voltage source file (PLOC) is
created. To perform a block level analysis, in a block-level GSR file, specify the
corresponding ploc file generated from the top-level dynamic run. Then run block
level static or dynamic analysis.
\end{lstlisting}

\subsubsection{\gsr{EZ_MERGE_NON_RECT_WIRE}}
\paragraph*{EZ_MERGE_NON_RECT_WIRE}

GSR 关键字 EZ_MERGE_NON_RECT_WIRE：用于设置与 EZ、MERGE、NON、RECT、WIRE 相关的 RedHawk 运行参数。（默认：0 (off)）

\textbf{语法：}
\begin{lstlisting}
EZ_MERGE_NON_RECT_WIRE [ 0 | 1 ]
\end{lstlisting}

\subsubsection{\gsr{EZ_MERGE_NON_RECT_WIRE_MAX_LENGTH}}
\paragraph*{EZ_MERGE_NON_RECT_WIRE_MAX_LENGTH}

GSR 关键字 EZ_MERGE_NON_RECT_WIRE_MAX_LENGTH：用于设置与 EZ、MERGE、NON、RECT、WIRE、MAX、LENGTH 相关的 RedHawk 运行参数。（默认：maximum length 100 um）

\textbf{语法：}
\begin{lstlisting}
EZ_MERGE_NON_RECT_WIRE_MAX_LENGTH <max_length_um>
\end{lstlisting}

\subsubsection{\gsr{FIND_ABUTTED_NONSTD_INSTS_PININSTS}}
\paragraph*{FIND_ABUTTED_NONSTD_INSTS_PININSTS}

GSR 关键字 FIND_ABUTTED_NONSTD_INSTS_PININSTS：用于设置与 FIND、ABUTTED、NONSTD、INSTS、PININSTS 相关的 RedHawk 运行参数。（默认：0 (off)）

\textbf{语法：}
\begin{lstlisting}
FIND_ABUTTED_NONSTD_INSTS_PININSTS [ 0 | 1 ]
\end{lstlisting}

\subsubsection{\gsr{IGNORE_DUMMY_PGNET}}
\paragraph*{IGNORE_DUMMY_PGNET}

在数据输入或分析流程中忽略与 DUMMY、PGNET 相关的对象或检查。（可选；默认：0；兼容 DMP）

\textbf{语法：}
\begin{lstlisting}
IGNORE_DUMMY_PGNET [ 0 | 1 ]
\end{lstlisting}

\subsubsection{\gsr{IGNORE_MACROEEQ}}
\paragraph*{IGNORE_MACROEEQ}

在数据输入或分析流程中忽略与 MACROEEQ 相关的对象或检查。（可选；默认：0；兼容 DMP）

\textbf{语法：}
\begin{lstlisting}
IGNORE_MACROEEQ [ 0 | 1 ]
\end{lstlisting}

\subsubsection{\gsr{IGNORE_OPC_METAL}}
\paragraph*{IGNORE_OPC_METAL}

在数据输入或分析流程中忽略与 OPC、金属 相关的对象或检查。（默认：0 (off - do not ignore any wires)）

\textbf{语法：}
\begin{lstlisting}
IGNORE_OPC_METAL <threshold_dimension>
\end{lstlisting}

\subsubsection{\gsr{IGNORE_PLOC_ON_OBS}}
\paragraph*{IGNORE_PLOC_ON_OBS}

在数据输入或分析流程中忽略与 焊盘位置、ON、OBS 相关的对象或检查。（默认：0）

\textbf{语法：}
\begin{lstlisting}
IGNORE_PLOC_ON_OBS [ 0 | 1 ]
\end{lstlisting}

\subsubsection{\gsr{IGNORE_THICKNESS_VARIATION}}
\paragraph*{IGNORE_THICKNESS_VARIATION}

在数据输入或分析流程中忽略与 厚度、变化 相关的对象或检查。（可选；默认：0）

\textbf{语法：}
\begin{lstlisting}
IGNORE_THICKNESS_VARIATION [ 0 | 1 ]
\end{lstlisting}

\subsubsection{\gsr{INST_CURRENT_DIST_MODE}}
\paragraph*{INST_CURRENT_DIST_MODE}

GSR 关键字 INST_CURRENT_DIST_MODE：用于设置与 INST、CURRENT、DIST、模式 相关的 RedHawk 运行参数。（可选；默认：-1；兼容 DMP）

\textbf{语法：}
\begin{lstlisting}
INST_CURRENT_DIST_MODE [ -1 | 0 | 1 ]
\end{lstlisting}

\subsubsection{\gsr{INTERNAL_CONNECT_PIN_CELLS_FILES}}
\paragraph*{INTERNAL_CONNECT_PIN_CELLS_FILES}

GSR 关键字 INTERNAL_CONNECT_PIN_CELLS_FILES：用于设置与 INTERNAL、CONNECT、引脚、单元、文件 相关的 RedHawk 运行参数。（可选；默认：None；兼容 DMP）

\textbf{语法：}
\begin{lstlisting}
INTERNAL_CONNECT_PIN_CELLS_FILES <cell_list_file>
\end{lstlisting}

\textbf{示例：}
\begin{lstlisting}
INTERNAL_CONNECT_PIN_CELLS_FILES Int_conn_cells
\end{lstlisting}

\subsubsection{\gsr{IR_REPORT_STACKVIA_METAL}}
\paragraph*{IR_REPORT_STACKVIA_METAL}

GSR 关键字 IR_REPORT_STACKVIA_METAL：用于设置与 IR、REPORT、STACKVIA、金属 相关的 RedHawk 运行参数。（可选；默认：0）

\textbf{语法：}
\begin{lstlisting}
IR_REPORT_STACKVIA_METAL [1|0]
\end{lstlisting}

\subsubsection{\gsr{LONG_WIRE_RES_CALC}}
\paragraph*{LONG_WIRE_RES_CALC}

GSR 关键字 LONG_WIRE_RES_CALC：用于设置与 LONG、WIRE、RES、CALC 相关的 RedHawk 运行参数。（可选；默认：Off (0)；兼容 DMP）

\textbf{语法：}
\begin{lstlisting}
LONG_WIRE_RES_CALC [ 0 | 1 ]
\end{lstlisting}

\subsubsection{\gsr{LOWEST_METAL}}
\paragraph*{LOWEST_METAL}

GSR 关键字 LOWEST_METAL：用于设置与 LOWEST、金属 相关的 RedHawk 运行参数。（可选；默认：None；兼容 DMP）

\textbf{语法：}
\begin{lstlisting}
LOWEST_METAL        <metal_layer_name>
\end{lstlisting}

\textbf{示例：}
\begin{lstlisting}
LOWEST_METAL Metal1
\end{lstlisting}

\subsubsection{\gsr{MACRO_IRDROP}}
\paragraph*{MACRO_IRDROP}

GSR 关键字 MACRO_IRDROP：用于设置与 宏、IRDROP 相关的 RedHawk 运行参数。（默认：0；兼容 DMP）

\textbf{语法：}
\begin{lstlisting}
MACRO_IRDROP [ 0 | 1 ]
\end{lstlisting}

\subsubsection{\gsr{MEM_RC_MODEL}}
\paragraph*{MEM_RC_MODEL}

GSR 关键字 MEM_RC_MODEL：用于设置与 MEM、RC、MODEL 相关的 RedHawk 运行参数。（可选；默认：0）

\subsubsection{\gsr{MERGE_ABUTTED_ASYM_CUTS}}
\paragraph*{MERGE_ABUTTED_ASYM_CUTS}

GSR 关键字 MERGE_ABUTTED_ASYM_CUTS：用于设置与 MERGE、ABUTTED、ASYM、CUTS 相关的 RedHawk 运行参数。（可选；默认：Off）

\textbf{语法：}
\begin{lstlisting}
MERGE_ABUTTED_ASYM_CUT [ 0 | 1 ]
\end{lstlisting}

\subsubsection{\gsr{MERGE_WIRE}}
\paragraph*{MERGE_WIRE}

GSR 关键字 MERGE_WIRE：用于设置与 MERGE、WIRE 相关的 RedHawk 运行参数。（可选；默认：(unless there are 45-degree wires in the；兼容 DMP）

\textbf{语法：}
\begin{lstlisting}
MERGE_WIRE [ 0 | 1 ]
\end{lstlisting}

\subsubsection{\gsr{MESH_VIAS_FILE}}
\paragraph*{MESH_VIAS_FILE}

GSR 关键字 MESH_VIAS_FILE：用于设置与 MESH、VIAS、文件 相关的 RedHawk 运行参数。（可选；默认：none）

\textbf{语法：}
\begin{lstlisting}
MESH_VIAS_FILE ecoRouting.tcl
\end{lstlisting}

\subsubsection{\gsr{METAL_DENSITY_BOUNDS}}
\paragraph*{METAL_DENSITY_BOUNDS}

GSR 关键字 METAL_DENSITY_BOUNDS：用于设置与 金属、DENSITY、BOUNDS 相关的 RedHawk 运行参数。（可选；默认：none）

\textbf{语法：}
\begin{lstlisting}
METAL_DENSITY_BOUNDS {
<layer> <min_density> <max_density >
}
\end{lstlisting}

\textbf{示例：}
\begin{lstlisting}
METAL_DENSITY_BOUNDS {
M4 0.1 0.3
}
In this example the computed metal area density for M4 layer metal is restricted
between 0.1 and 0.3.
\end{lstlisting}

\subsubsection{\gsr{MINWIDTH_FROM_LEF}}
\paragraph*{MINWIDTH_FROM_LEF}

GSR 关键字 MINWIDTH_FROM_LEF：用于设置与 MINWIDTH、从、LEF 相关的 RedHawk 运行参数。（可选；默认：1 (On)）

\textbf{语法：}
\begin{lstlisting}
MINWIDTH_FROM_LEF [ 0 | 1 ]
\end{lstlisting}

\subsubsection{\gsr{MIN_WIRE_DIMENSION}}
\paragraph*{MIN_WIRE_DIMENSION}

GSR 关键字 MIN_WIRE_DIMENSION：用于设置与 MIN、WIRE、DIMENSION 相关的 RedHawk 运行参数。（可选；默认：12；兼容 DMP）

\textbf{语法：}
\begin{lstlisting}
MIN_WIRE_DIMENSION <resolution_in_internal_db_unit>
\end{lstlisting}

\textbf{示例：}
\begin{lstlisting}
MIN_WIRE_DIMENSION 12
means a resolution of 0.006 um, since the internal DB unit is 2000 and 12/2000
=0.006.
\end{lstlisting}

\subsubsection{\gsr{MPR_MODE}}
\paragraph*{MPR_MODE}

GSR 关键字 MPR_MODE：用于设置与 MPR、模式 相关的 RedHawk 运行参数。（可选；默认：) 1 - uses pattern recognition, but no n；兼容 DMP）

\textbf{语法：}
\begin{lstlisting}
MPR_MODE [ 0 | 1 | 2 | 3 ]
\end{lstlisting}

\subsubsection{\gsr{MPR_PARTITION_REDUCTION}}
\paragraph*{MPR_PARTITION_REDUCTION}

GSR 关键字 MPR_PARTITION_REDUCTION：用于设置与 MPR、PARTITION、REDUCTION 相关的 RedHawk 运行参数。（可选；默认：off）

\textbf{语法：}
\begin{lstlisting}
MPR_PARTITION_REDUCTION [0|1]
\end{lstlisting}

\subsubsection{\gsr{MPR_POWER_LIMIT_FOR_RED}}
\paragraph*{MPR_POWER_LIMIT_FOR_RED}

GSR 关键字 MPR_POWER_LIMIT_FOR_RED：用于设置与 MPR、电源、LIMIT、FOR、RED 相关的 RedHawk 运行参数。（可选；默认：None）

\textbf{语法：}
\begin{lstlisting}
Syntax
MPR_POWER_LIMIT_FOR_RED <power_W>
\end{lstlisting}

\subsubsection{\gsr{NODE_REDUCTION_MODE}}
\paragraph*{NODE_REDUCTION_MODE}

GSR 关键字 NODE_REDUCTION_MODE：用于设置与 节点、REDUCTION、模式 相关的 RedHawk 运行参数。（可选；默认：ALP_ACCURATE (no node reduction)）

\textbf{语法：}
\begin{lstlisting}
Syntax
NODE_REDUCTION_MODE            [ 0 | ALP_ACCURATE ]
\end{lstlisting}

\subsubsection{\gsr{PIN_SLICE_CELL_LIST_FILE}}
\paragraph*{PIN_SLICE_CELL_LIST_FILE}

GSR 关键字 PIN_SLICE_CELL_LIST_FILE：用于设置与 引脚、SLICE、单元、LIST、文件 相关的 RedHawk 运行参数。（可选；默认：none）

\textbf{语法：}
\begin{lstlisting}
PIN_SLICE_CELL_LIST_FILE <cell_list_file>
\end{lstlisting}

\subsubsection{\gsr{PIN_SLICE_LIMIT}}
\paragraph*{PIN_SLICE_LIMIT}

GSR 关键字 PIN_SLICE_LIMIT：用于设置与 引脚、SLICE、LIMIT 相关的 RedHawk 运行参数。（可选；默认：none）

\textbf{语法：}
\begin{lstlisting}
PIN_SLICE_LIMIT <spacing_um>
\end{lstlisting}

\textbf{示例：}
\begin{lstlisting}
PIN_SLICE_LIMIT 5
This setting creates pin nodes every 5 microns.
\end{lstlisting}

\subsubsection{\gsr{PPI_STD_IGNORE_TOUCH_VIA_MET}}
\paragraph*{PPI_STD_IGNORE_TOUCH_VIA_MET}

GSR 关键字 PPI_STD_IGNORE_TOUCH_VIA_MET：用于设置与 PPI、STD、忽略、TOUCH、过孔、MET 相关的 RedHawk 运行参数。（可选；默认：Off；兼容 DMP）

\textbf{语法：}
\begin{lstlisting}
PPI_STD_IGNORE_TOUCH_VIA_MET [0 | 1 ]
\end{lstlisting}

\subsubsection{\gsr{PPI_CELL_EDGE_MAX_NM_THRESHOLD}}
\paragraph*{PPI_CELL_EDGE_MAX_NM_THRESHOLD}

GSR 关键字 PPI_CELL_EDGE_MAX_NM_THRESHOLD：用于设置与 PPI、单元、EDGE、MAX、NM、THRESHOLD 相关的 RedHawk 运行参数。（默认：is 50 nm)）

\textbf{语法：}
\begin{lstlisting}
PPI_CELL_EDGE_MAX_NM_THRESHOLD <max_length>
\end{lstlisting}

\textbf{示例：}
\begin{lstlisting}
Example
(See example for next keyword)
\end{lstlisting}

\subsubsection{\gsr{PPI_CELL_EDGE_THRESHOLD_PERCENT}}
\paragraph*{PPI_CELL_EDGE_THRESHOLD_PERCENT}

GSR 关键字 PPI_CELL_EDGE_THRESHOLD_PERCENT：用于设置与 PPI、单元、EDGE、THRESHOLD、PERCENT 相关的 RedHawk 运行参数。（可选；默认：1%）

\textbf{语法：}
\begin{lstlisting}
PPI_CELL_EDGE_THRESHOLD_PERCENT <percent>
\end{lstlisting}

\textbf{示例：}
\begin{lstlisting}
Example
Assume a cell with two pin instance geometries in the cell rail direction of:
a. 800 nm
b. 900 nm.
PPI_CELL_EDGE_THRESHOLD_PERCENT 20
PPI_CELL_EDGE_MAX_NM_THRESHOLD 150
So 20% of 1000 nm = 200 nm => 1000 nm – 200 nm = 800nm
So both pin instance geometries are > 800nm and are satisfying the
PPI_CELL_EDGE_THRESHOLD_PERCENT criteria.
In addition, the PPI_CELL_EDGE_MAX_NM_THRESHOLD needs to be met.
Cell length 1000nm – 150 nm = 850 nm
So only pin instance geometries b.) 900 that are > 850 nm will be pushed out.
\end{lstlisting}

\subsubsection{\gsr{PRIMARY_OUTPUT_LOAD_RC_MODEL}}
\paragraph*{PRIMARY_OUTPUT_LOAD_RC_MODEL}

GSR 关键字 PRIMARY_OUTPUT_LOAD_RC_MODEL：用于设置与 PRIMARY、OUTPUT、LOAD、RC、MODEL 相关的 RedHawk 运行参数。（可选；默认：none）

\textbf{语法：}
\begin{lstlisting}
PRIMARY_OUTPUT_LOAD_RC_MODEL {
<net1> <r1> <c1>
<net2> <r2> <c2>
\end{lstlisting}

\subsubsection{\gsr{PUSH_PININST_CELLS_FILES}}
\paragraph*{PUSH_PININST_CELLS_FILES}

GSR 关键字 PUSH_PININST_CELLS_FILES：用于设置与 PUSH、PININST、单元、文件 相关的 RedHawk 运行参数。

\textbf{语法：}
\begin{lstlisting}
Original format:
<cellName> [ All |<pinName1> { <layer1> <layer2> ...} <pinName2> {...} ...]
Newer format:
<cellName>[All|PIN {<pinName1> ...} LAYER {<layerName1> ...} {BBOX <bbox1>}]
...
<cellName>[All|PIN {<pinName1> ...} LAYER {<layerName1> ...} {BBOX <bboxN>}]
For both types of syntax above, wildcards can be used to specify both <cellName>
and <pinName>, using the single-character wildcard “?”, and the multiple-character
wildcard “*”. The <bbox> is specified as the x,y coordinates of the lower left and
upper right corners of the region. BBoxes and layers are the same for all wildcard
members. Multiple lines with the same cellname and different boundary boxes can
be specified. The layername and bbox parameters are optional. A sample push
pininst cells file is shown below:
AB2 PIN { VDDS } LAYER { MET1 } BBOX {0 0 5 5}
AB2 PIN { VDDS } LAYER { MET1 } BBOX {10 10 20 20}
In the above example, all the nodes in the region {0 0 5 5} for cell AB2 on layer
MET1 for pin VDDS are shorted together. Similarly, for the region {10 10 20 20} for
the cell AB2 on layer MET5 for pin VDDS, all nodes are shorted together. However,
these two regions are not shorted. NOTE: Only pins of the same layer and of the
same net can be shorted together. DMP compatible. Optional. Default: None.
PUSH_PININST_CELLS_FILES <ppi_cell_file>
\end{lstlisting}

\subsubsection{\gsr{PUSH_PG_PININST}}
\paragraph*{PUSH_PG_PININST}

GSR 关键字 PUSH_PG_PININST：用于设置与 PUSH、PG、PININST 相关的 RedHawk 运行参数。（可选；默认：behavior is that standard cell pins not；兼容 DMP）

\textbf{语法：}
\begin{lstlisting}
PUSH_PG_PININST [ 0 | 1 ]
\end{lstlisting}

\subsubsection{\gsr{PUSH_SIGNAL_PININST}}
\paragraph*{PUSH_SIGNAL_PININST}

GSR 关键字 PUSH_SIGNAL_PININST：用于设置与 PUSH、SIGNAL、PININST 相关的 RedHawk 运行参数。（可选；默认：0, off；兼容 DMP）

\textbf{语法：}
\begin{lstlisting}
PUSH_SIGNAL_PININST[ 0 | 1 ]
\end{lstlisting}

\subsubsection{\gsr{QUICK_MESH_WIRE_MERGE}}
\paragraph*{QUICK_MESH_WIRE_MERGE}

GSR 关键字 QUICK_MESH_WIRE_MERGE：用于设置与 QUICK、MESH、WIRE、MERGE 相关的 RedHawk 运行参数。（可选；默认：0, off）

\textbf{语法：}
\begin{lstlisting}
QUICK_MESH_WIRE_MERGE [ 0 | 1 ]
\end{lstlisting}

\subsubsection{\gsr{REPORT_ALL_UNCONNECT_PORTS}}
\paragraph*{REPORT_ALL_UNCONNECT_PORTS}

GSR 关键字 REPORT_ALL_UNCONNECT_PORTS：用于设置与 REPORT、ALL、UNCONNECT、PORTS 相关的 RedHawk 运行参数。（可选；默认：Off；兼容 DMP）

\textbf{语法：}
\begin{lstlisting}
REPORT_ALL_UNCONNECT_PORTS [ off | all | macro_only |
macro_detail | macro_std_all ]
\end{lstlisting}

\textbf{参数说明：}
\begin{itemize}
\item \texttt{off} — GSR 关键字 REPORT\_ALL\_UNCONNECT\_PORTS：用于设置与 REPORT、ALL、UNCONNECT、PORTS 相关的 RedHawk 运行参数。（默认：)）
\end{itemize}

\subsubsection{\gsr{SAVE_CONSOLIDATED_R}}
\paragraph*{SAVE_CONSOLIDATED_R}

GSR 关键字 SAVE_CONSOLIDATED_R：用于设置与 SAVE、CONSOLIDATED、R 相关的 RedHawk 运行参数。（默认：lower limit for any junction are saved）

\textbf{语法：}
\begin{lstlisting}
SAVE_CONSOLIDATED_R [ 0 | 1 ]
\end{lstlisting}

\subsubsection{\gsr{SCANLINE_MERGE_LAYERS}}
\paragraph*{SCANLINE_MERGE_LAYERS}

GSR 关键字 SCANLINE_MERGE_LAYERS：用于设置与 SCANLINE、MERGE、金属层 相关的 RedHawk 运行参数。（可选；默认：none；兼容 DMP）

\textbf{语法：}
\begin{lstlisting}
SCANLINE_MERGE_LAYERS {
[ ALL | <layer1> <layer2> ...]
}
\end{lstlisting}

\subsubsection{\gsr{SINGLE_CUT_SLICE_LAYERS}}
\paragraph*{SINGLE_CUT_SLICE_LAYERS}

GSR 关键字 SINGLE_CUT_SLICE_LAYERS：用于设置与 SINGLE、CUT、SLICE、金属层 相关的 RedHawk 运行参数。（可选；默认：none）

\textbf{语法：}
\begin{lstlisting}
SINGLE_CUT_SLICE_LAYERS {
<layer1>
<layer2>
...
}
\end{lstlisting}

\subsubsection{\gsr{SIZE_BASED_PININST_CURRENT_DISTRIBUTION}}
\paragraph*{SIZE_BASED_PININST_CURRENT_DISTRIBUTION}

GSR 关键字 SIZE_BASED_PININST_CURRENT_DISTRIBUTION：用于设置与 SIZE、BASED、PININST、CURRENT、DISTRIBUTION 相关的 RedHawk 运行参数。（可选；默认：1；兼容 DMP）

\textbf{语法：}
\begin{lstlisting}
SIZE_BASED_PININST_CURRENT_DISTRIBUTION [ 0 | 1 ]
\end{lstlisting}

\subsubsection{\gsr{SKIP_CONNECTIVITY_CHECK}}
\paragraph*{SKIP_CONNECTIVITY_CHECK}

GSR 关键字 SKIP_CONNECTIVITY_CHECK：用于设置与 SKIP、CONNECTIVITY、检查 相关的 RedHawk 运行参数。（可选；默认：0）

\textbf{语法：}
\begin{lstlisting}
SKIP_CONNECTIVITY_CHECK [ 0 | 1 ]
\end{lstlisting}

\subsubsection{\gsr{SKIP_LAYERS_FILE}}
\paragraph*{SKIP_LAYERS_FILE}

GSR 关键字 SKIP_LAYERS_FILE：用于设置与 SKIP、金属层、文件 相关的 RedHawk 运行参数。（可选；默认：none）

\textbf{语法：}
\begin{lstlisting}
SKIP_LAYERS_FILE <filename>
\end{lstlisting}

\subsubsection{\gsr{SPECIAL_SML}}
\paragraph*{SPECIAL_SML}

GSR 关键字 SPECIAL_SML：用于设置与 SPECIAL、SML 相关的 RedHawk 运行参数。（默认：none）

\textbf{语法：}
\begin{lstlisting}
SPECIAL_SML {
Layer1
Layer2
...
}
\end{lstlisting}

\subsubsection{\gsr{SPLIT_SPARSE_VIA_ARRAY}}
\paragraph*{SPLIT_SPARSE_VIA_ARRAY}

GSR 关键字 SPLIT_SPARSE_VIA_ARRAY：用于设置与 SPLIT、SPARSE、过孔、ARRAY 相关的 RedHawk 运行参数。（可选；默认：none；兼容 DMP）

\textbf{语法：}
\begin{lstlisting}
SPLIT_SPARSE_VIA_ARRAY <max_length_microns>
\end{lstlisting}

\subsubsection{\gsr{SPLIT_VIA_ARRAY}}
\paragraph*{SPLIT_VIA_ARRAY}

GSR 关键字 SPLIT_VIA_ARRAY：用于设置与 SPLIT、过孔、ARRAY 相关的 RedHawk 运行参数。（可选；默认：none；兼容 DMP）

\textbf{语法：}
\begin{lstlisting}
SPLIT_VIA_ARRAY {
<layer1> <split_threshold1>
<layer2> <split_threshold2>
...
}
\end{lstlisting}

\subsubsection{\gsr{SPLIT_VIA_ARRAY_CELL}}
\paragraph*{SPLIT_VIA_ARRAY_CELL}

GSR 关键字 SPLIT_VIA_ARRAY_CELL：用于设置与 SPLIT、过孔、ARRAY、单元 相关的 RedHawk 运行参数。（可选；默认：none；兼容 DMP）

\textbf{语法：}
\begin{lstlisting}
SPLIT_VIA_ARRAY_CELL {
<cellname1> <split_threshold1_u>
<cellname2> <split_threshold2_u>
...
\end{lstlisting}

\subsubsection{\gsr{STATIC_IR_HIDE_DISCON_INST}}
\paragraph*{STATIC_IR_HIDE_DISCON_INST}

GSR 关键字 STATIC_IR_HIDE_DISCON_INST：用于设置与 静态、IR、HIDE、DISCON、INST 相关的 RedHawk 运行参数。（可选；默认：0；兼容 DMP）

\textbf{语法：}
\begin{lstlisting}
STATIC_IR_HIDE_DISCON_INST [ 0 | 1 ]
\end{lstlisting}

\subsubsection{\gsr{STATIC_REDUCTION}}
\paragraph*{STATIC_REDUCTION}

GSR 关键字 STATIC_REDUCTION：用于设置与 静态、REDUCTION 相关的 RedHawk 运行参数。（可选；默认：1 (On)）

\textbf{语法：}
\begin{lstlisting}
STATIC_REDUCTION [ 0 | 1 ]
\end{lstlisting}

\subsubsection{\gsr{TSV_MODEL_FILE}}
\paragraph*{TSV_MODEL_FILE}

GSR 关键字 TSV_MODEL_FILE：用于设置与 TSV、MODEL、文件 相关的 RedHawk 运行参数。（可选；默认：none）

\textbf{语法：}
\begin{lstlisting}
TSV_MODEL_FILE <Sentinel_model_file>
\end{lstlisting}

\subsubsection{\gsr{USE_INTERP_ETCH_FOR_R}}
\paragraph*{USE_INTERP_ETCH_FOR_R}

GSR 关键字 USE_INTERP_ETCH_FOR_R：用于设置与 USE、INTERP、ETCH、FOR、R 相关的 RedHawk 运行参数。（可选；默认：1）

\textbf{语法：}
\begin{lstlisting}
USE_INTERP_ETCH_FOR_R [ 0 | 1 ]
\end{lstlisting}

\subsubsection{\gsr{USE_MVM_PIN_MODEL}}
\paragraph*{USE_MVM_PIN_MODEL}

GSR 关键字 USE_MVM_PIN_MODEL：用于设置与 USE、MVM、引脚、MODEL 相关的 RedHawk 运行参数。（可选；默认：none；兼容 DMP）

\textbf{语法：}
\begin{lstlisting}
USE_MVM_PIN_MODEL [ 0 | 1 ]
\end{lstlisting}

\subsubsection{\gsr{VIA_MAX_SPACE_FOR_COMP}}
\paragraph*{VIA_MAX_SPACE_FOR_COMP}

GSR 关键字 VIA_MAX_SPACE_FOR_COMP：用于设置与 过孔、MAX、SPACE、FOR、COMP 相关的 RedHawk 运行参数。（可选；默认：none；兼容 DMP）

\textbf{语法：}
\begin{lstlisting}
VIA_MAX_SPACE_FOR_COMP <max_space_um>
\end{lstlisting}

\subsubsection{\gsr{VIAMODEL_NAME_CHECK}}
\paragraph*{VIAMODEL_NAME_CHECK}

GSR 关键字 VIAMODEL_NAME_CHECK：用于设置与 过孔模型、NAME、检查 相关的 RedHawk 运行参数。（可选；默认：1；兼容 DMP）

\subsubsection{\gsr{WIRE_R_COMPENSATION}}
\paragraph*{WIRE_R_COMPENSATION}

GSR 关键字 WIRE_R_COMPENSATION：用于设置与 WIRE、R、COMPENSATION 相关的 RedHawk 运行参数。（可选；默认：extraction always includes the resistanc） 注意：GSR 关键字 WIRE_R_COMPENSATION：用于设置与 WIRE、R、COMPENSATION 相关的 RedHawk 运行参数。

\textbf{语法：}
\begin{lstlisting}
WIRE_R_COMPENSATION [ 0 | 1 ]
\end{lstlisting}

\begin{noteBox}
GSR 关键字 WIRE\_R\_COMPENSATION：用于设置与 WIRE、R、COMPENSATION 相关的 RedHawk 运行参数。
\end{noteBox}

\subsubsection{\gsr{WIRE_SLICE_MIN_DIM}}
\paragraph*{WIRE_SLICE_MIN_DIM}

GSR 关键字 WIRE_SLICE_MIN_DIM：用于设置与 WIRE、SLICE、MIN、DIM 相关的 RedHawk 运行参数。

\subsubsection{\gsr{WIRE_SLICE_WIDTH}}
\paragraph*{WIRE_SLICE_WIDTH}

GSR 关键字 WIRE_SLICE_WIDTH：用于设置与 WIRE、SLICE、WIDTH 相关的 RedHawk 运行参数。（可选；默认：no slicing performed）

\textbf{语法：}
\begin{lstlisting}
WIRE_SLICE_MIN_DIM <min_dimension>
WIRE_SLICE_WIDTH <width>
\end{lstlisting}

\subsection{表征关键字}
\subsubsection*{Characterization Keywords}

本类共 13 个 GSR 关键字。以下逐条给出中文用途、语法、默认值与示例。

\subsubsection{\gsr{APL_FILES}}
\paragraph*{APL_FILES}

GSR 关键字 APL_FILES：用于设置与 APL、文件 相关的 RedHawk 运行参数。（可选；默认：none；兼容 DMP）

\textbf{语法：}
\begin{lstlisting}
APL_FILES {
<input_file> [ current | cdev | pwcap | avm |
current_avm | cap_avm]
...
? <input_dir> [ current | cdev | pwcap ]?
...
}
\end{lstlisting}

\textbf{示例：}
\begin{lstlisting}
APL_FILES {
aaa.spcurrent current
bbb.cdev cdev
ccc.pwcdev pwcap
avm.config avm
vmemory.current current_avm
vmemory.cdev cap_avm
aaa_current_dir current
bbb_cdev_dir cdev
ccc_pwcap_dir pwcap                                     (EQ 1)
}
\end{lstlisting}

\subsubsection{\gsr{APLINT_NUMBER_THREADS}}
\paragraph*{APLINT_NUMBER_THREADS}

GSR 关键字 APLINT_NUMBER_THREADS：用于设置与 APLINT、NUMBER、THREADS 相关的 RedHawk 运行参数。（可选）

\textbf{语法：}
\begin{lstlisting}
APLINT_NUMBER_THREADS <num_threads>
\end{lstlisting}

\subsubsection{\gsr{DYNAMIC_APL_RESAMPLE}}
\paragraph*{DYNAMIC_APL_RESAMPLE}

GSR 关键字 DYNAMIC_APL_RESAMPLE：用于设置与 动态、APL、RESAMPLE 相关的 RedHawk 运行参数。（可选；默认：0）

\textbf{语法：}
\begin{lstlisting}
DYNAMIC_APL_RESAMPLE [ 0 | 1 ]
\end{lstlisting}

\subsubsection{\gsr{FAST_IMPORT_APL_MODE}}
\paragraph*{FAST_IMPORT_APL_MODE}

GSR 关键字 FAST_IMPORT_APL_MODE：用于设置与 FAST、导入、APL、模式 相关的 RedHawk 运行参数。（可选；默认：0）

\textbf{语法：}
\begin{lstlisting}
FAST_IMPORT_APL_MODE [ 0 | 1 ]
\end{lstlisting}

\subsubsection{\gsr{IGNORE_APL_CHECK}}
\paragraph*{IGNORE_APL_CHECK}

在数据输入或分析流程中忽略与 APL、检查 相关的对象或检查。（可选；默认：0 (do not ignore)；兼容 DMP）

\textbf{语法：}
\begin{lstlisting}
IGNORE_APL_CHECK [ 0 | 1 ]
\end{lstlisting}

\subsubsection{\gsr{IGNORE_APL_CHECK_SWITCH}}
\paragraph*{IGNORE_APL_CHECK_SWITCH}

在数据输入或分析流程中忽略与 APL、检查、SWITCH 相关的对象或检查。（可选；默认：0 (do not ignore)；兼容 DMP）

\textbf{语法：}
\begin{lstlisting}
IGNORE_APL_CHECK_SWITCH [ 0 | 1 ]
\end{lstlisting}

\subsubsection{\gsr{IGNORE_APL_PROCESS_CORNER}}
\paragraph*{IGNORE_APL_PROCESS_CORNER}

在数据输入或分析流程中忽略与 APL、PROCESS、CORNER 相关的对象或检查。（可选；默认：0 (off)；兼容 DMP） 注意：在数据输入或分析流程中忽略与 APL、PROCESS、CORNER 相关的对象或检查。

\textbf{语法：}
\begin{lstlisting}
Syntax
IGNORE_APL_PROCESS_CORNER [0|1]
may create incorrect RedHawk results.
\end{lstlisting}

\begin{noteBox}
在数据输入或分析流程中忽略与 APL、PROCESS、CORNER 相关的对象或检查。
\end{noteBox}

\subsubsection{\gsr{IGNORE_MULTISTATE_ERROR}}
\paragraph*{IGNORE_MULTISTATE_ERROR}

在数据输入或分析流程中忽略与 MULTISTATE、错误 相关的对象或检查。（可选；默认：0；兼容 DMP）

\textbf{语法：}
\begin{lstlisting}
Syntax
IGNORE_MULTISTATE_ERROR [1 | 0]
\end{lstlisting}

\subsubsection{\gsr{IP_MODEL_CELL_MAP}}
\paragraph*{IP_MODEL_CELL_MAP}

GSR 关键字 IP_MODEL_CELL_MAP：用于设置与 IP、MODEL、单元、MAP 相关的 RedHawk 运行参数。（可选；默认：view of the model is to be used for all；兼容 DMP）

\textbf{语法：}
\begin{lstlisting}
Syntax
IP_MODEL_CELL_MAP {
<cell name> <default view name>
}
\end{lstlisting}

\subsubsection{\gsr{IP_MODEL_INST_MAP}}
\paragraph*{IP_MODEL_INST_MAP}

GSR 关键字 IP_MODEL_INST_MAP：用于设置与 IP、MODEL、INST、MAP 相关的 RedHawk 运行参数。（可选；默认：view, to expand the usage of keyword IP_；兼容 DMP）

\textbf{语法：}
\begin{lstlisting}
Syntax
IP_MODEL_INST_MAP {
<inst_name> <cellname> <specified_instance_view>
...
}
\end{lstlisting}

\subsubsection{\gsr{IP_MODELS}}
\paragraph*{IP_MODELS}

GSR 关键字 IP_MODELS：用于设置与 IP、MODELS 相关的 RedHawk 运行参数。（可选；默认：None；兼容 DMP）

\textbf{语法：}
\begin{lstlisting}
Syntax
IP_MODELS {
<IP name> <IP_model_directory>
}
\end{lstlisting}

\subsubsection{\gsr{LIB2AVM}}
\paragraph*{LIB2AVM}

GSR 关键字 LIB2AVM：用于设置与 LIB2AVM 相关的 RedHawk 运行参数。（可选；默认：(1), a triangular-shaped profile is prov；兼容 DMP）

\textbf{语法：}
\begin{lstlisting}
Syntax
LIB2AVM [ 0 | 1 | 2 | 3 | 4]
\end{lstlisting}

\textbf{示例：}
\begin{lstlisting}
LIB2AVM 0
\end{lstlisting}

\subsubsection{\gsr{PS_IN_FLOW_APL}}
\paragraph*{PS_IN_FLOW_APL}

GSR 关键字 PS_IN_FLOW_APL：用于设置与 PS、IN、FLOW、APL 相关的 RedHawk 运行参数。（可选；默认：1(On)；兼容 DMP）

\textbf{语法：}
\begin{lstlisting}
Syntax
\end{lstlisting}

\subsection{时序关键字}
\subsubsection*{Timing Keywords}

本类共 21 个 GSR 关键字。以下逐条给出中文用途、语法、默认值与示例。

\subsubsection{\gsr{ATE_CONSTRAINT_FILES}}
\paragraph*{ATE_CONSTRAINT_FILES}

GSR 关键字 ATE_CONSTRAINT_FILES：用于设置与 ATE、CONSTRAINT、文件 相关的 RedHawk 运行参数。（可选；默认：none；兼容 DMP）

\textbf{语法：}
\begin{lstlisting}
ATE_CONSTRAINT_FILES {
<SDC_file1>
<SDC_file2>
...
USER_OPTIONS_FILE <TCL_file_name>
FREQ_OF_MISSING_INSTANCES <freq-Hertz>
}
\end{lstlisting}

\textbf{参数说明：}
\begin{itemize}
\item \texttt{<TCL\_file\_name>} — GSR 关键字 ATE\_CONSTRAINT\_FILES：用于设置与 ATE、CONSTRAINT、文件 相关的 RedHawk 运行参数。
\end{itemize}

\subsubsection{\gsr{ATE_USE_REDHAWK_DB}}
\paragraph*{ATE_USE_REDHAWK_DB}

GSR 关键字 ATE_USE_REDHAWK_DB：用于设置与 ATE、USE、REDHAWK、DB 相关的 RedHawk 运行参数。（可选；默认：1）

\textbf{语法：}
\begin{lstlisting}
ATE_USE_REDHAWK_DB [0|1]
\end{lstlisting}

\subsubsection{\gsr{BOUND_SLEW_TO_MAX_TRANSITION}}
\paragraph*{BOUND_SLEW_TO_MAX_TRANSITION}

GSR 关键字 BOUND_SLEW_TO_MAX_TRANSITION：用于设置与 BOUND、SLEW、TO、MAX、TRANSITION 相关的 RedHawk 运行参数。（可选；默认：1 (On)；兼容 DMP）

\textbf{语法：}
\begin{lstlisting}
BOUND_SLEW_TO_MAX_TRANSITION [ 0 | 1 ]
\end{lstlisting}

\subsubsection{\gsr{CLOCK_ROOTS}}
\paragraph*{CLOCK_ROOTS}

GSR 关键字 CLOCK_ROOTS：用于设置与 CLOCK、ROOTS 相关的 RedHawk 运行参数。（可选；默认：None；兼容 DMP）

\subsubsection{\gsr{CLOCK_ROOTS}}
\paragraph*{CLOCK_ROOTS}

GSR 关键字 CLOCK_ROOTS：用于设置与 CLOCK、ROOTS 相关的 RedHawk 运行参数。

\subsubsection{\gsr{CLOCK_ROOTS}}
\paragraph*{CLOCK_ROOTS}

GSR 关键字 CLOCK_ROOTS：用于设置与 CLOCK、ROOTS 相关的 RedHawk 运行参数。

\subsubsection{\gsr{ENABLE_ATE}}
\paragraph*{ENABLE_ATE}

启用 ATE 相关功能或检查。（可选；默认：RedHawk spef parser) and hence provide t；兼容 DMP）

\textbf{语法：}
\begin{lstlisting}
ENABLE_ATE [ 0 | 1 ]
\end{lstlisting}

\subsubsection{\gsr{INACTIVE_NETS}}
\paragraph*{INACTIVE_NETS}

GSR 关键字 INACTIVE_NETS：用于设置与 INACTIVE、网络 相关的 RedHawk 运行参数。（可选；默认：None；兼容 DMP）

\subsubsection{\gsr{INACTIVE_NETS}}
\paragraph*{INACTIVE_NETS}

GSR 关键字 INACTIVE_NETS：用于设置与 INACTIVE、网络 相关的 RedHawk 运行参数。

\subsubsection{\gsr{INACTIVE_NETS}}
\paragraph*{INACTIVE_NETS}

GSR 关键字 INACTIVE_NETS：用于设置与 INACTIVE、网络 相关的 RedHawk 运行参数。

\subsubsection{\gsr{INPUT_TRANSITION}}
\paragraph*{INPUT_TRANSITION}

GSR 关键字 INPUT_TRANSITION：用于设置与 INPUT、TRANSITION 相关的 RedHawk 运行参数。（可选；默认：10% of the value of the inverse of the f；兼容 DMP）

\textbf{语法：}
\begin{lstlisting}
INPUT_TRANSITION <value in seconds>
\end{lstlisting}

\textbf{示例：}
\begin{lstlisting}
INPUT_TRANSITION 0.2e-9 (or 0.2ns)
\end{lstlisting}

\subsubsection{\gsr{JITTER_ENABLE}}
\paragraph*{JITTER_ENABLE}

GSR 关键字 JITTER_ENABLE：用于设置与 JITTER、启用 相关的 RedHawk 运行参数。（可选；默认：0）

\textbf{语法：}
\begin{lstlisting}
JITTER_ENABLE [ 0 | 1 ]
\end{lstlisting}

\subsubsection{\gsr{PSI_SPICE_CELL_NETLIST_FILE}}
\paragraph*{PSI_SPICE_CELL_NETLIST_FILE}

GSR 关键字 PSI_SPICE_CELL_NETLIST_FILE：用于设置与 PSI、SPICE、单元、NETLIST、文件 相关的 RedHawk 运行参数。

\subsubsection{\gsr{PSI_SPICE_CELL_NETLIST_FILE}}
\paragraph*{PSI_SPICE_CELL_NETLIST_FILE}

GSR 关键字 PSI_SPICE_CELL_NETLIST_FILE：用于设置与 PSI、SPICE、单元、NETLIST、文件 相关的 RedHawk 运行参数。

\subsubsection{\gsr{PSI_SPICE_CELL_NETLIST_FILE}}
\paragraph*{PSI_SPICE_CELL_NETLIST_FILE}

GSR 关键字 PSI_SPICE_CELL_NETLIST_FILE：用于设置与 PSI、SPICE、单元、NETLIST、文件 相关的 RedHawk 运行参数。（可选；默认：s: minimum - 5e-12, maximum - 1e-9；兼容 DMP）

\textbf{语法：}
\begin{lstlisting}
SLEW_MIN_TRANSITION <min_rise_sec>
SLEW_MAX_TRANSITION <max_rise_sec>
\end{lstlisting}

\subsubsection{\gsr{SPARAM_CHECK_LOWEST_FREQ}}
\paragraph*{SPARAM_CHECK_LOWEST_FREQ}

GSR 关键字 SPARAM_CHECK_LOWEST_FREQ：用于设置与 SPARAM、检查、LOWEST、FREQ 相关的 RedHawk 运行参数。（默认：is off）

\subsubsection{\gsr{SPARAM_CHECK_REFERENCE_R}}
\paragraph*{SPARAM_CHECK_REFERENCE_R}

GSR 关键字 SPARAM_CHECK_REFERENCE_R：用于设置与 SPARAM、检查、REFERENCE、R 相关的 RedHawk 运行参数。（默认：is off）

\textbf{语法：}
\begin{lstlisting}
SPARAM_CHECK_REFERENCE_R [0|1]
\end{lstlisting}

\subsubsection{\gsr{SPARAM_EXACT_DC}}
\paragraph*{SPARAM_EXACT_DC}

GSR 关键字 SPARAM_EXACT_DC：用于设置与 SPARAM、EXACT、DC 相关的 RedHawk 运行参数。（默认：off）

\textbf{语法：}
\begin{lstlisting}
SPARAM_EXACT_DC [0|1]
\end{lstlisting}

\subsubsection{\gsr{SPARAM_HANDLING}}
\paragraph*{SPARAM_HANDLING}

GSR 关键字 SPARAM_HANDLING：用于设置与 SPARAM、HANDLING 相关的 RedHawk 运行参数。（可选；默认：) 1 - enables limited passivity enforcem；兼容 DMP）

\textbf{语法：}
\begin{lstlisting}
SPARAM_HANDLING [ 0 | 1 | 2 ]
\end{lstlisting}

\subsubsection{\gsr{STA_SLEW_SCALING}}
\paragraph*{STA_SLEW_SCALING}

GSR 关键字 STA_SLEW_SCALING：用于设置与 STA、SLEW、SCALING 相关的 RedHawk 运行参数。（默认：1）

\textbf{语法：}
\begin{lstlisting}
Syntax
STA_SLEW_SCALING <scaling factor>
\end{lstlisting}

\subsubsection{\gsr{USE_CCS_PIN_CAPS}}
\paragraph*{USE_CCS_PIN_CAPS}

GSR 关键字 USE_CCS_PIN_CAPS：用于设置与 USE、CCS、引脚、CAPS 相关的 RedHawk 运行参数。（默认：0）

\textbf{语法：}
\begin{lstlisting}
Syntax
USE_CCS_PIN_CAPS [ 0 | 1 ]
\end{lstlisting}

\subsection{仿真关键字}
\subsubsection*{Simulation Keywords}

本类共 59 个 GSR 关键字。以下逐条给出中文用途、语法、默认值与示例。

\subsubsection{\gsr{BLOCK_PAR}}
\paragraph*{BLOCK_PAR}

GSR 关键字 BLOCK_PAR：用于设置与 模块、PAR 相关的 RedHawk 运行参数。（可选；默认：1；兼容 DMP）

\subsubsection{\gsr{BLOCK_PAR}}
\paragraph*{BLOCK_PAR}

GSR 关键字 BLOCK_PAR：用于设置与 模块、PAR 相关的 RedHawk 运行参数。

\subsubsection{\gsr{BLOCK_PAR}}
\paragraph*{BLOCK_PAR}

GSR 关键字 BLOCK_PAR：用于设置与 模块、PAR 相关的 RedHawk 运行参数。

\subsubsection{\gsr{CONNECTIVITY_RES_THRESHOLD}}
\paragraph*{CONNECTIVITY_RES_THRESHOLD}

GSR 关键字 CONNECTIVITY_RES_THRESHOLD：用于设置与 CONNECTIVITY、RES、THRESHOLD 相关的 RedHawk 运行参数。（可选；默认：1e6）

\textbf{语法：}
\begin{lstlisting}
CONNECTIVITY_RES_THRESHOLD <threshold>
\end{lstlisting}

\subsubsection{\gsr{CONSISTENT_SCENARIO}}
\paragraph*{CONSISTENT_SCENARIO}

GSR 关键字 CONSISTENT_SCENARIO：用于设置与 CONSISTENT、SCENARIO 相关的 RedHawk 运行参数。（可选；默认：), this keyword insures that the same or；兼容 DMP）

\textbf{语法：}
\begin{lstlisting}
CONSISTENT_SCENARIO [ 0 | 1 ]
\end{lstlisting}

\subsubsection{\gsr{DVD_GLITCH_FILTER}}
\paragraph*{DVD_GLITCH_FILTER}

GSR 关键字 DVD_GLITCH_FILTER：用于设置与 DvD、GLITCH、FILTER 相关的 RedHawk 运行参数。（可选；默认：none；兼容 DMP） 注意：GSR 关键字 DVD_GLITCH_FILTER：用于设置与 DvD、GLITCH、FILTER 相关的 RedHawk 运行参数。

\textbf{语法：}
\begin{lstlisting}
DVD_GLITCH_FILTER {
VOLTAGE_LEVEL <Meas_V>
[ GLITCH_AREA <minA_psv> | GLITCH_WIDTH <minW_sec>]
CELL_LIST_FILE <filename>
}
\end{lstlisting}

\textbf{参数说明：}
\begin{itemize}
\item \texttt{VOLTAGE\_LEVEL <Meas\_V>} — GSR 关键字 DVD\_GLITCH\_FILTER：用于设置与 DvD、GLITCH、FILTER 相关的 RedHawk 运行参数。
\end{itemize}

\textbf{示例：}
\begin{lstlisting}
DVD_GLITCH_FILTER {
VOLTAGE_LEVEL 1.15
GLITCH_WIDTH 150ps
CELL_LIST_FILE DvD_glitch
}
\end{lstlisting}

\begin{noteBox}
GSR 关键字 DVD\_GLITCH\_FILTER：用于设置与 DvD、GLITCH、FILTER 相关的 RedHawk 运行参数。
\end{noteBox}

\subsubsection{\gsr{DYNAMIC_64BIT_SOLVER}}
\paragraph*{DYNAMIC_64BIT_SOLVER}

GSR 关键字 DYNAMIC_64BIT_SOLVER：用于设置与 动态、64BIT、SOLVER 相关的 RedHawk 运行参数。（可选；默认：0 (off)；兼容 DMP）

\textbf{语法：}
\begin{lstlisting}
DYNAMIC_64BIT_SOLVER [0 |1 ]
\end{lstlisting}

\subsubsection{\gsr{DYNAMIC_ARC_LOW_POWER}}
\paragraph*{DYNAMIC_ARC_LOW_POWER}

GSR 关键字 DYNAMIC_ARC_LOW_POWER：用于设置与 动态、ARC、LOW、电源 相关的 RedHawk 运行参数。（可选；默认：0）

\textbf{语法：}
\begin{lstlisting}
DYNAMIC_ARC_LOW_POWER [0 |1 ]
\end{lstlisting}

\subsubsection{\gsr{DYNAMIC_BYPASS_SHORT}}
\paragraph*{DYNAMIC_BYPASS_SHORT}

GSR 关键字 DYNAMIC_BYPASS_SHORT：用于设置与 动态、BYPASS、短路 相关的 RedHawk 运行参数。（可选；默认：0 (no bypass)）

\textbf{语法：}
\begin{lstlisting}
DYNAMIC_BYPASS_SHORT [0 |1 ]
\end{lstlisting}

\subsubsection{\gsr{DYNAMIC_CELL_CROSS_CHECK}}
\paragraph*{DYNAMIC_CELL_CROSS_CHECK}

GSR 关键字 DYNAMIC_CELL_CROSS_CHECK：用于设置与 动态、单元、CROSS、检查 相关的 RedHawk 运行参数。（默认：simulation errors out with an “inconsist）

\textbf{语法：}
\begin{lstlisting}
DYNAMIC_CELL_CROSS_CHECK [0 |1 ]
\end{lstlisting}

\subsubsection{\gsr{DYNAMIC_CLOCK_SCALE}}
\paragraph*{DYNAMIC_CLOCK_SCALE}

GSR 关键字 DYNAMIC_CLOCK_SCALE：用于设置与 动态、CLOCK、缩放 相关的 RedHawk 运行参数。（可选；默认：0 (no scaling)；兼容 DMP）

\textbf{语法：}
\begin{lstlisting}
DYNAMIC_CLOCK_SCALE [ 0| 1 | 2 ]
\end{lstlisting}

\subsubsection{\gsr{DYNAMIC_DETACHED_POSTPROCESS}}
\paragraph*{DYNAMIC_DETACHED_POSTPROCESS}

GSR 关键字 DYNAMIC_DETACHED_POSTPROCESS：用于设置与 动态、DETACHED、POSTPROCESS 相关的 RedHawk 运行参数。（可选；默认：0）

\textbf{语法：}
\begin{lstlisting}
DYNAMIC_DETACHED_POSTPROCESS [ 0| 1 ]
\end{lstlisting}

\subsubsection{\gsr{DYNAMIC_MBFF_TW_MODE}}
\paragraph*{DYNAMIC_MBFF_TW_MODE}

GSR 关键字 DYNAMIC_MBFF_TW_MODE：用于设置与 动态、MBFF、TW、模式 相关的 RedHawk 运行参数。（可选；默认：0）

\textbf{语法：}
\begin{lstlisting}
DYNAMIC_MBFF_TW_MODE[ 0| 1 ]
\end{lstlisting}

\subsubsection{\gsr{DYNAMIC_PRESIM_VCD}}
\paragraph*{DYNAMIC_PRESIM_VCD}

GSR 关键字 DYNAMIC_PRESIM_VCD：用于设置与 动态、PRESIM、VCD 相关的 RedHawk 运行参数。（默认：1）

\textbf{语法：}
\begin{lstlisting}
DYNAMIC_PRESIM_VCD [ 0 | 1 ]
\end{lstlisting}

\subsubsection{\gsr{DYNAMIC_PRESIM_DCINIT_SCALE}}
\paragraph*{DYNAMIC_PRESIM_DCINIT_SCALE}

GSR 关键字 DYNAMIC_PRESIM_DCINIT_SCALE：用于设置与 动态、PRESIM、DCINIT、缩放 相关的 RedHawk 运行参数。（可选；默认：0）

\textbf{语法：}
\begin{lstlisting}
DYNAMIC_PRESIM_DCINIT_SCALE <scale_value>
\end{lstlisting}

\subsubsection{\gsr{DYNAMIC_DISABLE_NEW_WFEXTRACT}}
\paragraph*{DYNAMIC_DISABLE_NEW_WFEXTRACT}

GSR 关键字 DYNAMIC_DISABLE_NEW_WFEXTRACT：用于设置与 动态、禁用、NEW、WFEXTRACT 相关的 RedHawk 运行参数。（可选；默认：the Vdd files with waveforms are split i；兼容 DMP）

\textbf{语法：}
\begin{lstlisting}
DYNAMIC_DISABLE_NEW_WFEXTRACT [0 | 1]
\end{lstlisting}

\subsubsection{\gsr{DYNAMIC_EXTEND_VCD}}
\paragraph*{DYNAMIC_EXTEND_VCD}

GSR 关键字 DYNAMIC_EXTEND_VCD：用于设置与 动态、EXTEND、VCD 相关的 RedHawk 运行参数。（可选；默认：if the presimulation period is not fully；兼容 DMP）

\textbf{语法：}
\begin{lstlisting}
DYNAMIC_EXTEND_VCD [0 | 1 | 2]
\end{lstlisting}

\subsubsection{\gsr{DYNAMIC_FF_ADJUST_NTRIG}}
\paragraph*{DYNAMIC_FF_ADJUST_NTRIG}

GSR 关键字 DYNAMIC_FF_ADJUST_NTRIG：用于设置与 动态、FF、ADJUST、NTRIG 相关的 RedHawk 运行参数。（可选；默认：1）

\subsubsection{\gsr{DYNAMIC_FRAME_SIZE}}
\paragraph*{DYNAMIC_FRAME_SIZE}

GSR 关键字 DYNAMIC_FRAME_SIZE：用于设置与 动态、FRAME、SIZE 相关的 RedHawk 运行参数。（默认：1/FREQ；兼容 DMP）

\textbf{语法：}
\begin{lstlisting}
DYNAMIC_FRAME_SIZE [ <frame_size_sec> | -1 ]
\end{lstlisting}

\subsubsection{\gsr{DYNAMIC_FREQUENCY_AWARE}}
\paragraph*{DYNAMIC_FREQUENCY_AWARE}

GSR 关键字 DYNAMIC_FREQUENCY_AWARE：用于设置与 动态、频率、AWARE 相关的 RedHawk 运行参数。（默认：none）

\textbf{语法：}
\begin{lstlisting}
DYNAMIC_FREQUENCY_AWARE              <freq_Hz>
\end{lstlisting}

\textbf{示例：}
\begin{lstlisting}
DYNAMIC_FREQUENCY_AWARE                300e6
In the example, the Ipwr is modulated at 300 MHz.
\end{lstlisting}

\subsubsection{\gsr{DYNAMIC_GROUP_WIRECAP}}
\paragraph*{DYNAMIC_GROUP_WIRECAP}

GSR 关键字 DYNAMIC_GROUP_WIRECAP：用于设置与 动态、GROUP、WIRECAP 相关的 RedHawk 运行参数。（默认：0 (off)）

\textbf{语法：}
\begin{lstlisting}
DYNAMIC_GROUP_WIRECAP [1 | 0 ]
\end{lstlisting}

\subsubsection{\gsr{DYNAMIC_MBFF_MODE}}
\paragraph*{DYNAMIC_MBFF_MODE}

GSR 关键字 DYNAMIC_MBFF_MODE：用于设置与 动态、MBFF、模式 相关的 RedHawk 运行参数。（默认：off）

\textbf{语法：}
\begin{lstlisting}
DYNAMIC_MBFF_MODE [ 0 | 1 | 2]
\end{lstlisting}

\subsubsection{\gsr{DYNAMIC_MCYC_TW}}
\paragraph*{DYNAMIC_MCYC_TW}

GSR 关键字 DYNAMIC_MCYC_TW：用于设置与 动态、MCYC、TW 相关的 RedHawk 运行参数。（可选；默认：0 (Off)；兼容 DMP）

\textbf{语法：}
\begin{lstlisting}
DYNAMIC_MCYC_TW [ 0 | 1 ]
\end{lstlisting}

\subsubsection{\gsr{DYNAMIC_MIXED_CONSISTENT_SCENARIO}}
\paragraph*{DYNAMIC_MIXED_CONSISTENT_SCENARIO}

GSR 关键字 DYNAMIC_MIXED_CONSISTENT_SCENARIO：用于设置与 动态、MIXED、CONSISTENT、SCENARIO 相关的 RedHawk 运行参数。（可选；默认：0；兼容 DMP）

\textbf{语法：}
\begin{lstlisting}
DYNAMIC_MIXED_CONSISTENT_SCENARIO [0 | 1]
\end{lstlisting}

\subsubsection{\gsr{DYNAMIC_MSTATE_FILTER}}
\paragraph*{DYNAMIC_MSTATE_FILTER}

GSR 关键字 DYNAMIC_MSTATE_FILTER：用于设置与 动态、MSTATE、FILTER 相关的 RedHawk 运行参数。（可选；默认：1 (on)；兼容 DMP）

\textbf{语法：}
\begin{lstlisting}
DYNAMIC_MSTATE_FILTER [0 | 1]
\end{lstlisting}

\subsubsection{\gsr{DYNAMIC_PEAK_CURRENT_AWARE}}
\paragraph*{DYNAMIC_PEAK_CURRENT_AWARE}

GSR 关键字 DYNAMIC_PEAK_CURRENT_AWARE：用于设置与 动态、峰值、CURRENT、AWARE 相关的 RedHawk 运行参数。（可选；默认：none；兼容 DMP）

\textbf{语法：}
\begin{lstlisting}
DYNAMIC_PEAK_CURRENT_AWARE <cell_list>
\end{lstlisting}

\subsubsection{\gsr{DYNAMIC_PGARC_REPORT_BASE}}
\paragraph*{DYNAMIC_PGARC_REPORT_BASE}

GSR 关键字 DYNAMIC_PGARC_REPORT_BASE：用于设置与 动态、PGARC、REPORT、BASE 相关的 RedHawk 运行参数。（可选；默认：avgTW)；兼容 DMP）

\textbf{语法：}
\begin{lstlisting}
DYNAMIC_PGARC_REPORT_BASE [avgTW | maxTW | minTW | minWC]
\end{lstlisting}

\textbf{示例：}
\begin{lstlisting}
DYNAMIC_PGARC_REPORT_BASE minTW
\end{lstlisting}

\subsubsection{\gsr{DYNAMIC_POST_BATCH}}
\paragraph*{DYNAMIC_POST_BATCH}

GSR 关键字 DYNAMIC_POST_BATCH：用于设置与 动态、POST、BATCH 相关的 RedHawk 运行参数。（可选；默认：RedHawk selects the appropriate number o；必需）

\textbf{语法：}
\begin{lstlisting}
DYNAMIC_POST_BATCH <num>
\end{lstlisting}

\subsubsection{\gsr{DYNAMIC_PRECHECK}}
\paragraph*{DYNAMIC_PRECHECK}

GSR 关键字 DYNAMIC_PRECHECK：用于设置与 动态、PRECHECK 相关的 RedHawk 运行参数。（可选；默认：On）

\textbf{语法：}
\begin{lstlisting}
DYNAMIC_PRECHECK [ 0| 1]
\end{lstlisting}

\subsubsection{\gsr{DYNAMIC_PRESIM_TIME}}
\paragraph*{DYNAMIC_PRESIM_TIME}

GSR 关键字 DYNAMIC_PRESIM_TIME：用于设置与 动态、PRESIM、TIME 相关的 RedHawk 运行参数。（可选；默认：presim time set by RedHawk；兼容 DMP）

\textbf{语法：}
\begin{lstlisting}
DYNAMIC_PRESIM_TIME [<presim_time_sec>| -1] ?<TSM> ?
?<TSM_fraction>?
\end{lstlisting}

\textbf{示例：}
\begin{lstlisting}
DYNAMIC_PRESIM_TIME 10e-09 5 0.9
In this example the presim time is set at 10e-09 seconds, the TSM is 5 and
TSM_fraction is 0.9. This specifies a time step 5 times the length of the specified
time step for the first 90% of the presim time, and a regular time step for the
remaining 10% of presimulation.
\end{lstlisting}

\subsubsection{\gsr{DYNAMIC_RELAX_CONTROL_PIN_CONSTRAINT}}
\paragraph*{DYNAMIC_RELAX_CONTROL_PIN_CONSTRAINT}

GSR 关键字 DYNAMIC_RELAX_CONTROL_PIN_CONSTRAINT：用于设置与 动态、RELAX、CONTROL、引脚、CONSTRAINT 相关的 RedHawk 运行参数。（可选；默认：Off, meaning that simulation ignores the）

\textbf{语法：}
\begin{lstlisting}
DYNAMIC_RELAX_CONTROL_PIN_CONSTRAINT [ 0 | 1 ]
\end{lstlisting}

\subsubsection{\gsr{DYNAMIC_REPORT_CLOCK_EVDD}}
\paragraph*{DYNAMIC_REPORT_CLOCK_EVDD}

GSR 关键字 DYNAMIC_REPORT_CLOCK_EVDD：用于设置与 动态、REPORT、CLOCK、EVDD 相关的 RedHawk 运行参数。（可选；默认：flow, all the effective Vdd values are c）

\textbf{语法：}
\begin{lstlisting}
DYNAMIC_REPORT_CLOCK_EVDD [ 0 | 1 ]
\end{lstlisting}

\subsubsection{\gsr{DYNAMIC_REPORT_DECAP}}
\paragraph*{DYNAMIC_REPORT_DECAP}

GSR 关键字 DYNAMIC_REPORT_DECAP：用于设置与 动态、REPORT、去耦电容 相关的 RedHawk 运行参数。（可选；默认：0 (off)；兼容 DMP）

\textbf{语法：}
\begin{lstlisting}
DYNAMIC_REPORT_DECAP [ 0 | 1 | 2]
\end{lstlisting}

\textbf{示例：}
\begin{lstlisting}
DYNAMIC_REPORT_DECAP 1
\end{lstlisting}

\subsubsection{\gsr{DYNAMIC_REPORT_DVD}}
\paragraph*{DYNAMIC_REPORT_DVD}

GSR 关键字 DYNAMIC_REPORT_DVD：用于设置与 动态、REPORT、DvD 相关的 RedHawk 运行参数。（可选；默认：, for basic vectorless analysis, RedHawk）

\textbf{语法：}
\begin{lstlisting}
DYNAMIC_REPORT_DVD [ 0 | 1 | 2 ]
\end{lstlisting}

\subsubsection{\gsr{DYNAMIC_REPORT_PRESIM_STATE}}
\paragraph*{DYNAMIC_REPORT_PRESIM_STATE}

GSR 关键字 DYNAMIC_REPORT_PRESIM_STATE：用于设置与 动态、REPORT、PRESIM、STATE 相关的 RedHawk 运行参数。（默认：0）

\textbf{语法：}
\begin{lstlisting}
DYNAMIC_REPORT_PRESIM_STATE [ 1 | 0 ]
\end{lstlisting}

\subsubsection{\gsr{DYNAMIC_SAVE_WAVEFORM}}
\paragraph*{DYNAMIC_SAVE_WAVEFORM}

GSR 关键字 DYNAMIC_SAVE_WAVEFORM：用于设置与 动态、SAVE、WAVEFORM 相关的 RedHawk 运行参数。（可选；默认：1）

\textbf{语法：}
\begin{lstlisting}
DYNAMIC_SAVE_WAVEFORM [ 0 | 1 | 2 ]
\end{lstlisting}

\subsubsection{\gsr{DYNAMIC_SELECTIVE_SAVE}}
\paragraph*{DYNAMIC_SELECTIVE_SAVE}

GSR 关键字 DYNAMIC_SELECTIVE_SAVE：用于设置与 动态、SELECTIVE、SAVE 相关的 RedHawk 运行参数。（可选；默认：(1) saves only waveforms for pad and pin）

\textbf{语法：}
\begin{lstlisting}
DYNAMIC_SELECTIVE_SAVE [1 | 0 ]
\end{lstlisting}

\subsubsection{\gsr{DYNAMIC_SIMULATION_TIME}}
\paragraph*{DYNAMIC_SIMULATION_TIME}

GSR 关键字 DYNAMIC_SIMULATION_TIME：用于设置与 动态、仿真、TIME 相关的 RedHawk 运行参数。（可选；默认：1/Frequency；兼容 DMP）

\textbf{语法：}
\begin{lstlisting}
DYNAMIC_SIMULATION_TIME ?<start_time_sec>? <end_time_sec>
\end{lstlisting}

\textbf{示例：}
\begin{lstlisting}
DYNAMIC_SIMULATION_TIME 3200e-12
is equal to
DYNAMIC_SIMULATION_TIME 0 3200e-12
\end{lstlisting}

\subsubsection{\gsr{DYNAMIC_SOLVER_MODE}}
\paragraph*{DYNAMIC_SOLVER_MODE}

GSR 关键字 DYNAMIC_SOLVER_MODE：用于设置与 动态、SOLVER、模式 相关的 RedHawk 运行参数。（可选；默认：an assumption of power-ground symmetry i；兼容 DMP）

\textbf{语法：}
\begin{lstlisting}
DYNAMIC_SOLVER_MODE [ 0 | 1]
\end{lstlisting}

\subsubsection{\gsr{DYNAMIC_SORT_BY_PERCENTAGE}}
\paragraph*{DYNAMIC_SORT_BY_PERCENTAGE}

GSR 关键字 DYNAMIC_SORT_BY_PERCENTAGE：用于设置与 动态、SORT、BY、PERCENTAGE 相关的 RedHawk 运行参数。（可选；默认：0 (off)）

\textbf{语法：}
\begin{lstlisting}
DYNAMIC_SORT_BY_PERCENTAGE [ 0 | 1]
\end{lstlisting}

\subsubsection{\gsr{DYNAMIC_TIME_STEP}}
\paragraph*{DYNAMIC_TIME_STEP}

GSR 关键字 DYNAMIC_TIME_STEP：用于设置与 动态、TIME、STEP 相关的 RedHawk 运行参数。（可选；默认：10 psec；兼容 DMP）

\textbf{语法：}
\begin{lstlisting}
DYNAMIC_TIME_STEP <time in sec>
\end{lstlisting}

\textbf{示例：}
\begin{lstlisting}
DYNAMIC_TIME_STEP 20e-12
\end{lstlisting}

\subsubsection{\gsr{DYNAMIC_USE_ZLIB}}
\paragraph*{DYNAMIC_USE_ZLIB}

GSR 关键字 DYNAMIC_USE_ZLIB：用于设置与 动态、USE、ZLIB 相关的 RedHawk 运行参数。（默认：1）

\textbf{语法：}
\begin{lstlisting}
DYNAMIC_USE_ZLIB          [ 0 | 1 ]
\end{lstlisting}

\subsubsection{\gsr{DYNAMIC_VOLTAGE_CHECK}}
\paragraph*{DYNAMIC_VOLTAGE_CHECK}

GSR 关键字 DYNAMIC_VOLTAGE_CHECK：用于设置与 动态、VOLTAGE、检查 相关的 RedHawk 运行参数。（可选；默认：, voltage checking for node voltages gre；兼容 DMP）

\textbf{语法：}
\begin{lstlisting}
DYNAMIC_VOLTAGE_CHECK [ 0 | 1]
\end{lstlisting}

\subsubsection{\gsr{EFFECTIVE_VDD_WINDOW}}
\paragraph*{EFFECTIVE_VDD_WINDOW}

GSR 关键字 EFFECTIVE_VDD_WINDOW：用于设置与 EFFECTIVE、VDD、WINDOW 相关的 RedHawk 运行参数。（可选；默认：the sliding window period starts when th）

\textbf{语法：}
\begin{lstlisting}
Syntax
EFFECTIVE_VDD_WINDOW {
? IN2OUT <inputV_fraction> <outputV_fraction> ?
? IN2IN <inputV_fraction1> <inputV_fraction2> ?
? MINIMUM <minWidth_ps> ?
}
\end{lstlisting}

\textbf{示例：}
\begin{lstlisting}
EFFECTIVE_VDD_WINDOW {
IN2OUT 0.2 0.7
MINIMUM 10
}
In the above example, the sliding window is the time interval between when the
instance’s input signal is at the 20% level to the time the instance’s output signal is
at the 70% level, with a minimum sliding window duration of 10ps.
\end{lstlisting}

\subsubsection{\gsr{GROUND_CURRENT_DISTRIBUTION}}
\paragraph*{GROUND_CURRENT_DISTRIBUTION}

GSR 关键字 GROUND_CURRENT_DISTRIBUTION：用于设置与 GROUND、CURRENT、DISTRIBUTION 相关的 RedHawk 运行参数。（可选；默认：0；兼容 DMP）

\textbf{语法：}
\begin{lstlisting}
Syntax
GROUND_CURRENT_DISTRIBUTION [ 0 | 1 ]
\end{lstlisting}

\subsubsection{\gsr{HIERARCHY_CONSISTENT_SCENARIO}}
\paragraph*{HIERARCHY_CONSISTENT_SCENARIO}

GSR 关键字 HIERARCHY_CONSISTENT_SCENARIO：用于设置与 HIERARCHY、CONSISTENT、SCENARIO 相关的 RedHawk 运行参数。（可选；默认：none）

\textbf{语法：}
\begin{lstlisting}
Syntax
HIERARCHY_CONSISTENT_SCENARIO <config_file>
\end{lstlisting}

\subsubsection{\gsr{ITERATIVE_SOLVER}}
\paragraph*{ITERATIVE_SOLVER}

GSR 关键字 ITERATIVE_SOLVER：用于设置与 ITERATIVE、SOLVER 相关的 RedHawk 运行参数。（可选；默认：-1；兼容 DMP） 注意：GSR 关键字 ITERATIVE_SOLVER：用于设置与 ITERATIVE、SOLVER 相关的 RedHawk 运行参数。

\textbf{语法：}
\begin{lstlisting}
Syntax
ITERATIVE_SOLVER [ -1 | 0 | 1 ]
\end{lstlisting}

\begin{noteBox}
GSR 关键字 ITERATIVE\_SOLVER：用于设置与 ITERATIVE、SOLVER 相关的 RedHawk 运行参数。
\end{noteBox}

\subsubsection{\gsr{MMX_ADAPTIVE_SAMPLING}}
\paragraph*{MMX_ADAPTIVE_SAMPLING}

GSR 关键字 MMX_ADAPTIVE_SAMPLING：用于设置与 MMX、ADAPTIVE、SAMPLING 相关的 RedHawk 运行参数。（可选；默认：On）

\textbf{语法：}
\begin{lstlisting}
Syntax
MMX_ADAPTIVE_SAMPLING [ 0 | 1 ]
\end{lstlisting}

\subsubsection{\gsr{MULTI_CYCLE_SCENARIO}}
\paragraph*{MULTI_CYCLE_SCENARIO}

GSR 关键字 MULTI_CYCLE_SCENARIO：用于设置与 MULTI、CYCLE、SCENARIO 相关的 RedHawk 运行参数。（可选；默认：Off；兼容 DMP）

\textbf{语法：}
\begin{lstlisting}
Syntax
MULTI_CYCLE_SCENARIO [ 0 | 1 ]
\end{lstlisting}

\subsubsection{\gsr{NX_SIM}}
\paragraph*{NX_SIM}

GSR 关键字 NX_SIM：用于设置与 NX、SIM 相关的 RedHawk 运行参数。（可选；默认：0；兼容 DMP） 注意：GSR 关键字 NX_SIM：用于设置与 NX、SIM 相关的 RedHawk 运行参数。

\textbf{语法：}
\begin{lstlisting}
NX_SIM [ 0 | 1 | 2 | 3]
\end{lstlisting}

\begin{noteBox}
GSR 关键字 NX\_SIM：用于设置与 NX、SIM 相关的 RedHawk 运行参数。
\end{noteBox}

\subsubsection{\gsr{NX_VECTORLESS}}
\paragraph*{NX_VECTORLESS}

GSR 关键字 NX_VECTORLESS：用于设置与 NX、VECTORLESS 相关的 RedHawk 运行参数。（可选；默认：On, as shown in Figure C-2；兼容 DMP）

\textbf{语法：}
\begin{lstlisting}
NX_VECTORLESS [ 0 | 1 ]
Figure C-2       Vectorless scenario generation
\end{lstlisting}

\subsubsection{\gsr{PIECEWISE_SWITCH_INPUT}}
\paragraph*{PIECEWISE_SWITCH_INPUT}

GSR 关键字 PIECEWISE_SWITCH_INPUT：用于设置与 PIECEWISE、SWITCH、INPUT 相关的 RedHawk 运行参数。

\textbf{语法：}
\begin{lstlisting}
PIECEWISE_SWITCH_INPUT <file_name>
\end{lstlisting}

\textbf{参数说明：}
\begin{itemize}
\item \texttt{Optional; default} — GSR 关键字 PIECEWISE\_SWITCH\_INPUT：用于设置与 PIECEWISE、SWITCH、INPUT 相关的 RedHawk 运行参数。
\end{itemize}

\subsubsection{\gsr{PLOC_SHORT_INTERNAL_NET}}
\paragraph*{PLOC_SHORT_INTERNAL_NET}

GSR 关键字 PLOC_SHORT_INTERNAL_NET：用于设置与 焊盘位置、短路、INTERNAL、网络 相关的 RedHawk 运行参数。（可选；默认：none；兼容 DMP；必需）

\textbf{语法：}
\begin{lstlisting}
PLOC_SHORT_INTERNAL_NET {
<Net1> ? <R1> <L1> <C1> ?
<Net2> ? <R2> <L2> <C2> ?
...
}
\end{lstlisting}

\subsubsection{\gsr{POWERUP_RANDOM_TOGGLE}}
\paragraph*{POWERUP_RANDOM_TOGGLE}

GSR 关键字 POWERUP_RANDOM_TOGGLE：用于设置与 POWERUP、RANDOM、翻转 相关的 RedHawk 运行参数。（可选；默认：0ff (no switching)；兼容 DMP）

\textbf{语法：}
\begin{lstlisting}
Syntax
POWERUP_RANDOM_TOGGLE [0|1]
\end{lstlisting}

\subsubsection{\gsr{PROBE_NODE_FILE}}
\paragraph*{PROBE_NODE_FILE}

GSR 关键字 PROBE_NODE_FILE：用于设置与 PROBE、节点、文件 相关的 RedHawk 运行参数。（可选；默认：none；兼容 DMP）

\textbf{语法：}
\begin{lstlisting}
PROBE_NODE_FILE <path_to_file>
\end{lstlisting}

\subsubsection{\gsr{SCAN_LAUNCH_CAPTURE_MODE}}
\paragraph*{SCAN_LAUNCH_CAPTURE_MODE}

GSR 关键字 SCAN_LAUNCH_CAPTURE_MODE：用于设置与 SCAN、LAUNCH、CAPTURE、模式 相关的 RedHawk 运行参数。（可选；默认：the shift procedure shifts test patterns；兼容 DMP）

\textbf{语法：}
\begin{lstlisting}
SCAN_LAUNCH_CAPTURE_MODE [ 0 | 1 ]
\end{lstlisting}

\subsubsection{\gsr{SIMULATION_CACHE_DIRECTORY}}
\paragraph*{SIMULATION_CACHE_DIRECTORY}

GSR 关键字 SIMULATION_CACHE_DIRECTORY：用于设置与 仿真、CACHE、DIRECTORY 相关的 RedHawk 运行参数。（可选；默认：none）

\textbf{语法：}
\begin{lstlisting}
SIMULATION_CACHE_DIRECTORY <sim_dir_path>
\end{lstlisting}

\subsubsection{\gsr{RECOMMENDED_SIMULATION_FACTORS}}
\paragraph*{RECOMMENDED_SIMULATION_FACTORS}

GSR 关键字 RECOMMENDED_SIMULATION_FACTORS：用于设置与 RECOMMENDED、仿真、FACTORS 相关的 RedHawk 运行参数。

\textbf{语法：}
\begin{lstlisting}
RECOMMENDED_SIMULATION_FACTORS {
SIMULATION_TIME_FACTOR <scale_factor1>
FRAME_SIZE_FACTOR <scale_factor2>
}
\end{lstlisting}

\textbf{示例：}
\begin{lstlisting}
RECOMMENDED_SIMULATION_FACTORS {
SIMULATION_TIME_FACTOR 4.0
FRAME_SIZE_FACTOR      2.0
}
and redhawk.log contents:
Recommended dynamic simulation time, 2560psec ,to include
97.0370% of total power for DYNAMIC_SIMULATION_TIME in GSR.
Here, DYNAMIC_SIMULATION_TIME is set to 10240ps (that is 2560ps x 4.0) and
DYNAMIC_FRAME_SIZE is set to 5120ps (2560ps x 2.0), where recommended
dynamic simulation time is 2560psec from redhawk.log file.
\end{lstlisting}

\subsubsection{\gsr{STATIC_IR_HIGHLIGHT_PART_CON_INST}}
\paragraph*{STATIC_IR_HIGHLIGHT_PART_CON_INST}

GSR 关键字 STATIC_IR_HIGHLIGHT_PART_CON_INST：用于设置与 静态、IR、HIGHLIGHT、PART、CON、INST 相关的 RedHawk 运行参数。（默认：0）

\textbf{语法：}
\begin{lstlisting}
STATIC_IR_HIGHLIGHT_PART_CON_INST [1 | 0]
\end{lstlisting}

\subsection{DMP 分布式关键字}
\subsubsection*{DMP Keywords}

本类共 5 个 GSR 关键字。以下逐条给出中文用途、语法、默认值与示例。

\subsubsection{\gsr{DMP_EP_FAST_MODE}}
\paragraph*{DMP_EP_FAST_MODE}

GSR 关键字 DMP_EP_FAST_MODE：用于设置与 DMP、EP、FAST、模式 相关的 RedHawk 运行参数。（可选；默认：, the number of shared events is the num）

\textbf{语法：}
\begin{lstlisting}
DMP_EP_FAST_MODE [ 0 | 1 ]
\end{lstlisting}

\subsubsection{\gsr{DMP_SETUP_BOUNDARY}}
\paragraph*{DMP_SETUP_BOUNDARY}

GSR 关键字 DMP_SETUP_BOUNDARY：用于设置与 DMP、SETUP、BOUNDARY 相关的 RedHawk 运行参数。

\textbf{语法：}
\begin{lstlisting}
DMP_SETUP_BOUNDARY {
bound1_x
bound2_x
...
}
\end{lstlisting}

\textbf{示例：}
\begin{lstlisting}
The below settings in GSR define the boundaries for a 4-way DMP run
DMP_SETUP_BOUNDARY {
286.440002
586.440002
886.440002
}
For cases where FDR is not able to create balanced partitioning (primarily due to
via merging resulting in non-uniform merging during flattening design), the tool will
dump out the coordinates that user can use for manual boundary setup in next runs
in the log. A sample below:
dmpBalanceInfo: if this run has node balance problem, please
set gsr KW,DMP_SETUP_BOUNDARY {
1770
3310
}
\end{lstlisting}

\subsubsection{\gsr{DMP_SIM_KEEP_LAYERS}}
\paragraph*{DMP_SIM_KEEP_LAYERS}

GSR 关键字 DMP_SIM_KEEP_LAYERS：用于设置与 DMP、SIM、KEEP、金属层 相关的 RedHawk 运行参数。（默认：, DMP will do RC network reduction on al）

\textbf{语法：}
\begin{lstlisting}
DMP_SIM_KEEP_LAYERS {
<LAYER_NAME> VDD(VSS)_NET_NAME
}
\end{lstlisting}

\textbf{示例：}
\begin{lstlisting}
DMP_SIM_KEEP_LAYERS {
RDL VDD18
METAL9 VDD18
RDL VSS
METAL9 VSS
}
"ALL" can be used either in layer name or net name, such as:
DMP_SIM_KEEP_LAYERS {
METAL9 ALL
}
\end{lstlisting}

\subsubsection{\gsr{DMP_SP_FAST_MODE}}
\paragraph*{DMP_SP_FAST_MODE}

GSR 关键字 DMP_SP_FAST_MODE：用于设置与 DMP、SP、FAST、模式 相关的 RedHawk 运行参数。（可选；默认：0）

\textbf{语法：}
\begin{lstlisting}
DMP_SP_FAST_MODE [ 0 | 1 ]
\end{lstlisting}

\subsubsection{\gsr{PS_DMP_PERFORMANCE_MODE}}
\paragraph*{PS_DMP_PERFORMANCE_MODE}

GSR 关键字 PS_DMP_PERFORMANCE_MODE：用于设置与 PS、DMP、PERFORMANCE、模式 相关的 RedHawk 运行参数。（默认：0）

\textbf{语法：}
\begin{lstlisting}
PS_DMP_PERFORMANCE_MODE [ 0 | 1 ]
\end{lstlisting}

\subsection{FAO 通用关键字}
\subsubsection*{FAO General Keywords}

本类共 12 个 GSR 关键字。以下逐条给出中文用途、语法、默认值与示例。

\subsubsection{\gsr{FAO_ACCURATE_VOLTAGE}}
\paragraph*{FAO_ACCURATE_VOLTAGE}

GSR 关键字 FAO_ACCURATE_VOLTAGE：用于设置与 FAO、ACCURATE、VOLTAGE 相关的 RedHawk 运行参数。（可选；默认：, performs accurate voltage difference c）

\textbf{语法：}
\begin{lstlisting}
FAO_ACCURATE_VOLTAGE[ 0 | 1 ]
\end{lstlisting}

\subsubsection{\gsr{FAO_ADD_STACK_VIA}}
\paragraph*{FAO_ADD_STACK_VIA}

GSR 关键字 FAO_ADD_STACK_VIA：用于设置与 FAO、添加、STACK、过孔 相关的 RedHawk 运行参数。（可选；默认：), turns on stack via insertion for the）

\textbf{语法：}
\begin{lstlisting}
Syntax
FAO_ADD_STACK_VIA [0|1]
\end{lstlisting}

\subsubsection{\gsr{FAO_BYPASS_CUTSIZE_CHECK}}
\paragraph*{FAO_BYPASS_CUTSIZE_CHECK}

GSR 关键字 FAO_BYPASS_CUTSIZE_CHECK：用于设置与 FAO、BYPASS、CUTSIZE、检查 相关的 RedHawk 运行参数。（可选；默认：On）

\textbf{语法：}
\begin{lstlisting}
Syntax
FAO_BYPASS_CUTSIZE_CHECK [ 0 | 1 ]
\end{lstlisting}

\subsubsection{\gsr{FAO_HOLD_LIC}}
\paragraph*{FAO_HOLD_LIC}

GSR 关键字 FAO_HOLD_LIC：用于设置与 FAO、HOLD、LIC 相关的 RedHawk 运行参数。（可选；默认：0)）

\textbf{语法：}
\begin{lstlisting}
Syntax
FAO_HOLD_LIC [ 0 | 1 ]
\end{lstlisting}

\subsubsection{\gsr{FAO_IGNORE_COMPRESS_VIAS}}
\paragraph*{FAO_IGNORE_COMPRESS_VIAS}

GSR 关键字 FAO_IGNORE_COMPRESS_VIAS：用于设置与 FAO、忽略、COMPRESS、VIAS 相关的 RedHawk 运行参数。（可选；默认：0）

\textbf{语法：}
\begin{lstlisting}
FAO_IGNORE_COMPRESS_VIAS [0|1]
\end{lstlisting}

\subsubsection{\gsr{FAO_MISVIA_ONE_FILE}}
\paragraph*{FAO_MISVIA_ONE_FILE}

GSR 关键字 FAO_MISVIA_ONE_FILE：用于设置与 FAO、MISVIA、ONE、文件 相关的 RedHawk 运行参数。（可选；默认：, lists are combined into one file；兼容 DMP）

\textbf{语法：}
\begin{lstlisting}
Syntax
FAO_MISVIA_ONE_FILE [ 0 | 1 ]
\end{lstlisting}

\subsubsection{\gsr{FAO_OBJ}}
\paragraph*{FAO_OBJ}

GSR 关键字 FAO_OBJ：用于设置与 FAO、OBJ 相关的 RedHawk 运行参数。（可选；默认：DVD） 注意：GSR 关键字 FAO_OBJ：用于设置与 FAO、OBJ 相关的 RedHawk 运行参数。

\textbf{语法：}
\begin{lstlisting}
Syntax
FAO_OBJ [ DVD | TIMING ?<slack threshold>? |
PATH ?<slack threshold>? ]
\end{lstlisting}

\textbf{参数说明：}
\begin{itemize}
\item \texttt{DVD} — GSR 关键字 FAO\_OBJ：用于设置与 FAO、OBJ 相关的 RedHawk 运行参数。（默认：value of the slack threshold is 0）
\end{itemize}

\textbf{示例：}
\begin{lstlisting}
FAO_OBJ DVD
\end{lstlisting}

\begin{noteBox}
GSR 关键字 FAO\_OBJ：用于设置与 FAO、OBJ 相关的 RedHawk 运行参数。
\end{noteBox}

\subsubsection{\gsr{FAO_REGION}}
\paragraph*{FAO_REGION}

GSR 关键字 FAO_REGION：用于设置与 FAO、区域 相关的 RedHawk 运行参数。（可选；默认：whole chip；兼容 DMP）

\textbf{语法：}
\begin{lstlisting}
fao_region          { <x1> <y1> ... <xn> <yn>}
\end{lstlisting}

\textbf{参数说明：}
GSR 关键字 FAO\_REGION：用于设置与 FAO、区域 相关的 RedHawk 运行参数。

\textbf{示例：}
\begin{lstlisting}
FAO_REGION {140 140 140 220 220 220 220 200 160 200 160 140 }
\end{lstlisting}

\subsubsection{\gsr{FAO_TURBO_MODE}}
\paragraph*{FAO_TURBO_MODE}

GSR 关键字 FAO_TURBO_MODE：用于设置与 FAO、TURBO、模式 相关的 RedHawk 运行参数。（可选；默认：), or for the most accurate mode, which）

\textbf{语法：}
\begin{lstlisting}
fao_turbo_mode [ 1 | 0 ]
\end{lstlisting}

\subsubsection{\gsr{FIX_WINDOW}}
\paragraph*{FIX_WINDOW}

GSR 关键字 FIX_WINDOW：用于设置与 FIX、WINDOW 相关的 RedHawk 运行参数。（可选；默认：100um x 100um）

\textbf{语法：}
\begin{lstlisting}
Syntax
fix_window          { x_dim y_dim }
\end{lstlisting}

\textbf{示例：}
\begin{lstlisting}
fix_window { 400 400 }
\end{lstlisting}

\subsubsection{\gsr{NOISE_LIMIT}}
\paragraph*{NOISE_LIMIT}

GSR 关键字 NOISE_LIMIT：用于设置与 NOISE、LIMIT 相关的 RedHawk 运行参数。（可选；默认：0；必需）

\textbf{语法：}
\begin{lstlisting}
noise_limit <percent_voltage_drop>
\end{lstlisting}

\textbf{示例：}
\begin{lstlisting}
noise_limit 5
\end{lstlisting}

\subsubsection{\gsr{NOISE_REDUCTION}}
\paragraph*{NOISE_REDUCTION}

GSR 关键字 NOISE_REDUCTION：用于设置与 NOISE、REDUCTION 相关的 RedHawk 运行参数。

\textbf{语法：}
\begin{lstlisting}
noise_reduction <percentage reduction>
\end{lstlisting}

\textbf{示例：}
\begin{lstlisting}
example, for an existing voltage drop of 100mV, to get an improvement (reduction)
of 20mV requires setting this GSR keyword to 20. Use with commands: mesh
optimize, mesh fix, decap optimize. Optional; default: 0.
noise_reduction 10
which targets a 10% reduction in present voltage drop
noise_reduction -10
which targets a 10% relaxation (increase) in present voltage drop
\end{lstlisting}

\subsection{网格修复与优化关键字}
\subsubsection*{Grid Fixing and Optimization Keywords}

本类共 16 个 GSR 关键字。以下逐条给出中文用途、语法、默认值与示例。

\subsubsection{\gsr{FAO_DRC_DROP_RATIO}}
\paragraph*{FAO_DRC_DROP_RATIO}

GSR 关键字 FAO_DRC_DROP_RATIO：用于设置与 FAO、DRC、DROP、RATIO 相关的 RedHawk 运行参数。（可选；默认：1 (100% of max width)）

\textbf{语法：}
\begin{lstlisting}
fao_drc_drop_ratio <fraction_of_max_width>
\end{lstlisting}

\textbf{示例：}
\begin{lstlisting}
FAO_DRC_DROP_RATIO 0.8
\end{lstlisting}

\subsubsection{\gsr{FAO_DVD_TYPE}}
\paragraph*{FAO_DVD_TYPE}

GSR 关键字 FAO_DVD_TYPE：用于设置与 FAO、DvD、TYPE 相关的 RedHawk 运行参数。（可选；默认：1）

\textbf{语法：}
\begin{lstlisting}
FAO_DVD_TYPE <integer 1-4>
\end{lstlisting}

\textbf{示例：}
\begin{lstlisting}
FAO_DVD_TYPE 2
\end{lstlisting}

\subsubsection{\gsr{FAO_DYNAMIC_MODE}}
\paragraph*{FAO_DYNAMIC_MODE}

GSR 关键字 FAO_DYNAMIC_MODE：用于设置与 FAO、动态、模式 相关的 RedHawk 运行参数。（可选；默认：0 (off)）

\textbf{语法：}
\begin{lstlisting}
FAO_DYNAMIC_MODE [ 0 | 1 ]
\end{lstlisting}

\subsubsection{\gsr{FAO_LAYERS}}
\paragraph*{FAO_LAYERS}

GSR 关键字 FAO_LAYERS：用于设置与 FAO、金属层 相关的 RedHawk 运行参数。（必需） 注意：GSR 关键字 FAO_LAYERS：用于设置与 FAO、金属层 相关的 RedHawk 运行参数。（必需）

\textbf{语法：}
\begin{lstlisting}
FAO_LAYERS { { <layer name> ?L ?<direction>?
?<wire_size_constraints> }...} ?
\end{lstlisting}

\textbf{参数说明：}
\begin{itemize}
\item \texttt{<layer name>} — GSR 关键字 FAO\_LAYERS：用于设置与 FAO、金属层 相关的 RedHawk 运行参数。（默认：, width- based constraint is available）
\end{itemize}

\textbf{示例：}
\begin{lstlisting}
FAO_LAYERS { { met4 ver > 1 < 3 } { met5 hor = 12 }}
With Length based constraint
FAO_LAYERS { { met4 ver > 1 < 3 } { met5 L hor = 12 }}
\end{lstlisting}

\begin{noteBox}
GSR 关键字 FAO\_LAYERS：用于设置与 FAO、金属层 相关的 RedHawk 运行参数。（必需）
\end{noteBox}

\subsubsection{\gsr{FAO_MISVIA_DISTANCE}}
\paragraph*{FAO_MISVIA_DISTANCE}

GSR 关键字 FAO_MISVIA_DISTANCE：用于设置与 FAO、MISVIA、DISTANCE 相关的 RedHawk 运行参数。（可选；默认：, all missing vias are reported；兼容 DMP）

\textbf{语法：}
\begin{lstlisting}
FAO_MISVIA_DISTANCE <value_um>
\end{lstlisting}

\subsubsection{\gsr{FAO_MISVIA_RPT_SHORT}}
\paragraph*{FAO_MISVIA_RPT_SHORT}

GSR 关键字 FAO_MISVIA_RPT_SHORT：用于设置与 FAO、MISVIA、RPT、短路 相关的 RedHawk 运行参数。（可选；默认：0 (off)）

\textbf{语法：}
\begin{lstlisting}
FAO_MISVIA_RPT_SHORT [ 1 | 0 ]
\end{lstlisting}

\subsubsection{\gsr{FAO_NETS}}
\paragraph*{FAO_NETS}

GSR 关键字 FAO_NETS：用于设置与 FAO、网络 相关的 RedHawk 运行参数。（必需）

\textbf{语法：}
\begin{lstlisting}
FAO_NETS [ <net1>| {
<net1>
<net2>
...
} ]
\end{lstlisting}

\textbf{示例：}
\begin{lstlisting}
FAO_NETS VDD
FAO_NETS {
\end{lstlisting}

\subsubsection{\gsr{VDD2}}
\paragraph*{VDD2}

GSR 关键字 VDD2：用于设置与 VDD2 相关的 RedHawk 运行参数。

\subsubsection{\gsr{VSS2}}
\paragraph*{VSS2}

GSR 关键字 VSS2：用于设置与 VSS2 相关的 RedHawk 运行参数。

\subsubsection{\gsr{FAO_PRECISE_VOLTAGE}}
\paragraph*{FAO_PRECISE_VOLTAGE}

GSR 关键字 FAO_PRECISE_VOLTAGE：用于设置与 FAO、PRECISE、VOLTAGE 相关的 RedHawk 运行参数。（可选；默认：1）

\textbf{语法：}
\begin{lstlisting}
Syntax
FAO_PRECISE_VOLTAGE [ 0 | 1]
\end{lstlisting}

\subsubsection{\gsr{FAO_RANGE}}
\paragraph*{FAO_RANGE}

GSR 关键字 FAO_RANGE：用于设置与 FAO、RANGE 相关的 RedHawk 运行参数。（必需） 注意：GSR 关键字 FAO_RANGE：用于设置与 FAO、RANGE 相关的 RedHawk 运行参数。（必需）

\textbf{语法：}
\begin{lstlisting}
FAO_RANGE { { <layer> <min_width> <max_width>
? <extend_dir> ?} ... }
\end{lstlisting}

\textbf{参数说明：}
\begin{itemize}
\item \texttt{<extend\_dir>} — GSR 关键字 FAO\_RANGE：用于设置与 FAO、RANGE 相关的 RedHawk 运行参数。
\end{itemize}

\textbf{示例：}
\begin{lstlisting}
FAO_RANGE { { met5 10 30 bot} { met6 15 40 } }
\end{lstlisting}

\begin{noteBox}
GSR 关键字 FAO\_RANGE：用于设置与 FAO、RANGE 相关的 RedHawk 运行参数。（必需）
\end{noteBox}

\subsubsection{\gsr{FAO_ROW_VDD_SITE}}
\paragraph*{FAO_ROW_VDD_SITE}

GSR 关键字 FAO_ROW_VDD_SITE：用于设置与 FAO、ROW、VDD、SITE 相关的 RedHawk 运行参数。

\textbf{语法：}
\begin{lstlisting}
FAO_ROW_VDD_SITE {
<power_domain_name> <sitename1> <sitename2> ...
}
\end{lstlisting}

\textbf{示例：}
\begin{lstlisting}
FAO_ROW_VDD_SITE {
VDD1 CORE
VDD2 CORE2
}
In the above example, DEF has two ROW sites defined: CORE that belongs to net
VDD1 and CORE2 that belongs to VDD2. FAO adds decaps only in the domain(s)
and sites defined by the keyword. This is currently supported only for the FAO
command 'decap fill'.
\end{lstlisting}

\subsubsection{\gsr{FAO_SUB_GRID_NETS}}
\paragraph*{FAO_SUB_GRID_NETS}

GSR 关键字 FAO_SUB_GRID_NETS：用于设置与 FAO、SUB、GRID、网络 相关的 RedHawk 运行参数。（可选；默认：none）

\textbf{语法：}
\begin{lstlisting}
fao_sub_grid_nets { <list of power/ground nets> }
\end{lstlisting}

\textbf{参数说明：}
GSR 关键字 FAO\_SUB\_GRID\_NETS：用于设置与 FAO、SUB、GRID、网络 相关的 RedHawk 运行参数。

\textbf{示例：}
\begin{lstlisting}
fao_sub_grid_nets { VDD VSS },
\end{lstlisting}

\subsubsection{\gsr{FAO_SUB_GRID_SPEC}}
\paragraph*{FAO_SUB_GRID_SPEC}

GSR 关键字 FAO_SUB_GRID_SPEC：用于设置与 FAO、SUB、GRID、SPEC 相关的 RedHawk 运行参数。（可选；默认：none Global System Requirements File (*） 注意：GSR 关键字 FAO_SUB_GRID_SPEC：用于设置与 FAO、SUB、GRID、SPEC 相关的 RedHawk 运行参数。（必需）

\textbf{语法：}
\begin{lstlisting}
FAO_SUB_GRID_SPEC { { <layer> <dir> <space> <pitch>
<top_via_layer> <bot_via_layer> } ... }
\end{lstlisting}

\textbf{参数说明：}
GSR 关键字 FAO\_SUB\_GRID\_SPEC：用于设置与 FAO、SUB、GRID、SPEC 相关的 RedHawk 运行参数。

\begin{noteBox}
GSR 关键字 FAO\_SUB\_GRID\_SPEC：用于设置与 FAO、SUB、GRID、SPEC 相关的 RedHawk 运行参数。（必需）
\end{noteBox}

\subsubsection{\gsr{FAO_WIDTH_CNSTR}}
\paragraph*{FAO_WIDTH_CNSTR}

GSR 关键字 FAO_WIDTH_CNSTR：用于设置与 FAO、WIDTH、CNSTR 相关的 RedHawk 运行参数。（可选；默认：no specification） 注意：GSR 关键字 FAO_WIDTH_CNSTR：用于设置与 FAO、WIDTH、CNSTR 相关的 RedHawk 运行参数。

\textbf{语法：}
\begin{lstlisting}
fao_width_cnstr { { layer1*<value> -/+ layer2*<value>
>/</>=/<= <value> } ...}
\end{lstlisting}

\textbf{示例：}
\begin{lstlisting}
fao_width_cnstr { { met4 - met3*2 > 0 }
{ met5 - met4*3 > 10 } }
required.
\end{lstlisting}

\begin{noteBox}
GSR 关键字 FAO\_WIDTH\_CNSTR：用于设置与 FAO、WIDTH、CNSTR 相关的 RedHawk 运行参数。
\end{noteBox}

\subsubsection{\gsr{NUM_HOTSPOT}}
\paragraph*{NUM_HOTSPOT}

GSR 关键字 NUM_HOTSPOT：用于设置与 NUM、HOTSPOT 相关的 RedHawk 运行参数。（可选；默认：1）

\textbf{语法：}
\begin{lstlisting}
num_hotspot <int>
\end{lstlisting}

\textbf{示例：}
\begin{lstlisting}
num_hotspot 5
\end{lstlisting}

\subsection{去耦电容优化关键字}
\subsubsection*{Decap Optimization Keywords}

本类共 17 个 GSR 关键字。以下逐条给出中文用途、语法、默认值与示例。

\subsubsection{\gsr{CAP_LIMIT}}
\paragraph*{CAP_LIMIT}

GSR 关键字 CAP_LIMIT：用于设置与 CAP、LIMIT 相关的 RedHawk 运行参数。（可选；默认：unlimited）

\textbf{语法：}
\begin{lstlisting}
cap_limit <cap_pF>
\end{lstlisting}

\textbf{示例：}
\begin{lstlisting}
cap_limit 1000
which defines 1000 pf as the maximum added decap.
\end{lstlisting}

\subsubsection{\gsr{DECAP_CELL}}
\paragraph*{DECAP_CELL}

GSR 关键字 DECAP_CELL：用于设置与 去耦电容、单元 相关的 RedHawk 运行参数。（可选；默认：none；兼容 DMP；必需）

\textbf{语法：}
\begin{lstlisting}
DECAP_CELL {
<decap_cellname> ? <width_um>? ?<height_um>? ?<C_pF>?
?<R_ohm>? ?<metal_layer_name>? ?<leakage_A>?
...
}
\end{lstlisting}

\textbf{示例：}
\begin{lstlisting}
DECAP_CELL {
decap1
decapAB 0.25 0.75 0.01 1000 MET1 1e-9
decapXY 0.75 0.75 0.03 1500 MET1 2e-9
}
\end{lstlisting}

\subsubsection{\gsr{DECAP_CELL_FILES}}
\paragraph*{DECAP_CELL_FILES}

允许 importing multiple 文件s that contain a list of 去耦电容 单元. 兼容 DMP 分布式流程。 可选; 默认: None.（可选；默认：None；兼容 DMP）

\textbf{语法：}
\begin{lstlisting}
DECAP_CELL_FILES {
<decap_file1>
...
}
\end{lstlisting}

\textbf{示例：}
\begin{lstlisting}
DECAP_CELL_FILES {
user_data/demo1.dec
user_data/demo2.dec
}
\end{lstlisting}

\subsubsection{\gsr{DECAP_DENSITY}}
\paragraph*{DECAP_DENSITY}

GSR 关键字 DECAP_DENSITY：用于设置与 去耦电容、DENSITY 相关的 RedHawk 运行参数。（可选；默认：1）

\textbf{语法：}
\begin{lstlisting}
decap_density <integer multiple>
\end{lstlisting}

\textbf{示例：}
\begin{lstlisting}
decap_density 3
\end{lstlisting}

\subsubsection{\gsr{DECAP_TILE_MAX}}
\paragraph*{DECAP_TILE_MAX}

GSR 关键字 DECAP_TILE_MAX：用于设置与 去耦电容、TILE、MAX 相关的 RedHawk 运行参数。（可选；默认：0 (off)）

\textbf{语法：}
\begin{lstlisting}
DECAP_TILE_MAX <value_pF>
\end{lstlisting}

\textbf{示例：}
\begin{lstlisting}
DECAP_TILE_MAX 100
\end{lstlisting}

\subsubsection{\gsr{FAO_DECAP_FILL_ALG}}
\paragraph*{FAO_DECAP_FILL_ALG}

GSR 关键字 FAO_DECAP_FILL_ALG：用于设置与 FAO、去耦电容、FILL、ALG 相关的 RedHawk 运行参数。（可选；默认：0 (off)）

\textbf{语法：}
\begin{lstlisting}
FAO_DECAP_FILL_ALG [ 0 | 1 | 2]
\end{lstlisting}

\subsubsection{\gsr{FAO_DECAP_FILL_NEW_FLOW}}
\paragraph*{FAO_DECAP_FILL_NEW_FLOW}

GSR 关键字 FAO_DECAP_FILL_NEW_FLOW：用于设置与 FAO、去耦电容、FILL、NEW、FLOW 相关的 RedHawk 运行参数。（可选；默认：0 (off)）

\textbf{语法：}
\begin{lstlisting}
FAO_DECAP_FILL_NEW_FLOW [ 0 | 1 ]
\end{lstlisting}

\subsubsection{\gsr{FAO_DECAP_OVERLAP}}
\paragraph*{FAO_DECAP_OVERLAP}

GSR 关键字 FAO_DECAP_OVERLAP：用于设置与 FAO、去耦电容、OVERLAP 相关的 RedHawk 运行参数。（可选；默认：false）

\textbf{语法：}
\begin{lstlisting}
fao_decap_overlap [true/1 | false/0]
\end{lstlisting}

\textbf{示例：}
\begin{lstlisting}
fao_decap_overlap true
\end{lstlisting}

\subsubsection{\gsr{FAO_DRC_PL_OBS}}
\paragraph*{FAO_DRC_PL_OBS}

GSR 关键字 FAO_DRC_PL_OBS：用于设置与 FAO、DRC、PL、OBS 相关的 RedHawk 运行参数。（可选；默认：1 (on)）

\textbf{语法：}
\begin{lstlisting}
fao_drc_pl_obs [0 | 1]
\end{lstlisting}

\subsubsection{\gsr{FAO_DRC_OBS}}
\paragraph*{FAO_DRC_OBS}

GSR 关键字 FAO_DRC_OBS：用于设置与 FAO、DRC、OBS 相关的 RedHawk 运行参数。（可选；默认：0 (off) Global System Requirements File）

\subsubsection{\gsr{FAO_ECO_NAME}}
\paragraph*{FAO_ECO_NAME}

GSR 关键字 FAO_ECO_NAME：用于设置与 FAO、ECO、NAME 相关的 RedHawk 运行参数。（可选；默认：instance name pattern for decap cells）

\textbf{语法：}
\begin{lstlisting}
fao_eco_name <decap_cell>_<eco_string><number>
\end{lstlisting}

\textbf{示例：}
\begin{lstlisting}
fao_eco_name decapABC_XY3
\end{lstlisting}

\subsubsection{\gsr{FAO_MAX_SHIFT}}
\paragraph*{FAO_MAX_SHIFT}

GSR 关键字 FAO_MAX_SHIFT：用于设置与 FAO、MAX、SHIFT 相关的 RedHawk 运行参数。（可选；默认：5 um）

\textbf{语法：}
\begin{lstlisting}
fao_max_shift <max dist to move cells>
\end{lstlisting}

\textbf{示例：}
\begin{lstlisting}
fao_max_shift 3
\end{lstlisting}

\subsubsection{\gsr{FAO_PLACE_GRID}}
\paragraph*{FAO_PLACE_GRID}

GSR 关键字 FAO_PLACE_GRID：用于设置与 FAO、PLACE、GRID 相关的 RedHawk 运行参数。（可选；默认：none）

\textbf{语法：}
\begin{lstlisting}
fao_place_grid <grid_adjustment>
\end{lstlisting}

\textbf{示例：}
\begin{lstlisting}
fao_place_grid 0.2
\end{lstlisting}

\subsubsection{\gsr{FAO_ROW_SITE}}
\paragraph*{FAO_ROW_SITE}

GSR 关键字 FAO_ROW_SITE：用于设置与 FAO、ROW、SITE 相关的 RedHawk 运行参数。（可选；默认：, without using FAO_ROW_SITE, FAO takes）

\textbf{语法：}
\begin{lstlisting}
FAO_ROW_SITE {
<LEF_site_name>
...
}
\end{lstlisting}

\textbf{示例：}
\begin{lstlisting}
For a DEF file ROW section containing several row definitions as follows:
ROW CORE_ROW_37549 bonuscoreaxg 3547600 2680600 N DO 7793 BY 1 STEP 690 0 ;
ROW CORE_ROW_52855 core 8927530 2680600 N DO 2054 BY 1 STEP 230 0 ;
ROW CORE_ROW_72819 bonuscore 8927530 2680600 N DO 684 BY 1 STEP 690 0 ;
Using ‘FAO_ROW_SITE core’, only rows defined by the ‘core’ site name are
created for FAO decap placement.
\end{lstlisting}

\subsubsection{\gsr{FAO_WIRE_WIDTH_LOW_LIMIT}}
\paragraph*{FAO_WIRE_WIDTH_LOW_LIMIT}

GSR 关键字 FAO_WIRE_WIDTH_LOW_LIMIT：用于设置与 FAO、WIRE、WIDTH、LOW、LIMIT 相关的 RedHawk 运行参数。（可选；默认：0）

\textbf{语法：}
\begin{lstlisting}
FAO_WIRE_WIDTH_LOW_LIMIT <tolerance_u>
\end{lstlisting}

\subsubsection{\gsr{LEAKAGE_LIMIT}}
\paragraph*{LEAKAGE_LIMIT}

GSR 关键字 LEAKAGE_LIMIT：用于设置与 LEAKAGE、LIMIT 相关的 RedHawk 运行参数。（可选；默认：infinity）

\textbf{语法：}
\begin{lstlisting}
leakage_limit <max decap power leakage>
\end{lstlisting}

\textbf{示例：}
\begin{lstlisting}
leakage_limit 10e-6
\end{lstlisting}

\subsubsection{\gsr{NUM_HOTINST}}
\paragraph*{NUM_HOTINST}

GSR 关键字 NUM_HOTINST：用于设置与 NUM、HOTINST 相关的 RedHawk 运行参数。（可选；默认：1）

\textbf{语法：}
\begin{lstlisting}
num_hotinst <integer>
\end{lstlisting}

\textbf{示例：}
\begin{lstlisting}
num_hotinst 3
\end{lstlisting}

\subsection{低功耗设计关键字}
\subsubsection*{Low Power Design Keywords}

本类共 22 个 GSR 关键字。以下逐条给出中文用途、语法、默认值与示例。

\subsubsection{\gsr{BLOCK_POWERUP_FILES}}
\paragraph*{BLOCK_POWERUP_FILES}

GSR 关键字 BLOCK_POWERUP_FILES：用于设置与 模块、POWERUP、文件 相关的 RedHawk 运行参数。（可选；默认：time offset =0；兼容 DMP）

\subsubsection{\gsr{BLOCK_POWERUP_FILES}}
\paragraph*{BLOCK_POWERUP_FILES}

GSR 关键字 BLOCK_POWERUP_FILES：用于设置与 模块、POWERUP、文件 相关的 RedHawk 运行参数。

\textbf{示例：}
\begin{lstlisting}
BLOCK_POWERUP_FILES block_powerup.file
\end{lstlisting}

\subsubsection{\gsr{CHARGE_SWITCH}}
\paragraph*{CHARGE_SWITCH}

GSR 关键字 CHARGE_SWITCH：用于设置与 CHARGE、SWITCH 相关的 RedHawk 运行参数。（可选；默认：none）

\textbf{语法：}
\begin{lstlisting}
CHARGE_SWITCH {
RAMPUP_THRESHOLD <Volts>
<cell_name1> ? <Volts>?
<cell_name2> <charge_control_pin> <funct_control pin> ?<Volts>?
...
}
\end{lstlisting}

\textbf{示例：}
\begin{lstlisting}
CHARGE_SWITCH {
RAMPUP_THRESHOLD 0.8
Charge_Sw_AB 0.7
Charge_Sw_BCD CCpin3 FCpin7
}
\end{lstlisting}

\subsubsection{\gsr{CHECK_SWITCH_POWERON}}
\paragraph*{CHECK_SWITCH_POWERON}

GSR 关键字 CHECK_SWITCH_POWERON：用于设置与 检查、SWITCH、POWERON 相关的 RedHawk 运行参数。（可选；默认：None）

\textbf{语法：}
\begin{lstlisting}
CHECK_SWITCH_POWERON {
[ THRESHOLD <threshold_all_V> |
<cellname> <cell_threshold_V> ]
...
}
\end{lstlisting}

\textbf{示例：}
\begin{lstlisting}
CHECK_SWITCH_POWERON {
SW_ABC 0.2
SW_CDE 0.1
SW_xyz 0.15
\end{lstlisting}

\subsubsection{\gsr{DECOUPLE_LDO_GROUND}}
\paragraph*{DECOUPLE_LDO_GROUND}

GSR 关键字 DECOUPLE_LDO_GROUND：用于设置与 DECOUPLE、LDO、GROUND 相关的 RedHawk 运行参数。（可选；默认：0）

\textbf{语法：}
\begin{lstlisting}
DECOUPLE_LDO_GROUND [ 0 | 1 ]
\end{lstlisting}

\subsubsection{\gsr{DYNAMIC_ADAPTIVE_RON}}
\paragraph*{DYNAMIC_ADAPTIVE_RON}

GSR 关键字 DYNAMIC_ADAPTIVE_RON：用于设置与 动态、ADAPTIVE、RON 相关的 RedHawk 运行参数。（可选；默认：switch model is used, RedHawk uses inter）

\textbf{语法：}
\begin{lstlisting}
DYNAMIC_ADAPTIVE_RON [0|1] <variation_threshold_%>
\end{lstlisting}

\subsubsection{\gsr{DYNAMIC_EXTEND_GSC}}
\paragraph*{DYNAMIC_EXTEND_GSC}

GSR 关键字 DYNAMIC_EXTEND_GSC：用于设置与 动态、EXTEND、GSC 相关的 RedHawk 运行参数。（可选；默认：1）

\textbf{语法：}
\begin{lstlisting}
DYNAMIC_EXTEND_GSC [0|1|2|3]
\end{lstlisting}

\textbf{参数说明：}
\begin{itemize}
\item \texttt{0} — GSR 关键字 DYNAMIC\_EXTEND\_GSC：用于设置与 动态、EXTEND、GSC 相关的 RedHawk 运行参数。
\end{itemize}

\subsubsection{\gsr{DYNAMIC_GSC_CHECK}}
\paragraph*{DYNAMIC_GSC_CHECK}

GSR 关键字 DYNAMIC_GSC_CHECK：用于设置与 动态、GSC、检查 相关的 RedHawk 运行参数。（可选；默认：1 (on)）

\textbf{语法：}
\begin{lstlisting}
DYNAMIC_GSC_CHECK [ 0 | 1 ]
\end{lstlisting}

\subsubsection{\gsr{EXTRACT_INTERNAL_NET}}
\paragraph*{EXTRACT_INTERNAL_NET}

GSR 关键字 EXTRACT_INTERNAL_NET：用于设置与 提取、INTERNAL、网络 相关的 RedHawk 运行参数。（可选；默认：1；兼容 DMP）

\textbf{语法：}
\begin{lstlisting}
EXTRACT_INTERNAL_NET [ 0 | 1 ]
\end{lstlisting}

\subsubsection{\gsr{PIECEWISE_CAP_FILES}}
\paragraph*{PIECEWISE_CAP_FILES}

GSR 关键字 PIECEWISE_CAP_FILES：用于设置与 PIECEWISE、CAP、文件 相关的 RedHawk 运行参数。（可选；默认：none）

\textbf{语法：}
\begin{lstlisting}
PIECEWISE_CAP_FILES <file name>
\end{lstlisting}

\textbf{示例：}
\begin{lstlisting}
PIECEWISE_CAP_FILES pwl
\end{lstlisting}

\subsubsection{\gsr{POWERUP_SAVE}}
\paragraph*{POWERUP_SAVE}

GSR 关键字 POWERUP_SAVE：用于设置与 POWERUP、SAVE 相关的 RedHawk 运行参数。（可选；默认：none）

\textbf{语法：}
\begin{lstlisting}
POWERUP_SAVE <FilePathName>
\end{lstlisting}

\textbf{示例：}
\begin{lstlisting}
POWERUP_SAVE powerup_save.file
\end{lstlisting}

\subsubsection{\gsr{POWERUP_OUTPUT_HIGH_PROB}}
\paragraph*{POWERUP_OUTPUT_HIGH_PROB}

GSR 关键字 POWERUP_OUTPUT_HIGH_PROB：用于设置与 POWERUP、OUTPUT、HIGH、PROB 相关的 RedHawk 运行参数。（可选；默认：0）

\textbf{语法：}
\begin{lstlisting}
POWERUP_OUTPUT_HIGH_PROB <probability>
\end{lstlisting}

\textbf{示例：}
\begin{lstlisting}
POWERUP_OUTPUT_HIGH_PROB 0.7
\end{lstlisting}

\subsubsection{\gsr{RAMPUP_OFFSTATE_VOLTAGE}}
\paragraph*{RAMPUP_OFFSTATE_VOLTAGE}

GSR 关键字 RAMPUP_OFFSTATE_VOLTAGE：用于设置与 RAMPUP、OFFSTATE、VOLTAGE 相关的 RedHawk 运行参数。（可选；默认：this keyword is not set, and simulation；兼容 DMP）

\textbf{语法：}
\begin{lstlisting}
RAMPUP_OFFSTATE_VOLTAGE <initial_voltage>
\end{lstlisting}

\textbf{示例：}
\begin{lstlisting}
RAMPUP_OFFSTATE_VOLTAGE 0.315
\end{lstlisting}

\subsubsection{\gsr{RAMPUP_SW_REPORT}}
\paragraph*{RAMPUP_SW_REPORT}

GSR 关键字 RAMPUP_SW_REPORT：用于设置与 RAMPUP、SW、REPORT 相关的 RedHawk 运行参数。（默认：0）

\textbf{语法：}
\begin{lstlisting}
RAMPUP_SW_REPORT [ 1| 0]
\end{lstlisting}

\subsubsection{\gsr{SWITCH_MODEL_FILES}}
\paragraph*{SWITCH_MODEL_FILES}

GSR 关键字 SWITCH_MODEL_FILES：用于设置与 SWITCH、MODEL、文件 相关的 RedHawk 运行参数。（可选；默认：none；兼容 DMP；必需）

\subsubsection{\gsr{SWITCH_MODEL_FILES}}
\paragraph*{SWITCH_MODEL_FILES}

GSR 关键字 SWITCH_MODEL_FILES：用于设置与 SWITCH、MODEL、文件 相关的 RedHawk 运行参数。

\subsubsection{\gsr{SWITCH_MODEL_FILES}}
\paragraph*{SWITCH_MODEL_FILES}

GSR 关键字 SWITCH_MODEL_FILES：用于设置与 SWITCH、MODEL、文件 相关的 RedHawk 运行参数。

\subsubsection{\gsr{SW_MOD_AB}}
\paragraph*{SW_MOD_AB}

GSR 关键字 SW_MOD_AB：用于设置与 SW、MOD、AB 相关的 RedHawk 运行参数。

\subsubsection{\gsr{SWITCH_MODEL_XTR_FILES}}
\paragraph*{SWITCH_MODEL_XTR_FILES}

GSR 关键字 SWITCH_MODEL_XTR_FILES：用于设置与 SWITCH、MODEL、XTR、文件 相关的 RedHawk 运行参数。

\textbf{语法：}
\begin{lstlisting}
SWITCH_MODEL_XTR_FILES {
<file1>
...
}
\end{lstlisting}

\textbf{参数说明：}
\begin{itemize}
\item \texttt{ext\_xN, ext\_yN} — GSR 关键字 SWITCH\_MODEL\_XTR\_FILES：用于设置与 SWITCH、MODEL、XTR、文件 相关的 RedHawk 运行参数。（可选；默认：none；兼容 DMP）
\end{itemize}

\subsubsection{\gsr{SWITCH_RES_FILES}}
\paragraph*{SWITCH_RES_FILES}

GSR 关键字 SWITCH_RES_FILES：用于设置与 SWITCH、RES、文件 相关的 RedHawk 运行参数。（可选；默认：none）

\textbf{语法：}
\begin{lstlisting}
SWITCH_RES_FILES {
<file name>
\end{lstlisting}

\subsubsection{\gsr{USE_MF_SWITCH_MODEL}}
\paragraph*{USE_MF_SWITCH_MODEL}

GSR 关键字 USE_MF_SWITCH_MODEL：用于设置与 USE、MF、SWITCH、MODEL 相关的 RedHawk 运行参数。（可选；默认：0；兼容 DMP）

\textbf{语法：}
\begin{lstlisting}
USE_MF_SWITCH_MODEL [ 0|1]
\end{lstlisting}

\subsubsection{\gsr{VP_CONTROL}}
\paragraph*{VP_CONTROL}

GSR 关键字 VP_CONTROL：用于设置与 VP、CONTROL 相关的 RedHawk 运行参数。（可选；默认：none） 注意：GSR 关键字 VP_CONTROL：用于设置与 VP、CONTROL 相关的 RedHawk 运行参数。

\textbf{语法：}
\begin{lstlisting}
VP_CONTROL {
<filename>
}
\end{lstlisting}

\textbf{参数说明：}
\begin{itemize}
\item \texttt{<daisy\_chain\_delay>} — GSR 关键字 VP\_CONTROL：用于设置与 VP、CONTROL 相关的 RedHawk 运行参数。
\end{itemize}

\textbf{示例：}
\begin{lstlisting}
VP_CONTROL {
vpcf_abcd4
}
Information in the VPC file has the following format (all times and delays in seconds):
<macro name> {
POWERUP <block_ctrl_pin_1> <int_domain1> <pwrup_time> ?<daisy_chain_delay>?
...
POWERUP <block_ctrl_pin_n> <int_domainN> <pwrup_time> ?<daisy_chain_delay>?
POWERON <block ctrl pin_1> <int_domain1> ?<poweron_time>?
...
POWERON <block ctrl pin_n> <int_domainN> ?<poweron_time>?
SWITCH_RAMPUP_TIMING <sw_inst_name> <pwrup_time> ?<poweron_time>?
}
Mem64x1024 {
POWERUP PON VDDINT 100e-12
POWERON PGOOD VDDINT 150e-12
}
The first line of the definition means that ‘PON’ is the name of controlling pin for the
switches between the external net and the internal net ‘VDDINT’, and 100e-12 is the
delay in seconds from PON to all the controlling pins of switches inside the memory
macro for that virtual domain. The keyword ‘POWERUP’ means that PON is the
charging pin and maps to the POWERUP state in the switch model. ‘POWERON’
means that ‘PGOOD’ is the controlling pin for the ON state in the switch model.
Mem64x1024 {
POWERUP PON VDDEXT VDDINT 100e-12 20e-12
}
If there are three switch instances in Mem64x1024, namely adsU2, adsU3 and
adsU4, then the controlling pin ‘PON’ of adsU2, adsU3 and adsU4 is assigned delay
numbers of 100ps, 120ps and 140ps respectively, depending on the switch instance
sequence in the block DEF file.
Mem64x1024 {
POWERUP PON VDDINT 50e-12
SWITCH_RAMPUP_TIMING adsU2 100e-12
SWITCH_RAMPUP_TIMING adsU3 120e-12
SWITCH_RAMPUP_TIMING adsU4 140e-12
}
The ‘POWERUP’ statement must be specified. The SWITCH_RAMPUP_TIMING
value overwrites the global delay assignment in POWERUP. For example,
Mem64x1024 has switch instances adsU2, adsU3, adsU4 and adsU5. Switch adsU2,
adsU3 and adsU4 are assigned delay numbers of 100ps, 120ps and 140ps respec-
\end{lstlisting}

\begin{noteBox}
GSR 关键字 VP\_CONTROL：用于设置与 VP、CONTROL 相关的 RedHawk 运行参数。
\end{noteBox}

\subsection{ESD 关键字}
\subsubsection*{ESD Keywords}

本类共 12 个 GSR 关键字。以下逐条给出中文用途、语法、默认值与示例。

\subsubsection{\gsr{ESD_CLAMP_FILE}}
\paragraph*{ESD_CLAMP_FILE}

GSR 关键字 ESD_CLAMP_FILE：用于设置与 ESD、CLAMP、文件 相关的 RedHawk 运行参数。（可选；默认：none）

\textbf{语法：}
\begin{lstlisting}
ESD_CLAMP_FILE <clamp_file>
or
ESD_CLAMP_FILE {
<clamp_fileName1>
\end{lstlisting}

\subsubsection{\gsr{ESD_CLAMP_PIN_FILE}}
\paragraph*{ESD_CLAMP_PIN_FILE}

GSR 关键字 ESD_CLAMP_PIN_FILE：用于设置与 ESD、CLAMP、引脚、文件 相关的 RedHawk 运行参数。

\textbf{语法：}
\begin{lstlisting}
ESD_CLAMP_PIN_FILES <file_name>
or
ESD_CLAMP_PIN_FILES {
<file_name1>
<file_name2>
...
}
\end{lstlisting}

\textbf{参数说明：}
\begin{itemize}
\item \texttt{all} — GSR 关键字 ESD\_CLAMP\_PIN\_FILE：用于设置与 ESD、CLAMP、引脚、文件 相关的 RedHawk 运行参数。（可选；默认：none）
\end{itemize}

\subsubsection{\gsr{ESD_CLAMP_PIN_NODE_DISTANCE}}
\paragraph*{ESD_CLAMP_PIN_NODE_DISTANCE}

GSR 关键字 ESD_CLAMP_PIN_NODE_DISTANCE：用于设置与 ESD、CLAMP、引脚、节点、DISTANCE 相关的 RedHawk 运行参数。（可选；默认：0 (no nodes are added)）

\textbf{语法：}
\begin{lstlisting}
ESD_CLAMP_PIN_NODE_DISTANCE <distance>
\end{lstlisting}

\subsubsection{\gsr{ESD_DEF_IGNORE_LAYER}}
\paragraph*{ESD_DEF_IGNORE_LAYER}

GSR 关键字 ESD_DEF_IGNORE_LAYER：用于设置与 ESD、DEF、忽略、层 相关的 RedHawk 运行参数。（可选；默认：none）

\textbf{语法：}
\begin{lstlisting}
ESD_DEF_IGNORE_LAYER {
<layer1>
<layer2>
...
}
\end{lstlisting}

\subsubsection{\gsr{ESD_EXTRACT_CLAMP_NET}}
\paragraph*{ESD_EXTRACT_CLAMP_NET}

GSR 关键字 ESD_EXTRACT_CLAMP_NET：用于设置与 ESD、提取、CLAMP、网络 相关的 RedHawk 运行参数。（可选；默认：, and also whether to add a ploc for the）

\textbf{语法：}
\begin{lstlisting}
ESD_EXTRACT_CLAMP_NET [ 0 | 1 | 2 ]
\end{lstlisting}

\subsubsection{\gsr{ESD_GSR_OPTIMIZE}}
\paragraph*{ESD_GSR_OPTIMIZE}

GSR 关键字 ESD_GSR_OPTIMIZE：用于设置与 ESD、GSR、OPTIMIZE 相关的 RedHawk 运行参数。（可选；默认：0）

\textbf{语法：}
\begin{lstlisting}
ESD_GSR_OPTIMIZE [ 0 | 1 | 2 | 3 ]
\end{lstlisting}

\textbf{参数说明：}
\begin{itemize}
\item \texttt{0 (or off)} — GSR 关键字 ESD\_GSR\_OPTIMIZE：用于设置与 ESD、GSR、OPTIMIZE 相关的 RedHawk 运行参数。（必需）
\end{itemize}

\subsubsection{\gsr{ESD_LOCAL_NETS}}
\paragraph*{ESD_LOCAL_NETS}

GSR 关键字 ESD_LOCAL_NETS：用于设置与 ESD、LOCAL、网络 相关的 RedHawk 运行参数。（可选；默认：none）

\textbf{语法：}
\begin{lstlisting}
ESD_LOCAL_NETS {
<netName1>
<netName2>
...
}
\end{lstlisting}

\subsubsection{\gsr{ESD_MERGE_NETS}}
\paragraph*{ESD_MERGE_NETS}

指定 网络 to be merged during ESD 分析.可选。默认：无。（可选；默认：none）

\textbf{语法：}
\begin{lstlisting}
ESD_MERGE_NETS {
<prime_net> {
<sub_net1>
<sub_net2>
...
}
....
Here, RedHawk will merge the sub_nets into the prime_net. The net
names are in terms of flattened net list similar to with VDD_NETS and
GND_NETS.
\end{lstlisting}

\subsubsection{\gsr{ESD_RULE_FILE}}
\paragraph*{ESD_RULE_FILE}

GSR 关键字 ESD_RULE_FILE：用于设置与 ESD、规则、文件 相关的 RedHawk 运行参数。（可选；默认：none）

\textbf{语法：}
\begin{lstlisting}
ESD_RULE_FILE <rule_filename>
or
\end{lstlisting}

\subsubsection{\gsr{ESD_SHORT_CLAMP_PIN}}
\paragraph*{ESD_SHORT_CLAMP_PIN}

GSR 关键字 ESD_SHORT_CLAMP_PIN：用于设置与 ESD、短路、CLAMP、引脚 相关的 RedHawk 运行参数。（可选；默认：-1）

\textbf{语法：}
\begin{lstlisting}
ESD_SHORT_CLAMP_PIN [ -1 | 0 | 1 | 2 |layer ]
\end{lstlisting}

\textbf{参数说明：}
\begin{itemize}
\item \texttt{-1} — GSR 关键字 ESD\_SHORT\_CLAMP\_PIN：用于设置与 ESD、短路、CLAMP、引脚 相关的 RedHawk 运行参数。
\end{itemize}

\subsubsection{\gsr{ESD_SIGNAL_NETS}}
\paragraph*{ESD_SIGNAL_NETS}

GSR 关键字 ESD_SIGNAL_NETS：用于设置与 ESD、SIGNAL、网络 相关的 RedHawk 运行参数。（可选；默认：none）

\textbf{语法：}
\begin{lstlisting}
ESD_SIGNAL_NETS {
<netname1> <cell_name1>
<netname2> <cell_name2>
...
}
\end{lstlisting}

\subsubsection{\gsr{ESD_SIGNAL_NET_FILE}}
\paragraph*{ESD_SIGNAL_NET_FILE}

GSR 关键字 ESD_SIGNAL_NET_FILE：用于设置与 ESD、SIGNAL、网络、文件 相关的 RedHawk 运行参数。（可选；默认：none；必需）

\textbf{语法：}
\begin{lstlisting}
ESD_SIGNAL_NET_FILE <filename>
\end{lstlisting}

\textbf{参数说明：}
GSR 关键字 ESD\_SIGNAL\_NET\_FILE：用于设置与 ESD、SIGNAL、网络、文件 相关的 RedHawk 运行参数。

\subsection{名称映射关键字}
\subsubsection*{Name Mapping Keywords}

本类共 8 个 GSR 关键字。以下逐条给出中文用途、语法、默认值与示例。

\subsubsection{\gsr{DEFER_VIA_CREATION}}
\paragraph*{DEFER_VIA_CREATION}

GSR 关键字 DEFER_VIA_CREATION：用于设置与 DEFER、过孔、CREATION 相关的 RedHawk 运行参数。（可选；默认：0；兼容 DMP）

\textbf{语法：}
\begin{lstlisting}
DEFER_VIA_CREATION [ 0 | 1 ]
\end{lstlisting}

\subsubsection{\gsr{NAME_CASE_SENSITIVE}}
\paragraph*{NAME_CASE_SENSITIVE}

GSR 关键字 NAME_CASE_SENSITIVE：用于设置与 NAME、CASE、SENSITIVE 相关的 RedHawk 运行参数。（可选；默认：1；兼容 DMP）

\textbf{语法：}
\begin{lstlisting}
NAME_CASE_SENSITIVE [ 0 | 1 ]
\end{lstlisting}

\subsubsection{\gsr{BUS_DELIMITER}}
\paragraph*{BUS_DELIMITER}

GSR 关键字 BUS_DELIMITER：用于设置与 BUS、DELIMITER 相关的 RedHawk 运行参数。（可选；默认：[ ]；兼容 DMP）

\textbf{语法：}
\begin{lstlisting}
BUS_DELIMITER <delim>
\end{lstlisting}

\textbf{示例：}
\begin{lstlisting}
BUS_DELIMITER ()
\end{lstlisting}

\subsubsection{\gsr{PIN_DELIMITER}}
\paragraph*{PIN_DELIMITER}

GSR 关键字 PIN_DELIMITER：用于设置与 引脚、DELIMITER 相关的 RedHawk 运行参数。（可选；默认：；兼容 DMP）

\textbf{语法：}
\begin{lstlisting}
PIN_DELIMITER <delim>
\end{lstlisting}

\textbf{示例：}
\begin{lstlisting}
PIN_DELIMITER @
\end{lstlisting}

\subsubsection{\gsr{BUS_DELIMITER_STA}}
\paragraph*{BUS_DELIMITER_STA}

GSR 关键字 BUS_DELIMITER_STA：用于设置与 BUS、DELIMITER、STA 相关的 RedHawk 运行参数。（可选；默认：[ ]；兼容 DMP）

\textbf{语法：}
\begin{lstlisting}
BUS_DELIMITER_STA <delim>
\end{lstlisting}

\textbf{示例：}
\begin{lstlisting}
BUS_DELIMITER_STA ()
\end{lstlisting}

\subsubsection{\gsr{PIN_DELIMITER_STA}}
\paragraph*{PIN_DELIMITER_STA}

GSR 关键字 PIN_DELIMITER_STA：用于设置与 引脚、DELIMITER、STA 相关的 RedHawk 运行参数。（可选；默认：/ Global System Requirements File (*；兼容 DMP）

\textbf{语法：}
\begin{lstlisting}
PIN_DELIMITER_STA <delim>
\end{lstlisting}

\textbf{示例：}
\begin{lstlisting}
PIN_DELIMITER_STA :
\end{lstlisting}

\subsubsection{\gsr{HIER_DIVIDER}}
\paragraph*{HIER_DIVIDER}

GSR 关键字 HIER_DIVIDER：用于设置与 HIER、DIVIDER 相关的 RedHawk 运行参数。（可选；默认：/；兼容 DMP）

\textbf{语法：}
\begin{lstlisting}
HIER_DIVIDER <divider>
\end{lstlisting}

\textbf{示例：}
\begin{lstlisting}
HIER_DIVIDER /
\end{lstlisting}

\subsubsection{\gsr{HIER_DIVIDER_STA}}
\paragraph*{HIER_DIVIDER_STA}

GSR 关键字 HIER_DIVIDER_STA：用于设置与 HIER、DIVIDER、STA 相关的 RedHawk 运行参数。（可选；默认：/；兼容 DMP）

\textbf{语法：}
\begin{lstlisting}
HIER_DIVIDER_STA <divider>
\end{lstlisting}

\textbf{示例：}
\begin{lstlisting}
HIER_DIVIDER_STA /
\end{lstlisting}

\subsection{警告与错误消息关键字}
\subsubsection*{Warning and Error Message Keywords}

本类共 4 个 GSR 关键字。以下逐条给出中文用途、语法、默认值与示例。

\subsubsection{\gsr{WARNING_COUNTS}}
\paragraph*{WARNING_COUNTS}

限制每类 Warning 写入 stdout 与 apache.log 的最大条数。（可选；默认：3；兼容 DMP）

\textbf{语法：}
\begin{lstlisting}
WARNING_COUNTS [ <number_mess> | 0 | -1 ]
\end{lstlisting}

\textbf{参数说明：}
\begin{itemize}
\item \texttt{<number\_mess>} — 限制每类 Warning 写入 stdout 与 apache.log 的最大条数。
\end{itemize}

\textbf{示例：}
\begin{lstlisting}
WARNING_COUNTS 10
\end{lstlisting}

\subsubsection{\gsr{WARNING_LOG_COUNTS}}
\paragraph*{WARNING_LOG_COUNTS}

限制写入 adsRpt/redhawk.warn 的每类 Warning 最大条数。（可选；默认：100；兼容 DMP）

\textbf{语法：}
\begin{lstlisting}
WARNING_LOG_COUNTS [ <number_mess> | 0 | -1 ]
\end{lstlisting}

\textbf{参数说明：}
\begin{itemize}
\item \texttt{<number\_mess>} — 限制写入 adsRpt/redhawk.warn 的每类 Warning 最大条数。
\end{itemize}

\textbf{示例：}
\begin{lstlisting}
WARNING_LOG_COUNTS 10
\end{lstlisting}

\subsubsection{\gsr{ERROR_COUNTS}}
\paragraph*{ERROR_COUNTS}

限制每类 Error 写入 stdout 的最大条数。（可选；默认：3；兼容 DMP）

\textbf{语法：}
\begin{lstlisting}
ERROR_COUNTS <number_mess>
\end{lstlisting}

\textbf{示例：}
\begin{lstlisting}
ERROR_COUNTS 10
\end{lstlisting}

\subsubsection{\gsr{ERROR_LOG_COUNTS}}
\paragraph*{ERROR_LOG_COUNTS}

限制写入 adsRpt/redhawk.err 的每类 Error 最大条数。（可选；默认：100；兼容 DMP）

\textbf{语法：}
\begin{lstlisting}
ERROR_LOG_COUNTS [ <number_mess> | 0 | -1 ]
\end{lstlisting}

\textbf{参数说明：}
\begin{itemize}
\item \texttt{<number\_mess>} — 限制写入 adsRpt/redhawk.err 的每类 Error 最大条数。
\end{itemize}

\textbf{示例：}
\begin{lstlisting}
ERROR_LOG_COUNTS 100
\end{lstlisting}

\subsection{忽略功能关键字}
\subsubsection*{Ignore Function Keywords}

本类共 40 个 GSR 关键字。以下逐条给出中文用途、语法、默认值与示例。

\subsubsection{\gsr{IGNORE_CELLS}}
\paragraph*{IGNORE_CELLS}

导入 DEF 时忽略列表中的单元（不写入设计数据库）。（可选；默认：none；兼容 DMP）

\subsubsection{\gsr{IGNORE_CELLS}}
\paragraph*{IGNORE_CELLS}

导入 DEF 时忽略列表中的单元（不写入设计数据库）。

\subsubsection{\gsr{IGNORE_CELLS}}
\paragraph*{IGNORE_CELLS}

导入 DEF 时忽略列表中的单元（不写入设计数据库）。

\subsubsection{\gsr{IGNORE_CELLS_FILES}}
\paragraph*{IGNORE_CELLS_FILES}

在数据输入或分析流程中忽略与 单元、文件 相关的对象或检查。（可选；默认：none；兼容 DMP）

\textbf{语法：}
\begin{lstlisting}
IGNORE_CELLS_FILES <Exclude_cells_file_pathName>
\end{lstlisting}

\textbf{示例：}
\begin{lstlisting}
IGNORE_CELLS_FILES /home/users/customerA/cells_exclude
\end{lstlisting}

\subsubsection{\gsr{IGNORE_DEF_ERROR}}
\paragraph*{IGNORE_DEF_ERROR}

在数据输入或分析流程中忽略与 DEF、错误 相关的对象或检查。（可选；默认：0 (perform check)；兼容 DMP）

\subsubsection{\gsr{IGNORE_DEF_ERROR}}
\paragraph*{IGNORE_DEF_ERROR}

在数据输入或分析流程中忽略与 DEF、错误 相关的对象或检查。

\subsubsection{\gsr{IGNORE_ERROR_POPUP}}
\paragraph*{IGNORE_ERROR_POPUP}

在数据输入或分析流程中忽略与 错误、弹窗 相关的对象或检查。（可选；默认：0 (Off)；兼容 DMP）

\textbf{语法：}
\begin{lstlisting}
IGNORE_ERROR_POPUP [ 0 | 1 ]
\end{lstlisting}

\subsubsection{\gsr{IGNORE_ESCAPE_CHAR}}
\paragraph*{IGNORE_ESCAPE_CHAR}

控制是否忽略名称中的转义字符（escape character，通常为反斜杠 \texttt{\textbackslash}）。可选；兼容 DMP。

\textbf{语法：}
\begin{lstlisting}
IGNORE_ESCAPE_CHAR [ true | false ]
\end{lstlisting}

\subsubsection{\gsr{IGNORE_FILE_PREPARSE}}
\paragraph*{IGNORE_FILE_PREPARSE}

在数据输入或分析流程中忽略与 文件、PREPARSE 相关的对象或检查。（可选；默认：0 (do not ignore preparse)；兼容 DMP）

\textbf{语法：}
\begin{lstlisting}
IGNORE_FILE_PREPARSE [ 0 | 1 ]
\end{lstlisting}

\subsubsection{\gsr{IGNORE_FILLER_DECAP_CELL_REPORT}}
\paragraph*{IGNORE_FILLER_DECAP_CELL_REPORT}

在数据输入或分析流程中忽略与 FILLER、去耦电容、单元、REPORT 相关的对象或检查。（可选；默认：0 (do not ignore)；兼容 DMP）

\textbf{语法：}
\begin{lstlisting}
IGNORE_FILLER_DECAP_CELL_REPORT [ 0 | 1 ]
\end{lstlisting}

\subsubsection{\gsr{IGNORE_FLOATING_INSTANCE_MISSING_TW_CHECK}}
\paragraph*{IGNORE_FLOATING_INSTANCE_MISSING_TW_CHECK}

在数据输入或分析流程中忽略与 悬空、实例、MISSING、TW、检查 相关的对象或检查。（可选；默认：0 (off)；兼容 DMP）

\textbf{语法：}
\begin{lstlisting}
Syntax
IGNORE_FLOATING_INSTANCE_MISSING_TW_CHECK [ 0 | 1 ]
\end{lstlisting}

\subsubsection{\gsr{IGNORE_FLOATING_PAD}}
\paragraph*{IGNORE_FLOATING_PAD}

在数据输入或分析流程中忽略与 悬空、焊盘 相关的对象或检查。（默认：0）

\textbf{语法：}
\begin{lstlisting}
IGNORE_FLOATING_PAD [ 1 | 0 ]
\end{lstlisting}

\subsubsection{\gsr{IGNORE_GDSMEM_ERROR}}
\paragraph*{IGNORE_GDSMEM_ERROR}

在数据输入或分析流程中忽略与 GDSMEM、错误 相关的对象或检查。（可选；默认：0 (perform check)；兼容 DMP）

\textbf{语法：}
\begin{lstlisting}
IGNORE_GDSMEM_ERROR [ 0 | 1 ]
\end{lstlisting}

\subsubsection{\gsr{IGNORE_GDS2DEF_UNCONNECTS}}
\paragraph*{IGNORE_GDS2DEF_UNCONNECTS}

在数据输入或分析流程中忽略与 GDS2DEF、UNCONNECTS 相关的对象或检查。（可选；默认：0）

\textbf{语法：}
\begin{lstlisting}
IGNORE_GDS2DEF_UNCONNECTS [ 0 | 1]
\end{lstlisting}

\subsubsection{\gsr{IGNORE_INSTANCES}}
\paragraph*{IGNORE_INSTANCES}

数据输入阶段忽略指定实例列表。（可选；默认：none Global System Requirements File (*；兼容 DMP）

\textbf{语法：}
\begin{lstlisting}
IGNORE_INSTANCES {
<exact_DEF_instance_name1> ? <exact DEF_FILES path1> ?
...<<hierarchical path to the ignore instance of a block>
hierarchy
}
\end{lstlisting}

\textbf{示例：}
\begin{lstlisting}
IGNORE_INSTANCES {
\end{lstlisting}

\subsubsection{\gsr{IGNORE_INSTANCES_FILES}}
\paragraph*{IGNORE_INSTANCES_FILES}

在数据输入或分析流程中忽略与 实例、文件 相关的对象或检查。（可选；默认：none；兼容 DMP）

\textbf{语法：}
\begin{lstlisting}
IGNORE_INSTANCES_FILES <exclude_instances_file_pathname>
\end{lstlisting}

\textbf{示例：}
\begin{lstlisting}
IGNORE_INSTANCES_FILES /home/users/customerA/nets_exclude
\end{lstlisting}

\subsubsection{\gsr{IGNORE_LEF_DEF_MISMATCH}}
\paragraph*{IGNORE_LEF_DEF_MISMATCH}

在数据输入或分析流程中忽略与 LEF、DEF、MISMATCH 相关的对象或检查。（可选；默认：0 (do not ignore)；兼容 DMP）

\textbf{语法：}
\begin{lstlisting}
IGNORE_LEF_DEF_MISMATCH [ 0 | 1 ]
\end{lstlisting}

\subsubsection{\gsr{IGNORE_LEF_CELL_PIN_FILE}}
\paragraph*{IGNORE_LEF_CELL_PIN_FILE}

在数据输入或分析流程中忽略与 LEF、单元、引脚、文件 相关的对象或检查。

\textbf{语法：}
\begin{lstlisting}
IGNORE_LEF_CELL_PIN_FILE              <file_name>
File format:
<cellName>         < pinName>
and2x                    VBP
and2x                    VBN
*                    VDDR
\end{lstlisting}

\subsubsection{\gsr{IGNORE_LEF_PIN_DIRECTION}}
\paragraph*{IGNORE_LEF_PIN_DIRECTION}

在数据输入或分析流程中忽略与 LEF、引脚、DIRECTION 相关的对象或检查。（可选；默认：all P/G pins with direction OUT defined）

\textbf{语法：}
\begin{lstlisting}
IGNORE_LEF_PIN_DIRECTION [ 0 | 1 ]
\end{lstlisting}

\subsubsection{\gsr{IGNORE_IO_POWER}}
\paragraph*{IGNORE_IO_POWER}

在数据输入或分析流程中忽略与 IO、电源 相关的对象或检查。（可选；默认：0 (do not ignore)；兼容 DMP）

\textbf{语法：}
\begin{lstlisting}
IGNORE_IO_POWER [ 0 | 1 ]
\end{lstlisting}

\subsubsection{\gsr{IGNORE_LEF_MACRO}}
\paragraph*{IGNORE_LEF_MACRO}

在数据输入或分析流程中忽略与 LEF、宏 相关的对象或检查。（可选；默认：none；兼容 DMP）

\subsubsection{\gsr{IGNORE_LEF_MACRO}}
\paragraph*{IGNORE_LEF_MACRO}

在数据输入或分析流程中忽略与 LEF、宏 相关的对象或检查。

\subsubsection{\gsr{IGNORE_LEF_MACRO}}
\paragraph*{IGNORE_LEF_MACRO}

在数据输入或分析流程中忽略与 LEF、宏 相关的对象或检查。

\subsubsection{\gsr{IGNORE_LIB_CHECK}}
\paragraph*{IGNORE_LIB_CHECK}

在数据输入或分析流程中忽略与 LIB、检查 相关的对象或检查。（可选；默认：0 (perform check)；兼容 DMP）

\textbf{语法：}
\begin{lstlisting}
IGNORE_LIB_CHECK [ 0 | 1 ]
\end{lstlisting}

\subsubsection{\gsr{IGNORE_NETS}}
\paragraph*{IGNORE_NETS}

在数据输入或分析流程中忽略与 网络 相关的对象或检查。（可选；默认：none；兼容 DMP）

\subsubsection{\gsr{IGNORE_NETS}}
\paragraph*{IGNORE_NETS}

在数据输入或分析流程中忽略与 网络 相关的对象或检查。

\subsubsection{\gsr{IGNORE_NETS}}
\paragraph*{IGNORE_NETS}

在数据输入或分析流程中忽略与 网络 相关的对象或检查。

\subsubsection{\gsr{IGNORE_NETS_FILES}}
\paragraph*{IGNORE_NETS_FILES}

在数据输入或分析流程中忽略与 网络、文件 相关的对象或检查。（可选；默认：none；兼容 DMP）

\textbf{语法：}
\begin{lstlisting}
IGNORE_NETS_FILES <ignoreNets_file_pathName>
\end{lstlisting}

\textbf{示例：}
\begin{lstlisting}
IGNORE_NETS_FILES /home/users/customerA/nets_exclude
\end{lstlisting}

\subsubsection{\gsr{IGNORE_PRECHECK_ERROR}}
\paragraph*{IGNORE_PRECHECK_ERROR}

在数据输入或分析流程中忽略与 PRECHECK、错误 相关的对象或检查。（可选；默认：off）

\textbf{语法：}
\begin{lstlisting}
Syntax
IGNORE_PRECHECK_ERROR [ 0 | 1]
\end{lstlisting}

\subsubsection{\gsr{IGNORE_PRIMARY_LOAD_DECAP}}
\paragraph*{IGNORE_PRIMARY_LOAD_DECAP}

在数据输入或分析流程中忽略与 PRIMARY、LOAD、去耦电容 相关的对象或检查。（可选；默认：0 (off)）

\textbf{语法：}
\begin{lstlisting}
IGNORE_PRIMARY_LOAD_DECAP [ 0 | 1 ]
\end{lstlisting}

\subsubsection{\gsr{IGNORE_ROUTE}}
\paragraph*{IGNORE_ROUTE}

在数据输入或分析流程中忽略与 布线 相关的对象或检查。（可选；默认：None；兼容 DMP）

\textbf{语法：}
\begin{lstlisting}
IGNORE_ROUTE {
<def_file_path_name>           <layer_name>
}
\end{lstlisting}

\subsubsection{\gsr{IGNORE_ROUTE}}
\paragraph*{IGNORE_ROUTE}

在数据输入或分析流程中忽略与 布线 相关的对象或检查。

\subsubsection{\gsr{IGNORE_SHORT}}
\paragraph*{IGNORE_SHORT}

在数据输入或分析流程中忽略与 短路 相关的对象或检查。（可选；默认：1 (on)；兼容 DMP）

\textbf{语法：}
\begin{lstlisting}
IGNORE_SHORT [ 1 | 0 ]
\end{lstlisting}

\subsubsection{\gsr{IGNORE_SIMTIME_CHECK}}
\paragraph*{IGNORE_SIMTIME_CHECK}

在数据输入或分析流程中忽略与 SIMTIME、检查 相关的对象或检查。（可选；默认：0 (off)；兼容 DMP）

\textbf{语法：}
\begin{lstlisting}
IGNORE_SIMTIME_CHECK [ 1 | 0 ]
\end{lstlisting}

\subsubsection{\gsr{IGNORE_TECH_ERROR}}
\paragraph*{IGNORE_TECH_ERROR}

在数据输入或分析流程中忽略与 工艺、错误 相关的对象或检查。（可选；默认：0 (stop for dielectric data failure)；兼容 DMP）

\textbf{语法：}
\begin{lstlisting}
IGNORE_TECH_ERROR [ 0 | 1 ]
\end{lstlisting}

\subsubsection{\gsr{IGNORE_UNDEFINED_LAYER}}
\paragraph*{IGNORE_UNDEFINED_LAYER}

在数据输入或分析流程中忽略与 未定义、层 相关的对象或检查。（可选；默认：0 (off)；兼容 DMP）

\textbf{语法：}
\begin{lstlisting}
IGNORE_UNDEFINED_LAYER [ 0 | 1 ]
\end{lstlisting}

\subsubsection{\gsr{IGNORE_UNPLACED_INSTANCE}}
\paragraph*{IGNORE_UNPLACED_INSTANCE}

在数据输入或分析流程中忽略与 未放置、实例 相关的对象或检查。（可选；默认：0 (off)）

\textbf{语法：}
\begin{lstlisting}
IGNORE_UNPLACED_INSTANCE [ 0 | 1 ]
\end{lstlisting}

\subsubsection{\gsr{IGNORE_UPF_PGARC}}
\paragraph*{IGNORE_UPF_PGARC}

在数据输入或分析流程中忽略与 UPF、PGARC 相关的对象或检查。（可选；默认：off)；兼容 DMP）

\textbf{语法：}
\begin{lstlisting}
IGNORE_UPF_PGARC [ 0 | 1 ]
To ignore/honor specific UPF libraries, include honor_upf_pgarc 0/1 in custom lib
file. Note that this setting works with IGNORE_UPF_PGARC set to 1.
library <library_name> {
nomVoltage <voltage value>
nomTemperature <temp value>
honor_upf_pgarc [ 0 | 1 ]
}
\end{lstlisting}

\subsubsection{\gsr{MISSING_VIA_CHECK_IGNORE_CELLS}}
\paragraph*{MISSING_VIA_CHECK_IGNORE_CELLS}

GSR 关键字 MISSING_VIA_CHECK_IGNORE_CELLS：用于设置与 MISSING、过孔、检查、忽略、单元 相关的 RedHawk 运行参数。（可选；默认：None）

\textbf{语法：}
\begin{lstlisting}
MISSING_VIA_CHECK_IGNORE_CELLS {
<cellname_1>
<cellname_2>
...
}
\end{lstlisting}

\subsubsection{\gsr{IGNORE_VP_CONTROL_ERROR}}
\paragraph*{IGNORE_VP_CONTROL_ERROR}

在数据输入或分析流程中忽略与 VP、CONTROL、错误 相关的对象或检查。（可选；默认：Off）

\textbf{语法：}
\begin{lstlisting}
IGNORE_VP_CONTROL_ERROR [ 0 | 1 ]
\end{lstlisting}

\subsection{GSR 宏关键字}
\subsubsection*{GSR Macro Keywords}

本类共 2 个 GSR 关键字。以下逐条给出中文用途、语法、默认值与示例。

\subsubsection{\gsr{DEFINE}}
\paragraph*{DEFINE}

定义可在多次运行间复用的变量或完整文件路径。（可选；默认：none）

\textbf{语法：}
\begin{lstlisting}
DEFINE <user_variable_name> <value>
\end{lstlisting}

\textbf{示例：}
\begin{lstlisting}
Example
DEFINE AAA /my/path/design_data/lef
DEFINE BBB /my/path/design_data/lef
LEF_FILES {
$AAA/a1.lef
$AAA/a2.lef
$BBB/b1.lef }
Or, using a command file definition:
run.tcl
gsr define AAA /my/path/design_data/lef
gsr define BBB /my/path/design_data/lef
\end{lstlisting}

\subsubsection{\gsr{INCLUDE}}
\paragraph*{INCLUDE}

将其他 GSR 文件内容嵌入当前 GSR，等效于直接写入。（可选；默认：none）

\textbf{语法：}
\begin{lstlisting}
INCLUDE <GSR_filename>
\end{lstlisting}

\textbf{示例：}
\begin{lstlisting}
Example
If a GSR file has the contents:
INCLUDE <GSR_fileA>
<gsr_keyword_1> <val_1>
<gsr_keyword_2> <val_2>
...
\end{lstlisting}

\section{焊盘、电源/地与 I/O 定义文件}
\subsection*{Pad, Power/Ground and I/O Definition Files}

\subsection{统一焊盘输入格式}
\subsubsection*{Unified Pad Input File Format}

GSR 关键字 ：用于设置与  相关的 RedHawk 运行参数。

GSR 关键字 ：用于设置与  相关的 RedHawk 运行参数。

\begin{lstlisting}
                         *PLOC
                               <.ploc file data>
\end{lstlisting}

GSR 关键字 ：用于设置与  相关的 RedHawk 运行参数。

\begin{lstlisting}
                         *PCELL
                               <.pcell file data>
\end{lstlisting}

GSR 关键字 ：用于设置与  相关的 RedHawk 运行参数。

\begin{lstlisting}
                         *PAD
                               <.pad file data>
\end{lstlisting}

GSR 关键字 ：用于设置与  相关的 RedHawk 运行参数。

\begin{lstlisting}
                        import pad unified_pad.txt
\end{lstlisting}

GSR 关键字 ：用于设置与  相关的 RedHawk 运行参数。

\begin{lstlisting}
                         PAD_FILES {
\end{lstlisting}

GSR 关键字 ：用于设置与  相关的 RedHawk 运行参数。

\begin{lstlisting}
                         }
\end{lstlisting}

GSR 关键字 ：用于设置与  相关的 RedHawk 运行参数。

\begin{lstlisting}
                    *PLOC_PSS
                    <pin_name> <x_loc> <y_loc> <layer> [POWER|GROUND]
                              <PAD_R_ohm> <PAD_L_pH> <PAD_C_pF>
\end{lstlisting}

\begin{lstlisting}
                   *PAD_PSS
                   <pad_inst name> <inst_pin name> <PAD_R_ohm>
                            <PAD_L_pH> <PAD_C_pF> <pkg_node>
\end{lstlisting}

GSR 关键字 ：用于设置与  相关的 RedHawk 运行参数。

GSR 关键字 ：用于设置与  相关的 RedHawk 运行参数。

GSR 关键字 ：用于设置与  相关的 RedHawk 运行参数。

\begin{lstlisting}
              ------------------------------------------------------------------------------------
              ## Format for .ploc files:
              *PLOC
              # For original .ploc format with 5 items
              # <pin_name> <x_loc> <y_loc> <layer> [POWER|GROUND]
\end{lstlisting}

GSR 关键字 ：用于设置与  相关的 RedHawk 运行参数。

\begin{lstlisting}
              # For original .ploc format with P/G_net name specified, 5 items
              # <pin_name> <x_loc> <y_loc> <layer> <PG_net name>
\end{lstlisting}

GSR 关键字 ：用于设置与  相关的 RedHawk 运行参数。

\begin{lstlisting}
              # For block pin voltage flow runs with 6 items
              # <pin_name> <x_loc> <y_loc> <layer> [POWER|GROUND] <pin voltage>
\end{lstlisting}

GSR 关键字 ：用于设置与  相关的 RedHawk 运行参数。

\begin{lstlisting}
              # For specifying bump L, R, C data - 8 items
              # <pin_name> <x_loc> <y_loc> <layer> <power|ground> <L_pH> <R_Ohms> <C_pF>
\end{lstlisting}

GSR 关键字 ：用于设置与  相关的 RedHawk 运行参数。

\begin{lstlisting}
              ## Format for .pcell files:
              *PCELL
\end{lstlisting}

\begin{lstlisting}
              # Original .pcell file with 1 item
              # <pad_cell macro name>
\end{lstlisting}

GSR 关键字 ：用于设置与  相关的 RedHawk 运行参数。

\begin{lstlisting}
              # Formats that specify which pins are used in a pad macro - 2 items
              # <pad_cell macro name> <macro pin name>
                  VDD_PAD_MASTER vdd
                  VDD_PAD_MASTER gnd
\end{lstlisting}

GSR 关键字 ：用于设置与  相关的 RedHawk 运行参数。

\begin{lstlisting}
              # Formats that specify connection locations and layers within
              # a pad macro: 1-2 line format
              # 1st : <pad_cell macro name>
              # <x_macro_loc> <y_macro_loc> <layer> [ Power|Ground ]
\end{lstlisting}

GSR 关键字 ：用于设置与  相关的 RedHawk 运行参数。

GSR 关键字 ：用于设置与  相关的 RedHawk 运行参数。

GSR 关键字 ：用于设置与  相关的 RedHawk 运行参数。

\begin{lstlisting}
              ## Format for .pad files:
              *PAD
\end{lstlisting}

\begin{lstlisting}
              # Original .pad file format - 1 item
              # <pad_instance name>
\end{lstlisting}

焊盘_INST1 焊盘_INST2

\begin{lstlisting}
              # To specify which pins are used in a pad instance - 2 items
              # <pad_inst name> <inst_pin name>
                  VDD_PAD_INST1 vdd
                  VDD_PAD_INST1 gnd
                  VDD_PAD_INST8 vdd
                  VDD_PAD_INST5 gnd
\end{lstlisting}

\begin{lstlisting}
              # To specify connection loc and layers for a pad instance: 1-2 lines
              # 1st: <pad_inst_name> <x_macro_loc> <y_macro_loc> <layer> [PWR|GRND]
              # 2nd: <x_macro_loc> <y_macro_loc> <layer> [PWR|GRND], if defined
              # for the same pad instance name
\end{lstlisting}

GSR 关键字 ：用于设置与  相关的 RedHawk 运行参数。

\begin{lstlisting}
              ## Format for .ploc files with a Package Spice Subckt (PSS)
              # definition, in a separate input file
              *PLOC_PSS
              # Package-Spice-Subckt(PSS) with 6 items
              # <pin_name> <x_loc> <y_loc> <layer> <PG_net name> <subckt node name>
\end{lstlisting}

GSR 关键字 ：用于设置与  相关的 RedHawk 运行参数。

\begin{lstlisting}
              ---------------------------------------------------------------------------------------------------
\end{lstlisting}

GSR 关键字 ：用于设置与  相关的 RedHawk 运行参数。

\begin{lstlisting}
                    setup pss -pad_file <filename> -subckt <pss_ckt_name>
\end{lstlisting}

GSR 关键字 ：用于设置与  相关的 RedHawk 运行参数。

\begin{lstlisting}
                        import pad <filename>
\end{lstlisting}

GSR 关键字 ：用于设置与  相关的 RedHawk 运行参数。

GSR 关键字 ：用于设置与  相关的 RedHawk 运行参数。

GSR 关键字 ：用于设置与  相关的 RedHawk 运行参数。（必需）

GSR 关键字 ：用于设置与  相关的 RedHawk 运行参数。

GSR 关键字 ：用于设置与  相关的 RedHawk 运行参数。

\begin{lstlisting}
              power and ground pads by using the Edit> Add Pad command.
              Note that you can use Edit -> Undo/Redo on Add/Delete Pad operations, and multiple
\end{lstlisting}

GSR 关键字 ：用于设置与  相关的 RedHawk 运行参数。

GSR 关键字 ：用于设置与  相关的 RedHawk 运行参数。

\begin{lstlisting}
              a <design>.pcell file:
                         <pad_cell_name_1> ?[<pin_name>| <pin_name> <layer_name>]?
\end{lstlisting}

GSR 关键字 ：用于设置与  相关的 RedHawk 运行参数。

\begin{lstlisting}
                         <pad_cell_name_n> ?[<pin_name>| <pin_name> <layer_name>]?
\end{lstlisting}

GSR 关键字 ：用于设置与  相关的 RedHawk 运行参数。

GSR 关键字 ：用于设置与  相关的 RedHawk 运行参数。

\begin{lstlisting}
                        <master cell name>
                        <x source loc> <y source loc> <layer> <P/G pad type>
\end{lstlisting}

GSR 关键字 ：用于设置与  相关的 RedHawk 运行参数。

GSR 关键字 ：用于设置与  相关的 RedHawk 运行参数。

\begin{lstlisting}
                      <master cell name>: specifies the master cell name (LEF macro name) indicating
\end{lstlisting}

GSR 关键字 ：用于设置与  相关的 RedHawk 运行参数。

\begin{lstlisting}
                      <x source loc> <y source loc> : x,y location of the master cell (referred to the
\end{lstlisting}

GSR 关键字 ：用于设置与  相关的 RedHawk 运行参数。

\begin{lstlisting}
                      <layer> : source layer name
\end{lstlisting}

GSR 关键字 ：用于设置与  相关的 RedHawk 运行参数。

\begin{lstlisting}
                      <P/G pad type> : the source P/G pad type
\end{lstlisting}

GSR 关键字 ：用于设置与  相关的 RedHawk 运行参数。

\begin{lstlisting}
              the instance location. A file called adsRpt/<design_name>.NEW.ploc reports all such
\end{lstlisting}

GSR 关键字 ：用于设置与  相关的 RedHawk 运行参数。

GSR 关键字 ：用于设置与  相关的 RedHawk 运行参数。

\begin{lstlisting}
              and the pins used in a <design>.pad file.
                   <pad_instance_name> ?[<pin_name> | <pin_name> <layer_name>]?
\end{lstlisting}

GSR 关键字 ：用于设置与  相关的 RedHawk 运行参数。

\begin{lstlisting}
                   <pad_instance_name> ? [<pin_name> | <pin_name> <layer_name>]?
\end{lstlisting}

GSR 关键字 ：用于设置与  相关的 RedHawk 运行参数。

\begin{lstlisting}
              file ) in a <design>.pad file:
\end{lstlisting}

GSR 关键字 ：用于设置与  相关的 RedHawk 运行参数。

GSR 关键字 ：用于设置与  相关的 RedHawk 运行参数。

\begin{lstlisting}
              <design_name>.NEW.ploc is also generated.
\end{lstlisting}

GSR 关键字 ：用于设置与  相关的 RedHawk 运行参数。

\begin{lstlisting}
              <design>.ploc file, which is formatted as follows:
              <pad_name> <X> <Y> <metal name-DEF> [POWER|GROUND] <L-pH> <R-Ohms> <C-pF>
\end{lstlisting}

GSR 关键字 ：用于设置与  相关的 RedHawk 运行参数。

GSR 关键字 ：用于设置与  相关的 RedHawk 运行参数。

GSR 关键字 ：用于设置与  相关的 RedHawk 运行参数。

GSR 关键字 ：用于设置与  相关的 RedHawk 运行参数。

GSR 关键字 ：用于设置与  相关的 RedHawk 运行参数。

\begin{lstlisting}
                        setup pad -power -r 30 -ground -r 30
                        setup package -power -r 70 -ground -r 70
                        setup wirebond -power -r 50 -ground -r 50
\end{lstlisting}

GSR 关键字 ：用于设置与  相关的 RedHawk 运行参数。

\begin{lstlisting}
                ** Pad/WireBond/Package RLC Setting Summary **
\end{lstlisting}

GSR 关键字 ：用于设置与  相关的 RedHawk 运行参数。

\begin{lstlisting}
                    *PIN<VDDC_1> will use the following R, I, C values.
\end{lstlisting}

GSR 关键字 ：用于设置与  相关的 RedHawk 运行参数。

\begin{lstlisting}
                     *PIN<VDDC_2> will use the following R, I, C values.
\end{lstlisting}

GSR 关键字 ：用于设置与  相关的 RedHawk 运行参数。

\begin{lstlisting}
                     *PIN<VSSC_1> will use the following R, I, C values.
\end{lstlisting}

GSR 关键字 ：用于设置与  相关的 RedHawk 运行参数。

GSR 关键字 ：用于设置与  相关的 RedHawk 运行参数。

GSR 关键字 ：用于设置与  相关的 RedHawk 运行参数。

\subsection{分散焊盘文件}
\subsubsection*{Individual Pad File Specification}

GSR 关键字 ：用于设置与  相关的 RedHawk 运行参数。

GSR 关键字 ：用于设置与  相关的 RedHawk 运行参数。

\begin{lstlisting}
              power and ground pads by using the Edit> Add Pad command.
              Note that you can use Edit -> Undo/Redo on Add/Delete Pad operations, and multiple
\end{lstlisting}

GSR 关键字 ：用于设置与  相关的 RedHawk 运行参数。

GSR 关键字 ：用于设置与  相关的 RedHawk 运行参数。

\begin{lstlisting}
              a <design>.pcell file:
                         <pad_cell_name_1> ?[<pin_name>| <pin_name> <layer_name>]?
\end{lstlisting}

GSR 关键字 ：用于设置与  相关的 RedHawk 运行参数。

\begin{lstlisting}
                         <pad_cell_name_n> ?[<pin_name>| <pin_name> <layer_name>]?
\end{lstlisting}

GSR 关键字 ：用于设置与  相关的 RedHawk 运行参数。

GSR 关键字 ：用于设置与  相关的 RedHawk 运行参数。

\begin{lstlisting}
                        <master cell name>
                        <x source loc> <y source loc> <layer> <P/G pad type>
\end{lstlisting}

GSR 关键字 ：用于设置与  相关的 RedHawk 运行参数。

GSR 关键字 ：用于设置与  相关的 RedHawk 运行参数。

\begin{lstlisting}
                      <master cell name>: specifies the master cell name (LEF macro name) indicating
\end{lstlisting}

GSR 关键字 ：用于设置与  相关的 RedHawk 运行参数。

\begin{lstlisting}
                      <x source loc> <y source loc> : x,y location of the master cell (referred to the
\end{lstlisting}

GSR 关键字 ：用于设置与  相关的 RedHawk 运行参数。

\begin{lstlisting}
                      <layer> : source layer name
\end{lstlisting}

GSR 关键字 ：用于设置与  相关的 RedHawk 运行参数。

\begin{lstlisting}
                      <P/G pad type> : the source P/G pad type
\end{lstlisting}

GSR 关键字 ：用于设置与  相关的 RedHawk 运行参数。

\begin{lstlisting}
              the instance location. A file called adsRpt/<design_name>.NEW.ploc reports all such
\end{lstlisting}

GSR 关键字 ：用于设置与  相关的 RedHawk 运行参数。

GSR 关键字 ：用于设置与  相关的 RedHawk 运行参数。

\begin{lstlisting}
              and the pins used in a <design>.pad file.
                   <pad_instance_name> ?[<pin_name> | <pin_name> <layer_name>]?
\end{lstlisting}

GSR 关键字 ：用于设置与  相关的 RedHawk 运行参数。

\begin{lstlisting}
                   <pad_instance_name> ? [<pin_name> | <pin_name> <layer_name>]?
\end{lstlisting}

GSR 关键字 ：用于设置与  相关的 RedHawk 运行参数。

\begin{lstlisting}
              file ) in a <design>.pad file:
\end{lstlisting}

GSR 关键字 ：用于设置与  相关的 RedHawk 运行参数。

GSR 关键字 ：用于设置与  相关的 RedHawk 运行参数。

\begin{lstlisting}
              <design_name>.NEW.ploc is also generated.
\end{lstlisting}

GSR 关键字 ：用于设置与  相关的 RedHawk 运行参数。

\begin{lstlisting}
              <design>.ploc file, which is formatted as follows:
              <pad_name> <X> <Y> <metal name-DEF> [POWER|GROUND] <L-pH> <R-Ohms> <C-pF>
\end{lstlisting}

GSR 关键字 ：用于设置与  相关的 RedHawk 运行参数。

GSR 关键字 ：用于设置与  相关的 RedHawk 运行参数。

GSR 关键字 ：用于设置与  相关的 RedHawk 运行参数。

GSR 关键字 ：用于设置与  相关的 RedHawk 运行参数。

GSR 关键字 ：用于设置与  相关的 RedHawk 运行参数。

\begin{lstlisting}
                        setup pad -power -r 30 -ground -r 30
                        setup package -power -r 70 -ground -r 70
                        setup wirebond -power -r 50 -ground -r 50
\end{lstlisting}

GSR 关键字 ：用于设置与  相关的 RedHawk 运行参数。

\begin{lstlisting}
                ** Pad/WireBond/Package RLC Setting Summary **
\end{lstlisting}

GSR 关键字 ：用于设置与  相关的 RedHawk 运行参数。

\begin{lstlisting}
                    *PIN<VDDC_1> will use the following R, I, C values.
\end{lstlisting}

GSR 关键字 ：用于设置与  相关的 RedHawk 运行参数。

\begin{lstlisting}
                     *PIN<VDDC_2> will use the following R, I, C values.
\end{lstlisting}

GSR 关键字 ：用于设置与  相关的 RedHawk 运行参数。

\begin{lstlisting}
                     *PIN<VSSC_1> will use the following R, I, C values.
\end{lstlisting}

GSR 关键字 ：用于设置与  相关的 RedHawk 运行参数。

GSR 关键字 ：用于设置与  相关的 RedHawk 运行参数。

GSR 关键字 ：用于设置与  相关的 RedHawk 运行参数。

\section{库工艺与设计网表文件}
\subsection*{Library Technology and Design Netlist Files}

\subsection{库工艺文件（LEF）}
\subsubsection*{Library Technology Files}

GSR 关键字 ：用于设置与  相关的 RedHawk 运行参数。

GSR 关键字 ：用于设置与  相关的 RedHawk 运行参数。（必需）

\begin{lstlisting}
                    *PAD_PSS
                    #<pad inst name> <P/G pin name> <package subcircuit port name>
\end{lstlisting}

GSR 关键字 ：用于设置与  相关的 RedHawk 运行参数。

\begin{lstlisting}
              The above file can have any extension and must be instantiated within the PAD_FILES { }
\end{lstlisting}

GSR 关键字 ：用于设置与  相关的 RedHawk 运行参数。

GSR 关键字 ：用于设置与  相关的 RedHawk 运行参数。

GSR 关键字 ：用于设置与  相关的 RedHawk 运行参数。

设计 网络list (DEF) 文件s

GSR 关键字 ：用于设置与  相关的 RedHawk 运行参数。

\begin{lstlisting}
                         <path>/<blockname>.def block/top
\end{lstlisting}

GSR 关键字 ：用于设置与  相关的 RedHawk 运行参数。

GSR 关键字 ：用于设置与  相关的 RedHawk 运行参数。

GSR 关键字 ：用于设置与  相关的 RedHawk 运行参数。

GSR 关键字 ：用于设置与  相关的 RedHawk 运行参数。

GSR 关键字 ：用于设置与  相关的 RedHawk 运行参数。

GSR 关键字 ：用于设置与  相关的 RedHawk 运行参数。

GSR 关键字 ：用于设置与  相关的 RedHawk 运行参数。（可选）

\begin{lstlisting}
              pgarc {
                    <vdd_pin_name> <gnd_pin_name>
\end{lstlisting}

GSR 关键字 ：用于设置与  相关的 RedHawk 运行参数。

\begin{lstlisting}
                    }
\end{lstlisting}

\begin{lstlisting}
              pin <pin_name>
                  {
                      capacitance <pF>
\end{lstlisting}

GSR 关键字 ：用于设置与  相关的 RedHawk 运行参数。

\begin{lstlisting}
                      function “<function_type>”
\end{lstlisting}

GSR 关键字 ：用于设置与  相关的 RedHawk 运行参数。

\begin{lstlisting}
                      vector “<inputPinName> <clockPinName> : <pinName> [0 | 1] ...”
                  }
\end{lstlisting}

\begin{lstlisting}
              cell <cell_name>
                  {
\end{lstlisting}

GSR 关键字 ：用于设置与  相关的 RedHawk 运行参数。

\begin{lstlisting}
                      pin <pin_name>
                      {
                           capacitance <pF>
\end{lstlisting}

GSR 关键字 ：用于设置与  相关的 RedHawk 运行参数。

\begin{lstlisting}
                           function “<function_type>”
\end{lstlisting}

GSR 关键字 ：用于设置与  相关的 RedHawk 运行参数。

\begin{lstlisting}
                           vector “<inputPinName> <clockPinName> : <pinName> [0 | 1] ...”
                      }
                  }
\end{lstlisting}

\begin{lstlisting}
              library <library_name> {
                  nomVoltage <volts>
                  nomTemperature <temp>
                  pin <pin_name> {
\end{lstlisting}

GSR 关键字 ：用于设置与  相关的 RedHawk 运行参数。

\begin{lstlisting}
                      }
                  }
\end{lstlisting}

GSR 关键字 ：用于设置与  相关的 RedHawk 运行参数。

\begin{lstlisting}
                  {
\end{lstlisting}

引脚 CK

\begin{lstlisting}
                       {
\end{lstlisting}

GSR 关键字 ：用于设置与  相关的 RedHawk 运行参数。

GSR 关键字 ：用于设置与  相关的 RedHawk 运行参数。

\begin{lstlisting}
                        }
                   }
\end{lstlisting}

GSR 关键字 ：用于设置与  相关的 RedHawk 运行参数。

\begin{lstlisting}
                  cell <cell_name>
                        {
                        type <comb | latch | ff | seq | memory>
                        subtype <input_pat | output_pad | inout_pad | clockgating | mux>
                        pin <pin_name>
                        {
                            capacitance <pF> |
                            direction <input | output | inout> |
                            function <function_type> |
                            type <clock | scan> |
                            vector “<inputPinName> <clockPinName> : <pinName> [0 | 1]”
                        }
                  }
\end{lstlisting}

GSR 关键字 ：用于设置与  相关的 RedHawk 运行参数。

\begin{lstlisting}
                  {
\end{lstlisting}

引脚 CK

\begin{lstlisting}
                       {
\end{lstlisting}

GSR 关键字 ：用于设置与  相关的 RedHawk 运行参数。

\begin{lstlisting}
                       }
                  }
\end{lstlisting}

GSR 关键字 ：用于设置与  相关的 RedHawk 运行参数。

\begin{lstlisting}
                  {
\end{lstlisting}

引脚 C

\begin{lstlisting}
                       {
\end{lstlisting}

GSR 关键字 ：用于设置与  相关的 RedHawk 运行参数。

\begin{lstlisting}
                       }
                  }
\end{lstlisting}

GSR 关键字 ：用于设置与  相关的 RedHawk 运行参数。

\begin{lstlisting}
                  pin <pin_name>
                  {
                        capacitance <pF> |
                        direction <input | output | inout> |
                        function <function_type> |
                        type <clock | scan> |
                        vector “<inputPinName> <clockPinName> : <pinName> [0 | 1]”
                  }
\end{lstlisting}

GSR 关键字 ：用于设置与  相关的 RedHawk 运行参数。

\begin{lstlisting}
                  {
\end{lstlisting}

GSR 关键字 ：用于设置与  相关的 RedHawk 运行参数。

GSR 关键字 ：用于设置与  相关的 RedHawk 运行参数。

\begin{lstlisting}
                   }
\end{lstlisting}

GSR 关键字 ：用于设置与  相关的 RedHawk 运行参数。

\begin{lstlisting}
              pgarc {          <- this statement covers all cells
\end{lstlisting}

GSR 关键字 ：用于设置与  相关的 RedHawk 运行参数。

\begin{lstlisting}
                }
              library LibABC {
                  cell ram_mvdd { <- this statement is cell-specific
                     pgarc {
\end{lstlisting}

GSR 关键字 ：用于设置与  相关的 RedHawk 运行参数。

\begin{lstlisting}
                     }
                  }
              }
\end{lstlisting}

GSR 关键字 ：用于设置与  相关的 RedHawk 运行参数。

GSR 关键字 ：用于设置与  相关的 RedHawk 运行参数。

GSR 关键字 ：用于设置与  相关的 RedHawk 运行参数。

\begin{lstlisting}
                         ###### Apache Design Solutions #####
                         # File Type                     : Timing Information
                         # Design Name              : <design_name>
\end{lstlisting}

GSR 关键字 ：用于设置与  相关的 RedHawk 运行参数。

\begin{lstlisting}
                         #################################
\end{lstlisting}

\begin{lstlisting}
                          # bus_naming_style <bus_open_close_chars>
                          # hierarchy_separator <separator_char>
                          # pin_delimiter <delimiter_char>
\end{lstlisting}

GSR 关键字 ：用于设置与  相关的 RedHawk 运行参数。

\begin{lstlisting}
                       CLOCK <rise> <fall> <period> <root> <idx>
\end{lstlisting}

GSR 关键字 ：用于设置与  相关的 RedHawk 运行参数。

GSR 关键字 ：用于设置与  相关的 RedHawk 运行参数。

GSR 关键字 ：用于设置与  相关的 RedHawk 运行参数。

\begin{lstlisting}
                        #CLOCK <clk_name> INVALID_PERIOD_OR_WAVEFORM
\end{lstlisting}

GSR 关键字 ：用于设置与  相关的 RedHawk 运行参数。

\begin{lstlisting}
                        #TIMESCALE 1e<digits>
\end{lstlisting}

GSR 关键字 ：用于设置与  相关的 RedHawk 运行参数。

\begin{lstlisting}
                      #TIMESCALE 1e-9
\end{lstlisting}

GSR 关键字 ：用于设置与  相关的 RedHawk 运行参数。

GSR 关键字 ：用于设置与  相关的 RedHawk 运行参数。

\begin{lstlisting}
                         #NAMEMAP (total_entries)
                         <id> <instance_hierarchy>
\end{lstlisting}

GSR 关键字 ：用于设置与  相关的 RedHawk 运行参数。

\begin{lstlisting}
                         #END NAMEMAP
\end{lstlisting}

GSR 关键字 ：用于设置与  相关的 RedHawk 运行参数。

\begin{lstlisting}
                      #NAMEMAP (2124)
\end{lstlisting}

GSR 关键字 ：用于设置与  相关的 RedHawk 运行参数。

\begin{lstlisting}
                      #END NAMEMAP
\end{lstlisting}

GSR 关键字 ：用于设置与  相关的 RedHawk 运行参数。

\begin{lstlisting}
                   #INSTANCE (total_instances)
                   I $<id>/<instance_name>
                   S <pin_name> <min_r_slew> <max_r_slew> <min_f_slew> <max_f_slew>
\end{lstlisting}

GSR 关键字 ：用于设置与  相关的 RedHawk 运行参数。

\begin{lstlisting}
                   L <pin_name> <C1> <R> <C2> <type> <Cpin>
                   C <pin_name> <1|0>
\end{lstlisting}

GSR 关键字 ：用于设置与  相关的 RedHawk 运行参数。

\begin{lstlisting}
                   T <pin_name> <is_clock> <min_r> <max_r> <min_f> <max_f> <clk_idx>
\end{lstlisting}

GSR 关键字 ：用于设置与  相关的 RedHawk 运行参数。

\begin{lstlisting}
                   D <pin_name> <is_clock> <min_r> <max_r> <min_f> <max_f> <clk_idx>
\end{lstlisting}

GSR 关键字 ：用于设置与  相关的 RedHawk 运行参数。

\begin{lstlisting}
                   #END INSTANCE
\end{lstlisting}

GSR 关键字 ：用于设置与  相关的 RedHawk 运行参数。

\begin{lstlisting}
                            do not have hierarchy, and are not prefixed with $<id>/.
\end{lstlisting}

GSR 关键字 ：用于设置与  相关的 RedHawk 运行参数。

GSR 关键字 ：用于设置与  相关的 RedHawk 运行参数。

\begin{lstlisting}
                      <type> is an integer that describes the signal load, with these values:
\end{lstlisting}

GSR 关键字 ：用于设置与  相关的 RedHawk 运行参数。

\begin{lstlisting}
                      <is_clock= 1 or 0> : indicates whether the pin is the output of an instance along
\end{lstlisting}

GSR 关键字 ：用于设置与  相关的 RedHawk 运行参数。

\begin{lstlisting}
                      <clk_idx> : specifies the index of the associated clock listed in the CLOCKS
\end{lstlisting}

GSR 关键字 ：用于设置与  相关的 RedHawk 运行参数。

GSR 关键字 ：用于设置与  相关的 RedHawk 运行参数。

\begin{lstlisting}
                      #INSTANCE (177763)
\end{lstlisting}

GSR 关键字 ：用于设置与  相关的 RedHawk 运行参数。

\begin{lstlisting}
                      #END INSTANCE
\end{lstlisting}

GSR 关键字 ：用于设置与  相关的 RedHawk 运行参数。

GSR 关键字 ：用于设置与  相关的 RedHawk 运行参数。

GSR 关键字 ：用于设置与  相关的 RedHawk 运行参数。

GSR 关键字 ：用于设置与  相关的 RedHawk 运行参数。

GSR 关键字 ：用于设置与  相关的 RedHawk 运行参数。

\begin{lstlisting}
                        <pin_name> L <C1> <R> <C2> <type> <Cpin>
\end{lstlisting}

GSR 关键字 ：用于设置与  相关的 RedHawk 运行参数。

\begin{lstlisting}
                   <pin_name> SL <min_r_slew> <max_r_slew> <min_f_slew> <max_f_slew>
\end{lstlisting}

GSR 关键字 ：用于设置与  相关的 RedHawk 运行参数。

\begin{lstlisting}
                   <pin_name> TW <is_clock> <min_r> <max_r> <min_f> <max_f> <clk_idx>
\end{lstlisting}

GSR 关键字 ：用于设置与  相关的 RedHawk 运行参数。

\begin{lstlisting}
                   <pin_name> IDEL <is_clock> <min_r> <max_r> <min_f> <max_f> <clk_idx>
\end{lstlisting}

GSR 关键字 ：用于设置与  相关的 RedHawk 运行参数。

\begin{lstlisting}
                   <pin_name> CONST [1|0]
\end{lstlisting}

GSR 关键字 ：用于设置与  相关的 RedHawk 运行参数。

\begin{lstlisting}
                   #<pin_name> DANGLING
\end{lstlisting}

GSR 关键字 ：用于设置与  相关的 RedHawk 运行参数。

\begin{lstlisting}
                   #<pin_name> NO_TW {}
\end{lstlisting}

GSR 关键字 ：用于设置与  相关的 RedHawk 运行参数。

\begin{lstlisting}
                   #<pin_name> NO_VALID_CLOCK
\end{lstlisting}

GSR 关键字 ：用于设置与  相关的 RedHawk 运行参数。

GSR 关键字 ：用于设置与  相关的 RedHawk 运行参数。

\subsection{设计网表文件（DEF）}
\subsubsection*{Design Netlist Files}

\subsection{Synopsys 库文件（LIB）}
\subsubsection*{Synopsys Library Files}

GSR 关键字 ：用于设置与  相关的 RedHawk 运行参数。

GSR 关键字 ：用于设置与  相关的 RedHawk 运行参数。

GSR 关键字 ：用于设置与  相关的 RedHawk 运行参数。

GSR 关键字 ：用于设置与  相关的 RedHawk 运行参数。

GSR 关键字 ：用于设置与  相关的 RedHawk 运行参数。（可选）

\begin{lstlisting}
              pgarc {
                    <vdd_pin_name> <gnd_pin_name>
\end{lstlisting}

GSR 关键字 ：用于设置与  相关的 RedHawk 运行参数。

\begin{lstlisting}
                    }
\end{lstlisting}

\begin{lstlisting}
              pin <pin_name>
                  {
                      capacitance <pF>
\end{lstlisting}

GSR 关键字 ：用于设置与  相关的 RedHawk 运行参数。

\begin{lstlisting}
                      function “<function_type>”
\end{lstlisting}

GSR 关键字 ：用于设置与  相关的 RedHawk 运行参数。

\begin{lstlisting}
                      vector “<inputPinName> <clockPinName> : <pinName> [0 | 1] ...”
                  }
\end{lstlisting}

\begin{lstlisting}
              cell <cell_name>
                  {
\end{lstlisting}

GSR 关键字 ：用于设置与  相关的 RedHawk 运行参数。

\begin{lstlisting}
                      pin <pin_name>
                      {
                           capacitance <pF>
\end{lstlisting}

GSR 关键字 ：用于设置与  相关的 RedHawk 运行参数。

\begin{lstlisting}
                           function “<function_type>”
\end{lstlisting}

GSR 关键字 ：用于设置与  相关的 RedHawk 运行参数。

\begin{lstlisting}
                           vector “<inputPinName> <clockPinName> : <pinName> [0 | 1] ...”
                      }
                  }
\end{lstlisting}

\begin{lstlisting}
              library <library_name> {
                  nomVoltage <volts>
                  nomTemperature <temp>
                  pin <pin_name> {
\end{lstlisting}

GSR 关键字 ：用于设置与  相关的 RedHawk 运行参数。

\begin{lstlisting}
                      }
                  }
\end{lstlisting}

GSR 关键字 ：用于设置与  相关的 RedHawk 运行参数。

\begin{lstlisting}
                  {
\end{lstlisting}

引脚 CK

\begin{lstlisting}
                       {
\end{lstlisting}

GSR 关键字 ：用于设置与  相关的 RedHawk 运行参数。

GSR 关键字 ：用于设置与  相关的 RedHawk 运行参数。

\begin{lstlisting}
                        }
                   }
\end{lstlisting}

GSR 关键字 ：用于设置与  相关的 RedHawk 运行参数。

\begin{lstlisting}
                  cell <cell_name>
                        {
                        type <comb | latch | ff | seq | memory>
                        subtype <input_pat | output_pad | inout_pad | clockgating | mux>
                        pin <pin_name>
                        {
                            capacitance <pF> |
                            direction <input | output | inout> |
                            function <function_type> |
                            type <clock | scan> |
                            vector “<inputPinName> <clockPinName> : <pinName> [0 | 1]”
                        }
                  }
\end{lstlisting}

GSR 关键字 ：用于设置与  相关的 RedHawk 运行参数。

\begin{lstlisting}
                  {
\end{lstlisting}

引脚 CK

\begin{lstlisting}
                       {
\end{lstlisting}

GSR 关键字 ：用于设置与  相关的 RedHawk 运行参数。

\begin{lstlisting}
                       }
                  }
\end{lstlisting}

GSR 关键字 ：用于设置与  相关的 RedHawk 运行参数。

\begin{lstlisting}
                  {
\end{lstlisting}

引脚 C

\begin{lstlisting}
                       {
\end{lstlisting}

GSR 关键字 ：用于设置与  相关的 RedHawk 运行参数。

\begin{lstlisting}
                       }
                  }
\end{lstlisting}

GSR 关键字 ：用于设置与  相关的 RedHawk 运行参数。

\begin{lstlisting}
                  pin <pin_name>
                  {
                        capacitance <pF> |
                        direction <input | output | inout> |
                        function <function_type> |
                        type <clock | scan> |
                        vector “<inputPinName> <clockPinName> : <pinName> [0 | 1]”
                  }
\end{lstlisting}

GSR 关键字 ：用于设置与  相关的 RedHawk 运行参数。

\begin{lstlisting}
                  {
\end{lstlisting}

GSR 关键字 ：用于设置与  相关的 RedHawk 运行参数。

\section{时序数据文件}
\subsection*{Timing Data File}

GSR 关键字 ：用于设置与  相关的 RedHawk 运行参数。

\begin{lstlisting}
                   }
\end{lstlisting}

GSR 关键字 ：用于设置与  相关的 RedHawk 运行参数。

\begin{lstlisting}
              pgarc {          <- this statement covers all cells
\end{lstlisting}

GSR 关键字 ：用于设置与  相关的 RedHawk 运行参数。

\begin{lstlisting}
                }
              library LibABC {
                  cell ram_mvdd { <- this statement is cell-specific
                     pgarc {
\end{lstlisting}

GSR 关键字 ：用于设置与  相关的 RedHawk 运行参数。

\begin{lstlisting}
                     }
                  }
              }
\end{lstlisting}

GSR 关键字 ：用于设置与  相关的 RedHawk 运行参数。

GSR 关键字 ：用于设置与  相关的 RedHawk 运行参数。

GSR 关键字 ：用于设置与  相关的 RedHawk 运行参数。

\begin{lstlisting}
                         ###### Apache Design Solutions #####
                         # File Type                     : Timing Information
                         # Design Name              : <design_name>
\end{lstlisting}

GSR 关键字 ：用于设置与  相关的 RedHawk 运行参数。

\begin{lstlisting}
                         #################################
\end{lstlisting}

\begin{lstlisting}
                          # bus_naming_style <bus_open_close_chars>
                          # hierarchy_separator <separator_char>
                          # pin_delimiter <delimiter_char>
\end{lstlisting}

GSR 关键字 ：用于设置与  相关的 RedHawk 运行参数。

\begin{lstlisting}
                       CLOCK <rise> <fall> <period> <root> <idx>
\end{lstlisting}

GSR 关键字 ：用于设置与  相关的 RedHawk 运行参数。

GSR 关键字 ：用于设置与  相关的 RedHawk 运行参数。

GSR 关键字 ：用于设置与  相关的 RedHawk 运行参数。

\begin{lstlisting}
                        #CLOCK <clk_name> INVALID_PERIOD_OR_WAVEFORM
\end{lstlisting}

GSR 关键字 ：用于设置与  相关的 RedHawk 运行参数。

\begin{lstlisting}
                        #TIMESCALE 1e<digits>
\end{lstlisting}

GSR 关键字 ：用于设置与  相关的 RedHawk 运行参数。

\begin{lstlisting}
                      #TIMESCALE 1e-9
\end{lstlisting}

GSR 关键字 ：用于设置与  相关的 RedHawk 运行参数。

GSR 关键字 ：用于设置与  相关的 RedHawk 运行参数。

\begin{lstlisting}
                         #NAMEMAP (total_entries)
                         <id> <instance_hierarchy>
\end{lstlisting}

GSR 关键字 ：用于设置与  相关的 RedHawk 运行参数。

\begin{lstlisting}
                         #END NAMEMAP
\end{lstlisting}

GSR 关键字 ：用于设置与  相关的 RedHawk 运行参数。

\begin{lstlisting}
                      #NAMEMAP (2124)
\end{lstlisting}

GSR 关键字 ：用于设置与  相关的 RedHawk 运行参数。

\begin{lstlisting}
                      #END NAMEMAP
\end{lstlisting}

GSR 关键字 ：用于设置与  相关的 RedHawk 运行参数。

\begin{lstlisting}
                   #INSTANCE (total_instances)
                   I $<id>/<instance_name>
                   S <pin_name> <min_r_slew> <max_r_slew> <min_f_slew> <max_f_slew>
\end{lstlisting}

GSR 关键字 ：用于设置与  相关的 RedHawk 运行参数。

\begin{lstlisting}
                   L <pin_name> <C1> <R> <C2> <type> <Cpin>
                   C <pin_name> <1|0>
\end{lstlisting}

GSR 关键字 ：用于设置与  相关的 RedHawk 运行参数。

\begin{lstlisting}
                   T <pin_name> <is_clock> <min_r> <max_r> <min_f> <max_f> <clk_idx>
\end{lstlisting}

GSR 关键字 ：用于设置与  相关的 RedHawk 运行参数。

\begin{lstlisting}
                   D <pin_name> <is_clock> <min_r> <max_r> <min_f> <max_f> <clk_idx>
\end{lstlisting}

GSR 关键字 ：用于设置与  相关的 RedHawk 运行参数。

\begin{lstlisting}
                   #END INSTANCE
\end{lstlisting}

GSR 关键字 ：用于设置与  相关的 RedHawk 运行参数。

\begin{lstlisting}
                            do not have hierarchy, and are not prefixed with $<id>/.
\end{lstlisting}

GSR 关键字 ：用于设置与  相关的 RedHawk 运行参数。

GSR 关键字 ：用于设置与  相关的 RedHawk 运行参数。

\begin{lstlisting}
                      <type> is an integer that describes the signal load, with these values:
\end{lstlisting}

GSR 关键字 ：用于设置与  相关的 RedHawk 运行参数。

\begin{lstlisting}
                      <is_clock= 1 or 0> : indicates whether the pin is the output of an instance along
\end{lstlisting}

GSR 关键字 ：用于设置与  相关的 RedHawk 运行参数。

\begin{lstlisting}
                      <clk_idx> : specifies the index of the associated clock listed in the CLOCKS
\end{lstlisting}

GSR 关键字 ：用于设置与  相关的 RedHawk 运行参数。

GSR 关键字 ：用于设置与  相关的 RedHawk 运行参数。

\begin{lstlisting}
                      #INSTANCE (177763)
\end{lstlisting}

GSR 关键字 ：用于设置与  相关的 RedHawk 运行参数。

\begin{lstlisting}
                      #END INSTANCE
\end{lstlisting}

GSR 关键字 ：用于设置与  相关的 RedHawk 运行参数。

GSR 关键字 ：用于设置与  相关的 RedHawk 运行参数。

\section{结果文件}
\subsection*{Result Files}

RedHawk 在 \texttt{adsRpt}、\texttt{adsPower} 等目录下生成功耗摘要、IR/DvD 结果、
EM 报告、电流波形与各类文本/图形报告。详见第~6~章「报告」。

\begin{seeAlsoBox}
结果文件类型与命名约定见第~6~章及各分析主题章节。
\end{seeAlsoBox}
